% =====================================================================
%  LIBRO ACADÉMICO: APLICACIONES WEB - EDICIÓN AMPLIADA
%  PARTE 1: Semanas 1 a 9
% =====================================================================
\documentclass[11pt,a4paper,openany]{book}

\usepackage[utf8]{inputenc}
\usepackage[T1]{fontenc}
\usepackage[spanish,es-tabla,es-noshorthands]{babel}
\usepackage{lmodern}
\usepackage{microtype}
\usepackage[a4paper,margin=2.5cm,headheight=14.5pt]{geometry}
\usepackage{setspace}\onehalfspacing
\usepackage{parskip}

\usepackage{xcolor}
\definecolor{uteqverde}{RGB}{0,102,51}
\definecolor{uteqdorado}{RGB}{204,153,0}
\definecolor{codebg}{RGB}{248,248,248}
\definecolor{codeborder}{RGB}{180,180,180}
\definecolor{codekey}{RGB}{0,0,180}
\definecolor{codestr}{RGB}{163,21,21}
\definecolor{codecom}{RGB}{0,128,0}

\usepackage{amsmath,amssymb}
\usepackage{graphicx}
\usepackage{tikz}
\usetikzlibrary{shapes,arrows,positioning,fit,backgrounds,calc}
\usetikzlibrary{babel}

\usepackage{fancyhdr}
\pagestyle{fancy}\fancyhf{}
\fancyhead[LE,RO]{\thepage}
\fancyhead[RE]{\nouppercase{\leftmark}}
\fancyhead[LO]{\nouppercase{\rightmark}}
\renewcommand{\headrulewidth}{0.4pt}

\usepackage[most]{tcolorbox}
\tcbuselibrary{breakable,skins,listings}

\newtcolorbox{cajadefinicion}[1]{colback=uteqverde!5,colframe=uteqverde,
  fonttitle=\bfseries\color{white},title=#1,breakable,enhanced}
\newtcolorbox{cajaejemplo}[1]{colback=uteqdorado!8,colframe=uteqdorado!80!black,
  fonttitle=\bfseries\color{white},title=#1,breakable,enhanced}
\newtcolorbox{cajaejercicio}[1]{colback=blue!4,colframe=blue!60!black,
  fonttitle=\bfseries\color{white},title=#1,breakable,enhanced}
\newtcolorbox{cajaadvertencia}[1]{colback=red!5,colframe=red!70!black,
  fonttitle=\bfseries\color{white},title=#1,breakable,enhanced}
\newtcolorbox{cajahistoria}[1]{colback=blue!3,colframe=blue!50!black,
  fonttitle=\bfseries\color{white},title=#1,breakable,enhanced}
\newtcolorbox{cajabuenas}[1]{colback=green!4,colframe=uteqverde!90!black,
  fonttitle=\bfseries\color{white},title=#1,breakable,enhanced}

\usepackage{listings}
\lstdefinelanguage{HTML5}{language=HTML,sensitive=true,
  morekeywords={article,header,nav,section,footer,aside,figure,figcaption,
    main,picture,video,audio,canvas,svg,details,summary,dialog,
    DOCTYPE,html,head,body,meta,link,title,div,span,form,input,
    label,button,select,option,textarea,table,thead,tbody,tr,th,td,
    p,a,img,ul,ol,li,h1,h2,h3,h4,h5,h6,script,style,em,strong,
    template,slot,source,track,output,progress,meter,datalist,
    fieldset,legend,time,mark,abbr,address,blockquote,cite,code,kbd}}
\lstdefinelanguage{JavaScript}{
  keywords={break,case,catch,continue,debugger,default,delete,do,else,
    finally,for,function,if,in,instanceof,new,return,switch,this,throw,
    try,typeof,var,void,while,with,let,const,of,yield,async,await,
    class,extends,super,static,import,export,from,as,true,false,null,
    undefined},
  ndkeywords={Array,Boolean,Date,Error,Function,JSON,Math,Number,Object,
    Promise,RegExp,String,Symbol,console,document,window,fetch,Map,Set,
    WeakMap,WeakSet,Proxy,Reflect},
  sensitive=true,comment=[l]{//},morecomment=[s]{/*}{*/},
  morestring=[b]',morestring=[b]",morestring=[b]`}
\lstdefinelanguage{PHP}{
  keywords={abstract,and,array,as,break,callable,case,catch,class,clone,
    const,continue,declare,default,do,echo,else,elseif,empty,enddeclare,
    endfor,endforeach,endif,endswitch,endwhile,enum,extends,final,
    finally,fn,for,foreach,function,global,goto,if,implements,include,
    include_once,instanceof,insteadof,interface,isset,list,match,
    namespace,new,null,or,print,private,protected,public,readonly,
    require,require_once,return,self,static,switch,throw,trait,true,
    false,try,unset,use,var,while,xor,yield,parent},
  sensitive=true,comment=[l]{//},morecomment=[l]{\#},
  morecomment=[s]{/*}{*/},morestring=[b]',morestring=[b]"}
\lstdefinelanguage{Perl5}{
  keywords={my,our,local,sub,if,elsif,else,unless,while,until,for,
    foreach,do,last,next,redo,return,use,no,require,package,qw,q,qq,
    defined,undef,eq,ne,lt,gt,le,ge,and,or,not,xor,print,printf,die,
    warn,exit},
  sensitive=true,comment=[l]{\#},morestring=[b]',morestring=[b]"}

\lstdefinestyle{codigobase}{basicstyle=\ttfamily\footnotesize,
  backgroundcolor=\color{codebg},frame=single,rulecolor=\color{codeborder},
  numbers=left,numberstyle=\tiny\color{gray},numbersep=8pt,stepnumber=1,
  breaklines=true,breakatwhitespace=true,showstringspaces=false,tabsize=2,
  keywordstyle=\color{codekey}\bfseries,stringstyle=\color{codestr},
  commentstyle=\color{codecom}\itshape,captionpos=b,
  xleftmargin=12pt,xrightmargin=4pt}
\lstdefinestyle{html}{style=codigobase,language=HTML5}
\lstdefinestyle{sass}{style=codigobase,
  morekeywords={include,extend,mixin,function,return,if,else,each,for,while,
  import,use,forward,debug,warn,error,at-root,content,
  background,color,padding,border-radius,font-size,margin,
  display,grid,flex,box-sizing}}
\lstdefinestyle{css}{style=codigobase,
  morekeywords={body,margin,padding,color,background,font-family,
  display,grid,flex,position,absolute,relative,fixed,width,height,
  border,border-radius,box-sizing,content-box,border-box,font-size,
  line-height,grid-template-columns,grid-template-rows,grid-template-areas,
  gap,justify-content,align-items,flex-direction,min-width,max-width,
  repeat,transition,transform,animation,important,root,var,calc,clamp,
  minmax,sticky,inherit,initial,unset,revert}}
\lstdefinestyle{js}{style=codigobase,language=JavaScript}
\lstdefinestyle{php}{style=codigobase,language=PHP}
\lstdefinestyle{java}{style=codigobase,language=Java}
\lstdefinestyle{perl}{style=codigobase,language=Perl5}
\lstdefinestyle{shell}{style=codigobase,language=bash}
\lstdefinestyle{xml}{style=codigobase,language=XML}
\lstdefinestyle{sql}{style=codigobase,language=SQL}

\usepackage{array}\usepackage{booktabs}\usepackage{longtable}
\usepackage{tabularx}\usepackage{multirow}

\usepackage[hidelinks,bookmarks=true,bookmarksnumbered=true]{hyperref}
\hypersetup{colorlinks=true,linkcolor=uteqverde,citecolor=uteqverde,
  urlcolor=uteqdorado!70!black,
  pdftitle={Aplicaciones Web - Libro Completo v3 (Semanas 1-18)},
  pdfauthor={UTEQ - Ingenieria de Software}}

\usepackage{enumitem}\setlist{itemsep=2pt,topsep=2pt}
\usepackage{epigraph}\setlength{\epigraphwidth}{0.7\textwidth}

\usepackage{titlesec}
\titleformat{\chapter}[display]
  {\normalfont\Huge\bfseries\color{uteqverde}}
  {\filleft\Large Semana \thechapter}
  {1ex}{\titlerule\vspace{1ex}\filright}
\titleformat{\section}
  {\normalfont\Large\bfseries\color{uteqverde}}{\thesection}{1em}{}
\titleformat{\subsection}
  {\normalfont\large\bfseries\color{uteqverde!85!black}}{\thesubsection}{1em}{}
\titleformat{\subsubsection}
  {\normalfont\normalsize\bfseries\color{uteqdorado!90!black}}{\thesubsubsection}{1em}{}

% Definicion del lenguaje C# para listings (semana 7 ASP.NET)
\lstdefinelanguage{CSharp}{sensitive=true,
  morekeywords={abstract,as,async,await,base,bool,break,byte,case,catch,char,checked,class,const,continue,decimal,default,delegate,do,double,else,enum,event,explicit,extern,false,finally,fixed,float,for,foreach,goto,if,implicit,in,int,interface,internal,is,lock,long,namespace,new,null,object,operator,out,override,params,private,protected,public,readonly,record,ref,return,sbyte,sealed,short,sizeof,stackalloc,static,string,struct,switch,this,throw,true,try,typeof,uint,ulong,unchecked,unsafe,ushort,using,var,virtual,void,volatile,when,where,while,yield},
  morecomment=[l]{//},morecomment=[s]{/*}{*/},morestring=[b]"}
\lstdefinestyle{cs}{style=codigobase,language=CSharp}

% =====================================================================
% Modificaciones del preámbulo para semanas 10-18 (Parte 2)
% =====================================================================

% --- Nuevas bibliotecas TikZ para diagramas arquitectónicos ---
\usetikzlibrary{shapes.geometric}
\usetikzlibrary{shapes.symbols}
\usetikzlibrary{shadows}
\usetikzlibrary{matrix}
\usetikzlibrary{decorations.markings}

% --- Soporte para imágenes H ---
\usepackage{float}

% --- Soporte para \sout (texto tachado) ---
% \usepackage[normalem]{ulem} (no disponible)

% --- Soporte para \citet, \citep ---
% \usepackage{natbib} (no disponible)

% --- Nuevo entorno: caja de decisión arquitectónica (ADR) ---
\newtcolorbox{cajaadr}[1]{colback=violet!4,colframe=violet!70!black,
  fonttitle=\bfseries\color{white},title=#1,breakable,enhanced}

% --- Nuevo entorno: caja de comparativa ---
\newtcolorbox{cajacompara}[1]{colback=cyan!5,colframe=cyan!60!black,
  fonttitle=\bfseries\color{white},title=#1,breakable,enhanced}

% --- Lenguaje YAML para configuraciones ---
\lstdefinelanguage{yaml}{
  keywords={true,false,null,yes,no,on,off},
  sensitive=true,
  comment=[l]{\#},
  morestring=[b]',
  morestring=[b]"
}

% --- Lenguaje Dockerfile ---
\lstdefinelanguage{dockerfile}{
  keywords={FROM,RUN,CMD,LABEL,MAINTAINER,EXPOSE,ENV,ADD,COPY,
    ENTRYPOINT,VOLUME,USER,WORKDIR,ARG,ONBUILD,STOPSIGNAL,
    HEALTHCHECK,SHELL,AS},
  sensitive=true,
  comment=[l]{\#},
  morestring=[b]',
  morestring=[b]"
}

% --- Estilos de listing nuevos ---
\lstdefinestyle{yaml}{style=codigobase,language=yaml,
  keywordstyle=\color{codekey}\bfseries}
\lstdefinestyle{docker}{style=codigobase,language=dockerfile,
  keywordstyle=\color{codekey}\bfseries}

% --- Comandos cortos para iconos C4 ---
\newcommand{\cfourPerson}[3]{%
  \node[circle,draw=blue!70!black,fill=blue!20,thick,
        minimum size=0.7cm,font=\tiny] (#1) at #2 {};%
  \node[below=0.05cm of #1,font=\scriptsize\bfseries,align=center] {#3};%
}
\newcommand{\cfourSystem}[3]{%
  \node[rectangle,draw=blue!70!black,fill=blue!10,thick,rounded corners,
        minimum width=2.5cm,minimum height=1.2cm,
        font=\scriptsize\bfseries,align=center] (#1) at #2 {#3};%
}
\newcommand{\cfourContainer}[3]{%
  \node[rectangle,draw=uteqverde!80!black,fill=uteqverde!15,thick,
        rounded corners,minimum width=2.5cm,minimum height=1.2cm,
        font=\scriptsize\bfseries,align=center] (#1) at #2 {#3};%
}
\newcommand{\cfourDb}[3]{%
  \node[cylinder,draw=uteqverde!80!black,fill=uteqverde!15,thick,
        shape border rotate=90,aspect=0.3,minimum width=1.4cm,
        minimum height=1.5cm,font=\scriptsize\bfseries,align=center,
        cylinder uses custom fill] (#1) at #2 {#3};%
}

\begin{document}

% --------------------- Portada -----------------------------------
\begin{titlepage}
\centering\vspace*{1cm}
{\color{uteqverde}\rule{\linewidth}{2pt}}\\[0.6cm]
{\Huge\bfseries\color{uteqverde} APLICACIONES WEB}\\[0.4cm]
{\Large\itshape Edición ampliada con historia, fundamentos detallados, prácticas paso a paso y buenas prácticas profesionales}\\[0.3cm]
{\color{uteqverde}\rule{\linewidth}{1pt}}\\[1.2cm]
{\Large\bfseries Libro de texto académico}\\[0.3cm]
{\large Parte 1 Ampliada --- Semanas 1 a 9}\\[0.3cm]
{\large Unidad I: \emph{Fundamentos y Diseño Web}}\\
{\large Unidad II: \emph{Herramientas de Programación Web del Servidor}}\\[2cm]
{\large\bfseries Universidad Técnica Estatal de Quevedo}\\
{\large Facultad de Ciencias de la Computación}\\
{\large Carrera de Ingeniería de Software (Rediseño)}\\[0.7cm]
{\large\itshape Quinto nivel --- Período Académico 2026--2027 PPA}\\[1.5cm]
{\large \textbf{Autor docente:} Dr. Gleiston Cicerón Guerrero Ulloa, Ph.D.}\\[2cm]
{\color{uteqdorado}\rule{\linewidth}{2pt}}\\[0.4cm]
{\large Quevedo, Los Ríos --- Ecuador}\\
{\large Mayo de 2026}
\end{titlepage}

\thispagestyle{empty}
\vspace*{2cm}
\noindent\textbf{APLICACIONES WEB --- Parte 1 Ampliada (Semanas 1 a 9)}\\[0.3cm]
\noindent Universidad Técnica Estatal de Quevedo, 2026.\\
Facultad de Ciencias de la Computación.\\
Carrera de Ingeniería de Software (Rediseño).\\[0.4cm]
\noindent Esta edición ampliada incorpora un tratamiento profundo de
la historia del desarrollo web, fundamentos teóricos extensos,
ejercicios paso a paso plenamente reproducibles en IntelliJ IDEA
Ultimate, tablas comparativas de versiones de tecnologías y secciones
dedicadas a buenas prácticas profesionales para cada tecnología
abordada. Su contenido se ha alineado con el SWEBOK Guide v4.0
\cite{washizaki2024swebok}, las directrices del W3C
\cite{w3c2018wcag}, la guía OWASP Top 10:2021 \cite{owasp2021top10}
y las especificaciones HTML Living Standard del WHATWG.\\[0.4cm]
\noindent\textbf{Uso académico no comercial.} Se permite la
reproducción total o parcial citando la fuente.
\newpage

\chapter*{Prefacio a la edición ampliada}
\addcontentsline{toc}{chapter}{Prefacio a la edición ampliada}

La presente edición ampliada de la Parte 1 del libro \emph{Aplicaciones
Web} responde a una reflexión pedagógica concreta: el desarrollo web
moderno no se domina memorizando sintaxis, se domina comprendiendo el
\textbf{por qué histórico} de cada decisión técnica, las
\textbf{razones científicas} detrás de cada práctica y los
\textbf{caminos prácticos} que llevan del concepto al artefacto
funcional.

A diferencia de la edición original, esta ampliación incorpora cuatro
elementos pedagógicos sistemáticos en cada capítulo:

\begin{enumerate}
\item \textbf{Cajas de historia} que sitúan cada tecnología en su
contexto histórico, mencionando autores, fechas, hitos y motivaciones
originales que rara vez aparecen en los manuales técnicos pero que
resultan indispensables para comprender por qué las cosas son como son.
\item \textbf{Tablas de versiones} que documentan la evolución de cada
tecnología con las mejoras sustantivas que cada revisión aportó.
\item \textbf{Cajas de buenas prácticas} que destilan en pautas
accionables el conocimiento profesional acumulado por la industria,
frecuentemente disperso en blogs y discusiones técnicas.
\item \textbf{Ejercicios paso a paso ampliados} con instrucciones tan
detalladas que cualquier estudiante con IntelliJ IDEA Ultimate puede
reproducir el resultado exacto descrito en el texto. Cada paso explica
\emph{qué} se hace, \emph{por qué} se hace y \emph{cómo verificar} que
el resultado es correcto.
\end{enumerate}

La elección de IntelliJ IDEA Ultimate como entorno de referencia
obedece a su soporte nativo de primera clase para HTML5, CSS3,
JavaScript moderno (ES2024), Sass, PHP, Perl, Java, Spring Boot y los
servidores Tomcat/Apache; a la integración nativa de Git, Maven,
Docker y bases de datos relacionales sin necesidad de plugins externos;
y a la disponibilidad de licencia educativa gratuita para estudiantes
universitarios.

\hfill\textit{El autor docente.}\\
\hfill Quevedo, mayo de 2026.

\tableofcontents
\newpage

% =====================================================================
%                     SEMANA 1: ENCUADRE Y FUNDAMENTOS
% =====================================================================

%%% ==== SEMANA 01 v3 ====
% =====================================================================
%   SEMANA 1: ENCUADRE Y FUNDAMENTOS — VERSIÓN MEJORADA v3
% =====================================================================
\chapter{Encuadre académico y fundamentos de las aplicaciones web}
\label{cap:semana1}

\epigraph{«La web no es un medio: es un ecosistema. Comprender sus
fundamentos es comprender la infraestructura digital de la civilización
contemporánea.»}{--- \textit{Tim Berners-Lee, inventor de la World Wide Web}}

\section*{Resultado de aprendizaje de la semana}
Al concluir la semana 1, el estudiante construye interfaces web
semánticas, accesibles y responsivas \textbf{en su nivel introductorio},
integrando HTML5, CSS3 y JavaScript para crear aplicaciones del lado
del cliente que cumplan con los estándares del W3C y los principios del
diseño web inclusivo (WCAG 2.1) \cite{w3c2018wcag}.

\section{Encuadre académico de la asignatura}

Antes de adentrarnos en los aspectos técnicos del desarrollo web,
conviene situar la asignatura dentro del plan de estudios y aclarar
los acuerdos que regirán el trabajo durante las dieciocho semanas.
Esta sección presenta tres aspectos institucionales fundamentales: la
ubicación curricular de la materia y los conocimientos previos que se
asumen; las modalidades formales de evaluación según la norma
académica de la UTEQ; y las políticas profesionales que se aplicarán
en clase (asistencia, entregas, integridad académica). La claridad
sobre estos acuerdos al inicio del curso evita malentendidos
posteriores y permite al estudiante planificar su esfuerzo con
realismo.

\subsection{Ubicación curricular y prerrequisitos}

La asignatura Aplicaciones Web pertenece al campo de estudio de
\emph{Diseño de Software (DES.dd)} en la Carrera de Ingeniería de
Software, ubicada en el quinto período académico de la malla curricular
vigente dentro de la Unidad de Organización Curricular Profesional.
Posee tres créditos académicos distribuidos en 144 horas totales, de
las cuales 48 corresponden a clases presenciales teóricas y prácticas,
48 al aprendizaje autónomo y 48 al contacto directo con el docente.

Sus prerrequisitos formales son \emph{Base de Datos} y \emph{Estructura
de Datos}, ambos de tercer nivel. El primero garantiza que el
estudiante maneje el modelo relacional y el lenguaje SQL como insumo
indispensable para el acceso a datos del bloque servidor; el segundo
asegura el dominio de listas, árboles, mapas hash y la complejidad
algorítmica $O(n\log n)$ y $O(n^2)$, elementos sin los cuales no es
posible diseñar APIs eficientes ni razonar sobre escalabilidad. La
ausencia de cualquiera de estos prerrequisitos compromete severamente
el aprovechamiento del curso.

\subsection{Articulación con el SWEBOK Guide v4.0}

El \emph{Software Engineering Body of Knowledge} en su cuarta versión
\cite{washizaki2024swebok} reconoce explícitamente al desarrollo web
como una de las áreas transversales de la práctica profesional. La
asignatura articula los siguientes núcleos del SWEBOK: requisitos de
software (KA02), diseño de software (KA03), construcción de software
(KA04), calidad del software (KA10) y seguridad de software (KA15). El
estudiante no construye únicamente código que «funciona»; construye
software que satisface estándares internacionales de calidad,
seguridad y accesibilidad.

\subsection{Competencias por desarrollar}

Esta asignatura forma al estudiante en cuatro frentes complementarios:
la \textbf{capacidad técnica} de implementar interfaces y servicios
web modernos; el \textbf{razonamiento arquitectónico} para tomar
decisiones de diseño justificadas; la \textbf{conciencia de calidad}
(seguridad, rendimiento, accesibilidad); y la \textbf{disciplina
profesional} (versionado, pruebas, documentación). La caja siguiente
recoge la formulación oficial de la competencia general.


\begin{cajadefinicion}{Competencia general de la asignatura}
Analizar, diseñar, implementar y evaluar aplicaciones web mediante
tecnologías de diseño, lenguajes de programación cliente y servidor,
patrones de arquitectura y servicios web, aplicando estándares de
calidad, accesibilidad e inclusión, con el fin de construir
\emph{soluciones informáticas robustas, escalables y pertinentes} que
respondan a las necesidades productivas, sociales y tecnológicas del
Ecuador, de la región y del mundo.
\end{cajadefinicion}

Esta competencia general se descompone en cuatro resultados de
aprendizaje unitarios (RAU):

\begin{description}[leftmargin=1.5cm]
\item[RAU 1.] Construir interfaces web semánticas, accesibles y
responsivas con HTML5, CSS3 y JavaScript siguiendo W3C y WCAG 2.1.
\item[RAU 2.] Desarrollar aplicaciones web dinámicas del lado del
servidor con PERL, PHP, ASP.NET y Java/JSP, implementando lógica de
negocio, procesamiento de formularios y mecanismos de seguridad.
\item[RAU 3.] Diseñar e implementar soluciones web que integren gestión
de estado, acceso a datos relacionales mediante SQL parametrizado y
ORM, y patrones arquitectónicos escalables (multicapa, SOA, EDA, P2P).
\item[RAU 4.] Construir una aplicación web completa bajo el patrón MVC
con un \emph{framework} moderno, exponiendo e integrando servicios web
REST y SOAP, y aplicando criterios de calidad, seguridad e
interoperabilidad.
\end{description}

\subsection{Sistema de evaluación}

La calificación final se construye a partir de tres tipos de
actividades ponderadas según la Norma Académica de la UTEQ:
\emph{Gestión del Aprendizaje (GA)}, \emph{Tareas y Proyectos
Autónomos (TA)} y \emph{Evaluaciones (EV)}. La asignatura organiza dos
cortes evaluativos (semanas 9 y 17) y un Proyecto Fin de Curso
entregado en tres hitos parciales (semanas 5, 9 y 12) y un sistema
final completo en la semana 16.

\section{Historia de las aplicaciones web}

Comprender la web actual exige conocer su origen: las decisiones
arquitectónicas tomadas en los años noventa (HTTP sin estado, URLs
opacas, hipertexto distribuido) condicionan todavía hoy la forma en
que diseñamos aplicaciones. La siguiente caja presenta los hitos
fundacionales que conviene tener presente como contexto histórico
permanente.


\begin{cajahistoria}{Los orígenes --- del hipertexto a la World Wide Web}
La idea del hipertexto precede a la web por décadas. Vannevar Bush
concibió en 1945 el sistema \emph{Memex} en el artículo \emph{«As We
May Think»} publicado en \emph{The Atlantic Monthly}, una máquina
capaz de almacenar libros y vincularlos mediante asociaciones
siguiendo el modo en que opera la mente humana \cite{bush1945memex}.
Ted Nelson acuñó el término \emph{hipertexto} en 1965 dentro de su
proyecto Xanadu, que aspiraba a crear un repositorio universal de
documentos enlazados \cite{nelson1965xanadu}.

El paso decisivo lo dio Tim Berners-Lee en marzo de 1989 cuando, en el
laboratorio del CERN en Ginebra, presentó el documento titulado
\emph{«Information Management: A Proposal»} \cite{bernerslee1989}. Su
intención era resolver un problema interno: la fuga de información
cuando los investigadores abandonaban el centro. La propuesta combinó
tres ingredientes existentes (el hipertexto, los protocolos de internet
y el sistema de nombres de dominio) en un sistema nuevo: la World Wide
Web.

A finales de 1990, Berners-Lee había desarrollado en una computadora
NeXT los tres pilares fundacionales de la web: el primer servidor
(\texttt{CERN httpd}), el primer navegador (\texttt{WorldWideWeb},
luego renombrado \texttt{Nexus} para evitar confusión con la red) y el
primer documento HTML. El 6 de agosto de 1991 se publicó el primer
sitio web público, alojado en \texttt{info.cern.ch}, que aún puede
visitarse en su versión preservada en
\url{http://info.cern.ch/hypertext/WWW/TheProject.html}
\cite{cern1991firstsite}.
\end{cajahistoria}

\subsection{Las tres eras de la web}

La World Wide Web no es un fenómeno monolítico: ha atravesado tres
grandes eras evolutivas que reflejan tanto cambios tecnológicos como
transformaciones sociales profundas en cómo las personas interactúan
con la información digital. Comprender estas eras es indispensable
para situar el desarrollo web contemporáneo dentro de una trayectoria
histórica significativa, y para anticipar hacia dónde se dirige.

Cada era se caracteriza por una combinación distintiva de:
\emph{(i)} el modelo de participación dominante (quién produce y quién
consume el contenido), \emph{(ii)} las tecnologías facilitadoras que
hicieron posible esa participación, \emph{(iii)} los modelos de
negocio que emergieron y \emph{(iv)} las consecuencias sociales y
políticas observables. Berners-Lee, en su revisión retrospectiva de
2019, describió esta evolución como un «doloroso aprendizaje»
porque ninguna era ha estado exenta de patologías propias: la Web 1.0
fue limitada en participación, la Web 2.0 concentró poder en
plataformas, y la Web 3.0 enfrenta el desafío de la
descentralización genuina \cite{bernerslee2019contract}.

\subsubsection{Web 1.0 --- la web de los documentos (1989--2004)}

Esta primera era se caracterizó por el modelo \emph{de solo lectura}
(\emph{read-only web}): una minoría de autores con conocimientos
técnicos publicaba contenido, y la mayoría lo consumía pasivamente.
Las páginas eran fundamentalmente estáticas, escritas a mano en HTML
sin interactividad significativa. O'Reilly y Battelle
\cite{oreilly2009web2} caracterizaron esta era como la del «sitio web
folleto»: una transposición digital del medio impreso.

Los navegadores dominantes evolucionaron desde el pionero Mosaic
(1993), liderado por Marc Andreessen y Eric Bina en el NCSA, que fue
el primero en mostrar imágenes \emph{inline} junto al texto, hasta
Netscape Navigator (1994), que popularizó la web entre el público
general. La respuesta de Microsoft fue Internet Explorer (1995),
desencadenando la \emph{primera guerra de los navegadores} que duró
hasta aproximadamente 2003 con la victoria temporal de IE6 (alcanzó
95\% de cuota de mercado en 2003) y la consiguiente década de
estancamiento en estándares web \cite{zeldman2010designing}.

\subsubsection{Web 2.0 --- la web social (2004--2014)}

El término \emph{Web 2.0} lo popularizó Tim O'Reilly durante una
conferencia homónima en octubre de 2004 \cite{oreilly2005what}. Su
característica distintiva: la \emph{participación}. Los usuarios
pasaron de consumidores a productores de contenido mediante blogs,
wikis, redes sociales y plataformas de video. O'Reilly resumió la
filosofía como «la web como plataforma», donde el valor reside en los
datos generados por usuarios y los efectos de red.

Tecnológicamente esta era se sostuvo en cuatro pilares:

\begin{itemize}
\item \textbf{AJAX} (Asynchronous JavaScript And XML), formalizado por
Jesse James Garrett en su artículo seminal de febrero de 2005
\cite{garrett2005ajax}, que permitió actualizar partes de una página
sin recargarla. Google Maps (lanzado ese mismo mes) fue la
demostración pública más espectacular del paradigma.
\item \textbf{Frameworks JavaScript de primera generación}: Prototype
(2005), MooTools (2006) y especialmente jQuery (2006), creado por
John Resig, que para 2013 estaba presente en más del 60\% de los
sitios web del mundo.
\item \textbf{Plataformas sociales}: Facebook (2004), YouTube (2005),
Twitter (2006), GitHub (2008), Instagram (2010), Pinterest (2010).
\item \textbf{APIs públicas}: Google Maps API (2005), Twitter API
(2006), Flickr API, abriendo el camino a los \emph{mashups} y al
ecosistema de aplicaciones de terceros.
\end{itemize}

Esta era también vio surgir las críticas: van Dijck
\cite{vandijck2013culture} analizó cómo las «plataformas sociales»
construyeron una nueva forma de capitalismo informacional basada en
la captura y monetización de datos personales, fenómeno que Zuboff
\cite{zuboff2019capitalismo} bautizó como \emph{capitalismo de
vigilancia}.

\subsubsection{Web 3.0 --- la web semántica, descentralizada e inteligente (2014--presente)}

La tercera era convive con tres significados parcialmente solapados:

\begin{enumerate}
\item La \emph{Semantic Web} originalmente propuesta por Berners-Lee,
Hendler y Lassila en \emph{Scientific American} en 2001
\cite{bernerslee2001semantic}, basada en RDF, OWL y SPARQL, que busca
representar el significado de la información para que las máquinas
razonen sobre ella. Su realización plena ha sido más lenta de lo
previsto, pero subyace en proyectos como \emph{schema.org} (consorcio
de Google, Microsoft, Yahoo y Yandex desde 2011), Wikidata y los
\emph{Knowledge Graphs} de los buscadores modernos.

\item La \emph{web descentralizada} (Web3 en sentido estricto),
emergente desde 2017, que reposa sobre blockchain, IPFS y contratos
inteligentes. Wood \cite{wood2014ethereum} acuñó el término con la
publicación del whitepaper de Ethereum. Sus aplicaciones (dApps)
incluyen finanzas descentralizadas (DeFi), tokens no fungibles (NFT)
y identidad soberana.

\item La \emph{web inteligente} potenciada por IA generativa,
emergente desde 2023 con ChatGPT, Claude y GitHub Copilot. Sus
implicaciones en el desarrollo web son tan profundas que algunos
autores hablan ya de una \emph{cuarta era} en gestación
\cite{kreutzer2024ai}.
\end{enumerate}

En paralelo, las aplicaciones SPA (Single Page Applications) y PWA
(Progressive Web Applications) desplazan progresivamente al modelo
\emph{server-rendered} clásico, y la integración de IA generativa
(2023+) está redefiniendo nuevamente el panorama del desarrollo
profesional.

\subsection{Evolución temporal de los hitos web}

La tabla \ref{tab:hitos-web} sintetiza los hitos cronológicos más
relevantes de la trayectoria de la web. No se trata de una lista
exhaustiva sino de los eventos que cualquier profesional del
desarrollo web debe conocer porque cada uno de ellos transformó la
manera de construir software para la web. Los hitos están organizados
en cinco grandes movimientos: \emph{(a)} la fundación
(1989--1995), \emph{(b)} la estandarización (1995--2004), \emph{(c)}
la era social (2004--2014), \emph{(d)} la madurez técnica
(2014--2022) y \emph{(e)} la era de la inteligencia generativa
(2023--presente).

\begin{table}[!htb]
\centering
\caption{Hitos seleccionados de la historia de las aplicaciones web.}
\label{tab:hitos-web}
\footnotesize
\begin{tabularx}{\textwidth}{lX}
\toprule
\textbf{Año} & \textbf{Hito}\\
\midrule
\multicolumn{2}{l}{\emph{(a) Fundación}}\\
1989 & Tim Berners-Lee propone la WWW en el CERN.\\
1990 & Primer servidor (\texttt{CERN httpd}) y navegador (\texttt{WorldWideWeb}).\\
1991 & Publicación del primer sitio web (\texttt{info.cern.ch}).\\
1993 & Mosaic populariza la navegación con imágenes inline.\\
1993 & Especificación CGI 1.0 para programación del servidor.\\
1994 & Fundación del W3C en MIT; nace Netscape Navigator.\\
\midrule
\multicolumn{2}{l}{\emph{(b) Estandarización}}\\
1995 & Brendan Eich crea JavaScript en 10 días para Netscape.\\
1995 & PHP/FI (Personal Home Page Tools) de Rasmus Lerdorf.\\
1996 & CSS1 publicado como recomendación W3C.\\
1996 & Macromedia Flash y los \emph{cookies} de Netscape.\\
1998 & XML 1.0 estandarizado; nace Google.\\
1999 & ECMAScript 3, base estable de JavaScript por más de una década.\\
2003 & WordPress y MediaWiki marcan la era de los CMS.\\
\midrule
\multicolumn{2}{l}{\emph{(c) Era social}}\\
2004 & Web 2.0 (O'Reilly); nace Facebook.\\
2005 & AJAX formalizado; lanzamiento de YouTube; Google Maps.\\
2006 & jQuery; Twitter; Amazon Web Services (S3, EC2).\\
2008 & Google Chrome con motor V8; Android SDK 1.0.\\
2009 & HTML5 borrador funcional WHATWG; Node.js (Ryan Dahl).\\
2010 & Responsive Web Design (Ethan Marcotte).\\
2011 & Laravel (Taylor Otwell); React (Jordan Walke, 2013).\\
\midrule
\multicolumn{2}{l}{\emph{(d) Madurez técnica}}\\
2014 & HTML5 oficial W3C; ECMAScript 6 anunciado (publicado en 2015).\\
2015 & ES6 (ES2015); Angular 2 reescrito desde AngularJS.\\
2018 & GDPR; Web Components estables; PWA en producción.\\
2019 & WebAssembly 1.0 estandarizado.\\
2020 & .NET 5 unifica .NET Framework y .NET Core.\\
2021 & OWASP Top 10:2021; HTTP/3 (RFC 9114).\\
2022 & ECMAScript 2022 (ES13); Angular 14 con standalone components.\\
\midrule
\multicolumn{2}{l}{\emph{(e) Era de la IA generativa}}\\
2023 & GitHub Copilot, ChatGPT, Claude transforman el desarrollo.\\
2024 & SWEBOK Guide v4.0; PHP 8.3; Spring Boot 3.2; .NET 8 LTS.\\
2026 & ECMAScript 2025; HTML Living Standard al día; IA agentic en IDEs.\\
\bottomrule
\end{tabularx}
\end{table}

\section{Fundamentos de las aplicaciones web}

Los fundamentos de las aplicaciones web no son meros tecnicismos
introductorios: constituyen el marco conceptual sobre el que se
edifica todo el resto del curso. Sin una comprensión sólida de qué
es exactamente una aplicación web, cómo se distingue de otros tipos
de software, qué arquitecturas la sostienen y cuáles son sus ventajas
y limitaciones, el aprendizaje posterior se vuelve mecánico y
superficial. Como observa Shklar y Rosen \cite{shklar2009web}, las
aplicaciones web son hoy «el género dominante del software»: la
mayoría de las personas interactúa más horas al día con aplicaciones
web (Gmail, redes sociales, banca en línea, sistemas educativos) que
con cualquier otro tipo de software.

Esta sección establece las definiciones operativas, la tipología y
los modelos arquitectónicos esenciales que el estudiante usará como
referencia durante las 18 semanas del curso.

\subsection{Definición y clasificación de aplicaciones web}

Una \textbf{aplicación web} es un programa de software cuya lógica se
ejecuta en uno o varios servidores remotos y cuya interfaz de usuario
se distribuye al cliente a través del Hypertext Transfer Protocol
(HTTP) y se renderiza en un navegador o agente de usuario. A diferencia
de una aplicación de escritorio, no requiere instalación local; a
diferencia de una página web estática, su contenido se genera
dinámicamente en respuesta a las acciones del usuario y al estado del
servidor \cite{shklar2009web}.

Siguiendo la clasificación de Ganzábal García
\cite{deitel2017java} y la actualización de Nixon
\cite{nixon2021learning}, las aplicaciones web modernas se agrupan en
seis tipologías principales que conviene distinguir:

\begin{enumerate}
\item \textbf{Aplicaciones web estáticas.} Sirven HTML, CSS y
JavaScript pre-compilado sin procesamiento de servidor más allá del
envío de archivos. Su contenido cambia solo cuando el desarrollador
publica una nueva versión. Ejemplo canónico: un currículum vitae
publicado en GitHub Pages, o un blog personal generado con Jekyll o
Hugo. \textbf{Ventaja:} rendimiento extremo (respuesta en
milisegundos desde un CDN). \textbf{Limitación:} incapacidad de
personalización por usuario.

\item \textbf{Aplicaciones web dinámicas.} Generan el HTML en el
servidor en cada petición. La lógica reside en lenguajes como PHP,
Java, Python, Ruby o C\#. Cada usuario ve una versión personalizada
del contenido según su sesión, sus permisos, los datos almacenados en
bases de datos. \textbf{Ejemplo:} el portal de matrícula de la UTEQ;
un panel de administración construido con Laravel; cualquier sitio
WordPress.

\item \textbf{Aplicaciones de página única (SPA).} Cargan un único
HTML inicial y a partir de ahí toda la interacción ocurre vía
JavaScript que consume APIs REST o GraphQL. La navegación entre
«páginas» ocurre sin recargar el navegador, gracias a la History API.
\textbf{Ejemplos:} Gmail, Trello, Notion. Frameworks habituales en
2026: React 19, Angular 19, Vue 3.5, Svelte 5.

\item \textbf{Aplicaciones web progresivas (PWA).} Combinan SPA con
capacidades nativas: instalación en pantalla de inicio,
notificaciones push, funcionamiento offline, acceso a sensores del
dispositivo. \textbf{Tecnologías:} Service Workers, Web App Manifest,
IndexedDB. Twitter Lite fue uno de los primeros casos de éxito; en
Latinoamérica, Pinterest PWA y MercadoLibre han demostrado mejoras
del 60\% en retención de usuarios \cite{richards2020fundamentals}.

\item \textbf{Aplicaciones de comercio electrónico.} Subgénero
especializado con flujos transaccionales: catálogo, carrito, pasarela
de pago (PayPal, Stripe, Datafast en Ecuador), gestión de pedidos,
post-venta. Plataformas dominantes: Shopify, WooCommerce,
Magento/Adobe Commerce.

\item \textbf{Portales corporativos y \emph{e-government}.}
Aplicaciones con altos requisitos de seguridad, trazabilidad e
integración con sistemas heredados. Ejemplos nacionales: el portal
del SRI, el sistema de compras públicas SOCE, la Ventanilla Única
Electrónica del SENAE. Suelen requerir autenticación con certificado
digital y firma electrónica.
\end{enumerate}

\subsection{Características de las aplicaciones web modernas}

Antes de profundizar en arquitectura y ventajas comparativas, conviene
sintetizar las características esenciales que distinguen a una
aplicación web de cualquier otro tipo de software. Estas
características no son meras curiosidades técnicas: cada una tiene
implicaciones directas en cómo se desarrollan, despliegan, usan,
escalan y aseguran las aplicaciones web. Newman \cite{newman2021building}
y Kleppmann \cite{kleppmann2017designing} dedican capítulos enteros a
discutir las consecuencias arquitectónicas de cada característica.

\begin{table}[htbp]
\centering
\caption{Características esenciales de las aplicaciones web modernas y sus implicaciones.}
\label{tab:caracteristicas-web}
\footnotesize
\begin{tabularx}{\textwidth}{lX}
\toprule
\textbf{Dimensión} & \textbf{Característica}\\
\midrule
\multicolumn{2}{l}{\emph{Desarrollo}}\\
Lenguajes & Múltiples coexistentes: HTML+CSS+JS en cliente; PHP/Java/C\#/Node en servidor.\\
Herramientas & IDE (IntelliJ IDEA Ultimate, VS Code), Git, Docker, gestores de paquetes.\\
Metodología & Predominantemente ágil (Scrum, Kanban) con CI/CD.\\
\midrule
\multicolumn{2}{l}{\emph{Despliegue}}\\
Empaquetado & Contenedores Docker, archivos JAR/WAR, paquetes \texttt{npm}.\\
Hosting & Cloud (AWS, Azure, GCP), VPS, on-premise; orquestación con Kubernetes.\\
Actualización & \emph{Rolling updates} sin interrumpir el servicio; \emph{blue-green deployment}.\\
\midrule
\multicolumn{2}{l}{\emph{Uso}}\\
Acceso & Universal vía URL, sin instalación local.\\
Multiplataforma & Mismo código accesible desde Windows, macOS, Linux, iOS, Android.\\
Conectividad & Requiere red; PWAs ofrecen modo offline limitado.\\
\midrule
\multicolumn{2}{l}{\emph{Operación}}\\
Monitoreo & Métricas (Prometheus), logs (ELK), traces (OpenTelemetry, Tempo).\\
Escalabilidad & Horizontal preferida sobre vertical en arquitecturas modernas.\\
Disponibilidad & Objetivo común 99.9\% (8.7 h/año de downtime aceptable).\\
\midrule
\multicolumn{2}{l}{\emph{Seguridad}}\\
Superficie de ataque & Pública en internet: requiere defensas multinivel.\\
Estándares & OWASP Top 10:2021, ISO 27001, GDPR, LOPDP (Ecuador).\\
Identidad & OAuth 2.0, OpenID Connect, SAML, JWT.\\
\bottomrule
\end{tabularx}
\end{table}

Las implicaciones prácticas de estas características son enormes. Por
ejemplo, la naturaleza multiplataforma elimina el costo histórico de
mantener versiones separadas para Windows, macOS y Linux, pero
introduce el reto de la consistencia entre navegadores. La actualización
\emph{rolling} elimina la fricción del «día de instalación» del
software clásico, pero exige diseñar el código para que dos versiones
coexistan durante el despliegue. La superficie pública de ataque obliga
a integrar la seguridad desde el día uno; no se puede «agregar
seguridad» al final como en un sistema interno.

\subsection{Arquitectura cliente-servidor en aplicaciones web}

La arquitectura cliente-servidor es el modelo de comunicación
distribuida que organiza prácticamente la totalidad del tráfico web
actual. Comprenderla a profundidad es indispensable porque cada
decisión técnica posterior (desde qué lenguaje usar en el navegador
hasta cómo escalar el servidor) se apoya en sus principios. En esta
sección revisaremos primero su concepto fundamental, después la
diferencia frecuentemente confundida entre \emph{capas} y
\emph{niveles}, y finalmente las arquitecturas de 2, 3 y N niveles
con ejemplos concretos y diagramas.

\subsubsection{Concepto fundamental}

La arquitectura cliente-servidor es el paradigma fundamental sobre el
que se erige la World Wide Web. En este modelo, un proceso
\emph{cliente} emite una solicitud (\emph{request}) a un proceso
\emph{servidor}, el cual procesa la solicitud y devuelve una respuesta
(\emph{response}). El protocolo que estandariza este intercambio es
HTTP/1.1 (RFC 7230--7235), HTTP/2 (RFC 7540) o HTTP/3 (RFC 9114), este
último ya dominante en 2026 gracias a su desempeño superior en redes
móviles.

\begin{cajadefinicion}{Anatomía de una petición HTTP}
\begin{lstlisting}[style=shell,numbers=none,caption={Ejemplo de Shell},label=lst:1.1]
GET /productos?categoria=cafe HTTP/1.1
Host: tienda.uteq.edu.ec
Accept: text/html
User-Agent: Mozilla/5.0 (Windows NT 10.0; Win64; x64) ...
Accept-Language: es-EC,es;q=0.9,en;q=0.8
Cookie: SESSID=ab39e1ff7e1c4
\end{lstlisting}
La primera línea es la \emph{línea de petición}, compuesta por método,
ruta y versión del protocolo. Las líneas siguientes son
\emph{cabeceras}; tras una línea en blanco puede ir el cuerpo (en
métodos POST, PUT y PATCH).
\end{cajadefinicion}

\subsubsection{Niveles vs capas: una distinción frecuentemente confundida}

Uno de los conceptos más confusos para el estudiante principiante es
la distinción entre \emph{capas} y \emph{niveles}. Aunque a menudo se
usan como sinónimos, en arquitectura de software profesional
\textbf{tienen significados diferentes} que conviene precisar desde el
inicio \cite{fowler2002patterns,richards2020fundamentals}:

\begin{cajadefinicion}{Capa (\emph{layer}) vs Nivel (\emph{tier})}
\textbf{Capa (\emph{layer}):} agrupación \emph{lógica} de
responsabilidades dentro del software. Es una abstracción
conceptual; varias capas pueden ejecutarse en el mismo proceso físico.
Ejemplo: en una aplicación monolítica, las capas «Presentación»,
«Negocio» y «Datos» pueden coexistir en el mismo \texttt{.jar}.

\textbf{Nivel (\emph{tier}):} unidad \emph{física} de despliegue. Cada
nivel corre típicamente en su propio servidor o contenedor, y los
niveles se comunican por red. Ejemplo: «servidor de aplicación» y
«servidor de BD» son dos niveles distintos aunque la app sea un
monolito de varias capas internas.

Un sistema puede tener \textbf{3 capas pero 2 niveles} (capas
presentación, negocio y datos en un mismo proceso que se conecta a una
BD en otro servidor), o \textbf{3 capas y 3 niveles} (cada capa en su
propio servidor).
\end{cajadefinicion}

\subsubsection{Arquitecturas N-niveles: panorama visual}

Las arquitecturas web se clasifican por la cantidad de
\emph{niveles físicos} (máquinas o procesos separados) en que se
despliega el sistema. Comprender visualmente esta diferencia entre
1, 2, 3 y N niveles es esencial antes de hablar de patrones más
complejos como microservicios o serverless. La siguiente figura
contrasta los cuatro casos canónicos.


\begin{figure}[htbp]
\centering
\begin{tikzpicture}[
  every node/.style={font=\scriptsize,align=center},
  nivel/.style={rectangle,draw=uteqverde,thick,fill=uteqverde!10,
                rounded corners=2pt,minimum width=2.4cm,
                minimum height=1cm},
  flecha/.style={<->,>=stealth,thick,uteqverde!70!black}]

  % Arquitectura 1-nivel
  \node[font=\bfseries] at (1.2,3.5) {1-nivel};
  \node[nivel] (a1) at (1.2,2.5) {Todo\\(monolito\\local)};
  \node[font=\tiny,align=center] at (1.2,1.3) {Ejemplo:\\app desktop\\con BD local};

  % Arquitectura 2-niveles
  \node[font=\bfseries] at (5,3.5) {2-niveles};
  \node[nivel] (b1) at (5,2.5) {Cliente\\(con UI\\y lógica)};
  \node[nivel,fill=uteqdorado!15,draw=uteqdorado] (b2) at (5,0.5)
       {Servidor\\BD};
  \draw[flecha] (b1) -- (b2);
  \node[font=\tiny,align=center] at (5,-0.8) {Ejemplo:\\cliente Excel\\con SQL Server};

  % Arquitectura 3-niveles
  \node[font=\bfseries] at (9,3.5) {3-niveles};
  \node[nivel] (c1) at (9,2.8) {Navegador\\(presentación)};
  \node[nivel,fill=uteqverde!25] (c2) at (9,1.4)
       {App Server\\(negocio)};
  \node[nivel,fill=uteqdorado!15,draw=uteqdorado] (c3) at (9,0)
       {BD};
  \draw[flecha] (c1) -- (c2);
  \draw[flecha] (c2) -- (c3);
  \node[font=\tiny,align=center] at (9,-1.2)
       {Ejemplo:\\Chrome $\to$ Tomcat\\$\to$ PostgreSQL};

  % Arquitectura N-niveles
  \node[font=\bfseries] at (13,3.5) {N-niveles};
  \node[nivel] (d1) at (13,3.0) {Navegador};
  \node[nivel,fill=uteqverde!18] (d2) at (13,2.1) {CDN};
  \node[nivel,fill=uteqverde!25] (d3) at (13,1.2) {LB / Gateway};
  \node[nivel,fill=uteqverde!30] (d4) at (13,0.3) {App Server};
  \node[nivel,fill=uteqdorado!10,draw=uteqdorado] (d5) at (13,-0.6) {Cache};
  \node[nivel,fill=uteqdorado!15,draw=uteqdorado] (d6) at (13,-1.5) {BD};
  \draw[flecha] (d1) -- (d2);
  \draw[flecha] (d2) -- (d3);
  \draw[flecha] (d3) -- (d4);
  \draw[flecha] (d4) -- (d5);
  \draw[flecha] (d4.east) to[bend left=40] (d6.east);

\end{tikzpicture}
\caption{Las cuatro variantes principales de arquitectura
cliente-servidor por número de niveles físicos. La elección depende
del tamaño, presupuesto, requisitos de escalabilidad y disponibilidad.
La mayoría de aplicaciones web modernas usan 3 o más niveles.}
\label{fig:arq-niveles}
\end{figure}

A continuación se describe cada variante:

\textbf{Arquitectura 1-nivel (monolito local).} Todo el software (UI,
lógica, datos) reside en una sola máquina sin red. No aplica a la
web. Ejemplos: una aplicación de escritorio con SQLite embebido,
Microsoft Access en modo monousuario.

\textbf{Arquitectura 2-niveles (cliente-servidor tradicional).} El
cliente contiene la UI y parte de la lógica de negocio; el servidor
expone solo la BD. Modelo dominante en los 90 con clientes
\emph{fat-client} (Visual Basic + Oracle, Delphi + Sybase). Para la
web no es óptimo: la lógica de negocio expuesta al cliente compromete
la seguridad. Ejemplos vigentes: algunas aplicaciones legadas de
contabilidad con cliente Windows + SQL Server.

\textbf{Arquitectura 3-niveles (la clásica web).} Separa
explícitamente presentación (navegador), lógica de negocio (servidor
de aplicación: Tomcat, IIS, Apache+PHP) y datos (SGBD). Es el modelo
canónico de la web desde 1998 y sigue siendo el más común. Ejemplo:
Chrome $\to$ Spring Boot $\to$ PostgreSQL.

\textbf{Arquitectura N-niveles.} Introduce niveles adicionales para
abordar requisitos no funcionales: CDN para contenido estático,
balanceador de carga para distribución, gateway API para
autenticación cruzada, caché distribuida (Redis), sistemas de
mensajería (Kafka). Es el modelo dominante en sistemas a gran escala
\cite{bass2021software,kleppmann2017designing}. Ejemplo: Netflix
opera con decenas de niveles distintos.

\subsubsection{Anatomía detallada del flujo de una petición web}

Una aplicación web real involucra al menos cinco actores físicos o
lógicos cuya comprensión es esencial para diagnosticar problemas en
producción:

\begin{enumerate}
\item \textbf{Navegador del usuario}: Chrome, Firefox, Safari, Edge.
Procesa HTML, CSS y JavaScript; gestiona cookies, caché y
almacenamiento local.
\item \textbf{Servidor DNS}: traduce el nombre de dominio
(\texttt{uteq.edu.ec}) a una dirección IP. Operación que dura
típicamente entre 20--80 ms.
\item \textbf{Balanceador o proxy inverso}: Nginx, HAProxy, AWS ELB.
Distribuye la carga entre servidores de aplicación, termina TLS,
sirve contenido estático.
\item \textbf{Servidor de aplicación}: Apache HTTPD, Tomcat, IIS,
Node.js, PHP-FPM. Ejecuta el código de la aplicación.
\item \textbf{Sistema gestor de bases de datos}: PostgreSQL, MySQL,
MariaDB, SQL Server, Oracle. Persiste el estado de la aplicación.
\end{enumerate}

\begin{figure}[htbp]
\centering
\begin{tikzpicture}[
  node distance=1.8cm and 1.8cm,
  every node/.style={font=\footnotesize},
  caja/.style={rectangle,draw=uteqverde,thick,minimum width=2.4cm,
               minimum height=1cm,rounded corners=2pt,fill=uteqverde!10},
  flecha/.style={->,>=stealth,thick,uteqverde!70!black}]
  \node[caja] (cli) {Navegador};
  \node[caja,right=of cli] (dns) {DNS};
  \node[caja,right=of dns] (lb)  {Proxy/LB};
  \node[caja,below=of lb] (app)  {App Server};
  \node[caja,below=of app] (db)  {BD};
  \draw[flecha] (cli) -- node[above]{1. resolver} (dns);
  \draw[flecha] (cli.east) to[bend right=20] node[below]{2. HTTPS} (lb.west);
  \draw[flecha] (lb)  -- node[right]{3. enrutar} (app);
  \draw[flecha] (app) -- node[right]{4. SQL} (db);
  \draw[flecha] (db.west) to[bend left=40] node[left]{5. rows} (app.west);
  \draw[flecha] (app.north west) to[bend left=15] node[above]{6. HTML/JSON} (cli.south east);
\end{tikzpicture}
\caption{Flujo simplificado de una petición web. Cada uno de estos seis
pasos puede convertirse en cuello de botella; la disciplina de
optimización web consiste en medir cada eslabón.}
\label{fig:flujo-web}
\end{figure}

\subsection{Ventajas y desventajas de las aplicaciones web}

Toda decisión arquitectónica implica trade-offs: las aplicaciones
web ofrecen ventajas operativas significativas frente a las nativas,
pero también limitaciones que debemos conocer antes de comprometernos
con la plataforma. La tabla siguiente resume el balance.


\begin{table}[htbp]
\centering
\caption{Comparativa entre aplicaciones web, escritorio y móviles nativas.}
\footnotesize
\begin{tabularx}{\textwidth}{lXXX}
\toprule
\textbf{Criterio} & \textbf{Web} & \textbf{Escritorio} & \textbf{Móvil nativa}\\
\midrule
Despliegue & Inmediato sin instalación & Empaquetado e instalador & Tienda de apps (revisión)\\
Compatibilidad & Multiplataforma & Sistema-dependiente & iOS y Android (separadas)\\
Acceso a hardware & Limitado vía Web APIs & Total & Total\\
Costo de desarrollo & Bajo (un \emph{stack}) & Medio & Alto (dos \emph{stacks})\\
Costo de mantenimiento & Muy bajo (deploy único) & Medio & Alto (dos versiones)\\
Rendimiento bruto & Medio & Alto & Alto\\
Modo \emph{offline} & Limitado (PWA) & Total & Total\\
Distribución & Instantánea (URL) & Descarga manual & Tienda de apps\\
Actualizaciones & Automáticas & Manuales o semi & Aceptación del usuario\\
\bottomrule
\end{tabularx}
\end{table}

\section{Diseño web inclusivo y accesibilidad}

La web nació con vocación universal: cualquier persona, desde
cualquier dispositivo, en cualquier lugar del mundo, debería poder
acceder a la información que contiene. Sin embargo, la realidad
muestra que millones de personas con discapacidad encuentran
barreras infranqueables cada día al navegar sitios que ignoraron sus
necesidades. La accesibilidad web no es un tema marginal ni un
añadido decorativo: es una responsabilidad profesional, una
obligación legal en cada vez más países (incluido Ecuador) y un
factor de competitividad comercial cuantificable. Esta sección
introduce el concepto, presenta las cuatro categorías de
discapacidad que afectan la experiencia web, expone las pautas WCAG
2.1 del W3C con sus principios POUR, y termina con los criterios de
éxito que todo desarrollador debe dominar.

\subsection{¿Qué es la accesibilidad web?}

La accesibilidad web es la práctica de diseñar y construir sitios y
aplicaciones que puedan ser percibidos, comprendidos, navegados y
utilizados por \emph{todas las personas}, independientemente de sus
capacidades físicas, sensoriales o cognitivas. La definición canónica
del W3C \cite{w3c2018wcag} la formula así: «accesibilidad significa
que personas con discapacidades pueden percibir, comprender, navegar
e interactuar con la Web, y que pueden contribuir a la Web».

Sin embargo, la accesibilidad web no es solo una buena causa moral:
es un imperativo legal, comercial y técnico. García-Holgado et al.\
\cite{henry2018accessibility} demuestran en un estudio sistemático sobre
instituciones de educación superior que solo el 23\% de los sitios
analizados cumplen el nivel mínimo de WCAG 2.1, lo que constituye un
incumplimiento normativo en la mayoría de jurisdicciones.

\subsubsection{Marco legal de la accesibilidad}

La accesibilidad no es solo una buena práctica: en muchas
jurisdicciones es una obligación legal con sanciones concretas. La
siguiente enumeración recopila los marcos normativos vigentes que un
desarrollador web debe conocer antes de publicar un sitio.


\begin{itemize}
\item \textbf{Ecuador.} La Ley Orgánica de Discapacidades (Registro
Oficial 2012) exige conformidad mínima AA en sitios públicos a partir
de su Reglamento de 2017. CONADIS supervisa el cumplimiento.
\item \textbf{Estados Unidos.} Section 508 del Rehabilitation Act
(enmienda de 1998, refundición de 2017) obliga a las agencias
federales a la accesibilidad. Casos como \emph{Robles v.\ Domino's
Pizza} (2019) extendieron la jurisprudencia al sector privado bajo el
ADA.
\item \textbf{Unión Europea.} Directiva 2016/2102 sobre accesibilidad
de los sitios y aplicaciones móviles de los organismos del sector
público. EN 301 549 V3.2.1 (2021) define los requisitos técnicos.
\item \textbf{España, México, Argentina, Brasil, Colombia, Chile.}
Todos cuentan con normativa específica derivada de los lineamientos
internacionales.
\end{itemize}

\subsubsection{Impacto comercial cuantificable}

Más allá del cumplimiento legal, la accesibilidad genera retorno
medible:

\begin{itemize}
\item \textbf{15\% de la población mundial} vive con alguna forma de
discapacidad según la OMS (2024).
\item Mejorar la accesibilidad mejora el SEO: Google penaliza sitios
con baja accesibilidad mediante Core Web Vitals.
\item Microsoft estimó que sus inversiones en accesibilidad
generaron un retorno de inversión del 1.000\% al ampliar mercados
\cite{microsoft2020inclusive}.
\end{itemize}

\subsection{Las cuatro categorías de discapacidad relevantes para la web}

Antes de profundizar en los criterios técnicos, es indispensable
comprender qué tipos de barreras enfrentan los usuarios con
discapacidad y qué tecnologías de asistencia emplean. Esta
comprensión es la base ética y empática de la accesibilidad: no se
trata de cumplir una norma abstracta sino de habilitar a personas
reales con necesidades específicas.

Las categorías que se presentan a continuación siguen la taxonomía
clásica de Henry \cite{henry2018accessibility} y el W3C WAI.
Importante: una misma persona puede tener \emph{múltiples}
discapacidades simultáneas (por ejemplo, baja visión y artritis), o
\emph{discapacidades temporales} (un brazo enyesado, una lesión
ocular) o \emph{situacionales} (usar el móvil bajo luz solar
intensa, manejar con manos ocupadas). El diseño accesible beneficia
a todos estos casos.

\begin{enumerate}
\item \textbf{Discapacidad visual.} Incluye ceguera total (1.5\% de
la población mundial), baja visión y daltonismo (8\% de hombres,
0.5\% de mujeres). Los usuarios con ceguera utilizan lectores de
pantalla (NVDA gratuito en Windows, JAWS comercial, VoiceOver en
Apple, TalkBack en Android) que recorren el árbol de accesibilidad
del DOM y convierten la información a voz. Los daltónicos requieren
que la información no se transmita exclusivamente por color (el rojo
y verde son indistinguibles para la mayoría).

\item \textbf{Discapacidad auditiva.} Sordera total o parcial (5\% de
la población mundial). Requiere subtítulos en contenido multimedia
(WCAG 1.2.2), transcripciones de podcasts, lengua de señas
intérprete en videos institucionales críticos, y notificaciones
visuales además de sonoras.

\item \textbf{Discapacidad motora.} Imposibilidad de usar mouse o
teclado de forma estándar. Causas incluyen parálisis, temblores,
artritis reumatoide, distrofia muscular, amputaciones. Requiere
navegación completa por teclado (sin depender del mouse), áreas de
clic amplias (mínimo 44$\times$44 px según WCAG 2.5.5), ausencia de
tiempos límite estrictos, soporte para entrada por voz, switches
adaptados.

\item \textbf{Discapacidad cognitiva.} La categoría más amplia y
diversa: incluye dislexia (10\% de la población), TDAH (5\%),
trastornos del espectro autista (1--2\%), discapacidad intelectual,
demencia. Requiere lenguaje claro y conciso, estructura predecible,
ausencia de elementos parpadeantes (riesgo de crisis epilépticas
fotosensibles entre 3--65 Hz), feedback claro de errores, opciones
para ajustar tipografía y espaciado.
\end{enumerate}

\subsection{Las Web Content Accessibility Guidelines 2.1 (WCAG 2.1)}

WCAG 2.1 \cite{w3c2018wcag} es la norma internacional de referencia,
publicada por el W3C en junio de 2018 y vigente. Está organizada en
torno a cuatro principios fundamentales conocidos por el acrónimo
\textbf{POUR} (Perceivable, Operable, Understandable, Robust). Cada
principio define directrices, cada directriz se materializa en
\emph{criterios de éxito} comprobables, y cada criterio se clasifica
en uno de tres niveles de conformidad: A (mínimo), AA (recomendado),
AAA (excepcional).

Esta sección no se limita a enunciar los principios: para cada uno
explica \emph{qué tecnologías} permiten implementarlo y \emph{cómo}
medir el cumplimiento de forma objetiva.

\subsubsection{Principio 1 --- Perceptible}

\textbf{Enunciado:} La información y los componentes de la interfaz
deben presentarse de modo que los usuarios puedan percibirlos.

\textbf{Cómo lograrlo en HTML:}
\begin{itemize}
\item Atributo \texttt{alt} descriptivo en toda \texttt{<img>}.
\item Subtítulos \texttt{<track kind=\char34 captions\char34 >} en \texttt{<video>}.
\item Transcripciones para podcasts publicadas junto al audio.
\item \texttt{aria-label} para iconos clicables sin texto.
\end{itemize}

\textbf{Cómo lograrlo en CSS:}
\begin{itemize}
\item Contraste mínimo 4.5:1 (texto normal) y 3:1 (texto grande).
\item No usar solo color para transmitir información crítica.
\item Diseño responsive que soporte zoom hasta 200\% sin pérdida
funcional.
\end{itemize}

\textbf{Cómo medir cumplimiento:}
\begin{itemize}
\item \textbf{WebAIM Contrast Checker} (\url{https://webaim.org/resources/contrastchecker/}):
ingresa colores hex y devuelve la ratio exacta.
\item \textbf{axe DevTools} (extensión Chrome/Firefox): detecta
imágenes sin \texttt{alt}, contraste insuficiente, audio sin
transcripción.
\item \textbf{Pa11y CLI} (\url{https://pa11y.org}): herramienta de
línea de comandos para integrar auditoría en CI/CD.
\item \textbf{Lighthouse Accessibility} (incluido en Chrome
DevTools): score 0--100 con recomendaciones priorizadas.
\end{itemize}

\subsubsection{Principio 2 --- Operable}

\textbf{Enunciado:} Los componentes y la navegación deben ser
operables por todos los usuarios.

\textbf{Cómo lograrlo en HTML:}
\begin{itemize}
\item Elementos interactivos nativos (\texttt{<button>},
\texttt{<a>}) en lugar de \texttt{<div onclick>}.
\item Atributo \texttt{tabindex} solo cuando estrictamente necesario;
nunca \texttt{tabindex=\char34 2\char34 } u otros valores positivos.
\item \texttt{skip-link} al inicio de cada página para saltar a
contenido principal.
\end{itemize}

\textbf{Cómo lograrlo en CSS:}
\begin{itemize}
\item \texttt{:focus-visible} con outline claramente visible.
\item \texttt{prefers-reduced-motion} respetado en animaciones.
\item Áreas de clic mínimo 44$\times$44 px.
\end{itemize}

\textbf{Cómo lograrlo en JavaScript:}
\begin{itemize}
\item Sin tiempo límite ajustable o que pueda extenderse.
\item Manejo de teclado en componentes custom: \texttt{Enter},
\texttt{Space}, \texttt{Escape}, flechas direccionales.
\item Sin contenido parpadeante en rango 3--55 Hz.
\end{itemize}

\textbf{Cómo medir cumplimiento:}
\begin{itemize}
\item Navegación completa solo con \texttt{Tab} y \texttt{Enter} sin
mouse.
\item Recorrido visual del foco siguiendo orden lógico.
\item Lighthouse y axe DevTools detectan tabbing roto.
\end{itemize}

\subsubsection{Principio 3 --- Comprensible}

\textbf{Enunciado:} La información y el manejo de la interfaz deben
ser comprensibles.

\textbf{Cómo lograrlo en HTML:}
\begin{itemize}
\item Atributo \texttt{lang=\char34 es-EC\char34 } en \texttt{<html>}; cambiar con
\texttt{<span lang=\char34 en\char34 >} en fragmentos en otros idiomas.
\item \texttt{<label for=\char34 id\char34 >} asociado a cada \texttt{<input>}.
\item Atributo \texttt{autocomplete} con valores estándares.
\item \texttt{aria-describedby} para descripción extendida.
\end{itemize}

\textbf{Cómo lograrlo en lenguaje natural:}
\begin{itemize}
\item Nivel de lectura grado 8--9 según escala Flesch-Kincaid.
\item Frases cortas (máx.\ 25 palabras).
\item Voz activa preferida a pasiva.
\item Glosario de términos técnicos accesible desde cada uso.
\end{itemize}

\textbf{Cómo medir cumplimiento:}
\begin{itemize}
\item \textbf{Hemingway Editor} (\url{https://hemingwayapp.com}):
evalúa legibilidad.
\item Pruebas de usabilidad con usuarios reales de baja
alfabetización digital.
\end{itemize}

\subsubsection{Principio 4 --- Robusto}

\textbf{Enunciado:} El contenido debe ser robusto suficiente para que
pueda ser interpretado por una amplia variedad de agentes de usuario,
incluyendo tecnologías asistivas.

\textbf{Cómo lograrlo en HTML:}
\begin{itemize}
\item HTML válido sin errores (verificar en
\url{https://validator.w3.org/nu/}).
\item Roles ARIA solo cuando los elementos semánticos nativos no
basten.
\item IDs únicos en todo el documento.
\item Nombres accesibles para todos los controles
(\texttt{aria-label} o texto visible asociado).
\end{itemize}

\textbf{Cómo medir cumplimiento:}
\begin{itemize}
\item W3C Markup Validator: 0 errores.
\item NVDA o VoiceOver: recorrido auditivo manual.
\item axe DevTools: 0 violaciones serias o críticas.
\end{itemize}

\subsubsection{Niveles de conformidad y herramientas integradas de auditoría}

WCAG 2.1 establece tres niveles de conformidad: \textbf{A} (mínimo
imprescindible), \textbf{AA} (recomendado para sector público y
privado serio) y \textbf{AAA} (excepcional, raro de alcanzar
íntegramente). En 2023 se publicó WCAG 2.2 que añadió nueve criterios
de éxito adicionales y previsiblemente WCAG 3.0 (\emph{Silver},
borrador desde 2021) será la próxima revolución mayor, reemplazando
el sistema de niveles A/AA/AAA por un puntaje continuo.

\subsection{Criterios de éxito clave de WCAG 2.1}

Esta subsección presenta los diez criterios de éxito más importantes
que todo desarrollador web debe interiorizar. Para cada criterio se
explica: \emph{(a)} qué exige exactamente la norma, \emph{(b)} cómo
implementarlo en código real, y \emph{(c)} un ejemplo concreto.

\begin{cajabuenas}{Diez criterios WCAG imprescindibles --- detallados}
\textbf{1.1.1 Contenido no textual (Nivel A).}\\
\emph{Qué exige:} todo contenido no textual tiene una alternativa
textual equivalente.\\
\emph{Cómo lograrlo:} \texttt{alt} en \texttt{<img>};
\texttt{aria-label} en iconos clicables; \texttt{<title>} en SVGs
significativas; \texttt{alt=\char34 \char34 } (vacío) si la imagen es decorativa.\\
\emph{Ejemplo:} \texttt{<img src=\char34 grafico.png\char34  alt=\char34 Ventas
trimestrales 2026: Q1 \$45k, Q2 \$52k, Q3 \$48k, Q4 \$61k\char34 >}.

\textbf{1.3.1 Información y relaciones (Nivel A).}\\
\emph{Qué exige:} la estructura semántica del HTML refleja la
jerarquía visual y las relaciones de la información.\\
\emph{Cómo lograrlo:} usar encabezados correctos
\texttt{<h1>}-\texttt{<h6>}, listas \texttt{<ul>}/\texttt{<ol>} para
listas, \texttt{<table>} con \texttt{<th scope>} para datos
tabulares, \texttt{<fieldset>+<legend>} para agrupar campos
relacionados.\\
\emph{Anti-ejemplo:} simular tabla con \texttt{<div>} y CSS Grid sin
roles ARIA, rompe lectores de pantalla.

\textbf{1.4.3 Contraste mínimo (Nivel AA).}\\
\emph{Qué exige:} ratio de contraste $\geq$ 4.5:1 para texto normal;
$\geq$ 3:1 para texto grande (18pt+ o 14pt+ negrita).\\
\emph{Cómo lograrlo:} verificar con WebAIM Contrast Checker; usar
paletas pre-validadas.\\
\emph{Ejemplo de combinaciones AA conformes sobre fondo blanco:}
\texttt{\#006633} (verde UTEQ) ratio 6.27:1; negro
\texttt{\#000000} ratio 21:1.

\textbf{1.4.10 Reflujo (Nivel AA).}\\
\emph{Qué exige:} el contenido reflowea sin pérdida de información ni
funcionalidad y sin scroll horizontal en 320 px de ancho.\\
\emph{Cómo lograrlo:} CSS responsive mobile-first;
\texttt{viewport meta}; unidades relativas (rem, em, \%).

\textbf{1.4.11 Contraste no textual (Nivel AA).}\\
\emph{Qué exige:} componentes UI (botones, iconos de estado, bordes
de campos) con contraste $\geq$ 3:1 con sus alrededores.\\
\emph{Cómo lograrlo:} bordes visibles en inputs; iconos de estado
con fondo o borde definido.

\textbf{2.1.1 Teclado (Nivel A).}\\
\emph{Qué exige:} toda funcionalidad disponible solo desde teclado.\\
\emph{Cómo lograrlo:} usar elementos interactivos nativos; manejar
\texttt{Enter} y \texttt{Space} en componentes custom; no atrapar el
foco.\\
\emph{Anti-ejemplo:} dropdown que solo abre con hover de mouse.

\textbf{2.4.7 Foco visible (Nivel AA).}\\
\emph{Qué exige:} indicador visible cuando un elemento recibe foco.\\
\emph{Cómo lograrlo:} no quitar el outline por defecto, o
reemplazarlo con \texttt{:focus-visible} de mayor visibilidad
($\geq$ 2 px, contraste $\geq$ 3:1).

\textbf{3.1.1 Idioma de la página (Nivel A).}\\
\emph{Qué exige:} elemento \texttt{<html>} tiene atributo
\texttt{lang} correcto.\\
\emph{Cómo lograrlo:} \texttt{<html lang=\char34 es-EC\char34 >}.

\textbf{3.3.1 Identificación de errores (Nivel A).}\\
\emph{Qué exige:} si el sistema detecta error en entrada del usuario,
lo identifica y describe en texto.\\
\emph{Cómo lograrlo:} \texttt{<input aria-invalid=\char34 true\char34 
aria-describedby=\char34 err-email\char34 >} y mensaje en
\texttt{<span id=\char34 err-email\char34 >Email inválido</span>}.

\textbf{4.1.2 Nombre, función, valor (Nivel A).}\\
\emph{Qué exige:} para todos los componentes de interfaz, nombre,
función y valor son determinables programáticamente.\\
\emph{Cómo lograrlo:} usar elementos nativos siempre que sea posible;
si se construye un componente custom, declarar
\texttt{role}, \texttt{aria-label}, \texttt{aria-checked} etc.
\end{cajabuenas}

\section{Configuración del entorno con IntelliJ IDEA Ultimate}

Antes de continuar con el resto de la asignatura, el estudiante
debe disponer de un entorno IDE funcional. El siguiente ejercicio
resuelto guía paso a paso la instalación y verificación de IntelliJ
IDEA Ultimate con la licencia educativa UTEQ.


\begin{cajaejercicio}{Ejercicio 1.1 (RESUELTO) --- Instalación y verificación de IntelliJ IDEA Ultimate}
\textbf{Objetivo:} configurar el entorno integrado de desarrollo que
se utilizará durante todo el semestre y verificar que los plugins
necesarios para HTML, CSS y JavaScript están operativos.

\textbf{Paso 1 --- Obtener la licencia educativa.} Ingresar a
\url{https://www.jetbrains.com/community/education} y crear una
cuenta con el correo institucional \texttt{<usuario>@uteq.edu.ec}. Al
verificar el dominio, JetBrains envía un correo de confirmación con
una clave educativa renovable anualmente. Sin esta clave, IntelliJ
IDEA Ultimate funciona únicamente 30 días en modo evaluación.

\textbf{Paso 2 --- Descargar e instalar.} Descargar IntelliJ IDEA
Ultimate desde \url{https://www.jetbrains.com/idea/download}. El
instalador pesa aproximadamente 850 MB.

\textbf{Paso 3 --- Activar la licencia.} En la primera ejecución
elegir «Activate License $\rightarrow$ JB Account». Verificar en
\texttt{Help $\rightarrow$ Register} que la licencia es de tipo
\emph{Educational License}.

\textbf{Paso 4 --- Crear el primer proyecto.} En la pantalla
\emph{Welcome} pulsar \texttt{New Project}. Seleccionar la plantilla
\texttt{HTML}, establecer \textbf{Name}: \texttt{appweb-semana01},
y pulsar \texttt{Create}.

\textbf{Paso 5 --- Verificar plugins esenciales} en
\texttt{File $\rightarrow$ Settings $\rightarrow$ Plugins}: HTML
Tools, CSS, JavaScript and TypeScript, Markdown, Live Edit.

\textbf{Paso 6 --- Probar.} Abrir el archivo \texttt{index.html} que
el IDE generó automáticamente. Posicionar el cursor sobre el código
y pulsar \texttt{Alt+F2}; elegir Chrome.

\textbf{Resultado esperado (verificación):}

\begin{center}
\fbox{\begin{minipage}{0.85\textwidth}
\footnotesize\ttfamily
\textbf{[ Ventana del navegador Chrome ]}\\
\textbf{URL:} \texttt{http://localhost:63342/appweb-semana01/index.html}\\[0.3em]
\textbf{[ Pestaña: index.html ]}\\[0.3em]
Hello, World!\\[0.3em]
\rmfamily\textit{(Página HTML mínima generada por IntelliJ con el texto
por defecto, servida desde el servidor de desarrollo embebido del IDE
en el puerto 63342.)}
\end{minipage}}
\end{center}

\textbf{Indicadores de éxito verificables:}
\begin{itemize}
\item La barra de URL muestra una dirección \texttt{localhost:63342/\dots}
\item El título de la pestaña coincide con \texttt{index.html}.
\item En IntelliJ, el panel \emph{Run} (parte inferior) muestra el
servidor en estado «Running».
\end{itemize}
\end{cajaejercicio}

\begin{cajaejercicio}{Ejercicio 1.2 (RESUELTO) --- Primera página HTML5 semántica y accesible}
\textbf{Objetivo:} crear una página HTML5 estructuralmente correcta,
semánticamente significativa y accesible. Verificar su validez en el
W3C Validator y la ausencia de violaciones críticas en axe DevTools.

\textbf{Código completo (Listado~\ref{lst:1.2}):}

\begin{lstlisting}[style=html,caption={index.html (Ejercicio 1.2)},label=lst:1.2]
<!DOCTYPE html>
<html lang="es-EC">
<head>
  <meta charset="UTF-8">
  <meta name="viewport" content="width=device-width, initial-scale=1.0">
  <title>Mi primera pagina semantica - UTEQ</title>
  <meta name="description" content="Pagina de practica
        para Aplicaciones Web - UTEQ 2026">
</head>
<body>
  <header role="banner">
    <h1>Universidad Tecnica Estatal de Quevedo</h1>
    <p>La primera universidad agropecuaria del Ecuador</p>
    <nav role="navigation" aria-label="Menu principal">
      <ul>
        <li><a href="#inicio">Inicio</a></li>
        <li><a href="#carreras">Carreras</a></li>
        <li><a href="#contacto">Contacto</a></li>
      </ul>
    </nav>
  </header>
  <main role="main">
    <article>
      <h2 id="inicio">Bienvenido al curso de Aplicaciones Web</h2>
      <p>Este es el primer ejercicio del semestre. La fecha de hoy es
      <time datetime="2026-05-04">4 de mayo de 2026</time>.</p>
      <figure>
        <img src="https://placehold.co/240x240/006633/ffffff?text=UTEQ"
             alt="Escudo de la UTEQ"
             width="240" height="240">
        <figcaption>Escudo institucional UTEQ</figcaption>
      </figure>
    </article>
    <aside aria-label="Informacion adicional">
      <h2>Datos de contacto</h2>
      <address>
        Av. Quito km. 1.5 via a Santo Domingo<br>
        Quevedo, Los Rios, Ecuador<br>
        Telefono: <a href="tel:+59353702220">(+593) 5 3702-220</a>
      </address>
    </aside>
  </main>
  <footer role="contentinfo">
    <small>&copy; 2026 UTEQ. Todos los derechos reservados.</small>
  </footer>
</body>
</html>
\end{lstlisting}

\textbf{Explicación condensada:}
\texttt{<!DOCTYPE html>} declara HTML5; \texttt{lang=\char34 es-EC\char34 } permite
al lector de pantalla elegir voz española; \texttt{<meta charset>}
evita corrupción de tildes; \texttt{<meta viewport>} habilita
responsive en móvil. Los roles ARIA (\texttt{banner},
\texttt{navigation}, \texttt{main}, \texttt{contentinfo}) refuerzan
la semántica para tecnologías asistivas antiguas. \texttt{<time
datetime>} hace la fecha legible por máquinas. \texttt{alt} en la
imagen cumple WCAG 1.1.1.

\textbf{Resultado esperado en el navegador:}

\begin{center}
\fbox{\begin{minipage}{0.88\textwidth}
\footnotesize
\textbf{[ Vista renderizada en Chrome - sin estilos CSS ]}\\[0.3em]
\rule{\textwidth}{0.4pt}\\
\textbf{\Large Universidad Tecnica Estatal de Quevedo}\\
La primera universidad agropecuaria del Ecuador\\[0.3em]
$\bullet$ Inicio\quad $\bullet$ Carreras\quad $\bullet$ Contacto\\
\rule{\textwidth}{0.4pt}\\[0.5em]
\textbf{\large Bienvenido al curso de Aplicaciones Web}\\
Este es el primer ejercicio del semestre. La fecha de hoy es 4 de
mayo de 2026.\\[0.3em]
\fbox{\begin{minipage}{2cm}\centering
\colorbox{uteqverde}{\textcolor{white}{\rule{0pt}{1.8cm}\hspace{1.8cm}}}\\
\textit{UTEQ}\end{minipage}}\\
\textit{Escudo institucional UTEQ}\\[0.5em]
\textbf{\large Datos de contacto}\\
Av. Quito km. 1.5 via a Santo Domingo\\
Quevedo, Los Rios, Ecuador\\
Telefono: \underline{(+593) 5 3702-220}\\
\rule{\textwidth}{0.4pt}\\
\textit{\textcopyright{} 2026 UTEQ. Todos los derechos reservados.}
\end{minipage}}
\end{center}

\textbf{Verificación obligatoria del éxito:}
\begin{enumerate}
\item \emph{W3C Validator} (\url{https://validator.w3.org/nu/}):
copiar el código completo, pegar en «Validate by direct input»,
pulsar Check. Debe aparecer en verde: «Document checking completed.
No errors or warnings to show.»
\item \emph{axe DevTools}: instalar la extensión desde Chrome Web
Store. F12 $\to$ pestaña axe DevTools $\to$ «Scan ALL of my page».
Resultado esperado: 0 violaciones críticas.
\item \emph{Navegación con teclado}: pulsar Tab repetidamente desde
la barra de URL. El foco debe recorrer en orden los tres enlaces del
menú y el enlace del teléfono.
\end{enumerate}
\end{cajaejercicio}

\begin{cajaejercicio}{Ejercicio 1.3 (PROPUESTO) --- Página de presentación personal accesible}
\textbf{Objetivo:} construir una página de presentación personal que
integre todos los conceptos vistos: estructura semántica,
accesibilidad, metadatos correctos, navegación con teclado.

\textbf{Especificaciones obligatorias:}
\begin{enumerate}
\item Crear un nuevo proyecto en IntelliJ:
\texttt{appweb-presentacion-<TuApellido>}.
\item El archivo \texttt{index.html} debe contener:
   \begin{itemize}
   \item DOCTYPE HTML5 y \texttt{lang=\char34 es-EC\char34 }.
   \item Metadatos: charset, viewport, description (máx.\ 160
   caracteres), author, keywords (5--7 términos).
   \item \texttt{<header>} con nombre completo (h1), foto con
   \texttt{alt} descriptivo, y un \texttt{<nav>} con cinco enlaces
   internos.
   \item Cinco \texttt{<section>} (cada una con su \texttt{<h2>}):
   «Sobre mí», «Estudios», «Habilidades técnicas», «Proyectos» y
   «Contacto».
   \item En «Habilidades técnicas», una tabla \texttt{<table>}
   con cuatro filas (HTML, CSS, JavaScript, otra de elección) y
   tres columnas (tecnología, nivel 1--5, años de experiencia).
   \item En «Proyectos», tres \texttt{<article>} con
   \texttt{<figure>+<figcaption>} cada uno.
   \item \texttt{<footer>} con copyright y fecha
   \texttt{<time datetime=\char34 2026\char34 >}.
   \end{itemize}
\item La página debe ser \emph{válida} en
\url{https://validator.w3.org/nu/} sin errores ni advertencias.
\item axe DevTools debe arrojar cero violaciones de cualquier
severidad.
\item La navegación por teclado debe permitir alcanzar todos los
enlaces y secciones.
\end{enumerate}

\textbf{Entrega:} \texttt{index.html} + captura del W3C Validator
mostrando «No errors» + captura de axe DevTools mostrando «No issues
found».

\textbf{Esta práctica forma parte de la \emph{Sumativa GA-Taller:
Encuadre académico}.}
\end{cajaejercicio}

\section{Buenas prácticas profesionales para HTML semántico}

Las buenas prácticas profesionales no son convenciones arbitrarias:
cada una tiene una justificación técnica, de mantenibilidad, de
accesibilidad o de SEO documentada en la literatura especializada
\cite{robbins2018learning,duckett2011html}. La adherencia disciplinada
a estas prácticas es lo que separa al código profesional del código
amateur.

\begin{cajabuenas}{Doce reglas de oro del HTML semántico profesional}
\begin{enumerate}
\item \textbf{Un solo \texttt{<h1>} por página}, que coincide con el
título principal y suele ser similar al del \texttt{<title>}.
\item \textbf{Jerarquía de encabezados sin saltos}:
\texttt{h1 $\rightarrow$ h2 $\rightarrow$ h3}, nunca
\texttt{h1 $\rightarrow$ h3} saltando \texttt{h2}.
\item \textbf{Ningún encabezado vacío y ningún encabezado meramente
decorativo}: si un texto necesita estilo de título pero no es título,
usar \texttt{<p>} con clase CSS.
\item \textbf{Listas para listas}: cualquier serie de elementos
similares debe ir en \texttt{<ul>}, \texttt{<ol>} o \texttt{<dl>},
nunca en una serie de \texttt{<div>}.
\item \textbf{Tablas \emph{solo} para datos tabulares}, nunca para
maquetación.
\item \textbf{Botones para acciones, enlaces para navegación}.
\item \textbf{Imágenes con \texttt{alt} descriptivo}; si la imagen
es puramente decorativa, \texttt{alt=\char34 \char34 } (vacío).
\item \textbf{Formularios con \texttt{<label>} asociado a cada
\texttt{<input>}}.
\item \textbf{Validar siempre} con \url{https://validator.w3.org/nu/}
antes de cada commit.
\item \textbf{Atributo \texttt{lang} obligatorio} en \texttt{<html>}.
\item \textbf{Evitar \texttt{<br>} para crear espacios}: usar CSS.
\item \textbf{Mantener IDs únicos}: dos elementos con el mismo
\texttt{id} son HTML inválido y rompen tecnologías asistivas.
\end{enumerate}
\end{cajabuenas}

\section{Tabla de referencia rápida: tags HTML5 más usados, atributos y valores}

La tabla \ref{tab:tags-html} compendia los tags HTML5 más frecuentes
en el desarrollo profesional, sus atributos principales, los valores
admitidos y el efecto esperado. Algunos atributos
(\texttt{id}, \texttt{class}, \texttt{lang}, \texttt{title},
\texttt{hidden}, \texttt{tabindex}, eventos \texttt{onclick} etc.)
son \textbf{globales}: aplican a casi todos los elementos. Estos
se listan una sola vez al inicio para evitar repetición.
\scriptsize
\setlength{\LTpre}{0pt}
\setlength{\LTpost}{0pt}
\begin{longtable}{@{}>{\raggedright\arraybackslash}p{2.6cm}
		>{\raggedright\arraybackslash}p{2.4cm}
		>{\raggedright\arraybackslash}p{6.0cm}
		>{\raggedright\arraybackslash}p{3.4cm}@{}}
	\caption{Tags HTML5 más usados: estructura, contenido y formularios.}
	\label{tab:tags-html}\\
	\toprule
	\textbf{Tag} & \textbf{Atributo} & \textbf{Valores / efecto} & \textbf{Ejemplo}\\
	\midrule
	\endfirsthead
	
	\multicolumn{4}{l}{\emph{(Continuación de la Tabla \ref{tab:tags-html})}}\\
	\toprule
	\textbf{Tag} & \textbf{Atributo} & \textbf{Valores / efecto} & \textbf{Ejemplo}\\
	\midrule
	\endhead
	
	\midrule
	\multicolumn{4}{r}{\emph{Continúa en la siguiente página\ldots}}\\
	\endfoot
	
	\bottomrule
	\endlastfoot
	
	\multicolumn{4}{l}{\emph{Estructura semántica (HTML5)}}\\
	\texttt{<header>}    & --- & Cabecera de página o sección & \texttt{<header>\dots</header>}\\
	\texttt{<nav>}       & --- & Bloque de navegación & \texttt{<nav>\dots</nav>}\\
	\texttt{<main>}      & --- & Contenido principal (único por página) & \texttt{<main>\dots</main>}\\
	\texttt{<article>}   & --- & Contenido autónomo reutilizable & \texttt{<article>\dots</article>}\\
	\texttt{<section>}   & --- & Sección temática con encabezado & \texttt{<section>\dots</section>}\\
	\texttt{<aside>}     & --- & Contenido relacionado pero secundario & \texttt{<aside>\dots</aside>}\\
	\texttt{<footer>}    & --- & Pie de página o sección & \texttt{<footer>\dots</footer>}\\
	\midrule
	\multicolumn{4}{l}{\emph{Contenido textual}}\\
	\texttt{<h1>}--\texttt{<h6>} & --- & Encabezados de niveles 1 a 6 & \texttt{<h1>Titulo</h1>}\\
	\texttt{<p>}         & --- & Párrafo & \texttt{<p>Texto.</p>}\\
	\texttt{<strong>}    & --- & Énfasis fuerte (importancia) & \texttt{<strong>Atencion</strong>}\\
	\texttt{<em>}        & --- & Énfasis (entonación) & \texttt{<em>realmente</em>}\\
	\texttt{<mark>}      & --- & Texto destacado contextualmente & \texttt{<mark>resaltado</mark>}\\
	\texttt{<time>}      & \texttt{datetime} & Fecha/hora ISO 8601 legible máquina & \texttt{<time datetime=\char34 2026-05-04\char34 >}\\
	\texttt{<code>}      & --- & Fragmento de código & \texttt{<code>let x=1</code>}\\
	\texttt{<blockquote>} & \texttt{cite} & URL fuente de la cita & \texttt{<blockquote cite=\char34 ...\char34 >}\\
	\midrule
	\multicolumn{4}{l}{\emph{Hipervínculos y enlaces}}\\
	\texttt{<a>} & \texttt{href} & URL destino & \texttt{href=\char34 https://uteq.edu.ec\char34 }\\
	& \texttt{target} & \texttt{\_blank}, \texttt{\_self}, \texttt{\_parent} & \texttt{target=\char34 \_blank\char34 }\\
	& \texttt{rel}  & \texttt{noopener}, \texttt{noreferrer}, \texttt{me} & \texttt{rel=\char34 noopener\char34 }\\
	& \texttt{download} & Fuerza descarga & \texttt{download=\char34 archivo.pdf\char34 }\\
	\midrule
	\multicolumn{4}{l}{\emph{Imágenes y multimedia}}\\
	\texttt{<img>} & \texttt{src} & URL de la imagen & \texttt{src=\char34 logo.png\char34 }\\
	& \texttt{alt} & Texto alternativo (obligatorio) & \texttt{alt=\char34 Logo UTEQ\char34 }\\
	& \texttt{width}, \texttt{height} & Dimensiones (px o \%) & \texttt{width=\char34 200\char34 }\\
	& \texttt{loading} & \texttt{lazy}, \texttt{eager} & \texttt{loading=\char34 lazy\char34 }\\
	& \texttt{srcset} & Múltiples resoluciones & \texttt{srcset=\char34 img-2x.png 2x\char34 }\\
	\texttt{<video>} & \texttt{src} & URL del video & \texttt{src=\char34 video.mp4\char34 }\\
	& \texttt{controls} & Booleano: controles nativos & \texttt{controls}\\
	& \texttt{poster} & Imagen previa & \texttt{poster=\char34 prev.jpg\char34 }\\
	& \texttt{preload} & \texttt{none}, \texttt{metadata}, \texttt{auto} & \texttt{preload=\char34 metadata\char34 }\\
	\midrule
	\multicolumn{4}{l}{\emph{Listas}}\\
	\texttt{<ul>}, \texttt{<ol>} & --- & Lista no/ordenada & \texttt{<ul><li>\dots</li></ul>}\\
	\texttt{<li>}       & --- & Ítem de lista & \texttt{<li>texto</li>}\\
	\texttt{<dl>}, \texttt{<dt>}, \texttt{<dd>} & --- & Lista de definiciones & ---\\
	\midrule
	\multicolumn{4}{l}{\emph{Tablas}}\\
	\texttt{<table>}    & --- & Datos tabulares & \texttt{<table>\dots</table>}\\
	\texttt{<thead>}, \texttt{<tbody>}, \texttt{<tfoot>} & --- & Secciones & ---\\
	\texttt{<tr>}       & --- & Fila & \texttt{<tr>\dots</tr>}\\
	\texttt{<th>}       & \texttt{scope} & \texttt{col}, \texttt{row}, \texttt{colgroup} & \texttt{<th scope=\char34 col\char34 >}\\
	\texttt{<td>}       & \texttt{colspan}, \texttt{rowspan} & N de celdas a abarcar & \texttt{colspan=\char34 2\char34 }\\
	\midrule
	\multicolumn{4}{l}{\emph{Formularios}}\\
	\texttt{<form>}     & \texttt{action} & URL del handler & \texttt{action=\char34 /api/users\char34 }\\
	& \texttt{method} & \texttt{GET}, \texttt{POST} & \texttt{method=\char34 POST\char34 }\\
	& \texttt{enctype} & MIME para uploads & \texttt{enctype=\char34 multipart/form-data\char34 }\\
	\texttt{<input>}    & \texttt{type} & \texttt{text}, \texttt{email}, \texttt{password}, \texttt{number}, \texttt{date}, \texttt{tel}, \texttt{url}, \texttt{checkbox}, \texttt{radio}, \texttt{file}, \texttt{hidden}, \texttt{range}, \texttt{color} & \texttt{type=\char34 email\char34 }\\
	& \texttt{name} & Nombre del campo & \texttt{name=\char34 email\char34 }\\
	& \texttt{value} & Valor inicial & \texttt{value=\char34 default\char34 }\\
	& \texttt{required} & Booleano: obligatorio & \texttt{required}\\
	& \texttt{pattern} & Regex de validación & \texttt{pattern=\char34{}[0-9]\{10\}\char34{}}\\
	& \texttt{min}/\texttt{max}/\texttt{step} & Rango numérico/fecha & \texttt{min=\char34 0\char34  max=\char34 100\char34 }\\
	& \texttt{minlength}/\texttt{maxlength} & Longitud texto & \texttt{maxlength=\char34 80\char34 }\\
	& \texttt{autocomplete} & \texttt{on}, \texttt{off}, \texttt{name}, \texttt{email}\dots & \texttt{autocomplete=\char34 email\char34 }\\
	\texttt{<label>}    & \texttt{for}  & ID del input asociado & \texttt{<label for=\char34 em\char34 >}\\
	\texttt{<select>}   & \texttt{multiple} & Booleano: múltiple selección & \texttt{multiple}\\
	\texttt{<option>}   & \texttt{value}, \texttt{selected} & Valor; preselección & \texttt{<option value=\char34 EC\char34  selected>}\\
	\texttt{<textarea>} & \texttt{rows}, \texttt{cols} & Tamaño visible & \texttt{rows=\char34 5\char34 }\\
	\texttt{<button>}   & \texttt{type} & \texttt{submit}, \texttt{button}, \texttt{reset} & \texttt{type=\char34 submit\char34 }\\
\end{longtable}

Esta tabla está organizada para servir como \emph{referencia rápida}
durante la práctica. Para profundizar en los atributos específicos de
formularios (tipos de input, validación nativa, accesibilidad de
controles) la semana 2 dedica una sección completa al tema.

\subsection{Atributos globales (aplican a casi todos los tags)}

Los atributos globales son aquellos que pueden aplicarse a
\emph{cualquier} elemento HTML5, no solo a unos pocos. Comprenderlos
en profundidad evita reaprenderlos para cada nuevo tag. La tabla
\ref{tab:atributos_globales} cataloga los más importantes con su efecto y un ejemplo
breve.


\begin{table}[htbp]
	\centering
	\caption{Atributos HTML5 globales. \label{tab:atributos_globales}}
	\footnotesize
	\begin{tabularx}{\textwidth}{lXl}
		\toprule
		\textbf{Atributo} & \textbf{Efecto / valores} & \textbf{Ejemplo}\\
		\midrule
		\texttt{id}    & Identificador único en el documento & \texttt{id=\char34 hero\char34 }\\
		\texttt{class} & Una o varias clases CSS separadas por espacio & \texttt{class=\char34 btn primario\char34 }\\
		\texttt{style} & CSS inline (uso desaconsejado) & \texttt{style=\char34 color:red\char34 }\\
		\texttt{title} & Tooltip flotante al pasar el mouse & \texttt{title=\char34 Cerrar\char34 }\\
		\texttt{lang}  & Idioma del contenido (ISO 639-1 y 3166-1) & \texttt{lang=\char34 es-EC\char34 }\\
		\texttt{dir}   & Dirección texto: \texttt{ltr}, \texttt{rtl}, \texttt{auto} & \texttt{dir=\char34 rtl\char34 }\\
		\texttt{hidden} & Booleano: oculta el elemento & \texttt{hidden}\\
		\texttt{tabindex} & Orden de tabulación; 0 (natural), -1 (programático) & \texttt{tabindex=\char34 0\char34 }\\
		\texttt{role}  & Rol ARIA explícito & \texttt{role=\char34 navigation\char34 }\\
		\texttt{aria-label} & Etiqueta accesible cuando no hay texto visible & \texttt{aria-label=\char34 Cerrar\char34 }\\
		\texttt{aria-hidden} & \texttt{true}/\texttt{false}: oculta para lectores & \texttt{aria-hidden=\char34 true\char34 }\\
		\texttt{data-*} & Datos custom para JavaScript & \texttt{data-id=\char34 42\char34 }\\
		\texttt{contenteditable} & \texttt{true}/\texttt{false} & \texttt{contenteditable=\char34 true\char34 }\\
		\texttt{draggable} & Booleano: arrastrable & \texttt{draggable=\char34 true\char34 }\\
		\bottomrule
	\end{tabularx}
\end{table}

\textbf{Ejemplo corto comentado (atributos globales en acción) (Listado~\ref{lst:1.3}):}
\begin{lstlisting}[language=HTML5,numbers=none,caption={Ejemplo de HTML},label=lst:1.3]
<!-- 'id' identifica univocamente al elemento (uno por documento) -->
<!-- 'class' agrupa elementos para aplicar el mismo estilo CSS    -->
<!-- 'title' muestra un tooltip al pasar el cursor               -->
<!-- 'lang' especifica el idioma del contenido (importante a11y)  -->
<!-- 'data-*' guarda datos personalizados accesibles desde JS     -->
<article id="post-42" class="card destacado"
         title="Noticia destacada" lang="es"
         data-autor="ana" data-categoria="tecnologia">
  Contenido del articulo
</article>
\end{lstlisting}

\subsection{Tags estructurales y de contenido}

HTML5 introduce un conjunto de tags semánticos pensados para
describir el \emph{significado} del contenido, no solo su apariencia.
La tabla \ref{tab:tags-html} cataloga los tags estructurales y de contenido más usados, con sus atributos específicos y un ejemplo de uso para cada uno.

\textbf{Ejemplo corto comentado (estructura semántica completa) (Listado~\ref{lst:1.4}):}

\begin{lstlisting}[language=HTML5,numbers=none,caption={Ejemplo de HTML},label=lst:1.4]
<!-- <header> contiene la cabecera del sitio o de una seccion       -->
<header>
  <!-- <nav> agrupa enlaces de navegacion principal                  -->
  <nav>
    <a href="/inicio">Inicio</a>
    <a href="/contacto">Contacto</a>
  </nav>
</header>
<!-- <main> envuelve el contenido unico de la pagina (uno por doc)   -->
<main>
  <!-- <article> es contenido autonomo (post, noticia, comentario)   -->
  <article>
    <!-- <h1>...<h6> jerarquia de encabezados (h1 = mas importante)  -->
    <h1>Titulo del articulo</h1>
    <!-- <section> agrupa contenido tematicamente relacionado         -->
    <section>
      <p>Parrafo de texto en la seccion.</p>
    </section>
  </article>
</main>
<!-- <footer> contiene pie de pagina (copyright, contacto, sitemap)  -->
<footer><p>(c) 2026 UTEQ</p></footer>
\end{lstlisting}

%%% ==== SEMANA 02 v3 ====
% =====================================================================
%   SEMANA 2: HTML5 AMPLIADA — VERSIÓN MEJORADA v3
% =====================================================================
\chapter{HTML y HTML5 para el desarrollo web}
\label{cap:semana2}

\epigraph{«HTML no es código fuente: es un contrato semántico entre el
autor del documento y el ecosistema de agentes que lo interpretan.»}
{--- \textit{Ian Hickson, editor de la especificación HTML5}}

\section*{Resultado de aprendizaje de la semana}
Al concluir la semana 2, el estudiante construye documentos HTML5
estructuralmente correctos, semánticamente significativos y validables
ante el W3C, integrando formularios con validación nativa, contenido
multimedia accesible, APIs nativas del lenguaje y entendiendo el lugar
de XML en el ecosistema de marcado.

\section{Historia y evolución de HTML}

El lenguaje HTML ha evolucionado durante tres décadas
adaptándose a las demandas cambiantes de la web: del documento
estático al \emph{single-page application} interactivo. Conocer las
etapas de esta evolución ayuda a entender por qué algunos tags
existen, por qué otros se deprecaron y por qué HTML5 adoptó las
formas actuales. La siguiente caja narra esa historia.


\begin{cajahistoria}{De SGML a HTML Living Standard}
La historia de HTML no comienza con HTML sino con SGML (Standard
Generalized Markup Language), estándar ISO 8879 publicado en 1986.
SGML era un meta-lenguaje complejo que permitía definir lenguajes de
marcado específicos para documentos de gobierno, ejército y empresa.
Tim Berners-Lee, al inventar la web en 1989, tomó SGML como punto de
partida pero buscó una versión radicalmente simplificada: HTML.

La primera descripción pública de HTML data de 1991 en una nota
informal titulada \emph{«HTML Tags»} que listaba apenas 18 etiquetas
\cite{bernerslee1991htmltags}. HTML 2.0 (1995, RFC 1866) fue el
primer estándar formal, manteniendo la simplicidad pero añadiendo
formularios. HTML 3.2 (W3C, 1997) incorporó tablas, applets y
\emph{scripting}. HTML 4.01 (1999) trajo la separación entre
estructura y presentación al recomendar el uso de hojas de estilo
CSS \cite{raggett1999html401}.

A finales de los noventa el W3C decidió reemplazar HTML por XHTML,
una versión XML-estricta del lenguaje. Esta decisión, aunque
técnicamente elegante, ignoró que la mayoría del HTML real en la web
era \emph{tag soup}: marcado mal formado pero que funcionaba en los
navegadores. La presión por estandarizar las prácticas reales llevó a
la formación, en 2004, del \emph{Web Hypertext Application Technology
Working Group} (WHATWG) por parte de Apple, Mozilla y Opera
\cite{hickson2004whatwg}. El WHATWG se desvinculó del proceso XHTML
2.0 y comenzó a desarrollar HTML5.

En 2014 el W3C publicó HTML5 como Recomendación oficial. En 2019
W3C y WHATWG firmaron un memorando reconociendo que el WHATWG sería
el único editor del HTML Living Standard
\cite{whatwgw3c2019memorandum}, abandonando la idea de versiones
numeradas. Desde entonces, HTML evoluciona como \emph{Living
Standard}: cambios incrementales continuos sin números de versión
formales. Por convención y comodidad, seguimos hablando de «HTML5»
para referirnos al HTML moderno.
\end{cajahistoria}

La evolución de HTML refleja la transformación de la web misma: lo
que comenzó como un sistema simple para compartir documentos
científicos hipertextuales en el CERN se convirtió en la plataforma
de software más ubicua del planeta. Cada versión documenta no solo
adiciones técnicas sino también cambios profundos en cómo se entendía
la naturaleza del documento web.

La tabla \ref{tab:evolucion-html} resume los hitos. Conviene
contextualizarla: las versiones \emph{numeradas} de HTML
(\texttt{HTML 2.0}, \texttt{3.2}, \texttt{4.01}) corresponden a una
etapa en la que el lenguaje se actualizaba mediante \emph{snapshots}
formales del W3C. A partir de HTML5, y especialmente desde el
memorando de 2019 entre WHATWG y W3C
\cite{whatwgw3c2019memorandum}, el modelo cambió a \emph{Living
Standard}: la especificación se actualiza continuamente sin saltos
de versión. Esto refleja la madurez del lenguaje: no necesita grandes
revisiones disruptivas, solo evoluciones incrementales.

\begin{table}[!htb]
\centering
\caption{Hitos de la evolución de HTML.}
\label{tab:evolucion-html}
\footnotesize
\begin{tabularx}{\textwidth}{lll}
\toprule
\textbf{Versión} & \textbf{Año} & \textbf{Aportes principales}\\
\midrule
HTML «Tags» & 1991 & 18 etiquetas iniciales (\texttt{<a>}, \texttt{<p>}, \texttt{<h1>}\dots).\\
HTML 2.0 & 1995 & Primera estandarización formal (RFC 1866); formularios.\\
HTML 3.2 & 1997 & Tablas, applets, fuentes, alineación.\\
HTML 4.01 & 1999 & Marcos, scripts, separación de presentación (CSS).\\
XHTML 1.0 & 2000 & Reformulación XML de HTML 4.01.\\
XHTML 1.1 & 2001 & Modularización; estricta sintaxis XML.\\
HTML5 borrador & 2008 & WHATWG inicia HTML5; nuevos elementos semánticos.\\
HTML5 W3C Rec. & 2014 & Recomendación oficial; \texttt{<video>}, \texttt{<canvas>}, APIs.\\
HTML 5.1 & 2016 & Mejoras menores (\texttt{<picture>}, \texttt{<details>}).\\
HTML 5.2 & 2017 & \texttt{<dialog>}, modal nativo.\\
HTML 5.3 & --- & Cancelado en favor del Living Standard.\\
HTML Living Standard & desde 2019 & Estado actual; cambios continuos por WHATWG.\\
\bottomrule
\end{tabularx}
\end{table}

La trayectoria muestra tres etapas claramente diferenciadas:
\emph{(a)} la etapa fundacional (1991--1999) donde HTML pasa de
prototipo a estándar formal; \emph{(b)} el periodo de transición
XML (2000--2008) en el que el W3C intentó migrar HTML al ecosistema
XML con XHTML 1.0 y 2.0, esfuerzo que finalmente fracasó porque la
web real era demasiado tolerante a errores
\cite{hickson2014html5}; y \emph{(c)} la era moderna de HTML5 (desde
2008) en la que el lenguaje recupera su pragmatismo original mientras
incorpora capacidades multimedia, APIs avanzadas y semántica rica.

\section{Fundamentos de HTML}

HTML (\emph{HyperText Markup Language}) es un lenguaje de marcado
declarativo cuyo propósito es describir la estructura semántica de
un documento hipertextual. A diferencia de los lenguajes
procedurales o imperativos, HTML no «hace» nada: describe
\emph{qué es} cada parte del contenido y deja al navegador la
responsabilidad de renderizarlo. Como argumentan Duckett
\cite{duckett2011html} y Robbins \cite{robbins2018learning}, esta
naturaleza declarativa es lo que ha permitido que HTML sobreviva 35
años casi sin cambios sintácticos: agentes muy diversos
(navegadores, lectores de pantalla, motores de búsqueda, crawlers,
extractores de datos, asistentes de IA) pueden interpretar el mismo
documento de maneras distintas pero coherentes.

Sus tres pilares son:

\begin{itemize}
\item \textbf{Etiquetas (tags).} Construyen la estructura. Vienen
en pares de apertura y cierre: \texttt{<p>texto</p>}. Algunas son
\emph{void} (no contienen contenido): \texttt{<img>}, \texttt{<br>},
\texttt{<input>}.
\item \textbf{Atributos.} Aportan información adicional sobre los
elementos: \texttt{class}, \texttt{id}, \texttt{lang},
\texttt{aria-*}. Se escriben dentro de la etiqueta de apertura como
pares \texttt{nombre=\char34 valor\char34 }.
\item \textbf{Contenido.} El texto y los elementos hijos que se
anidan dentro de las etiquetas.
\end{itemize}

\subsection{Estructura mínima de un documento HTML5 válido}

Todo documento HTML5 válido debe contener cinco elementos
imprescindibles: \emph{(i)} la declaración \texttt{<!DOCTYPE html>}
en la primera línea; \emph{(ii)} un elemento raíz \texttt{<html>} con
atributo \texttt{lang}; \emph{(iii)} una sección \texttt{<head>} con
metadatos (al menos \texttt{<meta charset>} y \texttt{<title>});
\emph{(iv)} una sección \texttt{<body>} con el contenido visible; y
\emph{(v)} idealmente, una estructura semántica que separe cabecera,
navegación, contenido principal y pie.

El siguiente código presenta una página HTML5 completa que ilustra
el uso de \textbf{todos} los tags principales: tanto los nuevos tags
semánticos de HTML5 (\texttt{<header>}, \texttt{<nav>},
\texttt{<main>}, \texttt{<section>}, \texttt{<article>},
\texttt{<aside>}, \texttt{<footer>}, \texttt{<figure>},
\texttt{<time>}) como los tags clásicos que siguen siendo
fundamentales (\texttt{<div>}, \texttt{<table>}, \texttt{<form>},
\texttt{<ul>}, \texttt{<a>}, \texttt{<img>}). Este código sirve como
\textbf{patrón canónico} de referencia para construir cualquier
página web.

\begin{cajaejemplo}{Ejemplo 2.1 --- Documento HTML5 completo con tags clásicos y semánticos}

\textbf{Enunciado.} Construir un documento HTML5 mínimo pero completo que incluya todos los tags semánticos esenciales (\texttt{<header>}, \texttt{<nav>}, \texttt{<main>}, \texttt{<article>}, \texttt{<section>}, \texttt{<aside>}, \texttt{<footer>}) y los tags clásicos de organización (\texttt{<div>}, \texttt{<table>}). Servirá como página de referencia para el resto de los ejercicios.

\begin{lstlisting}[style=html,caption={pagina-completa.html},label=lst:2.1]
<!DOCTYPE html>
<html lang="es-EC">
<head>
  <meta charset="UTF-8">
  <meta name="viewport" content="width=device-width, initial-scale=1">
  <meta name="description" content="Pagina demo con todos los tags clave">
  <meta name="author" content="UTEQ - Aplicaciones Web">
  <title>Pagina completa con tags clave - UTEQ</title>
</head>
<body>

  <!-- Tags HTML5 semanticos: estructura -->
  <header>
    <h1>Universidad Tecnica Estatal de Quevedo</h1>
    <p>La primera universidad agropecuaria del Ecuador</p>
    <nav aria-label="Menu principal">
      <ul>
        <li><a href="#inicio">Inicio</a></li>
        <li><a href="#carreras">Carreras</a></li>
        <li><a href="#contacto">Contacto</a></li>
      </ul>
    </nav>
  </header>

  <main>
    <article id="inicio">
      <h2>Bienvenidos al sitio web</h2>
      <p>Publicado el <time datetime="2026-05-04">4 de mayo
      de 2026</time>.</p>
      <figure>
        <img src="campus.jpg"
             alt="Vista panoramica del campus UTEQ"
             width="600" height="400">
        <figcaption>Campus central UTEQ - Quevedo</figcaption>
      </figure>
    </article>

    <section id="carreras">
      <h2>Carreras de la Facultad</h2>

      <!-- Tag clasico: tabla para datos tabulares -->
      <table>
        <caption>Carreras de Ciencias de la Computacion</caption>
        <thead>
          <tr><th scope="col">Carrera</th>
              <th scope="col">Duracion</th>
              <th scope="col">Titulo</th></tr>
        </thead>
        <tbody>
          <tr><td>Ingenieria de Software</td>
              <td>9 semestres</td>
              <td>Ing. de Software</td></tr>
          <tr><td>Ingenieria en Sistemas</td>
              <td>9 semestres</td>
              <td>Ing. de Sistemas</td></tr>
        </tbody>
      </table>
    </section>

    <section id="contacto">
      <h2>Contactenos</h2>

      <!-- Tag clasico: formulario -->
      <form action="/api/contacto" method="POST">
        <!-- Tag clasico: div como contenedor generico -->
        <div>
          <label for="nombre">Nombre:</label>
          <input id="nombre" name="nombre" type="text" required>
        </div>
        <div>
          <label for="email">Correo:</label>
          <input id="email" name="email" type="email" required>
        </div>
        <div>
          <label for="msg">Mensaje:</label>
          <textarea id="msg" name="msg" rows="4"></textarea>
        </div>
        <button type="submit">Enviar</button>
      </form>
    </section>

    <aside aria-label="Datos adicionales">
      <h2>Datos rapidos</h2>
      <ul>
        <li>Fundacion: <strong>1984</strong></li>
        <li>Estudiantes: <strong>13.000+</strong></li>
        <li>Programas: <strong>22</strong></li>
      </ul>
    </aside>
  </main>

  <footer>
    <small>&copy; 2026 UTEQ. Todos los derechos reservados.</small>
    <address>
      Av. Quito km 1.5 via a Santo Domingo<br>
      Telefono: <a href="tel:+59353702220">(+593) 5 3702-220</a>
    </address>
  </footer>

</body>
</html>
\end{lstlisting}

\textbf{Explicación de las decisiones tomadas:}
\begin{itemize}
\item \texttt{<header>}, \texttt{<nav>}, \texttt{<main>},
\texttt{<article>}, \texttt{<section>}, \texttt{<aside>},
\texttt{<footer>}: la estructura semántica de HTML5 reemplaza el uso
clásico de \texttt{<div class=\char34 header\char34 >}, \texttt{<div id=\char34 content\char34 >}.
\item \texttt{<table>}: usada \emph{solo} para datos tabulares, con
\texttt{<caption>}, \texttt{<thead>}, \texttt{<tbody>} y atributo
\texttt{scope} en los \texttt{<th>}, conforme a las prácticas
accesibles \cite{robbins2018learning}.
\item \texttt{<form>}, \texttt{<input>}, \texttt{<label>}: cada
\texttt{<input>} tiene su \texttt{<label for>} asociado (criterio
WCAG 1.3.1).
\item \texttt{<div>}: usado solo cuando no hay un elemento semántico
apropiado, como contenedor neutro para cada grupo \texttt{<label>+<input>}.
\item \texttt{<time datetime>}: fecha en formato ISO 8601 legible
por máquinas.
\end{itemize}
\end{cajaejemplo}

\subsection{Las cinco categorías de contenido en HTML5}

HTML5 abandona la dicotomía clásica \emph{block-level / inline}
(propia de HTML 4) y la sustituye por las \textbf{categorías de
contenido}, que se solapan y permiten razonar con mayor precisión
sobre qué puede contener qué \cite{hickson2014html5}. Este cambio
conceptual es significativo: antes, los desarrolladores aprendían
reglas mecánicas («un \texttt{<p>} no puede ir dentro de un
\texttt{<a>}»); ahora aprenden categorías («un elemento de
\emph{phrasing content} no admite \emph{flow content}»), lo que
generaliza mucho mejor.

A continuación se presentan las cinco categorías principales con
ejemplos de uso en una página web real:

\begin{enumerate}
\item \textbf{Flow content} (la mayoría del contenido del cuerpo):
texto, párrafos, listas, encabezados, imágenes, formularios,
secciones. Casi todo lo que puede ir en \texttt{<body>} es flow
content. \emph{Ejemplo de uso:} dentro de un \texttt{<section>}
podemos poner \texttt{<h2>}, \texttt{<p>}, \texttt{<ul>},
\texttt{<table>}, \texttt{<form>}: todos son flow content.

\item \textbf{Phrasing content} (texto y sus marcas en línea):
\texttt{<em>}, \texttt{<strong>}, \texttt{<a>}, \texttt{<span>},
\texttt{<code>}, \texttt{<time>}, \texttt{<mark>}. Es un subconjunto
de \emph{flow}. \emph{Ejemplo de uso:} dentro de un párrafo podemos
escribir \texttt{<p>El libro <em>1984</em> fue publicado en
<time datetime=\char34 1949\char34 >1949</time>.</p>}.

\item \textbf{Embedded content} (contenido externo o interactivo
incrustado): \texttt{<img>}, \texttt{<video>}, \texttt{<audio>},
\texttt{<canvas>}, \texttt{<iframe>}, \texttt{<object>},
\texttt{<svg>}. \emph{Ejemplo de uso:} dentro de un \texttt{<figure>}
incrustamos \texttt{<img>} y agregamos \texttt{<figcaption>}.

\item \textbf{Interactive content} (elementos con los que el usuario
interactúa): \texttt{<a>} con \texttt{href}, \texttt{<button>},
\texttt{<input>} (excepto \texttt{type=\char34 hidden\char34 }), \texttt{<select>},
\texttt{<textarea>}, \texttt{<details>}, \texttt{<label>}.
\emph{Ejemplo de uso:} dentro de un \texttt{<form>} ponemos
varios \texttt{<input>} interactivos.

\item \textbf{Sectioning content} (secciones del documento):
\texttt{<article>}, \texttt{<section>}, \texttt{<nav>},
\texttt{<aside>}. Cada uno define un nuevo «contexto de outline» que
afecta la jerarquía de encabezados. \emph{Ejemplo de uso:} la
página de la sección anterior tiene un \texttt{<article>} (sectioning
content) que contiene un \texttt{<h2>} de bienvenida.
\end{enumerate}

Existen además categorías menos comunes: \emph{metadata content}
(\texttt{<meta>}, \texttt{<link>}, \texttt{<style>}, \texttt{<title>}),
\emph{heading content} (\texttt{<h1>}-\texttt{<h6>}), \emph{form-associated
content} (\texttt{<input>}, \texttt{<button>}, etc.). Un mismo
elemento puede pertenecer a varias categorías simultáneamente: por
ejemplo, \texttt{<a>} es flow, phrasing e interactive.

\subsection{Modelo de contenido y reglas de anidamiento}

Cada elemento HTML define qué puede contener y qué no, mediante un
\textbf{modelo de contenido} declarado en la especificación. Estas
reglas no son arbitrarias: derivan de la semántica del elemento y de
limitaciones técnicas históricas. Comprenderlas evita errores
sutiles que pasan desapercibidos a primera vista pero rompen la
accesibilidad o el SEO.

\textbf{Reglas fundamentales con ejemplos:}

\begin{enumerate}
\item \textbf{\texttt{<p>} solo admite phrasing content.} Por eso un
párrafo no puede contener un \texttt{<div>} o un \texttt{<table>}.
\textbf{Anti-ejemplo (inválido) (Listado~\ref{lst:2.2}):}
\begin{lstlisting}[style=html,numbers=none,caption={Ejemplo de HTML},label=lst:2.2]
<p>Aqui hay una tabla:
  <table><tr><td>X</td></tr></table>
</p>
\end{lstlisting}
\textbf{Correcto:} cerrar el \texttt{<p>} antes de la tabla.

\item \textbf{\texttt{<a>} acepta cualquier contenido excepto
interactive content.} Por eso un enlace no puede contener otro
enlace, ni un \texttt{<button>}, ni un \texttt{<input>}.
\textbf{Anti-ejemplo (inválido) (Listado~\ref{lst:2.3}):}
\begin{lstlisting}[style=html,numbers=none,caption={Ejemplo de HTML},label=lst:2.3]
<a href="/x"><button>Click</button></a>
\end{lstlisting}
\textbf{Correcto:} usar solo \texttt{<a>} con texto y CSS para
darle apariencia de botón, o solo \texttt{<button>} con
\texttt{onclick} que redirija.

\item \textbf{\texttt{<button>} no admite interactive content.} Por
eso no se puede tener un \texttt{<a>} dentro de un \texttt{<button>}.

\item \textbf{\texttt{<ul>} y \texttt{<ol>} solo aceptan
\texttt{<li>}} (y opcionalmente \texttt{<script>}/\texttt{<template>}).
\textbf{Anti-ejemplo (inválido) (Listado~\ref{lst:2.4}):}
\begin{lstlisting}[style=html,numbers=none,caption={Ejemplo de HTML},label=lst:2.4]
<ul>
  <h2>Items</h2>
  <li>Uno</li>
  <li>Dos</li>
</ul>
\end{lstlisting}
\textbf{Correcto:} el encabezado va \emph{antes} de la \texttt{<ul>}.

\item \textbf{\texttt{<table>} requiere estructura específica:}
solo admite \texttt{<caption>}, \texttt{<colgroup>}, \texttt{<thead>},
\texttt{<tbody>}, \texttt{<tfoot>}, \texttt{<tr>} (este último solo
sin \texttt{<thead>}/\texttt{<tbody>}). Las filas \texttt{<tr>} solo
admiten \texttt{<td>} y \texttt{<th>}.

\item \textbf{\texttt{<form>} no puede contener otro \texttt{<form>}.}
Pero sí puede tener \texttt{<input form=\char34 otroForm\char34 >} que asocia un
input con un form externo (técnica avanzada).
\end{enumerate}

\textbf{¿Por qué el navegador es indulgente?} Los navegadores
intentan recuperarse de HTML inválido por compatibilidad histórica.
Sin embargo:
\begin{itemize}
\item Los \emph{validadores} (W3C) reportan errores.
\item Los \emph{lectores de pantalla} pueden anunciar contenido
incoherente.
\item Los \emph{motores de SEO} penalizan estructuras inválidas.
\item Las \emph{herramientas automatizadas} (extracción de datos,
asistentes de IA) producen resultados impredecibles.
\end{itemize}

Por estas razones, todo HTML profesional debe pasar el W3C Markup
Validator \cite{w3cvalidator} sin errores antes de ir a producción.

\section{Elementos semánticos de HTML5 en profundidad}

Antes de HTML5 la mayoría de la maquetación se hacía con \texttt{<div>}
genéricos diferenciados por clases CSS. Este enfoque
\emph{semánticamente ciego} tenía varios problemas documentados
\cite{lawson2011html5}: los motores de búsqueda no entendían la
estructura del contenido, los lectores de pantalla no podían
ofrecer navegación rápida entre regiones, y el código se volvía
ilegible para humanos. HTML5 introdujo un vocabulario semántico
explícito que mejora la accesibilidad, el SEO, la mantenibilidad y
la interoperabilidad con herramientas de \emph{reader mode},
extracción automatizada y asistentes de IA.

\subsection{Catálogo de elementos semánticos}

La tabla \ref{tab:semanticos} sintetiza los elementos semánticos
esenciales de HTML5 con su uso correcto e incorrecto. La razón
para incluir ambas columnas es pedagógica: en los entregables del
curso, los errores más frecuentes consisten en usar elementos
semánticos para fines genéricos (por ejemplo, \texttt{<section>}
como contenedor visual sin encabezado), lo que invalida su
semántica.

\begin{table}[htbp]
\centering
\caption{Elementos semánticos esenciales de HTML5 con su uso correcto e incorrecto.}
\label{tab:semanticos}
\footnotesize
\begin{tabularx}{\textwidth}{lXX}
\toprule
\textbf{Elemento} & \textbf{Uso correcto} & \textbf{Uso incorrecto}\\
\midrule
\texttt{<header>}  & Cabecera de página o sección & Como contenedor genérico\\
\texttt{<nav>}     & Bloque de navegación principal & Para listas de enlaces secundarios\\
\texttt{<main>}    & Contenido principal único & En múltiples instancias\\
\texttt{<article>} & Contenido autónomo y reutilizable & Para cualquier párrafo\\
\texttt{<section>} & Agrupación temática con encabezado & Sin \texttt{<h2>}/\texttt{<h3>} dentro\\
\texttt{<aside>}   & Contenido relacionado pero secundario & Para barras laterales decorativas\\
\texttt{<footer>}  & Pie de página o sección & Para contenido principal\\
\texttt{<figure>}  & Imagen, código o tabla autocontenida & Para cualquier imagen\\
\texttt{<figcaption>} & Pie de figura & Fuera de \texttt{<figure>}\\
\texttt{<time>}    & Fecha u hora marcada para máquinas & Sin atributo \texttt{datetime}\\
\texttt{<mark>}    & Texto destacado contextualmente & Como sinónimo de \texttt{<strong>}\\
\texttt{<details>}+\texttt{<summary>} & Bloque colapsable & Para acordeones complejos\\
\texttt{<dialog>}  & Modal nativo & Para tooltips simples\\
\texttt{<address>} & Datos de contacto & Para direcciones postales arbitrarias\\
\bottomrule
\end{tabularx}
\end{table}

Para cada elemento, la regla mnemotécnica es: \emph{«si puedo
explicar en una frase qué significa este bloque, debo poder elegir
el elemento semántico apropiado; si solo puedo describirlo por su
apariencia visual, probablemente debo usar un \texttt{<div>}»}.

\subsection{¿\texttt{<article>} o \texttt{<section>}? Guía decisional con ejemplo real}

La confusión más común es cuándo usar \texttt{<article>} y cuándo
\texttt{<section>}. Ambos son \emph{sectioning content}, ambos
crean un nuevo contexto de outline, pero tienen semánticas distintas
\cite{lawson2011html5,robbins2018learning}.

\textbf{Regla heurística canónica:}

\begin{itemize}
\item \texttt{<article>}: si el bloque de contenido podría
\emph{publicarse de manera independiente} en otro contexto y seguir
teniendo sentido, es un artículo. La «prueba RSS»: ¿tendría sentido
incluir este bloque en un feed RSS aislado de su contexto?

Ejemplos canónicos: una entrada de blog, una noticia, un comentario
en un foro, una tarjeta de producto, un \emph{tweet}, una reseña.

\item \texttt{<section>}: si el bloque es una agrupación temática
que solo tiene sentido como parte del documento mayor.

Ejemplos: la sección «Estudios» de un currículum, el capítulo de un
libro, la sección «Características» de una landing page.
\end{itemize}

\textbf{Ejemplo real con código completo:} blog con múltiples entradas.

\begin{cajaejemplo}{Ejemplo 2.2 --- Uso correcto de \texttt{<article>} y \texttt{<section>}}
\begin{lstlisting}[style=html,caption={blog.html (extracto significativo)},label=lst:2.5]
<main>
  <h1>Blog de Aplicaciones Web - UTEQ</h1>

  <!-- ARTICLE: entrada que puede publicarse en RSS independiente -->
  <article>
    <header>
      <h2>Por que HTML5 cambio el desarrollo web</h2>
      <p>Por <a href="/autor/jdoe" rel="author">Juan Doe</a>
         el <time datetime="2026-05-04">4 mayo 2026</time></p>
    </header>

    <p>Introduccion al articulo...</p>

    <!-- SECTION dentro de article: agrupacion tematica interna -->
    <section>
      <h3>Elementos semanticos nuevos</h3>
      <p>HTML5 introdujo header, nav, main...</p>
    </section>

    <section>
      <h3>APIs multimedia nativas</h3>
      <p>Las etiquetas video y audio reemplazaron Flash...</p>
    </section>

    <footer>
      <p>Etiquetas:
        <a href="/tag/html5">HTML5</a>,
        <a href="/tag/web">Web</a>
      </p>
    </footer>
  </article>

  <!-- Otro ARTICLE independiente -->
  <article>
    <header>
      <h2>CSS Grid: el sistema de maquetacion definitivo</h2>
      <p>Por <a href="/autor/asmith">Ana Smith</a>
         el <time datetime="2026-04-28">28 abril 2026</time></p>
    </header>
    <p>CSS Grid revoluciono la maquetacion web...</p>
  </article>

  <!-- SECTION: agrupacion que NO tendria sentido fuera del blog -->
  <section aria-label="Articulos relacionados">
    <h2>Tambien te puede interesar</h2>
    <ul>
      <li><a href="/post/1">Una guia rapida de Flexbox</a></li>
      <li><a href="/post/2">JavaScript moderno en 2026</a></li>
    </ul>
  </section>
</main>
\end{lstlisting}

\textbf{Razonamiento de las decisiones:}
\begin{itemize}
\item Cada entrada del blog es \texttt{<article>}: si se extrae y
publica en otro sitio o en un feed, tiene sentido por sí misma.
\item Dentro de cada artículo, las \texttt{<section>} agrupan
subtemas pero \emph{no} podrían publicarse independientemente: «APIs
multimedia nativas» sin el contexto del artículo carece de
significado.
\item «Artículos relacionados» es una \texttt{<section>} porque es
una agrupación temática (lista de enlaces) que pertenece a la página
del blog pero no es \emph{contenido} autónomo.
\end{itemize}
\end{cajaejemplo}

\section{Formularios HTML5 y validación nativa}

Los formularios constituyen el principal mecanismo de entrada de
datos del usuario hacia el servidor. HTML5 amplió drásticamente las
capacidades de validación nativa, reduciendo la dependencia de
JavaScript para tareas básicas y mejorando significativamente la
experiencia del usuario en dispositivos móviles
\cite{nixon2021learning}. Con HTML5 podemos confiar en el navegador
para verificar formato de email, rangos numéricos, longitudes
mínimas y máximas, y patrones complejos mediante expresiones
regulares.

Es importante destacar: \textbf{la validación HTML5 es para
usabilidad, no para seguridad}. Un atacante puede saltarse fácilmente
las validaciones del cliente; por eso la validación robusta debe
duplicarse en el servidor (semanas 5--9). Sin embargo, la
validación HTML5 mejora la experiencia ofreciendo retroalimentación
inmediata al usuario \emph{antes} de enviar el formulario.

\subsection{Tipos de input modernos}

A los tipos clásicos de input que existían desde HTML 2.0
(\texttt{text}, \texttt{password}, \texttt{checkbox}, \texttt{radio},
\texttt{submit}, \texttt{hidden}, \texttt{file}), HTML5 añadió
\textbf{trece tipos modernos} que tres ventajas operacionales
fundamentales: \emph{(i)} validación automática del formato (correo,
URL, número), \emph{(ii)} interfaces de entrada especializadas
(calendarios, selectores de color, sliders) y \emph{(iii)} en
móviles, invocan automáticamente el teclado virtual más apropiado
para el tipo de dato (numérico para \texttt{tel}, con \texttt{@}
para \texttt{email}, etc.).

La tabla \ref{tab:input-types} cataloga estos tipos modernos junto
con su uso recomendado y la validación nativa que cada uno
implementa. Conviene memorizar esta tabla: usar el tipo correcto de
input es una de las maneras más simples y efectivas de mejorar la
usabilidad de un formulario.

\begin{table}[htbp]
\centering
\caption{Tipos de \texttt{<input>} de HTML5 y su validación nativa.}
\label{tab:input-types}
\footnotesize
\begin{tabularx}{\textwidth}{lXX}
\toprule
\textbf{Tipo} & \textbf{Uso} & \textbf{Validación nativa}\\
\midrule
\texttt{email} & Correo electrónico & Formato \texttt{x@y.z}\\
\texttt{tel} & Teléfono & No valida formato (varía por país); usar \texttt{pattern}\\
\texttt{url} & URL & Formato URL absoluto\\
\texttt{number} & Número & Solo dígitos; \texttt{min}/\texttt{max}/\texttt{step}\\
\texttt{range} & Rango (slider) & Numérico; sin teclado de entrada\\
\texttt{date} & Fecha & Calendario nativo\\
\texttt{time} & Hora & Selector horario\\
\texttt{datetime-local} & Fecha y hora local & Calendario + reloj\\
\texttt{month} & Año-mes & Selector de mes\\
\texttt{week} & Semana del año & Selector de semana\\
\texttt{color} & Color & Selector de color\\
\texttt{search} & Búsqueda & Cosmética: borrar con X\\
\bottomrule
\end{tabularx}
\end{table}

En móviles, además, cada tipo de input invoca el teclado virtual
adecuado: \texttt{email} muestra el carácter \texttt{@};
\texttt{number} y \texttt{tel} muestran el teclado numérico;
\texttt{url} muestra \texttt{.com} y \texttt{/}. Esto reduce
drásticamente los errores de tipeo en pantallas pequeñas.

\subsection{Atributos de validación}

Los siguientes atributos permiten especificar restricciones que el
navegador valida automáticamente antes de enviar el formulario:

\begin{itemize}
\item \texttt{required}: el campo no puede enviarse vacío.
\item \texttt{minlength} / \texttt{maxlength}: rango de longitud de
texto.
\item \texttt{min} / \texttt{max}: rango numérico o temporal.
\item \texttt{pattern}: expresión regular (sintaxis JavaScript) que
debe cumplir el valor.
\item \texttt{step}: incremento permitido en numéricos y fechas
(\texttt{step=\char34 0.01\char34 } para dinero, \texttt{step=\char34 60\char34 } para minutos).
\item \texttt{autocomplete}: directrices al navegador (\texttt{on},
\texttt{off}, \texttt{name}, \texttt{email}, \texttt{tel},
\texttt{cc-number}, \texttt{street-address}, \texttt{country},
\texttt{bday}\dots).
\item \texttt{novalidate} (en \texttt{<form>}): deshabilita la
validación nativa.
\item \texttt{multiple}: permite múltiples valores
(\texttt{<input type=\char34 email\char34  multiple>}).
\item \texttt{accept}: tipos MIME permitidos en \texttt{<input
type=\char34 file\char34 >}.
\end{itemize}

\textbf{Ejemplo de formulario completo con validación nativa:}

\begin{cajaejemplo}{Ejemplo 2.3 --- Formulario de inscripción a curso con validación HTML5 completa}

\textbf{Enunciado.} Construir un formulario de inscripción a curso con validación HTML5 nativa (\texttt{required}, \texttt{type}, \texttt{pattern}, \texttt{min}/\texttt{max}) que rechace datos inválidos antes de enviarlos al servidor.

\begin{lstlisting}[style=html,caption={inscripcion.html (formulario funcional)},label=lst:2.6]
<form action="/api/inscripcion" method="POST">
  <fieldset>
    <legend>Datos personales</legend>

    <label for="cedula">Cedula (10 digitos):</label>
    <input id="cedula" name="cedula" type="text"
           pattern="[0-9]{10}" required
           autocomplete="off"
           aria-describedby="ayuda-cedula">
    <small id="ayuda-cedula">Sin guiones ni espacios</small>

    <label for="nombre">Nombre completo:</label>
    <input id="nombre" name="nombre" type="text"
           minlength="3" maxlength="80" required
           autocomplete="name">

    <label for="email">Correo institucional:</label>
    <input id="email" name="email" type="email"
           pattern=".+@uteq\.edu\.ec" required
           autocomplete="email">

    <label for="tel">Telefono celular:</label>
    <input id="tel" name="tel" type="tel"
           pattern="09[0-9]{8}" required
           autocomplete="tel">

    <label for="fnac">Fecha de nacimiento:</label>
    <input id="fnac" name="fnac" type="date"
           min="1980-01-01" max="2010-12-31" required
           autocomplete="bday">
  </fieldset>

  <fieldset>
    <legend>Curso</legend>

    <label for="nivel">Nivel academico (1-10):</label>
    <input id="nivel" name="nivel" type="number"
           min="1" max="10" step="1" required>

    <label for="prom">Promedio acumulado:</label>
    <input id="prom" name="prom" type="number"
           min="6.0" max="10.0" step="0.01" required>

    <label for="hrs">Horas semanales de estudio:</label>
    <input id="hrs" name="hrs" type="range"
           min="0" max="40" step="1" value="20">

    <label for="carrera">Carrera:</label>
    <select id="carrera" name="carrera" required>
      <option value="">-- Seleccione --</option>
      <option value="SW">Ingenieria de Software</option>
      <option value="SE">Ingenieria de Sistemas</option>
    </select>
  </fieldset>

  <button type="submit">Inscribirme</button>
  <button type="reset">Limpiar</button>
</form>
\end{lstlisting}

\textbf{Validaciones automáticas que aplica el navegador:}
\begin{itemize}
\item Si la cédula no tiene exactamente 10 dígitos: error nativo.
\item Si el email no termina en \texttt{@uteq.edu.ec}: error nativo.
\item Si el teléfono no comienza con \texttt{09}: error nativo.
\item Si el promedio está fuera de rango 6.0--10.0: error nativo.
\item Si el campo carrera no se selecciona: error nativo.
\end{itemize}
Todo esto sin escribir una sola línea de JavaScript.
\end{cajaejemplo}

\textbf{Nota didáctica:} las pseudo-clases CSS para dar
retroalimentación visual al estado de validación (\texttt{:valid},
\texttt{:invalid}, \texttt{:required}, \texttt{:user-invalid}) se
estudiarán en detalle en la semana 3, cuando dispongamos del marco
conceptual completo de CSS. Por ahora, basta con saber que existen y
que podemos darles estilo al campo según esté correctamente lleno o
no.

\section{Multimedia y APIs de HTML5}

HTML5 incorporó al estándar nativo cinco APIs de gran importancia
práctica que reemplazaron tecnologías propietarias como Flash y
Silverlight \cite{macdonald2018css}. Su importancia no es solo
técnica: políticamente representó la victoria de los estándares
abiertos sobre los plugins propietarios. Apple inició la batalla en
2010 cuando Steve Jobs publicó la carta abierta «Thoughts on
Flash» \cite{jobs2010flash} declarando que iOS nunca soportaría
Flash. La transición a HTML5 nativo se completó alrededor de 2020,
cuando Adobe descontinuó oficialmente Flash Player.

Las cinco APIs nativas son:

\begin{enumerate}
\item \texttt{<video>} y \texttt{<audio>}: reproducción multimedia
nativa.
\item \texttt{<canvas>}: gráficos 2D y 3D (con WebGL).
\item \textbf{Geolocation API}: posición geográfica del usuario con
consentimiento.
\item \textbf{Web Storage}: \texttt{localStorage} y
\texttt{sessionStorage}.
\item \textbf{Drag-and-Drop API}: arrastrar y soltar entre
componentes.
\end{enumerate}

\subsection{La etiqueta \texttt{<video>} con accesibilidad completa}

La etiqueta \texttt{<video>} es uno de los elementos más complejos
de HTML5 porque integra múltiples preocupaciones: formato del video,
controles de reproducción, accesibilidad para usuarios sordos
(subtítulos), accesibilidad para usuarios ciegos (audiodescripciones),
optimización de carga (preload, poster) y compatibilidad con navegadores
distintos (múltiples \texttt{<source>}).

Para construir un video accesible profesional necesitamos conocer
los siguientes atributos y elementos asociados:

\begin{itemize}
\item \texttt{controls}: muestra los controles nativos del navegador
(reproducir, pausa, volumen, pantalla completa).
\item \texttt{preload}: indica al navegador cuánto descargar antes de
reproducir. Valores: \texttt{none} (nada), \texttt{metadata}
(solo cabeceras: duración y dimensiones), \texttt{auto} (todo).
\item \texttt{poster}: imagen que se muestra antes de iniciar la
reproducción (mejora la percepción de carga).
\item \texttt{autoplay}: reproducción automática (la mayoría de
navegadores la bloquean si el video tiene audio, por buenas
razones).
\item \texttt{loop}: reproducción en bucle.
\item \texttt{muted}: silenciado por defecto (necesario para
autoplay).
\item Múltiples \texttt{<source>}: el navegador usa el primer formato
que reconoce. Buena práctica: WebM primero (Chrome/Firefox), MP4
después (Safari/Edge).
\item \texttt{<track>}: pistas de subtítulos, descripciones,
metadatos. Sus tipos (\texttt{kind}):
   \begin{itemize}
   \item \texttt{captions}: subtítulos para usuarios sordos en el
   idioma del video (criterio WCAG 1.2.2).
   \item \texttt{subtitles}: subtítulos traducidos a otro idioma.
   \item \texttt{descriptions}: descripciones de audio para usuarios
   ciegos (criterio WCAG 1.2.5).
   \item \texttt{chapters}: marcadores de capítulos.
   \item \texttt{metadata}: datos para scripts.
   \end{itemize}
\end{itemize}

\begin{cajaejemplo}{Ejemplo 2.4 --- Video accesible profesional}

\textbf{Enunciado.} Mostrar la etiqueta \texttt{<video>} con todos los atributos de accesibilidad necesarios en producción: \texttt{controls}, \texttt{preload}, \texttt{poster}, pistas de subtítulos con \texttt{<track>} y fuentes alternativas con \texttt{<source>}.

\begin{lstlisting}[style=html,caption={Ejemplo de HTML},label=lst:2.7]
<figure>
  <video controls width="640"
         preload="metadata"
         poster="poster-uteq.jpg"
         crossorigin="anonymous">
    <source src="bienvenida.webm" type="video/webm">
    <source src="bienvenida.mp4"  type="video/mp4">
    <track kind="captions" src="bienvenida.es.vtt"
           srclang="es" label="Espanol" default>
    <track kind="captions" src="bienvenida.en.vtt"
           srclang="en" label="English">
    <track kind="descriptions" src="bienvenida.desc.vtt"
           srclang="es" label="Audiodescripcion">
    <p>Su navegador no soporta video HTML5.
       <a href="bienvenida.mp4">Descargar el video</a>.</p>
  </video>
  <figcaption>Mensaje de bienvenida del Rector --- Mayo 2026</figcaption>
</figure>
\end{lstlisting}

\textbf{Aspectos clave de este código:}
\begin{itemize}
\item Múltiples \texttt{<source>}: el navegador elige el primero
soportado.
\item Subtítulos en español y en inglés (\texttt{default} marca el
predeterminado).
\item \texttt{<track kind=\char34 descriptions\char34 >}: audiodescripciones para
usuarios ciegos.
\item Texto fallback dentro de \texttt{<video>}: navegadores muy
antiguos.
\item Envuelto en \texttt{<figure>+<figcaption>} para semántica
correcta.
\end{itemize}

\textbf{Formato del archivo \texttt{.vtt}} (WebVTT) para subtítulos (Listado~\ref{lst:2.8}):
\begin{lstlisting}[style=shell,caption={bienvenida.es.vtt},label=lst:2.8]
WEBVTT

00:00:00.000 --> 00:00:04.000
Bienvenidos a la Universidad Tecnica Estatal de Quevedo.

00:00:04.500 --> 00:00:08.000
Soy el Rector y les hablo desde nuestro campus central.
\end{lstlisting}
\end{cajaejemplo}

\subsection{El elemento \texttt{<picture>} para imágenes responsivas}

El elemento \texttt{<picture>} resuelve un problema concreto: servir
\emph{la imagen correcta para cada dispositivo}. Antes de
\texttt{<picture>}, la web servía la misma imagen a todos los
dispositivos, lo que penalizaba a usuarios móviles con redes lentas:
una imagen optimizada para desktop (1200 px de ancho) cargaba
también en un teléfono que solo necesitaba 400 px, gastando datos
innecesarios y prolongando tiempos de carga.

\texttt{<picture>} permite al navegador elegir la imagen óptima
según: \emph{(i)} el tamaño del viewport (media queries), \emph{(ii)}
la densidad de pixels (Retina, etc.) y \emph{(iii)} el formato
soportado (WebP es 25--35\% más pequeño que JPEG, AVIF es 50\% más
pequeño aún, pero no todos los navegadores los soportan).

\begin{cajaejemplo}{Ejemplo 2.5 --- Imagen optimizada para múltiples dispositivos}

\textbf{Enunciado.} Implementar una imagen responsive con \texttt{<picture>} y \texttt{<source>} para servir el formato y tamaño más apropiados según el dispositivo y soporte del navegador.

\begin{lstlisting}[style=html,caption={Ejemplo de HTML},label=lst:2.9]
<picture>
  <!-- Para pantallas grandes con AVIF soportado -->
  <source media="(min-width: 1200px)"
          srcset="banner-grande.avif"
          type="image/avif">
  <source media="(min-width: 1200px)"
          srcset="banner-grande.webp"
          type="image/webp">
  <source media="(min-width: 1200px)"
          srcset="banner-grande.jpg"
          type="image/jpeg">

  <!-- Para tablets -->
  <source media="(min-width: 600px)"
          srcset="banner-medio.webp"
          type="image/webp">

  <!-- Fallback para todos -->
  <img src="banner-pequeno.jpg"
       alt="Vista panoramica del campus UTEQ"
       width="1200" height="600"
       loading="lazy">
</picture>
\end{lstlisting}
\textbf{Beneficios:}
\begin{itemize}
\item El navegador descarga \emph{solo} la versión adecuada al
viewport y al formato soportado.
\item AVIF/WebP reducen el peso 25--50\% respecto a JPEG con calidad
equivalente.
\item \texttt{loading=\char34 lazy\char34 } difiere la descarga hasta que la
imagen está cerca del viewport.
\item Las dimensiones explícitas (\texttt{width}, \texttt{height})
evitan el \emph{layout shift} durante la carga.
\end{itemize}
\end{cajaejemplo}

\section{XML y familias de lenguajes de marcado}

HTML pertenece a una familia mayor de lenguajes de marcado que
descienden de SGML. XML (\emph{eXtensible Markup Language}) es el
otro gran descendiente de SGML, estandarizado por el W3C en 1998
\cite{w3c1998xml}. La diferencia fundamental: HTML tiene un
vocabulario fijo (los tags están predefinidos); XML es
\emph{extensible}, el autor define su propio vocabulario.

Esta extensibilidad permitió que XML se convirtiera en el formato
de intercambio dominante en sistemas empresariales durante los
2000s. Aunque desde 2015 ha sido desplazado parcialmente por JSON
en APIs REST modernas, XML sigue siendo esencial en muchos
contextos: configuración de aplicaciones Java (Maven, Spring),
servicios web SOAP, sindicación de contenidos, formatos de
documentos (DOCX, ODT, EPUB son ZIP de archivos XML).

\subsection{Subdialectos de XML relevantes en aplicaciones web}

XML como tecnología generó múltiples subdialectos
especializados que siguen vigentes en distintas áreas. Conocer cuáles
están en uso en aplicaciones web reales permite reconocerlos cuando
aparecen en bibliotecas, configuraciones o servicios externos. La
lista siguiente enumera los más relevantes en 2026.


\begin{itemize}
\item \textbf{XHTML}: HTML reformulado como XML estricto. Hoy
relegado pero existe en casos específicos (correos electrónicos,
algunos sistemas legados).
\item \textbf{SVG} (Scalable Vector Graphics): gráficos vectoriales,
nativos en HTML5. Iconos, gráficos, diagramas que escalan sin
pérdida de calidad.
\item \textbf{MathML}: notación matemática. Usada en sitios
académicos y educativos.
\item \textbf{RSS} y \textbf{Atom}: sindicación de contenido.
\item \textbf{SOAP}: servicios web XML (semana 15).
\item \textbf{XSLT}: transformaciones XML a otros formatos.
\item \textbf{DocBook}: documentación técnica.
\item \textbf{XAML}: interfaces de WPF y WinUI.
\end{itemize}

\subsection{Formatos de intercambio de datos: XML, JSON, YAML, TOML y otros}

XML y JSON no son los únicos formatos de intercambio de datos
relevantes en 2026. Para que el desarrollador profesional pueda
elegir bien, conviene conocer el panorama completo. Cada formato
tiene fortalezas y casos de uso óptimos.

\begin{table}[htbp]
\centering
\caption{Formatos de intercambio de datos modernos.}
\label{tab:formatos-datos}
\footnotesize
\begin{tabularx}{\textwidth}{lXX}
\toprule
\textbf{Formato} & \textbf{Características} & \textbf{Casos de uso 2026}\\
\midrule
XML & Verboso; etiquetas anidadas con atributos; validación con XSD; comentarios soportados. & Maven \texttt{pom.xml}, SOAP, SVG, configuración legada (Spring XML).\\
JSON & Sintaxis tipo objeto JS; tipos básicos; sin comentarios; soporte nativo en JavaScript. & APIs REST, \texttt{package.json}, NoSQL (MongoDB).\\
YAML & Indentación significativa; muy legible; tipos inferidos; soporta comentarios. & Docker Compose, Kubernetes, GitHub Actions, Spring Boot \texttt{application.yml}.\\
TOML & Sintaxis INI mejorada; tipado explícito; secciones; pensado para configuración. & Rust (\texttt{Cargo.toml}), Python (\texttt{pyproject.toml}).\\
INI & Formato \texttt{clave=valor} con secciones; muy simple. & Configuración legada en Windows, PHP \texttt{php.ini}.\\
CSV & Tabular plano; valores separados por comas. & Exportación a Excel, ETL básico.\\
Protobuf & Binario; eficiente; requiere esquema \texttt{.proto}; soporte multi-lenguaje. & gRPC, comunicación servicio-servicio de alto rendimiento.\\
MessagePack & Binario tipo JSON; más compacto. & Comunicación IoT, caches binarias.\\
Avro & Binario; esquema JSON embebido. & Apache Kafka, big data, streaming.\\
Parquet & Columnar binario; comprimido. & Análisis de datos a escala (Spark, BigQuery).\\
GraphQL & Lenguaje de consulta; respuesta JSON. & APIs con clientes que necesitan flexibilidad.\\
\bottomrule
\end{tabularx}
\end{table}

\textbf{Comparativa XML vs JSON (la elección más frecuente en
desarrollo web):}

\begin{table}[htbp]
\centering
\caption{XML vs JSON --- comparación detallada.}
\footnotesize
\begin{tabularx}{\textwidth}{lXX}
\toprule
\textbf{Criterio} & \textbf{XML} & \textbf{JSON}\\
\midrule
Verbosidad & Alta (etiquetas de apertura y cierre) & Baja (sintaxis tipo objeto)\\
Tipos de datos & Solo strings, requiere XSD para tipado & Strings, números, booleanos, null, arrays, objetos\\
Validación & XSD, DTD, RelaxNG & JSON Schema\\
Soporte JavaScript & Requiere parsing manual & \texttt{JSON.parse()} nativo\\
Atributos & Sí (atributos en elementos) & No (todo son propiedades)\\
Comentarios & Sí (\texttt{<!-- -->}) & No (JSON5 permite)\\
Tamaño típico & 100\% (referencia) & 60--70\% del XML equivalente\\
Casos de uso 2026 & Sistemas legados, configuración (Maven, Spring XML) & APIs REST, configuración (\texttt{package.json})\\
\bottomrule
\end{tabularx}
\end{table}

\textbf{Ejemplo comparativo: la misma información en cuatro formatos.}

\textbf{XML (Listado~\ref{lst:2.10}):}
\begin{lstlisting}[style=xml,numbers=none,caption={Ejemplo de XML},label=lst:2.10]
<usuario id="42">
  <nombre>Maria Lopez</nombre>
  <edad>23</edad>
  <activo>true</activo>
  <roles>
    <rol>USER</rol>
    <rol>EDITOR</rol>
  </roles>
</usuario>
\end{lstlisting}

\textbf{JSON (Listado~\ref{lst:2.11}):}
\begin{lstlisting}[style=js,numbers=none,caption={Ejemplo de JavaScript},label=lst:2.11]
{
  "id": 42,
  "nombre": "Maria Lopez",
  "edad": 23,
  "activo": true,
  "roles": ["USER", "EDITOR"]
}
\end{lstlisting}

\textbf{YAML (Listado~\ref{lst:2.12}):}
\begin{lstlisting}[numbers=none,basicstyle=\ttfamily\footnotesize,frame=single,caption={Ejemplo de Código},label=lst:2.12]
id: 42
nombre: Maria Lopez
edad: 23
activo: true
roles:
  - USER
  - EDITOR
\end{lstlisting}

\textbf{TOML (Listado~\ref{lst:2.13}):}
\begin{lstlisting}[style=shell,numbers=none,caption={Ejemplo de Shell},label=lst:2.13]
id = 42
nombre = "Maria Lopez"
edad = 23
activo = true
roles = ["USER", "EDITOR"]
\end{lstlisting}

\section{Buenas prácticas profesionales en HTML5}

Conocer la sintaxis no basta para escribir HTML profesional:
hay convenciones, criterios de mantenibilidad y restricciones de
accesibilidad que distinguen el código de producción del código
estudiantil. La caja siguiente compila quince reglas que el
estudiante debe respetar al entregar HTML5 en cualquier proyecto.


\begin{cajabuenas}{Quince reglas profesionales de HTML5}
\begin{enumerate}
\item \textbf{Usa el doctype HTML5 simple} \texttt{<!DOCTYPE html>}.
\item \textbf{Define el charset al inicio del \texttt{<head>}}, en
los primeros 1024 bytes del archivo.
\item \textbf{Define el viewport} en toda página que pretenda ser
mobile-friendly.
\item \textbf{Cierra todas las etiquetas} (incluso las que HTML5
permite no cerrar).
\item \textbf{Usa comillas dobles} en valores de atributos por
consistencia.
\item \textbf{No abuses de \texttt{<div>}}: si existe un elemento
semántico, úsalo.
\item \textbf{Atributos booleanos sin valor}.
\item \textbf{Imágenes con \texttt{alt}} siempre.
\item \textbf{Imágenes con \texttt{width} y \texttt{height}} para
evitar \emph{layout shift}.
\item \textbf{Lazy loading} en imágenes \emph{below the fold}.
\item \textbf{Formularios con \texttt{<label>}} asociados.
\item \textbf{Atributo \texttt{autocomplete}} con valores estándares.
\item \textbf{Headings sin saltos jerárquicos}.
\item \textbf{Validar} en \url{https://validator.w3.org/nu/} antes
de cada commit.
\item \textbf{No uses CSS inline} excepto en correos electrónicos.
\end{enumerate}
\end{cajabuenas}

\section{Ejercicios}

Los siguientes ejercicios consolidan los temas de la semana 2.
El primero está completamente resuelto con código y captura de
pantalla del resultado; el segundo es propuesto para que el
estudiante lo desarrolle.


\begin{cajaejercicio}{Ejercicio 2.1 (RESUELTO) --- Formulario de registro completo y accesible}
\textbf{Objetivo:} construir un formulario completo con ocho tipos
diferentes de \texttt{<input>}, validación HTML5, accesibilidad y
feedback visual.

\textbf{Paso 1 --- Crear proyecto en IntelliJ.}
\texttt{New Project $\rightarrow$ HTML $\rightarrow$ Name:
appweb-formularios}. Crear \texttt{registro.html}.

\textbf{Paso 2 --- Código completo.}

\begin{lstlisting}[style=html,caption={registro.html},label=lst:2.14]
<!DOCTYPE html>
<html lang="es-EC">
<head>
  <meta charset="UTF-8">
  <meta name="viewport" content="width=device-width, initial-scale=1">
  <title>Registro completo - UTEQ</title>
  <style>
    body { font-family: system-ui, sans-serif; max-width: 600px;
           margin: 2rem auto; padding: 0 1rem; }
    fieldset { margin-bottom: 1rem; padding: 1rem;
               border: 1px solid #aaa; border-radius: 6px; }
    legend { font-weight: bold; color: #006633; }
    label { display: block; margin-top: .8rem; font-weight: 600; }
    input, select, textarea {
      width: 100%; padding: .5rem; box-sizing: border-box;
      border: 1px solid #aaa; border-radius: 4px;
      font-size: 1rem; margin-top: .25rem;
    }
    input:user-invalid { border-color: #b00020; }
    input:user-valid { border-color: #2e7d32; }
    small { display: block; color: #555; font-size: .85rem; }
    button { margin-top: 1rem; padding: .7rem 1.2rem;
             background: #006633; color: white; border: 0;
             border-radius: 4px; font-size: 1rem; cursor: pointer; }
  </style>
</head>
<body>
<main>
  <h1>Registro de estudiantes UTEQ</h1>
  <form action="/api/registro" method="POST">
    <fieldset>
      <legend>Datos personales</legend>

      <label for="cedula">Cedula ecuatoriana:</label>
      <input id="cedula" name="cedula" type="text"
             pattern="[0-9]{10}" required>
      <small>10 digitos sin guiones</small>

      <label for="email">Correo institucional:</label>
      <input id="email" name="email" type="email"
             pattern=".+@uteq\.edu\.ec" required>

      <label for="tel">Telefono celular:</label>
      <input id="tel" name="tel" type="tel"
             pattern="09[0-9]{8}" required>

      <label for="fnac">Fecha de nacimiento:</label>
      <input id="fnac" name="fnac" type="date"
             min="1980-01-01" max="2010-12-31" required>

      <label for="url">Sitio web (opcional):</label>
      <input id="url" name="url" type="url">
    </fieldset>

    <fieldset>
      <legend>Datos academicos</legend>

      <label for="prom">Promedio (6.0-10.0):</label>
      <input id="prom" name="prom" type="number"
             min="6.0" max="10.0" step="0.01" required>

      <label for="color">Color institucional preferido:</label>
      <input id="color" name="color" type="color" value="#006633">

      <label for="carrera">Carrera:</label>
      <select id="carrera" name="carrera" required>
        <option value="">-- Seleccione --</option>
        <option value="SW">Ingenieria de Software</option>
        <option value="SE">Ingenieria de Sistemas</option>
      </select>
    </fieldset>

    <button type="submit">Registrarme</button>
  </form>
</main>
</body>
</html>
\end{lstlisting}

\textbf{Resultado esperado en el navegador:}

\begin{center}
\fbox{\begin{minipage}{0.85\textwidth}
\footnotesize
\textbf{[ Vista renderizada en Chrome ]}\\[0.3em]
\textbf{\large Registro de estudiantes UTEQ}\\[0.4em]
\rule{\textwidth}{0.4pt}\\
\textit{Datos personales}\\[0.3em]
Cedula ecuatoriana:\\
\fbox{\rule{0pt}{1em}\hspace{12cm}}\\
\textit{\scriptsize 10 dígitos sin guiones}\\[0.3em]
Correo institucional:\\
\fbox{\rule{0pt}{1em}\hspace{12cm}}\\[0.3em]
Teléfono celular:\\
\fbox{\rule{0pt}{1em}\hspace{12cm}}\\[0.3em]
Fecha de nacimiento:\\
\fbox{\rule{0pt}{1em}\hspace{12cm}} \textit{(calendario nativo)}\\[0.3em]
\rule{\textwidth}{0.4pt}\\
\textit{Datos académicos}\\[0.3em]
Promedio (6.0-10.0):\\
\fbox{\rule{0pt}{1em}\hspace{12cm}} \textit{(controles arriba/abajo)}\\[0.3em]
Color preferido: \colorbox{uteqverde}{\hspace{1cm}} \textit{(selector)}\\[0.3em]
Carrera: \fbox{-- Seleccione -- \quad\quad $\vee$}\\[0.5em]
\colorbox{uteqverde}{\textcolor{white}{\hspace{0.5em}\textbf{Registrarme}\hspace{0.5em}}}
\end{minipage}}
\end{center}

\textbf{Verificación de validaciones nativas:}
\begin{enumerate}
\item Pulsar «Registrarme» con campos vacíos: el navegador detiene el
envío y enfoca el primer campo inválido mostrando tooltip.
\item Escribir email con dominio diferente a \texttt{uteq.edu.ec}:
el borde se vuelve rojo (gracias a \texttt{:user-invalid}).
\item En móvil (F12 $\to$ Toggle device toolbar): hacer clic en
campo teléfono $\to$ aparece teclado numérico; en email $\to$ aparece
tecla \texttt{@}; en fecha $\to$ aparece selector calendario.
\end{enumerate}
\end{cajaejercicio}

\begin{cajaejercicio}{Ejercicio 2.2 (PROPUESTO) --- Página de portafolio profesional}
Construir una página \texttt{portafolio.html} que demuestre dominio
integral de HTML5 semántico. Requisitos:

\begin{itemize}
\item Estructura semántica completa: \texttt{<header>}, \texttt{<nav>},
\texttt{<main>}, \texttt{<article>} (al menos 3), \texttt{<aside>},
\texttt{<footer>}.
\item Al menos un \texttt{<figure>+<figcaption>} con
\texttt{<picture>} sirviendo múltiples resoluciones.
\item Un \texttt{<video>} con al menos un \texttt{<track>} de
subtítulos (crear archivo \texttt{.vtt} con 3--4 líneas).
\item Un formulario de contacto con al menos 6 tipos diferentes de
\texttt{<input>} HTML5: \texttt{text}, \texttt{email}, \texttt{tel},
\texttt{date}, \texttt{number}, \texttt{url}.
\item Un elemento \texttt{<details>+<summary>} para una sección de
preguntas frecuentes.
\item Una tabla \texttt{<table>} con \texttt{<caption>},
\texttt{<thead>}, \texttt{<tbody>}, \texttt{<th scope>}.
\end{itemize}

\textbf{Validar:} cero errores en W3C, cero violaciones críticas en
axe DevTools, navegación completa por teclado, y todos los tipos de
input invocando teclado virtual correcto en móvil.

\textbf{Entrega:} archivo HTML + screenshots de validadores.
\end{cajaejercicio}

%%% ==== SEMANA 03 v3 ====
% =====================================================================
%   SEMANA 3 - VERSIÓN V3 MEJORADA
%   CSS para el diseño y la presentación web
%   Aplica directrices: historia CSS ampliada, referencias contrastadas,
%   ejercicios resueltos con imagen + propuestos, tablas
% =====================================================================

\chapter{CSS para el diseño y la presentación web}
\label{cap:semana3}

\epigraph{«El diseño no es lo que se ve, es lo que funciona. CSS es
el puente entre la intención del diseñador y la experiencia del
usuario.»}{--- \textit{Adaptado de Steve Jobs}}

\section*{Resultado de aprendizaje}
Al concluir la semana 3, el estudiante diseña hojas de estilo CSS3
modernas aplicando metodologías de organización (BEM, mobile-first),
sistemas de \emph{layout} (Flexbox, Grid), preprocesadores (Sass) y
buenas prácticas profesionales \cite{macdonald2018css,frain2020responsive}.

\section{Historia y evolución de CSS}

CSS pasó en 30 años de un lenguaje minimalista para colores y
fuentes a un sistema robusto de diseño y \emph{layout} con soporte
nativo para grids, container queries y propiedades personalizadas
(\texttt{--variables}). Conocer las etapas de su evolución revela
por qué hoy podemos prescindir de frameworks pesados como Bootstrap
en muchos proyectos. La caja siguiente narra esa historia.


\begin{cajahistoria}{De CSS1 al motor de los \emph{design systems} modernos}
En 1994, \textbf{Håkon Wium Lie}, trabajando junto a Tim Berners-Lee
en el CERN, propuso un mecanismo para separar la presentación del
contenido. La idea era radical en su simplicidad: que los autores y
los usuarios pudieran proponer estilos en \emph{cascada}, donde el
navegador combinara las reglas según prioridades. La propuesta se
convertiría en \textbf{CSS Level 1} (1996), el primer estándar W3C de
hojas de estilo.

Durante quince años CSS evolucionó lentamente. CSS2 (1998) añadió
posicionamiento absoluto y \emph{media types}. La explosión llegó con
\textbf{CSS3} (modular, desde 2011 en adelante), que introdujo
gradientes, transiciones, animaciones, transformaciones,
\texttt{border-radius}, \texttt{box-shadow} y dos sistemas
revolucionarios: \textbf{Flexbox} (2012) y \textbf{Grid} (2017). Hoy
CSS es el motor de los \emph{design systems} modernos: Material
Design (Google), Fluent (Microsoft), Tailwind CSS, Bootstrap.
\end{cajahistoria}

\subsection{Tabla evolutiva de CSS}

La especificación de CSS ha avanzado en módulos
independientes desde 2011, lo que significa que distintas partes del
estándar maduran a velocidades distintas. La tabla siguiente resume
los hitos más relevantes para el desarrollador web 2026.


\begin{table}[htbp]
\centering
\caption{Hitos en la evolución de CSS.}
\footnotesize
\begin{tabularx}{\textwidth}{llX}
\toprule
\textbf{Año} & \textbf{Versión} & \textbf{Aporte}\\
\midrule
1996 & CSS 1     & Selectores básicos, colores, fuentes, márgenes.\\
1998 & CSS 2     & Posicionamiento (absolute, relative), z-index.\\
2004 & CSS 2.1   & Revisión de CSS 2 que es la usada en navegadores antiguos.\\
2011 & CSS 3 (modular) & Transiciones, transformaciones, animaciones, gradientes.\\
2012 & Flexbox CR  & Layout unidimensional con propiedades de eje.\\
2017 & CSS Grid CR & Layout bidimensional con áreas nombradas.\\
2020 & Custom Properties & Variables CSS nativas (\texttt{-{}-color: \#006633}).\\
2022 & Container Queries & Estilos basados en el contenedor padre, no en el viewport.\\
2023 & \texttt{:has()} selector & Selector «padre» tras décadas de espera.\\
2024 & Cascade Layers & Capas declarativas para resolver especificidad.\\
2025 & CSS Nesting nativo & Anidamiento estilo Sass sin preprocesador.\\
\bottomrule
\end{tabularx}
\end{table}

\section{Fundamentos de CSS}

Antes de discutir layouts modernos, conviene revisar los tres
pilares de CSS: cómo se aplica el estilo, qué selectores existen y
cómo funciona la cascada y la especificidad. Estas bases son la
diferencia entre escribir CSS que mantiene el código bajo control y
escribir CSS que se vuelve ingobernable.


\subsection{Las tres formas de aplicar estilos}

Existen tres mecanismos para aplicar estilo a elementos HTML, en
orden de preferencia profesional. La enumeración siguiente los
expone con su justificación.

\textbf{Ejemplo corto comentado (las tres formas) (Listado~\ref{lst:3.1}):}
\begin{lstlisting}[language=HTML5,numbers=none,caption={Ejemplo de HTML},label=lst:3.1]
<!-- 1. EXTERNO: archivo .css separado. Es la forma profesional       -->
<!--    porque permite cache del navegador y separa contenido/forma. -->
<link rel="stylesheet" href="estilos.css">

<!-- 2. INTERNO: <style> dentro de <head>. Util para una sola pagina  -->
<!--    o para CSS critico que debe cargar sin peticion adicional.   -->
<style>
  p { color: navy; }   /* p = selector de tipo (todos los <p>) */
</style>

<!-- 3. INLINE: atributo style dentro del tag. Evitar en produccion   -->
<!--    porque pulveriza la separacion de responsabilidades.        -->
<p style="color: red">Texto rojo (inline, anti-patron).</p>
\end{lstlisting}


\begin{enumerate}
\item \textbf{Inline} (\texttt{style=\char34{}...\char34{}}): evitar en producción.
\item \textbf{Interno} (\texttt{<style>} en \texttt{<head>}): útil
para páginas pequeñas o casos especiales.
\item \textbf{Externo} (\texttt{<link rel=\char34{}stylesheet\char34{}>}): el modelo
profesional. Permite caché del navegador y separación de
responsabilidades.
\end{enumerate}

\subsection{Selectores principales}

Un \emph{selector} CSS es la expresión que decide a qué
elementos del DOM se les aplica un conjunto de propiedades. CSS
ofrece más de una docena de tipos de selectores con distinta
\emph{especificidad}. La tabla siguiente recoge los más usados.

\textbf{Ejemplo corto comentado (selectores principales) (Listado~\ref{lst:3.2}):}
\begin{lstlisting}[language=HTML5,numbers=none,caption={Ejemplo de HTML},label=lst:3.2]
/* Selector de TIPO: todos los parrafos                  */
p          { line-height: 1.6; }
/* Selector de CLASE: elementos con class="destacado"    */
.destacado { background: yellow; }
/* Selector de ID: el unico elemento con id="logo"       */
#logo      { width: 120px; }
/* Selector DESCENDIENTE: <a> dentro de <nav>            */
nav a      { text-decoration: none; }
/* Selector de ATRIBUTO: inputs de tipo email            */
input[type="email"] { border: 1px solid blue; }
/* PSEUDO-CLASE: estado del elemento                     */
a:hover    { color: red; }   /* Cuando el cursor pasa encima */
\end{lstlisting}


\begin{table}[htbp]
\centering
\caption{Selectores CSS más usados.}
\footnotesize
\begin{tabularx}{\textwidth}{lXl}
\toprule
\textbf{Selector} & \textbf{Aplica a} & \textbf{Especificidad}\\
\midrule
\texttt{*} & Todos los elementos & 0\\
\texttt{p} & Por tag & 1\\
\texttt{.clase} & Por clase & 10\\
\texttt{\#id} & Por ID & 100\\
\texttt{a:hover} & Pseudo-clase & 11\\
\texttt{p::first-letter} & Pseudo-elemento & 2\\
\texttt{ul > li} & Hijo directo & 2\\
\texttt{[type=\char34{}email\char34{}]} & Por atributo & 11\\
\texttt{:has(img)} & Padre con descendiente (2023) & varía\\
\bottomrule
\end{tabularx}
\end{table}

\subsection{Pseudo-clases CSS para feedback de validación}

Las pseudo-clases relacionadas con formularios permiten dar feedback
visual al usuario sin JavaScript:

\begin{itemize}
\item \texttt{:valid} --- el campo cumple la validación.
\item \texttt{:invalid} --- el campo no cumple.
\item \texttt{:user-valid} --- el usuario ya interactuó y es válido
(no marca el campo vacío al cargar).
\item \texttt{:user-invalid} --- el usuario interactuó y es inválido.
\item \texttt{:required} --- el campo es obligatorio.
\item \texttt{:placeholder-shown} --- mientras el placeholder es
visible.
\item \texttt{:focus-visible} --- foco visible (solo navegación con
teclado).
\end{itemize}

\textbf{Ejemplo aplicable a formularios HTML5} (los formularios fueron
vistos en la semana 2) (Listado~\ref{lst:3.3}):
\begin{lstlisting}[style=css,numbers=none,caption={Ejemplo de CSS},label=lst:3.3]
input:user-invalid { border-color: red; }
input:user-valid   { border-color: green; }
input:focus-visible { outline: 3px solid #0066CC; }
\end{lstlisting}

\section{Sistemas modernos de layout}

Históricamente, maquetar páginas con CSS fue una pesadilla:
floats, posicionamientos hack, márgenes negativos. Desde 2017,
Flexbox y CSS Grid resuelven la mayoría de los layouts de forma
limpia y declarativa. Las dos subsecciones siguientes los explican
con ejemplos visuales.


\subsection{Flexbox: layout unidimensional}

Flexbox organiza elementos en una dimensión (fila o columna) y resuelve
problemas históricos de CSS como centrado vertical, distribución de
espacio y orden visual independiente del HTML \cite{robbins2018learning}.

\begin{cajaejemplo}{Ejemplo 3.1 (RESUELTO) --- Navbar responsivo con Flexbox}

\textbf{Objetivo:} crear una barra de navegación que se adapte a
pantalla grande (horizontal) y pequeña (vertical).

\textbf{IDE:} IntelliJ IDEA Ultimate.

\textbf{Paso 1:} Crear archivo \texttt{navbar.html}. V\'ease el Listado~\ref{lst:3.4}.

\begin{lstlisting}[style=html,caption={navbar.html},label=lst:3.4]
<!DOCTYPE html>
<html lang="es-EC">
<head>
  <meta charset="UTF-8">
  <meta name="viewport" content="width=device-width, initial-scale=1.0">
  <title>Navbar Flexbox</title>
  <link rel="stylesheet" href="navbar.css">
</head>
<body>
  <nav class="navbar" aria-label="Principal">
    <div class="logo">UTEQ</div>
    <ul class="menu">
      <li><a href="#">Inicio</a></li>
      <li><a href="#">Carreras</a></li>
      <li><a href="#">Investigacion</a></li>
      <li><a href="#">Contacto</a></li>
    </ul>
  </nav>
  <main><p>Contenido de la pagina.</p></main>
</body>
</html>
\end{lstlisting}

\textbf{Paso 2:} Crear \texttt{navbar.css}. V\'ease el Listado~\ref{lst:3.5}.

\begin{lstlisting}[style=css,caption={navbar.css},label=lst:3.5]
* { margin: 0; padding: 0; box-sizing: border-box; }
body { font-family: system-ui, sans-serif; }

.navbar {
  /* activa Flexbox (layout unidimensional) */
  display: flex;                                            /* activa Flexbox (layout unidimensional) */
  /* alineacion principal en el eje horizontal */
  justify-content: space-between;                           /* alineacion principal en el eje horizontal */
  /* alineacion en el eje secundario (vertical) */
  align-items: center;                                      /* alineacion en el eje secundario (vertical) */
  background: #006633;
  padding: 1rem 2rem;
  /* permite saltos de linea en flex */
  flex-wrap: wrap;                                          /* permite saltos de linea en flex */
}

.logo {
  color: white;
  font-size: 1.5rem;
  font-weight: bold;
}

.menu {
  /* activa Flexbox (layout unidimensional) */
  display: flex;                                            /* activa Flexbox (layout unidimensional) */
  list-style: none;
  /* espacio entre elementos hijos */
  gap: 1.5rem;                                              /* espacio entre elementos hijos */
}

.menu a {
  color: white;
  text-decoration: none;
  padding: 0.5rem 0;
}

.menu a:hover { color: #CC9900; }

/* media query (responsive) */
@media (max-width: 600px) {                                 /* media query (responsive) */
  /* direccion principal del flex container */
  /* direccion principal del flex container */
  .menu { flex-direction: column; width: 100%; gap: 0.5rem; }
}

main { padding: 2rem; }
\end{lstlisting}

\textbf{Resultado esperado (renderizado en pantalla ancha):}

\fbox{\begin{minipage}{0.95\textwidth}
\vspace{0.4em}
\colorbox{uteqverde}{\parbox{0.94\textwidth}{
\textcolor{white}{\bfseries UTEQ}\hfill
\textcolor{white}{Inicio \quad Carreras \quad Investigación \quad Contacto}
}}\\[0.5em]
Contenido de la página.
\vspace{0.4em}
\end{minipage}}

\textbf{Resultado esperado (pantalla móvil $<$ 600px):}

\fbox{\begin{minipage}{0.45\textwidth}
\vspace{0.4em}
\colorbox{uteqverde}{\parbox{0.94\linewidth}{
\textcolor{white}{\bfseries UTEQ}\\
\textcolor{white}{Inicio}\\
\textcolor{white}{Carreras}\\
\textcolor{white}{Investigación}\\
\textcolor{white}{Contacto}
}}\\[0.3em]
Contenido de la página.
\vspace{0.4em}
\end{minipage}}
\end{cajaejemplo}

\subsection{CSS Grid: layout bidimensional}

Grid permite layouts en dos dimensiones simultáneas. Es ideal para
maquetar páginas completas, dashboards y galerías
\cite{macdonald2018css}.

\begin{cajaejemplo}{Ejemplo 3.2 (RESUELTO) --- Galería de productos con Grid}


\textbf{Enunciado.} Construir una galería de productos responsive con CSS Grid que se adapte automáticamente al ancho del viewport sin media queries explícitos.

\begin{lstlisting}[style=css,caption={galeria.css (fragmento)},label=lst:3.6]
.galeria {
  /* activa CSS Grid (layout bidimensional) */
  display: grid;                                            /* activa CSS Grid (layout bidimensional) */
  /* define columnas del grid */
  /* define columnas del grid */
  grid-template-columns: repeat(auto-fill, minmax(220px, 1fr));
  /* espacio entre elementos hijos */
  gap: 1rem;                                                /* espacio entre elementos hijos */
}
.producto {
  border: 1px solid #ddd;
  padding: 1rem;
  border-radius: 8px;
}
\end{lstlisting}

\textbf{Resultado esperado:} una cuadrícula que ajusta automáticamente
el número de columnas según el ancho del viewport (4 columnas en
desktop, 2 en tablet, 1 en móvil).

\fbox{\begin{minipage}{0.95\textwidth}
\vspace{0.4em}
\fbox{\parbox{0.22\linewidth}{\small Producto 1\\\$2.50}}
\fbox{\parbox{0.22\linewidth}{\small Producto 2\\\$3.20}}
\fbox{\parbox{0.22\linewidth}{\small Producto 3\\\$4.50}}
\fbox{\parbox{0.22\linewidth}{\small Producto 4\\\$1.80}}
\vspace{0.4em}
\end{minipage}}
\end{cajaejemplo}

\section{Sass: el preprocesador profesional}

Sass añade variables, mixins, anidamiento y módulos a CSS. Aunque
CSS moderno ya tiene variables nativas, Sass sigue siendo el estándar
en grandes proyectos por su capacidad de \emph{computación}. V\'ease el Listado~\ref{lst:3.7}.

\begin{lstlisting}[style=sass,caption={ejemplo.scss},numbers=none,label=lst:3.7]
$verde: #006633;
$dorado: #CC9900;

@mixin boton($bg) {                                         // definicion de mixin (codigo reutilizable)
  background: $bg;
  color: white;
  padding: 0.6rem 1.2rem;
  border: 0;
  border-radius: 4px;
  &:hover { filter: brightness(1.1); }                      // referencia al selector padre
}

.btn-primario { @include boton($verde); }
.btn-acento   { @include boton($dorado); }
\end{lstlisting}

\section{Ejercicios}

Los siguientes ejercicios consolidan Flexbox, Grid y Sass.
El estudiante debe entregar el código fuente más una captura del
resultado renderizado en el navegador.


\begin{cajaejercicio}{Ejercicio 3.1 (PROPUESTO) --- Card responsive con identidad UTEQ}

\textbf{Objetivo:} construir un componente \emph{card} reutilizable
con HTML semántico + CSS Flexbox + Grid.

\textbf{Requisitos:}
\begin{itemize}
\item Tres cards alineadas con CSS Grid.
\item Cada card tiene imagen (placeholder), título, descripción y
botón.
\item Colores institucionales (uteqverde \texttt{\#006633}, uteqdorado
\texttt{\#CC9900}).
\item Responsive: 3 columnas en desktop, 2 en tablet, 1 en móvil.
\item Estado \texttt{:hover} con elevación (box-shadow).
\item Score Lighthouse Accessibility $\geq$ 95.
\end{itemize}

\textbf{Entrega:} archivos \texttt{cards.html} + \texttt{cards.css}
con capturas en escritorio, tablet y móvil (3 imágenes).
\end{cajaejercicio}

\begin{cajaejercicio}{Ejercicio 3.2 (PROPUESTO) --- Dashboard básico con CSS Grid}


\textbf{Enunciado.} Construir un dashboard básico con CSS Grid de 3 columnas que se reorganice a 1 columna en móviles.

Crear un dashboard con cuatro tarjetas de estadísticas alineadas
con Grid, una gráfica simulada con \texttt{<div>}+CSS y un menú
lateral fijo. Aplicar mobile-first y validar accesibilidad.
\end{cajaejercicio}

\section{Buenas prácticas CSS profesional}

Conocer CSS no basta: la diferencia entre un proyecto
mantenible y uno caótico está en aplicar disciplina al escribir
estilos. La caja siguiente recopila doce reglas profesionales que
el estudiante debe seguir desde el inicio.


\begin{cajabuenas}{Doce reglas CSS profesionales 2026}
\begin{enumerate}
\item \texttt{box-sizing: border-box} global.
\item Mobile-first siempre.
\item Variables CSS para colores y tipografía (\texttt{-{}-primario}).
\item BEM o similar para nombres de clase
(\texttt{.card\_\_titulo}).
\item Unidades relativas (\texttt{rem}, \texttt{em}, \texttt{\%})
sobre \texttt{px}.
\item Sin \texttt{!important} excepto utilidades muy específicas.
\item No usar IDs como selectores (alta especificidad).
\item Flexbox para 1D, Grid para 2D.
\item Animaciones con \texttt{transform} y \texttt{opacity}
(no afectan layout).
\item Respetar \texttt{prefers-reduced-motion} para usuarios sensibles
a animaciones.
\item Lighthouse Performance + Accessibility $\geq$ 90.
\item Minificación y autoprefixer en pipeline de build.
\end{enumerate}
\end{cajabuenas}

%%% ==== SEMANA 04 v3 ====
% =====================================================================
%   SEMANA 4 - VERSIÓN V3 MEJORADA
%   JavaScript moderno y programación asíncrona
% =====================================================================

\chapter{JavaScript moderno y programación asíncrona}
\label{cap:semana4}

\epigraph{«Cualquier aplicación que pueda escribirse en JavaScript,
eventualmente se escribirá en JavaScript.»}{--- \textit{Jeff
Atwood, 2007 (Ley de Atwood)}}

\section*{Resultado de aprendizaje}
Al concluir la semana 4, el estudiante construye lógica interactiva
del lado del cliente con JavaScript moderno (ES2024+), domina el DOM,
los eventos, las APIs asíncronas (fetch, Promises, async/await), y
aplica las buenas prácticas de la industria
\cite{flanagan2020javascript,crockford2008javascript}.

\section{Historia de JavaScript}

JavaScript nació en 1995 como un lenguaje juguete creado en
diez días por Brendan Eich. Hoy es la lingua franca del navegador,
ejecuta servidores enteros (Node.js), corre en satélites y dispositivos
IoT. Conocer su historia ayuda a entender por qué tiene
particularidades sintácticas y por qué evoluciona tan rápido.


\begin{cajahistoria}{De 10 días de diseño a la lingua franca de la web}
En mayo de 1995, Netscape contrató a \textbf{Brendan Eich} para
crear un lenguaje de scripting que se ejecutara en el navegador.
Le dieron \emph{diez días}. El resultado, originalmente llamado
\textbf{Mocha}, luego \textbf{LiveScript}, y finalmente
\textbf{JavaScript} (1996, por razones de marketing aprovechando la
popularidad de Java), tenía rasgos memorables: tipado dinámico,
\emph{first-class functions}, prototipos en lugar de clases,
\texttt{undefined} y \texttt{null} como dos formas distintas de
«nada» \cite{crockford2008javascript}.

Microsoft respondió con JScript en Internet Explorer 3 (1996),
provocando una guerra de incompatibilidades que se cerró con la
estandarización ECMA en 1997 (\textbf{ECMAScript}). Durante una
década, JavaScript fue considerado un lenguaje de juguete. El cambio
llegó con \textbf{Google Chrome} y su motor V8 (2008), que ejecutaba
JavaScript hasta 100x más rápido. Ese mismo año Ryan Dahl creó
\textbf{Node.js} (2009), llevando JS al servidor. \textbf{ECMAScript
2015 (ES6)} formalizó clases, módulos, \texttt{let}, \texttt{const},
arrow functions, Promises. ECMAScript 2024 ya es estándar y 2026
\cite{ecma2024} incorpora características como \texttt{Array.fromAsync},
\texttt{Promise.try}, y \emph{records} y \emph{tuples} inmutables.
\end{cajahistoria}

\subsection{Tabla evolutiva de ECMAScript}

ECMAScript es el estándar formal que define el lenguaje
JavaScript. Desde 2015, evoluciona anualmente añadiendo
funcionalidades. La tabla siguiente lista las versiones más
significativas y sus aportes.


\begin{table}[htbp]
\centering
\caption{Versiones recientes de ECMAScript.}
\footnotesize
\begin{tabularx}{\textwidth}{lXX}
\toprule
\textbf{Versión} & \textbf{Año} & \textbf{Aportes principales}\\
\midrule
ES5  & 2009 & \texttt{strict mode}, JSON nativo, getters/setters.\\
ES6 (ES2015) & 2015 & Clases, módulos, \texttt{let}/\texttt{const}, arrow, Promises, destructuring.\\
ES2017 & 2017 & \texttt{async/await}, \texttt{Object.entries}.\\
ES2019 & 2019 & \texttt{Array.flat}, \texttt{Object.fromEntries}.\\
ES2020 & 2020 & Optional chaining (\texttt{?.}), Nullish coalescing (\texttt{??}), BigInt.\\
ES2022 & 2022 & Top-level await, campos privados (\texttt{\#}).\\
ES2024 & 2024 & \texttt{Object.groupBy}, ArrayBuffer transferible.\\
ES2026 \cite{ecma2024} & 2026 & \texttt{Array.fromAsync}, \texttt{Promise.try}, records y tuples.\\
\bottomrule
\end{tabularx}
\end{table}

\section{Fundamentos modernos}

Antes de ver el DOM y el código asíncrono, el estudiante debe
dominar los fundamentos modernos: declaración de variables con
\texttt{let}/\texttt{const}, funciones flecha, \emph{destructuring}
y \emph{spread}. Estos elementos aparecen en todo código JavaScript
profesional de 2026.


\subsection{Variables: por qué \texttt{let}/\texttt{const} y nunca \texttt{var}}

La palabra clave \texttt{var} tenía dos defectos graves:
\emph{hoisting} (las variables se declaraban silenciosamente al
inicio del scope) y \emph{function scope} en lugar de \emph{block
scope}. ES6 introdujo \texttt{let} y \texttt{const} para resolverlos. V\'ease el Listado~\ref{lst:4.1}.


\begin{lstlisting}[style=js,numbers=none,caption={Ejemplo de JavaScript},label=lst:4.1]
const PI = 3.14159;          // inmutable, scope de bloque
let contador = 0;            // mutable, scope de bloque
var legacy = "no usar";      // hoisted, scope de funcion - EVITAR
\end{lstlisting}

\subsection{Funciones flecha y funciones regulares}

Las funciones flecha (\texttt{=>}) introducidas en ES6 son
sintaxis más concisa para funciones, pero además tienen una
diferencia semántica importante: no enlazan su propio
\texttt{this}. Esto las hace ideales como callbacks pero
inadecuadas como métodos de objetos en algunos contextos. V\'ease el Listado~\ref{lst:4.2}.


\begin{lstlisting}[style=js,numbers=none,caption={Ejemplo de JavaScript},label=lst:4.2]
// Funcion regular
function sumar(a, b) { return a + b; }                      // declaracion de funcion

// Arrow function equivalente
const sumar = (a, b) => a + b;                              // constante (no reasignable)

// Arrow sin parametros
const ahora = () => new Date();                             // constante (no reasignable)

// Arrow con cuerpo
const dividir = (a, b) => {                                 // constante (no reasignable)
  if (b === 0) throw new Error("Division por cero");        // condicional
  return a / b;                                             // devuelve valor de la funcion
};
\end{lstlisting}

\subsection{Destructuring y spread}

\emph{Destructuring} permite extraer valores de objetos y
arrays con una sintaxis declarativa concisa; \emph{spread}
(\texttt{...}) permite expandir colecciones en distintos
contextos. Ambas técnicas son omnipresentes en código moderno. V\'ease el Listado~\ref{lst:4.3}.


\begin{lstlisting}[style=js,numbers=none,caption={Ejemplo de JavaScript},label=lst:4.3]
const usuario = { nombre: "Ana", edad: 22, rol: "admin" };  // constante (no reasignable)
const { nombre, rol } = usuario;          // destructuring objeto

const [primero, segundo] = [10, 20, 30];  // destructuring array

const copia = { ...usuario, edad: 23 };   // spread + override
\end{lstlisting}

\section{Manipulación del DOM}

El DOM (\emph{Document Object Model}) es la representación
en memoria de la página HTML, expuesta como un árbol de objetos que
JavaScript puede leer y modificar. Los siguientes ejemplos muestran
las operaciones más frecuentes.


\begin{cajaejemplo}{Ejemplo 4.1 (RESUELTO) --- Lista de tareas dinámica}

\textbf{Objetivo:} crear una pequeña aplicación de lista de tareas
con agregar y eliminar.

\textbf{Paso 1:} En IntelliJ, crear \texttt{tareas.html}. V\'ease el Listado~\ref{lst:4.4}.

\begin{lstlisting}[style=html,caption={tareas.html},label=lst:4.4]
<!DOCTYPE html>
<html lang="es-EC">
<head>
  <meta charset="UTF-8">
  <title>Lista de tareas</title>
  <style>
    body { font-family: system-ui; max-width: 500px;
           margin: 2rem auto; padding: 0 1rem; }
    h1 { color: #006633; }
    input, button { padding: 0.5rem; font-size: 1rem; }
    button { background: #006633; color: white;
             border: 0; cursor: pointer; }
    ul { list-style: none; padding: 0; }
    li { display: flex; justify-content: space-between;
         padding: 0.5rem; border-bottom: 1px solid #eee; }
  </style>
</head>
<body>
  <h1>Mis tareas</h1>
  <form id="form-tarea">
    <input type="text" id="texto-tarea" required
           placeholder="Nueva tarea" minlength="3">
    <button type="submit">Agregar</button>
  </form>
  <ul id="lista"></ul>
  <script src="tareas.js"></script>
</body>
</html>
\end{lstlisting}

\textbf{Paso 2:} Crear \texttt{tareas.js}. V\'ease el Listado~\ref{lst:4.5}.

\begin{lstlisting}[style=js,caption={tareas.js},label=lst:4.5]
"use strict";

const form  = document.getElementById("form-tarea");        // constante (no reasignable)
const input = document.getElementById("texto-tarea");       // constante (no reasignable)
const lista = document.getElementById("lista");             // constante (no reasignable)

form.addEventListener("submit", (e) => {                    // registra manejador de evento
  e.preventDefault();
  const texto = input.value.trim();                         // constante (no reasignable)
  if (texto.length < 3) return;                             // condicional

  const li = document.createElement("li");                  // constante (no reasignable)
  li.textContent = texto;                                   // cambia el texto (seguro vs XSS)

  const btn = document.createElement("button");             // constante (no reasignable)
  btn.textContent = "X";                                    // cambia el texto (seguro vs XSS)
  btn.addEventListener("click", () => li.remove());         // registra manejador de evento

  li.appendChild(btn);
  lista.appendChild(li);
  input.value = "";
  input.focus();
});
\end{lstlisting}

\textbf{Resultado esperado:}

\fbox{\begin{minipage}{0.85\textwidth}
\vspace{0.4em}
{\Large\bfseries\color{uteqverde} Mis tareas}\\[0.5em]
\fbox{\parbox{0.55\linewidth}{Nueva tarea}}
\colorbox{uteqverde}{\textcolor{white}{\bfseries\ Agregar\ }}\\[0.5em]
Estudiar JavaScript \hfill \fbox{X}\\
\hrule
Hacer ejercicios HTML \hfill \fbox{X}\\
\hrule
Practicar Flexbox \hfill \fbox{X}
\vspace{0.4em}
\end{minipage}}
\end{cajaejemplo}

\section{Programación asíncrona: fetch, Promises, async/await}

La programación asíncrona es ineludible en el navegador:
peticiones HTTP, lectura de archivos, temporizadores y eventos no
deben bloquear la interfaz. JavaScript provee tres mecanismos
encadenados históricamente: callbacks, Promises y \texttt{async/await}.
El ejemplo siguiente muestra el patrón moderno.


\begin{cajaejemplo}{Ejemplo 4.2 (RESUELTO) --- Consumir API pública con fetch}

\textbf{Enunciado.} Consumir una API REST pública con \texttt{fetch} usando \texttt{async/await} para mostrar los datos en la página, manejando correctamente los errores de red.

\begin{lstlisting}[style=js,caption={api.js},label=lst:4.6]
"use strict";

async function obtenerUsuarios() {                          // funcion asincrona (devuelve Promise)
  try {
    const res = await fetch("https://jsonplaceholder.typicode.com/users");
    if (!res.ok) throw new Error(`HTTP ${res.status}`);     // condicional
    const usuarios = await res.json();                      // constante (no reasignable)
    return usuarios;                                        // devuelve valor de la funcion
  } catch (e) {
    console.error("Error:", e.message);
    return [];                                              // devuelve valor de la funcion
  }
}

obtenerUsuarios().then(lista => {                           // callback ejecutado cuando la Promise se resuelve
  lista.forEach(u => console.log(u.name, "-", u.email));    // imprime en la consola del navegador
});
\end{lstlisting}

\textbf{Cómo probarlo:} abrir \texttt{tareas.html} (o cualquier
archivo con un \texttt{<script>}), abrir la consola del navegador
(F12 $\rightarrow$ Console), pegar el código. Aparecen 10 usuarios
con su nombre y correo en la consola.
\end{cajaejemplo}

\section{Ejercicios}

Los siguientes ejercicios consolidan DOM, fetch y patrones
asíncronos. El estudiante debe entregar código y capturas del
resultado en el navegador.


\begin{cajaejercicio}{Ejercicio 4.1 (PROPUESTO) --- Calculadora simple}


\textbf{Enunciado.} Implementar una calculadora simple en JavaScript puro con operaciones aritméticas básicas, capturando eventos de los botones.

Crear una calculadora HTML+CSS+JS con las cuatro operaciones básicas.
Validar división por cero. Mostrar el resultado en pantalla. Aplicar
buenas prácticas (\texttt{strict mode}, \texttt{const}/\texttt{let},
manejo de errores).
\end{cajaejercicio}

\begin{cajaejercicio}{Ejercicio 4.2 (PROPUESTO) --- Buscador de pokémon}


\textbf{Enunciado.} Construir un buscador de Pokémon que consuma la PokeAPI pública y muestre imagen, tipos y estadísticas del Pokémon buscado.

Consumir la API \url{https://pokeapi.co/api/v2/pokemon/\{nombre\}} con
\texttt{fetch} y \texttt{async/await}. Mostrar al usuario el nombre,
la imagen y los tipos del pokémon buscado. Manejar errores 404.
Aplicar la pseudo-clase \texttt{:focus-visible} aprendida en la
semana 3.
\end{cajaejercicio}

\section{Buenas prácticas JavaScript}

Escribir JavaScript profesional es más que conocer la
sintaxis: requiere disciplina con tipos, manejo de errores,
inmutabilidad y modularidad. La caja siguiente reúne doce reglas
profesionales que distinguen al desarrollador junior del senior.


\begin{cajabuenas}{Doce reglas profesionales 2026}
\begin{enumerate}
\item \texttt{\char34{}use strict\char34{}} en todo módulo.
\item \texttt{const} por defecto; \texttt{let} solo cuando muta;
\texttt{var} nunca.
\item \texttt{===} sobre \texttt{==} siempre (comparación estricta).
\item Manejo de errores con \texttt{try/catch} en código asíncrono.
\item Validar entrada de funciones públicas.
\item Funciones pequeñas, una responsabilidad.
\item Sin variables globales: usar módulos ES (\texttt{type=\char34{}module\char34{}}).
\item Linter (ESLint) con configuración del equipo.
\item Tests con Jest, Vitest o Playwright.
\item No mezclar \texttt{Promise.then()} y \texttt{await}: elegir uno.
\item Optional chaining (\texttt{?.}) para acceso seguro a objetos.
\item Tipado con JSDoc o TypeScript en proyectos serios.
\end{enumerate}
\end{cajabuenas}

%%% ==== SEMANA 05 v3 ====
% =====================================================================
%   SEMANA 5 - VERSIÓN V3 MEJORADA
%   Programación del lado del servidor: PERL y PHP
% =====================================================================

\chapter{Introducción a la programación del lado del servidor: PERL y PHP}
\label{cap:semana5}

\epigraph{«PHP creció orgánicamente para resolver problemas reales.
PERL fue diseñado para ser práctico. Esa pragmática es lo que hizo a
ambos sobrevivientes en una industria de modas.»}{--- \textit{Adaptado
de las palabras de Rasmus Lerdorf y Larry Wall}}

\section*{Resultado de aprendizaje}
Al concluir la semana 5, el estudiante desarrolla aplicaciones web
dinámicas del lado del servidor utilizando PERL/CGI y PHP 8.3,
implementa procesamiento de formularios y comprende el modelo
\emph{request-response} de HTTP \cite{nixon2021learning}.

\section{Programación del lado del servidor: conceptos}

El servidor web es el segundo gran protagonista del modelo
cliente-servidor. Comprender el ciclo \emph{request-response} y la
arquitectura del lado servidor es prerrequisito para programar en
cualquier lenguaje de backend.


\subsection{Modelo request-response}

El servidor web recibe peticiones HTTP del cliente y devuelve
respuestas. El \emph{servidor de aplicaciones} ejecuta código
(PHP, Java, .NET, Python) que genera dinámicamente la respuesta
según la petición, datos de sesión, base de datos, etc.

\subsection{Arquitectura: navegador $\to$ Apache/Nginx $\to$ intérprete}

Cuando el navegador hace una petición, esta atraviesa varias
capas hasta llegar al intérprete del lenguaje y volver con la
respuesta. La figura siguiente lo ilustra para el caso PHP.


\begin{figure}[htbp]
\centering
\begin{tikzpicture}[every node/.style={font=\scriptsize,align=center},
  caja/.style={rectangle,draw=uteqverde,thick,minimum width=2cm,
               minimum height=0.8cm,fill=uteqverde!10,rounded corners=2pt}]
  \node[caja] (a) {Navegador};
  \node[caja,right=1cm of a] (b) {Apache /\\Nginx};
  \node[caja,right=1cm of b] (c) {Interprete\\PHP/PERL};
  \node[cylinder,draw=uteqverde,fill=uteqverde!15,thick,
        shape border rotate=90,aspect=0.3,minimum width=1.2cm,
        minimum height=1.5cm,right=1cm of c] (d) {MySQL};
  \draw[->,>=stealth,thick] (a) -- node[above]{HTTP} (b);
  \draw[->,>=stealth,thick] (b) -- node[above]{exec} (c);
  \draw[->,>=stealth,thick] (c) -- node[above]{SQL} (d);
\end{tikzpicture}
\caption{Flujo típico de una petición a un servidor PHP/PERL: el
navegador envía HTTP a Apache, éste delega al intérprete PHP/PERL,
que consulta la BD MySQL y devuelve HTML.}
\end{figure}

\section{PERL para CGI: el pionero}

PERL fue el primer lenguaje en dominar el desarrollo web
servidor durante los años 90, gracias al CGI (\emph{Common Gateway
Interface}). Aunque hoy ha sido desplazado por PHP, Java y .NET,
estudiarlo brevemente da perspectiva histórica.


\begin{cajahistoria}{Larry Wall y la web de los 90}
\textbf{Larry Wall}, lingüista y programador, creó PERL (Practical
Extraction and Report Language) en 1987. En 1993, cuando la web
empezaba a hacerse dinámica, PERL+CGI fue \emph{la} combinación que
permitió procesar formularios, enviar correos y generar páginas
dinámicas. Empresas como Amazon, Slashdot y muchos primeros sitios
de e-commerce nacieron en PERL.

PERL es notable por su lema \emph{«There's More Than One Way To Do
It»} (TIMTOWTDI): hay más de una forma de hacerlo. Esto lo hizo
expresivo pero también difícil de mantener. En 2026 PERL sigue
vigente en \emph{sysadmin}, integración de sistemas y análisis de
texto, aunque su rol en la web ha quedado relegado a sistemas
legados.
\end{cajahistoria}

\subsection{Hola mundo en PERL+CGI}

El siguiente ejemplo presenta el \emph{Hola mundo} en
PERL/CGI ---el código mínimo viable--- con su salida HTTP y HTML.


\begin{cajaejemplo}{Ejemplo 5.1 (RESUELTO) --- Hola mundo en PERL/CGI}


\textbf{Enunciado.} Construir el \emph{Hola Mundo} canónico en PERL/CGI: un script que reciba un parámetro \texttt{nombre} por GET y devuelva una página HTML personalizada.

\textbf{IDE:} IntelliJ IDEA Ultimate con plugin \emph{Perl}.

\textbf{Paso 1:} Verificar Perl. En terminal: \texttt{perl -v}. Si no
está, instalarlo con \texttt{sudo apt install perl} (Linux) o desde
\url{https://strawberryperl.com} (Windows).

\textbf{Paso 2:} Crear \texttt{C:\textbackslash{}xampp\textbackslash{}cgi-bin\textbackslash{}hola.pl}
con el siguiente código (Listado~\ref{lst:5.1}):

\begin{lstlisting}[style=perl,caption={hola.pl},label=lst:5.1]
#!/usr/bin/perl
use strict;                                                 # importa modulo Perl
use warnings;                                               # importa modulo Perl

# imprime al output (al cliente en CGI)
print "Content-Type: text/html; charset=utf-8\n\n";         # imprime al output (al cliente en CGI)

my $hora = (localtime)[2];                                  # declara variable con alcance lexico
# declara variable con alcance lexico
my $saludo = $hora < 12 ? "Buenos dias" :                   # declara variable con alcance lexico
             $hora < 19 ? "Buenas tardes" : "Buenas noches";

print <<HTML;                                               # imprime al output (al cliente en CGI)
<!DOCTYPE html>
<html lang="es-EC">
<head><meta charset="UTF-8"><title>PERL CGI</title></head>
<body>
  <h1>$saludo desde PERL</h1>
  <p>La hora del servidor es: $hora</p>
</body>
</html>
HTML
\end{lstlisting}

\textbf{Paso 3:} Hacer ejecutable
(\texttt{chmod +x hola.pl} en Linux) y configurar Apache CGI.
Acceder a \url{http://localhost/cgi-bin/hola.pl}.

\textbf{Resultado esperado:}

\fbox{\begin{minipage}{0.85\textwidth}
\vspace{0.4em}
{\Large\bfseries Buenas tardes desde PERL}\\[0.4em]
La hora del servidor es: 15
\vspace{0.4em}
\end{minipage}}
\end{cajaejemplo}

\section{PHP: el lenguaje dominante de la web tradicional}

PHP es el lenguaje de servidor más extendido en la web
pública (\(\approx\) 75\% de sitios). Comenzó como un \emph{hobby
project} de Rasmus Lerdorf en 1994 y hoy potencia Wikipedia,
Facebook (en sus raíces), WordPress y miles de aplicaciones
empresariales.


\begin{cajahistoria}{Rasmus Lerdorf y el accidente que dominó la web}
En 1994, \textbf{Rasmus Lerdorf} creó un conjunto de scripts CGI en
PERL para gestionar su currículum online. Los llamó
\emph{Personal Home Page Tools}, abreviado \textbf{PHP}. Lo que
empezó como código personal evolucionó hasta convertirse en
\textbf{PHP/FI 2.0}, luego en PHP 3 (1998, reescrito por Andi
Gutmans y Zeev Suraski), PHP 4 (2000, con el motor Zend), PHP 5
(2004, OOP madura), PHP 7 (2015, salto de rendimiento), y PHP 8.x
(2020--2026) con tipado estricto, atributos y JIT compilation.

PHP corre el 76\% de los sitios web cuyo lenguaje servidor es
identificable (W3Techs 2026): WordPress, Wikipedia, Facebook
(parcialmente, vía Hack), Slack, Mailchimp. Aunque no es el lenguaje
más \emph{cool}, su pragmatismo y vasto ecosistema lo mantienen
dominante.
\end{cajahistoria}

\subsection{Tabla evolutiva de PHP}

PHP ha sufrido tres reescrituras de motor mayores. Conocer
su evolución importa porque el código heredado de PHP 5 sigue vivo
en muchos sitios y requiere migración. La tabla siguiente resume
las versiones.


\begin{table}[htbp]
\centering
\caption{Versiones recientes de PHP.}
\footnotesize
\begin{tabularx}{\textwidth}{lll}
\toprule
\textbf{Versión} & \textbf{Año} & \textbf{Aporte clave}\\
\midrule
PHP 5.6  & 2014 & Última de la rama 5, fin de soporte 2018.\\
PHP 7.0  & 2015 & Doble de rendimiento, scalar type hints.\\
PHP 7.4  & 2019 & Properties typed, arrow functions, FFI.\\
PHP 8.0  & 2020 & JIT, named arguments, attributes, match.\\
PHP 8.1  & 2021 & Readonly, enums, fibers, never type.\\
PHP 8.2  & 2022 & Readonly classes, DNF types.\\
PHP 8.3  & 2023 & Typed class constants, randomizer.\\
PHP 8.4  & 2024 & Property hooks, asymmetric visibility.\\
\bottomrule
\end{tabularx}
\end{table}

\subsection{Sintaxis básica}

PHP es C-like en su sintaxis pero con variables siempre
prefijadas por \texttt{\$}, tipado débil y un sistema de funciones
gigantesco. El siguiente fragmento muestra los elementos
fundamentales que el estudiante usará en el resto de la semana. V\'ease el Listado~\ref{lst:5.2}.


\begin{lstlisting}[style=php,numbers=none,caption={Ejemplo de PHP},label=lst:5.2]
<?php
declare(strict_types=1);

$nombre = "Mundo";
echo "Hola $nombre!";                                       // imprime al output del servidor

function sumar(int $a, int $b): int {
    return $a + $b;                                         // retorna valor
}

$total = sumar(3, 5);  // 8

// Arrays
$frutas = ["manzana", "pera", "uva"];
foreach ($frutas as $f) {                                   // itera sobre colecion
    echo $f . "\n";                                         // imprime al output del servidor
}

// Asociativo
$user = ["nombre" => "Ana", "edad" => 22];
echo $user["nombre"];                                       // imprime al output del servidor
\end{lstlisting}

\subsection{Procesamiento de formularios}

El procesamiento de formularios HTTP fue el caso de uso
original de PHP: leer datos POST/GET, validar, persistir, responder.
El siguiente ejemplo resuelto cubre el ciclo completo con buenas
prácticas de seguridad.


\begin{cajaejemplo}{Ejemplo 5.2 (RESUELTO) --- Formulario de contacto en PHP}


\textbf{Enunciado.} Construir un formulario de contacto en PHP que valide los datos del lado servidor, escape la salida contra XSS y muestre un mensaje de confirmación.

\textbf{Paso 1:} Iniciar XAMPP, abrir Apache. Crear carpeta
\texttt{C:\textbackslash{}xampp\textbackslash{}htdocs\textbackslash{}contacto}.

\textbf{Paso 2:} En IntelliJ, abrir esa carpeta y crear
\texttt{contacto.php} (Listado~\ref{lst:5.3}):

\begin{lstlisting}[style=php,caption={contacto.php},label=lst:5.3]
<?php
declare(strict_types=1);

$mensaje_exito = "";
$errores = [];

if ($_SERVER["REQUEST_METHOD"] === "POST") {                // condicional
    $nombre = trim($_POST["nombre"] ?? "");
    $email  = trim($_POST["email"]  ?? "");
    $msg    = trim($_POST["mensaje"] ?? "");

    // condicional
    // condicional
    if (strlen($nombre) < 3) $errores[] = "Nombre muy corto";
    if (!filter_var($email, FILTER_VALIDATE_EMAIL))         // valida/sanea entrada con filtro
        $errores[] = "Correo invalido";
    // condicional
    // condicional
    if (strlen($msg) < 10) $errores[] = "Mensaje muy corto";

    if (empty($errores)) {                                  // condicional
        // Aqui normalmente se envia el correo / guarda en BD
        $mensaje_exito = "Gracias $nombre, su mensaje fue recibido.";
    }
}
?>
<!DOCTYPE html>
<html lang="es-EC">
<head>
  <meta charset="UTF-8">
  <title>Contacto UTEQ</title>
  <style>
    body { font-family: system-ui; max-width: 500px;
           margin: 2rem auto; padding: 0 1rem; }
    h1 { color: #006633; }
    label { display: block; margin: 0.5rem 0 0.2rem; }
    input, textarea { width: 100%; padding: 0.4rem;
                      font-size: 1rem; }
    button { margin-top: 1rem; padding: 0.6rem 1.2rem;
             background: #006633; color: white; border: 0;
             cursor: pointer; }
    .error { color: red; }
    .exito { color: green; padding: 1rem;
             background: #efffef; border-left: 3px solid green; }
  </style>
</head>
<body>
  <h1>Formulario de contacto</h1>

  <?php if ($mensaje_exito): ?>
    // escapa caracteres HTML (defensa anti-XSS)
    // escapa caracteres HTML (defensa anti-XSS)
    <div class="exito"><?= htmlspecialchars($mensaje_exito) ?></div>
  <?php endif; ?>

  <?php foreach ($errores as $e): ?>
    <p class="error"><?= htmlspecialchars($e) ?></p>        // escapa caracteres HTML (defensa anti-XSS)
  <?php endforeach; ?>

  <form method="POST">
    <label for="nombre">Nombre:</label>
    <input type="text" id="nombre" name="nombre" required minlength="3">

    <label for="email">Correo:</label>
    <input type="email" id="email" name="email" required>

    <label for="mensaje">Mensaje:</label>
    <textarea id="mensaje" name="mensaje" rows="5"
              required minlength="10"></textarea>           // incluye otro archivo PHP (fatal si falla)

    <button type="submit">Enviar</button>
  </form>
</body>
</html>
\end{lstlisting}

\textbf{Paso 3:} Acceder a
\url{http://localhost/contacto/contacto.php}. Probar:
\begin{itemize}
\item Enviar formulario vacío $\to$ aparecen mensajes de error.
\item Llenar correctamente $\to$ aparece mensaje de éxito.
\end{itemize}

\textbf{Resultado esperado (envío exitoso):}

\fbox{\begin{minipage}{0.85\textwidth}
\vspace{0.4em}
{\Large\bfseries\color{uteqverde} Formulario de contacto}\\[0.5em]
\colorbox{green!15}{\parbox{0.92\linewidth}{
\textcolor{green!50!black}{Gracias Juan, su mensaje fue recibido.}
}}\\[0.5em]
Nombre:\\ \fbox{\parbox{0.92\linewidth}{Juan Pérez}}\\
Correo:\\ \fbox{\parbox{0.92\linewidth}{juan@uteq.edu.ec}}\\
Mensaje:\\ \fbox{\parbox{0.92\linewidth}{Necesito información sobre los cursos de software.}}\\[0.5em]
\colorbox{uteqverde}{\textcolor{white}{\bfseries\ Enviar\ }}
\vspace{0.4em}
\end{minipage}}

\textbf{Defensas implementadas:}
\begin{itemize}
\item \texttt{filter\_var(\$email, FILTER\_VALIDATE\_EMAIL)} para
formato.
\item \texttt{htmlspecialchars} al renderizar (anti-XSS).
\item Validación server-side (longitud mínima).
\item HTML5 \texttt{required}, \texttt{minlength} (UX, no seguridad).
\end{itemize}
\end{cajaejemplo}

\section{Ejercicios}

Los siguientes ejercicios consolidan PHP y formularios.
El estudiante debe entregar código fuente y capturas funcionando.


\begin{cajaejercicio}{Ejercicio 5.1 (PROPUESTO) --- Encuesta en PHP}


\textbf{Enunciado.} Construir un sistema de encuesta en PHP que recoja respuestas por POST y persista los resultados en un archivo o tabla.

Crear una encuesta de 5 preguntas con \texttt{radio} y \texttt{select}
HTML. Al enviar, mostrar las respuestas en pantalla. Sumar puntaje
según opciones y mostrar un mensaje personalizado.
\end{cajaejercicio}

\begin{cajaejercicio}{Ejercicio 5.2 (PROPUESTO) --- Conversor de monedas}


\textbf{Enunciado.} Construir un conversor de monedas en PHP que consulte una API de tipos de cambio en tiempo real.

Formulario que reciba un monto en USD y un destino (EUR, GBP, BRL,
COP). Calcular el monto convertido con tasas hardcodeadas. Aplicar
validación servidor y mostrar el resultado.
\end{cajaejercicio}

\section{Buenas prácticas PHP}

Escribir PHP profesional requiere disciplina específica:
PHP es flexible pero permisivo, y esa permisividad genera vulnerabilidades
si no se siguen prácticas estrictas. La caja siguiente reúne doce
reglas profesionales.


\begin{cajabuenas}{Doce reglas profesionales 2026}
\begin{enumerate}
\item \texttt{declare(strict\_types=1)} al inicio de cada archivo.
\item Type hints en parámetros y retornos.
\item NUNCA \texttt{echo \$\_GET['x']} sin
\texttt{htmlspecialchars}: anti-XSS.
\item NUNCA concatenar SQL: usar PDO con prepared statements.
\item Usar Composer para dependencias.
\item Frameworks modernos: Laravel, Symfony, Slim.
\item PSR-12 para estilo de código.
\item Sesiones siempre con flags
\texttt{HttpOnly}/\texttt{Secure}/\texttt{SameSite}.
\item PHPStan o Psalm para análisis estático.
\item Tests con PHPUnit o Pest.
\item OPcache activado en producción.
\item Logs estructurados con Monolog.
\end{enumerate}
\end{cajabuenas}

%%% ==== SEMANA 06 v3 ====
% =====================================================================
%   SEMANA 6 - VERSIÓN V3 MEJORADA
%   Java para aplicaciones web: JSP, Servlets y Spring Boot
% =====================================================================

\chapter{Java para aplicaciones web: JSP, Servlets y Spring Boot}
\label{cap:semana6}

\epigraph{«Write once, run anywhere. Java cumplió la promesa en el
servidor: 30 años después, sigue siendo la columna vertebral de la
banca, la salud y el gobierno digital del planeta.»}{--- \textit{Adaptado
de la promesa original de Sun Microsystems}}

\section*{Resultado de aprendizaje}
Al concluir la semana 6, el estudiante construye aplicaciones web
con Java moderno: comprende la trayectoria desde Servlets/JSP hasta
Spring Boot 3 \cite{walls2022spring,deitel2017java}, e implementa un
proyecto funcional con persistencia JPA.

\section{Historia de Java en la web}

Java cumple 30 años en 2025 y ha sido protagonista de la web
empresarial desde los Applets de 1995 hasta Spring Boot 3 en 2026.
Conocer su trayectoria explica por qué hay tantos modelos de
programación web Java vivos al mismo tiempo (Servlets, JSP, JSF,
Struts, Spring, Quarkus).


\begin{cajahistoria}{De los Applets a Spring Boot --- 30 años de Java web}
James Gosling y su equipo en Sun Microsystems crearon Java en 1995.
La idea original era usarlo en dispositivos embebidos (el famoso
Star7), pero pronto se reveló como el lenguaje ideal para la naciente
World Wide Web. Sun introdujo en 1996 los \textbf{Applets} (Java
ejecutándose en el navegador) y los \textbf{Servlets} (Java en el
servidor) en 1997.

En 1999 Sun lanzó \textbf{J2EE} (Java 2 Enterprise Edition), el
estándar empresarial. Era poderoso pero pesado: configurar un proyecto
requería decenas de archivos XML. La frustración generó dos
respuestas: \textbf{Spring Framework} (Rod Johnson, 2003), que
simplificó la inyección de dependencias; y posteriormente
\textbf{Spring Boot} (2014), que eliminó la configuración manual con
\emph{convention over configuration}. Hoy Spring Boot es el framework
Java dominante en aplicaciones web empresariales.
\end{cajahistoria}

\subsection{Tabla evolutiva}

La evolución de Java para web ha pasado por al menos cinco
generaciones tecnológicas. La tabla siguiente sintetiza esa
trayectoria.


\begin{table}[htbp]
\centering
\caption{Versiones de Java y Spring relevantes en 2026.}
\footnotesize
\begin{tabularx}{\textwidth}{lXX}
\toprule
\textbf{Año} & \textbf{Versión} & \textbf{Aporte}\\
\midrule
2014 & Java 8       & Lambdas, streams, Optional.\\
2014 & Spring Boot 1.0 & Convención sobre configuración.\\
2018 & Java 11 LTS  & Versión LTS estable.\\
2021 & Java 17 LTS  & Sealed classes, records, pattern matching.\\
2022 & Spring Boot 3.0 & Jakarta EE 9 (paquetes \texttt{jakarta.*}).\\
2023 & Java 21 LTS  & Virtual threads, sequenced collections.\\
2024 & Spring Boot 3.4 & Compatible con Java 21+, observabilidad.\\
2025 & Java 25 LTS  & Patrones de cadena en switch.\\
\bottomrule
\end{tabularx}
\end{table}

\section{Servlets y JSP: el modelo histórico}

El modelo Servlet/JSP es el cimiento sobre el que se
construyen todos los frameworks Java web modernos, incluido Spring
Boot. Comprenderlo aunque solo sea brevemente revela cómo funciona
el motor por debajo.


\begin{cajaejemplo}{Ejemplo 6.1 (RESUELTO) --- Servlet hola mundo en Tomcat 10}


\textbf{Enunciado.} Crear el primer Servlet en Java desplegado sobre Tomcat 10 que responda con \texttt{Hola Mundo} a peticiones GET.

\begin{lstlisting}[style=java,caption={HolaServlet.java},label=lst:6.1]
package ec.edu.uteq.appweb;                                 // paquete del archivo

import jakarta.servlet.*;                                   // importacion de clase externa
import jakarta.servlet.annotation.WebServlet;               // importacion de clase externa
import jakarta.servlet.http.*;                              // importacion de clase externa
import java.io.IOException;                                 // importacion de clase externa

@WebServlet("/hola")
public class HolaServlet extends HttpServlet {              // declaracion de clase publica
    @Override
    // metodo protegido
    // metodo protegido
    protected void doGet(HttpServletRequest req, HttpServletResponse resp)
            throws IOException {
        String nombre = req.getParameter("nombre");
        // condicional
        // condicional
        if (nombre == null || nombre.isBlank()) nombre = "Mundo";

        resp.setContentType("text/html; charset=utf-8");
        resp.getWriter().write("""
            <!DOCTYPE html>
            <html lang="es-EC">
            <head><meta charset="UTF-8"><title>Servlet</title></head>
            <body><h1>Hola, %s desde Servlet 6.0!</h1></body>
            </html>
            """.formatted(nombre));
    }
}
\end{lstlisting}

Acceder a \url{http://localhost:8080/app/hola?nombre=Ana} y aparece
«Hola, Ana desde Servlet 6.0!»

\textbf{Resultado esperado:}

\fbox{\begin{minipage}{0.85\textwidth}
\vspace{0.4em}
\hspace{0.5em}{\Large\bfseries Hola, Ana desde Servlet 6.0!}
\vspace{0.4em}
\end{minipage}}
\end{cajaejemplo}

\section{Spring Boot: el estándar moderno}

Spring Boot es el framework Java dominante para aplicaciones
web modernas: opinionado, con \emph{auto-configuration}, dependencias
preempaquetadas y un \emph{starter} para cada necesidad. El
siguiente ejemplo construye desde cero la primera aplicación.


\begin{cajaejemplo}{Ejemplo 6.2 (RESUELTO) --- Mi primera app Spring Boot}


\textbf{Enunciado.} Construir la primera aplicación Spring Boot con un controlador REST mínimo que devuelva JSON. Demostrará la productividad del framework respecto al Servlet puro.

\textbf{Paso 1:} En IntelliJ IDEA Ultimate:
\texttt{File $\rightarrow$ New $\rightarrow$ Project $\rightarrow$
Spring Boot}. Configurar:
\begin{itemize}
\item Java 21, Maven, Spring Boot 3.4.x.
\item Group: \texttt{ec.edu.uteq}; Artifact: \texttt{primera}.
\item Dependencies: \emph{Spring Web}, \emph{Spring Boot DevTools},
\emph{Thymeleaf}.
\end{itemize}

\textbf{Paso 2:} Crear el controller (Listado~\ref{lst:6.2}):

\begin{lstlisting}[style=java,caption={HolaController.java},label=lst:6.2]
package ec.edu.uteq.primera.controller;                     // paquete del archivo

import org.springframework.stereotype.Controller;           // importacion de clase externa
import org.springframework.ui.Model;                        // importacion de clase externa
import org.springframework.web.bind.annotation.*;           // importacion de clase externa

@Controller                                                 // controlador MVC server-side (devuelve vistas)
public class HolaController {                               // declaracion de clase publica

    @GetMapping("/")                                        // endpoint HTTP GET
    public String inicio(                                   // metodo publico
            // extrae parametro de query string
            // extrae parametro de query string
            @RequestParam(defaultValue = "Mundo") String nombre,
            Model model) {
        model.addAttribute("nombre", nombre);
        return "inicio";   // resuelve a templates/inicio.html
    }
}
\end{lstlisting}

\textbf{Paso 3:} Crear la vista
\texttt{src/main/resources/templates/inicio.html} (Listado~\ref{lst:6.3}):

\begin{lstlisting}[style=html,caption={inicio.html (Thymeleaf)},label=lst:6.3]
<!DOCTYPE html>
<html xmlns:th="http://www.thymeleaf.org" lang="es-EC">
<head>
  <meta charset="UTF-8">
  <title>Hola Spring Boot</title>
  <style>
    body { font-family: system-ui; max-width: 600px;
           margin: 2rem auto; padding: 0 1rem; }
    h1 { color: #006633; }
  </style>
</head>
<body>
  <h1 th:text="'Hola, ' + ${nombre} + ' desde Spring Boot!'">
    Hola, Mundo desde Spring Boot!
  </h1>
  <p>Servidor Tomcat embebido en puerto 8080.</p>
</body>
</html>
\end{lstlisting}

\textbf{Paso 4:} Pulsar el icono verde de \emph{play}
(\texttt{Shift+F10}). El servidor inicia en pocos segundos.

\textbf{Paso 5:} Abrir \url{http://localhost:8080/?nombre=UTEQ}.

\textbf{Resultado esperado:}

\fbox{\begin{minipage}{0.85\textwidth}
\vspace{0.4em}
{\Large\bfseries\color{uteqverde} Hola, UTEQ desde Spring Boot!}\\[0.4em]
Servidor Tomcat embebido en puerto 8080.
\vspace{0.4em}
\end{minipage}}
\end{cajaejemplo}

\section{Spring Boot con JPA y BD H2}

Spring Data JPA es la pieza que conecta Spring Boot con bases
de datos relacionales. El siguiente ejemplo construye un CRUD
completo con la base H2 embebida para que el estudiante pueda
ejecutarlo sin instalar PostgreSQL.


\begin{cajaejemplo}{Ejemplo 6.3 (RESUELTO) --- CRUD de productos completo}


\textbf{Enunciado.} Construir un CRUD completo de productos con Spring Boot, Spring Data JPA y H2 embebida: los cuatro endpoints REST (\texttt{GET}, \texttt{POST}, \texttt{PUT}, \texttt{DELETE}) con persistencia real.

\textbf{Dependencias adicionales en \texttt{pom.xml}:}
\emph{Spring Data JPA}, \emph{H2 Database}, \emph{Lombok},
\emph{Validation}.

\textbf{Entidad (Listado~\ref{lst:6.4}):}
\begin{lstlisting}[style=java,caption={Producto.java},label=lst:6.4]
package ec.edu.uteq.primera.model;                          // paquete del archivo

import jakarta.persistence.*;                               // importacion de clase externa
import jakarta.validation.constraints.*;                    // importacion de clase externa
import lombok.*;                                            // importacion de clase externa
import java.math.BigDecimal;                                // importacion de clase externa

@Entity @Table(name = "productos")                          // entidad JPA (mapea a tabla)
@Getter @Setter @NoArgsConstructor @AllArgsConstructor
public class Producto {                                     // declaracion de clase publica
    @Id @GeneratedValue(strategy = GenerationType.IDENTITY) // campo clave primaria
    private Long id;                                        // campo privado

    @NotBlank @Size(min = 3, max = 100)                     // validacion: no nulo y no vacio
    private String nombre;                                  // campo privado

    @NotNull @DecimalMin("0.01")                            // validacion: no nulo
    private BigDecimal precio;                              // campo privado

    @NotNull @Min(0)                                        // validacion: no nulo
    private Integer stock;                                  // campo privado
}
\end{lstlisting}

\textbf{Repository (Listado~\ref{lst:6.5}):}
\begin{lstlisting}[style=java,numbers=none,caption={Ejemplo de Java},label=lst:6.5]
public interface ProductoRepository                         // declaracion de interfaz
    extends JpaRepository<Producto, Long> {}
\end{lstlisting}

\textbf{Controller con CRUD completo (Listado~\ref{lst:6.6}):}
\begin{lstlisting}[style=java,caption={ProductoController.java},label=lst:6.6]
@Controller                                                 // controlador MVC server-side (devuelve vistas)
@RequestMapping("/productos")                               // ruta base del controlador
@RequiredArgsConstructor
public class ProductoController {                           // declaracion de clase publica
    private final ProductoRepository repo;                  // campo inmutable (mejor practica)

    @GetMapping                                             // endpoint HTTP GET
    public String listar(Model model) {                     // metodo publico
        model.addAttribute("productos", repo.findAll());
        return "productos/lista";                           // devuelve el resultado al llamador
    }

    @GetMapping("/nuevo")                                   // endpoint HTTP GET
    public String nuevo(Model model) {                      // metodo publico
        model.addAttribute("producto", new Producto());
        return "productos/form";                            // devuelve el resultado al llamador
    }

    @PostMapping                                            // endpoint HTTP POST (crear)
    public String guardar(@Valid Producto producto,         // metodo publico
                          BindingResult br,
                          RedirectAttributes ra) {
        if (br.hasErrors()) return "productos/form";        // condicional
        repo.save(producto);
        ra.addFlashAttribute("msg", "Producto guardado");
        return "redirect:/productos";                       // devuelve el resultado al llamador
    }

    @PostMapping("/{id}/eliminar")                          // endpoint HTTP POST (crear)
    public String eliminar(@PathVariable Long id,           // metodo publico
                           RedirectAttributes ra) {
        repo.deleteById(id);
        ra.addFlashAttribute("msg", "Producto eliminado");
        return "redirect:/productos";                       // devuelve el resultado al llamador
    }
}
\end{lstlisting}

\textbf{Resultado esperado (listado):}

\fbox{\begin{minipage}{0.85\textwidth}
\vspace{0.4em}
{\Large\bfseries\color{uteqverde} Productos UTEQ}\\[0.3em]
\colorbox{uteqverde}{\textcolor{white}{\bfseries\ Nuevo producto\ }}\\[0.5em]
\begin{tabular}{lll}
\textbf{Nombre} & \textbf{Precio} & \textbf{Stock} \\
\hline
Café galletti & \$4.50 & 30 \\
Cacao Pacari  & \$6.00 & 5  \\
Té orgánico   & \$3.20 & 12 \\
\end{tabular}
\vspace{0.4em}
\end{minipage}}
\end{cajaejemplo}

\section{Ejercicios}

Los siguientes ejercicios consolidan Spring Boot. El
estudiante debe entregar el código fuente y capturas funcionando
desde Postman o el navegador.


\begin{cajaejercicio}{Ejercicio 6.1 (PROPUESTO) --- API básica de tareas}


\textbf{Enunciado.} Construir una API REST básica de tareas en Spring Boot con las cuatro operaciones CRUD y persistencia en H2.

Crear una API REST en Spring Boot que exponga endpoints CRUD para
una entidad \texttt{Tarea} (id, titulo, completada, fechaCreacion).
Probar con Postman. Devolver JSON. Validar entrada con
\texttt{@Valid}. Score \emph{thumbs-up} si los endpoints siguen
convenciones REST.
\end{cajaejercicio}

\begin{cajaejercicio}{Ejercicio 6.2 (PROPUESTO) --- Sistema de inscripciones}


\textbf{Enunciado.} Construir un sistema de inscripciones a cursos en Spring Boot con relaciones JPA (un curso tiene muchos estudiantes).

Construir una pequeña aplicación Spring Boot + Thymeleaf donde un
estudiante pueda inscribirse a un curso (formulario), ver los cursos
disponibles (listado) y darse de baja (botón eliminar). BD H2 en
memoria.
\end{cajaejercicio}

\section{Buenas prácticas Spring Boot}

Escribir Spring Boot profesional requiere disciplina específica
del ecosistema. La caja siguiente reúne doce reglas que el
estudiante debe respetar en todo proyecto Spring.


\begin{cajabuenas}{Doce reglas profesionales 2026}
\begin{enumerate}
\item Controllers \emph{thin}: lógica en servicios.
\item Capas: controller $\to$ service $\to$ repository.
\item DTOs separados de entidades JPA.
\item \texttt{@Valid} en TODA entrada de controllers.
\item Manejo global de excepciones con
\texttt{@ControllerAdvice}.
\item Patrón PRG (Post-Redirect-Get).
\item Lombok para reducir boilerplate.
\item Slf4j + Logback con niveles correctos.
\item Tests con \texttt{@WebMvcTest},
\texttt{@DataJpaTest}.
\item Profiles separados (dev, test, prod).
\item Configuración externa con
\texttt{application-\{profile\}.yml}.
\item Health checks habilitados con Spring Boot Actuator.
\end{enumerate}
\end{cajabuenas}

%%% ==== SEMANA 07 v3 ====
% =====================================================================
%   SEMANA 7 - VERSIÓN V3 MEJORADA
%   ASP.NET Core: la plataforma .NET para aplicaciones web modernas
% =====================================================================

\chapter{ASP.NET Core: la plataforma .NET para aplicaciones web modernas}
\label{cap:semana7}

\epigraph{«La plataforma .NET pasó de ser propietaria y Windows-only a
ser open-source y multiplataforma. Es el mayor giro estratégico de
Microsoft en los últimos 20 años.»}{--- \textit{Adaptado de Scott
Hanselman, 2024}}

\section*{Resultado de aprendizaje}
Al concluir la semana 7, el estudiante construye aplicaciones web con
ASP.NET Core 8 y C\# 12, comprende su modelo de hosting Kestrel y
desarrolla un CRUD funcional con Entity Framework Core
\cite{lock2021asp,galloway2011mvc}.

\section{Historia de ASP.NET}

ASP.NET es la plataforma de Microsoft para desarrollo web.
Tras tres décadas de evoluciones (ASP clásico → WebForms → MVC →
.NET Core → .NET 8), hoy es multiplataforma, de código abierto y
con rendimiento comparable a Java/Go.


\begin{cajahistoria}{De ASP clásico a ASP.NET Core multiplataforma}
Microsoft lanzó \textbf{ASP} (Active Server Pages) en 1996 como
respuesta a PHP y JSP. Era propietario, Windows-only y basado en
VBScript/JScript. En 2002 evolucionó a \textbf{ASP.NET WebForms},
que usaba .NET Framework y un modelo de eventos similar a Windows
Forms. En 2009 llegó \textbf{ASP.NET MVC} \cite{galloway2011mvc},
inspirado en Ruby on Rails.

El giro mayor ocurrió en 2016 con \textbf{ASP.NET Core 1.0}: open
source, multiplataforma (Linux, macOS, Windows), modular y de alto
rendimiento. ASP.NET Core 8.0 (2023) es la versión actual estable;
ASP.NET Core 9 (noviembre 2024) y 10 (en preview en 2026) marcan la
hoja de ruta. Hoy ASP.NET Core compite con Spring Boot y Node.js en
las grandes empresas \cite{lock2021asp}.
\end{cajahistoria}

\section{Conceptos fundamentales}

Antes de implementar la Web API CRUD del ejemplo, el
estudiante debe conocer los dos conceptos básicos del ecosistema:
la línea de comandos \texttt{dotnet} y la estructura típica de un
proyecto ASP.NET Core.


\subsection{La línea de comandos \texttt{dotnet}}

La CLI \texttt{dotnet} es la herramienta universal para
crear, compilar y ejecutar proyectos .NET. El estudiante debe
familiarizarse con sus comandos esenciales.

\textbf{Ejemplo corto comentado (comandos esenciales \texttt{dotnet}):}
\begin{lstlisting}[language=bash,numbers=none,caption={Ejemplo de Shell},label=lst:7.1]
# 'new' crea un proyecto a partir de una plantilla
dotnet new webapi -n MiApi      # nueva Web API
dotnet new mvc    -n MiSitio    # nueva app MVC

# 'restore' descarga dependencias (NuGet packages)
dotnet restore

# 'build' compila el codigo
dotnet build

# 'run' ejecuta la aplicacion en modo desarrollo
dotnet run

# 'test' ejecuta los tests unitarios
dotnet test
\end{lstlisting}


\begin{lstlisting}[style=shell,numbers=none,caption={Ejemplo de Shell},label=lst:7.2]
dotnet --version              # ver version instalada
dotnet new webapi -n MiApi    # crear proyecto Web API
dotnet new mvc -n MiMvc       # crear proyecto MVC
dotnet run                    # ejecutar
dotnet watch run              # con hot reload
dotnet test                   # ejecutar tests
dotnet publish -c Release     # publicar para produccion
\end{lstlisting}

\subsection{Estructura típica de un proyecto}

Un proyecto ASP.NET Core sigue una convención de carpetas y
archivos que conviene reconocer antes de programar. La estructura
siguiente es la más habitual. V\'ease el Listado~\ref{lst:7.3}.


\begin{lstlisting}[style=shell,numbers=none,caption={Ejemplo de Shell},label=lst:7.3]
MiApp/
+-- Controllers/      # API o MVC controllers
+-- Models/           # Entidades y DTOs
+-- Views/            # Razor templates (solo MVC)
+-- wwwroot/          # Recursos estaticos
+-- Program.cs        # Punto de entrada
+-- appsettings.json  # Configuracion
`-- MiApp.csproj      # Definicion del proyecto
\end{lstlisting}

\section{Ejemplo completo: Web API CRUD}

Con los conceptos básicos cubiertos, el estudiante está
listo para construir una Web API CRUD completa con Entity Framework
Core y SQLite. El siguiente ejemplo resuelto la implementa paso a
paso.


\begin{cajaejemplo}{Ejemplo 7.1 (RESUELTO) --- API REST de productos en ASP.NET Core 8}


\textbf{Enunciado.} Construir una Web API CRUD completa con ASP.NET Core 8 y Entity Framework Core sobre SQLite, exponiendo Swagger UI para pruebas interactivas.

\textbf{IDE recomendado:} JetBrains Rider (mejor para .NET) o
IntelliJ IDEA Ultimate con plugin .NET.

\textbf{Paso 1:} Verificar .NET 8: \texttt{dotnet --version} debe
mostrar \texttt{8.0.x}. Si no, instalar desde
\url{https://dotnet.microsoft.com/download}.

\textbf{Paso 2:} Crear el proyecto en terminal (Listado~\ref{lst:7.4}):
\begin{lstlisting}[style=shell,numbers=none,caption={Ejemplo de Shell},label=lst:7.4]
dotnet new webapi -n ProductosApi
cd ProductosApi                                             # cambia directorio
dotnet add package Microsoft.EntityFrameworkCore.Sqlite
dotnet add package Microsoft.EntityFrameworkCore.Design
\end{lstlisting}

\textbf{Paso 3:} Abrir el proyecto en el IDE. Crear \texttt{Models/Producto.cs} (Listado~\ref{lst:7.5}):

\begin{lstlisting}[style=cs,caption={Models/Producto.cs},label=lst:7.5]
using System.ComponentModel.DataAnnotations;                // importa namespace

namespace ProductosApi.Models;                              // namespace del archivo

public class Producto                                       // declaracion de clase publica
{
    public long Id { get; set; }                            // propiedad auto-implementada

    [Required, MinLength(3), MaxLength(120)]                // validacion: campo obligatorio
    public string Nombre { get; set; } = "";                // propiedad auto-implementada

    [Range(0.01, 99999.99)]                                 // validacion: rango numerico
    public decimal Precio { get; set; }                     // propiedad auto-implementada

    [Range(0, int.MaxValue)]                                // validacion: rango numerico
    public int Stock { get; set; }                          // propiedad auto-implementada
}
\end{lstlisting}

\textbf{Paso 4:} Crear \texttt{Data/AppDbContext.cs} (Listado~\ref{lst:7.6}):

\begin{lstlisting}[style=cs,caption={Data/AppDbContext.cs},label=lst:7.6]
using Microsoft.EntityFrameworkCore;                        // importa namespace
using ProductosApi.Models;                                  // importa namespace

namespace ProductosApi.Data;                                // namespace del archivo

public class AppDbContext : DbContext                       // declaracion de clase publica
{
    public AppDbContext(DbContextOptions<AppDbContext> opt) // metodo o miembro publico
        : base(opt) {}

    public DbSet<Producto> Productos => Set<Producto>();    // tabla expuesta por Entity Framework
}
\end{lstlisting}

\textbf{Paso 5:} Configurar \texttt{Program.cs} (Listado~\ref{lst:7.7}):

\begin{lstlisting}[style=cs,caption={Program.cs},label=lst:7.7]
using Microsoft.EntityFrameworkCore;                        // importa namespace
using ProductosApi.Data;                                    // importa namespace

var builder = WebApplication.CreateBuilder(args);           // crea el builder de la app
builder.Services.AddControllers();                          // configuracion de servicios DI
builder.Services.AddEndpointsApiExplorer();                 // configuracion de servicios DI
builder.Services.AddSwaggerGen();                           // configuracion de servicios DI
builder.Services.AddDbContext<AppDbContext>(opt =>          // registra el DbContext en el DI container
    opt.UseSqlite("Data Source=productos.db"));

var app = builder.Build();                                  // construye la app a partir del builder
app.UseSwagger();
app.UseSwaggerUI();
app.UseHttpsRedirection();
app.MapControllers();                                       // mapea los controladores REST a rutas
app.Run();                                                  // arranca la aplicacion
\end{lstlisting}

\textbf{Paso 6:} Crear el controller \texttt{Controllers/ProductosController.cs} (Listado~\ref{lst:7.8}):

\begin{lstlisting}[style=cs,caption={ProductosController.cs},label=lst:7.8]
using Microsoft.AspNetCore.Mvc;                             // importa namespace
using Microsoft.EntityFrameworkCore;                        // importa namespace
using ProductosApi.Data;                                    // importa namespace
using ProductosApi.Models;                                  // importa namespace

namespace ProductosApi.Controllers;                         // namespace del archivo

[ApiController]                                             // controlador de Web API (validacion automatica)
[Route("api/[controller]")]                                 // ruta base del controlador
public class ProductosController : ControllerBase           // declaracion de clase publica
{
    private readonly AppDbContext _ctx;                     // metodo o miembro privado
    // metodo o miembro publico
    // metodo o miembro publico
    public ProductosController(AppDbContext ctx) => _ctx = ctx;

    [HttpGet]                                               // endpoint HTTP GET
    public async Task<IEnumerable<Producto>> Listar()       // metodo o miembro publico
        => await _ctx.Productos.ToListAsync();              // ejecuta consulta y devuelve lista (async)

    [HttpGet("{id:long}")]                                  // endpoint HTTP GET
    // metodo o miembro publico
    // metodo o miembro publico
    public async Task<ActionResult<Producto>> Obtener(long id)
    {
        var p = await _ctx.Productos.FindAsync(id);         // busca por clave primaria (async)
        return p is null ? NotFound() : Ok(p);              // devuelve el resultado al llamador
    }

    [HttpPost]                                              // endpoint HTTP POST (crear)
    // metodo o miembro publico
    // metodo o miembro publico
    public async Task<ActionResult<Producto>> Crear(Producto p)
    {
        _ctx.Productos.Add(p);
        await _ctx.SaveChangesAsync();                      // persiste cambios en la BD (async)
        // HTTP 201 (recurso creado)
        // HTTP 201 (recurso creado)
        return CreatedAtAction(nameof(Obtener), new { id = p.Id }, p);
    }

    [HttpDelete("{id:long}")]                               // endpoint HTTP DELETE (eliminar)
    public async Task<IActionResult> Eliminar(long id)      // metodo o miembro publico
    {
        var p = await _ctx.Productos.FindAsync(id);         // busca por clave primaria (async)
        if (p is null) return NotFound();                   // HTTP 404 (no encontrado)
        _ctx.Productos.Remove(p);
        await _ctx.SaveChangesAsync();                      // persiste cambios en la BD (async)
        return NoContent();                                 // HTTP 204 (sin cuerpo)
    }
}
\end{lstlisting}

\textbf{Paso 7:} Crear la BD con migraciones (Listado~\ref{lst:7.9}):
\begin{lstlisting}[style=shell,numbers=none,caption={Ejemplo de Shell},label=lst:7.9]
dotnet ef migrations add Inicial
dotnet ef database update
\end{lstlisting}

\textbf{Paso 8:} Ejecutar: \texttt{dotnet run}. Abrir
\url{https://localhost:7xxx/swagger}: aparece Swagger UI con todos
los endpoints.

\textbf{Resultado esperado (Swagger UI):}

\fbox{\begin{minipage}{0.85\textwidth}
\vspace{0.4em}
{\Large\bfseries ProductosApi --- Swagger UI}\\[0.3em]
\colorbox{green!20}{\textbf{GET}}\quad \texttt{/api/Productos} \quad \emph{Listar}\\[0.2em]
\colorbox{green!20}{\textbf{GET}}\quad \texttt{/api/Productos/\{id\}} \quad \emph{Obtener}\\[0.2em]
\colorbox{blue!20}{\textbf{POST}}\quad \texttt{/api/Productos} \quad \emph{Crear}\\[0.2em]
\colorbox{red!20}{\textbf{DELETE}}\quad \texttt{/api/Productos/\{id\}} \quad \emph{Eliminar}
\vspace{0.4em}
\end{minipage}}

\textbf{Probar con curl o Postman (Listado~\ref{lst:7.10}):}
\begin{lstlisting}[style=shell,numbers=none,caption={Ejemplo de Shell},label=lst:7.10]
# Crear
curl -X POST https://localhost:7000/api/productos \
  -H "Content-Type: application/json" \
  -d '{"nombre":"Cafe","precio":4.50,"stock":30}'

# Listar
curl https://localhost:7000/api/productos
\end{lstlisting}
\end{cajaejemplo}

\section{Razor Pages y MVC tradicional}

ASP.NET Core soporta tres estilos de aplicación web:
\begin{itemize}
\item \textbf{Web API}: solo JSON, sin vistas (lo que acabamos de
hacer).
\item \textbf{MVC}: con controllers, views Razor y models.
\item \textbf{Razor Pages}: cada página tiene su PageModel
(similar a Blazor).
\end{itemize}

Para el PFC del curso se recomienda \textbf{Web API + frontend
Angular separado}, que es el patrón profesional moderno.

\section{Ejercicios}

Los siguientes ejercicios consolidan ASP.NET Core. El
estudiante debe entregar el código y capturas de Swagger UI
funcionando.


\begin{cajaejercicio}{Ejercicio 7.1 (PROPUESTO) --- API de biblioteca}


\textbf{Enunciado.} Construir una API REST de biblioteca con ASP.NET Core 8, EF Core SQLite y Swagger UI.

Implementar una API REST de gestión de libros con ASP.NET Core 8 +
EF Core + SQLite. Entidad \texttt{Libro} (id, titulo, autor, isbn,
disponible). Endpoints CRUD completos. Documentar con Swagger.
\end{cajaejercicio}

\begin{cajaejercicio}{Ejercicio 7.2 (PROPUESTO) --- Comparativa Spring Boot vs ASP.NET}


\textbf{Enunciado.} Producir un informe comparativo entre Spring Boot y ASP.NET Core implementando el mismo CRUD en ambos y midiendo tiempos.

Implementar el mismo CRUD (Tarea) en Spring Boot 3 y en ASP.NET Core
8. Comparar:
\begin{itemize}
\item Líneas de código.
\item Tiempo de arranque del servidor.
\item Requests/segundo con Apache Bench (1000 req, 10 conc).
\item Productividad subjetiva (1--5).
\end{itemize}
Producir tabla comparativa en archivo \texttt{comparativa.md}.
\end{cajaejercicio}

\section{Buenas prácticas ASP.NET Core}

La caja siguiente reúne doce reglas profesionales específicas
del ecosistema .NET que el estudiante debe respetar en todo
proyecto ASP.NET Core.


\begin{cajabuenas}{Doce reglas profesionales 2026}
\begin{enumerate}
\item Inyección de dependencias siempre (no \texttt{new}).
\item \texttt{async/await} en todos los endpoints I/O.
\item DTOs separados de entidades EF.
\item Data Annotations + \texttt{ModelState.IsValid}.
\item Swagger documentando 100\% de endpoints.
\item Logs estructurados con Serilog.
\item HTTPS obligatorio (\texttt{UseHttpsRedirection}).
\item Health checks (\texttt{/health}).
\item Manejo global de excepciones con middleware.
\item Configuración por ambiente (\texttt{appsettings.\{env\}.json}).
\item Tests con xUnit y \texttt{WebApplicationFactory}.
\item EF Core con migraciones, jamás
\texttt{EnsureCreated} en producción.
\end{enumerate}
\end{cajabuenas}

%%% ==== SEMANA 08 v3 ====
% =====================================================================
%   SEMANA 8 - VERSIÓN V3 MEJORADA
%   Tutoría de refuerzo del primer corte
% =====================================================================

\chapter{Tutoría de refuerzo: consolidación y resolución de blockers}
\label{cap:semana8}

\epigraph{«La diferencia entre un buen y un excelente estudiante no
es lo que sabe: es la disciplina con que repasa lo que cree saber.»}
{--- \textit{Anónimo, atribuido al folklore académico}}

\section*{Naturaleza de la semana}
La semana 8 es una semana de \textbf{tutoría de refuerzo} previa a la
Evaluación Parcial Primer Corte (semana 9). Sus objetivos son tres:
\begin{enumerate}
\item Resolver los temas donde el grupo presentó mayor dificultad
en las semanas 1--7.
\item Consolidar el aprendizaje con ejercicios integradores.
\item Revisar el avance de la primera entrega del Proyecto Fin de
Curso (semana 5).
\end{enumerate}

\section{Mapa conceptual del primer corte}

Antes del repaso intensivo y los ejercicios integradores,
conviene visualizar globalmente qué se ha cubierto en las semanas
1--7. El siguiente mapa conceptual lo sintetiza.


\begin{figure}[htbp]
\centering
\begin{tikzpicture}[every node/.style={font=\scriptsize,align=center},
  unidad/.style={rectangle,draw=uteqverde,thick,fill=uteqverde!10,
                 rounded corners=4pt,minimum width=3cm,minimum height=1.5cm}]
  \node[unidad] (cli) {Cliente\\HTML5 + CSS3 + JS\\Sem 1--4};
  \node[unidad,right=1cm of cli] (srv) {Servidor\\PHP + Java + .NET\\Sem 5--7};
  \node[circle,draw=uteqdorado,thick,fill=uteqdorado!20,
        minimum size=1.5cm,below right=0.5cm and 0cm of cli,
        font=\bfseries\scriptsize] (ep1) {EP1\\Sem 9};
  \draw[->,thick] (cli) -- (srv);
  \draw[->,thick,dashed,uteqdorado!70!black] (cli) -- (ep1);
  \draw[->,thick,dashed,uteqdorado!70!black] (srv) -- (ep1);
\end{tikzpicture}
\caption{Estructura del primer corte: las semanas 1--4 cubren cliente,
las 5--7 servidor, y todas convergen en la EP1.}
\end{figure}

\section{Errores frecuentes detectados (sem 1--7)}

La calificación de prácticas y entregas de las semanas 1--7
permite identificar los errores más frecuentes del grupo. Atacarlos
ahora ---antes del EP1--- mejora significativamente el desempeño.


\begin{cajaadvertencia}{Diez errores que se debe haber superado}
\begin{enumerate}
\item No declarar \texttt{<!DOCTYPE html>} o usar \texttt{lang}
incorrecto.
\item Anidar elementos de bloque dentro de inline
(\texttt{<a>} dentro de \texttt{<span>}).
\item Saltar niveles de encabezado (\texttt{h1} $\to$ \texttt{h3}).
\item Falta de \texttt{alt} en imágenes (criterio WCAG 1.1.1).
\item Uso de \texttt{<table>} para maquetar (no es 2005).
\item CSS desktop-first en lugar de mobile-first.
\item \texttt{var} en lugar de \texttt{const}/\texttt{let}.
\item Comparación con \texttt{==} en lugar de \texttt{===}.
\item Concatenar SQL con datos de usuario en PHP.
\item No usar \texttt{declare(strict\_types=1)} en PHP 8.
\end{enumerate}
\end{cajaadvertencia}

\section{Ejercicios integradores}

Los siguientes ejercicios integradores combinan HTML, CSS,
JavaScript, PHP y Spring Boot. Sirven como simulación de la
profundidad esperada en el EP1.


\begin{cajaejercicio}{Ejercicio 8.1 (RESUELTO --- guía paso a paso) --- Mini-CRUD frontend+backend}

\textbf{Objetivo:} integrar HTML5 + CSS3 + JavaScript (cliente) con
PHP (servidor) en un mini-CRUD completo.

\textbf{Paso 1:} En XAMPP, crear carpeta
\texttt{htdocs/mini-crud}. Dentro, crear:
\begin{itemize}
\item \texttt{index.html}: formulario y lista.
\item \texttt{api.php}: endpoint que recibe POST y GET.
\item \texttt{datos.json}: archivo de persistencia simple.
\end{itemize}

\textbf{Paso 2:} \texttt{api.php} (Listado~\ref{lst:8.1}):
\begin{lstlisting}[style=php,numbers=none,caption={Ejemplo de PHP},label=lst:8.1]
<?php
declare(strict_types=1);
header('Content-Type: application/json');                   // envia cabecera HTTP al cliente

$archivo = __DIR__ . '/datos.json';
$datos = file_exists($archivo)
       ? json_decode(file_get_contents($archivo), true) : [];

if ($_SERVER['REQUEST_METHOD'] === 'POST') {                // condicional
    $body = json_decode(file_get_contents('php://input'), true);
    $nombre = trim($body['nombre'] ?? '');
    if (strlen($nombre) < 3) {                              // condicional
        http_response_code(422);
        echo json_encode(['error' => 'Nombre muy corto']);  // imprime al output del servidor
        exit;
    }
    $datos[] = ['id' => count($datos)+1, 'nombre' => $nombre];
    file_put_contents($archivo, json_encode($datos));
}
echo json_encode($datos);                                   // imprime al output del servidor
\end{lstlisting}

\textbf{Paso 3:} \texttt{index.html} con JavaScript que consume la
API.

\textbf{Resultado esperado:} formulario que añade nombres y los
muestra en lista sin recargar la página.

\fbox{\begin{minipage}{0.75\textwidth}
\vspace{0.4em}
\fbox{\parbox{0.6\linewidth}{Nuevo nombre}}
\colorbox{uteqverde}{\textcolor{white}{\bfseries\ Agregar\ }}\\[0.5em]
1. Ana\\
2. Luis\\
3. María
\vspace{0.4em}
\end{minipage}}
\end{cajaejercicio}

\begin{cajaejercicio}{Ejercicio 8.2 (PROPUESTO) --- Página de inscripción end-to-end}


\textbf{Enunciado.} Construir una página de inscripción end-to-end (cliente + servidor + base de datos) que pueda ejecutarse localmente.

Construir una página de inscripción a un curso que integre todo lo
aprendido:
\begin{itemize}
\item HTML5 semántico (\texttt{<form>}, \texttt{<label>}).
\item CSS responsive (mobile-first).
\item JavaScript con validación client-side.
\item PHP server-side validando y guardando en JSON o BD.
\item Mensaje de éxito tras envío correcto.
\end{itemize}
\end{cajaejercicio}

\section{Repaso intensivo de los temas más difíciles}

Tres temas concentran la mayor parte de los errores en EP1
históricamente. Reforzarlos explícitamente antes del examen es
prioritario.


\subsection{Tema 1: Diferencia GET vs POST}

La diferencia entre GET y POST no es trivial: tiene
implicaciones de seguridad, caché, idempotencia y semántica.

\textbf{Ejemplo corto comentado (Listado~\ref{lst:8.2}):}
\begin{lstlisting}[language=HTML5,numbers=none,caption={Ejemplo de HTML},label=lst:8.2]
<!-- GET: parametros en la URL, visibles, cacheables, idempotente   -->
<!-- Uso: BUSQUEDAS, consultas, navegacion. No modifica el servidor -->
<form action="/buscar" method="GET">
  <input name="q">               <!-- ?q=ingenieria             -->
</form>

<!-- POST: parametros en el cuerpo, no visibles, no cacheables       -->
<!-- Uso: CREAR recursos, enviar contrasenas, subir archivos        -->
<form action="/usuarios" method="POST">
  <input name="email">
  <input name="password" type="password">
</form>
\end{lstlisting}


\begin{cajadefinicion}{Resumen ejecutivo}
\textbf{GET}: idempotente, seguro, parámetros en URL,
cacheable. Para \emph{consultar} recursos.

\textbf{POST}: no idempotente, modifica estado, parámetros en cuerpo,
no cacheable. Para \emph{crear} recursos.
\end{cajadefinicion}

\subsection{Tema 2: Validación cliente vs servidor}

Validar solo en el cliente es como cerrar la puerta de la
calle pero dejar la ventana abierta: cualquier usuario puede saltarse
las validaciones del navegador con DevTools o Postman. La validación
del servidor es \emph{obligatoria}; la del cliente solo añade UX.

\textbf{Ejemplo corto comentado (Listado~\ref{lst:8.3}):}
\begin{lstlisting}[language=PHP,numbers=none,caption={Ejemplo de PHP},label=lst:8.3]
<?php
// CLIENTE (HTML5): mejora UX, no es seguridad
?>
<input type="email" required maxlength="100">

<?php
// SERVIDOR (PHP): obligatorio, es la unica defensa real
// valida/sanea entrada con filtro
// valida/sanea entrada con filtro
$email = filter_var($_POST['email'] ?? '', FILTER_VALIDATE_EMAIL);
if (!$email)          die('Email invalido');           // valido?
if (strlen($email) > 100) die('Demasiado largo');      // limite?
?>
\end{lstlisting}


\begin{cajadefinicion}{Resumen ejecutivo}
\textbf{Cliente (HTML5/JS)}: para UX rápida. NUNCA es la única
defensa.

\textbf{Servidor (PHP/Java/.NET)}: indispensable. Toda entrada debe
revalidarse aquí porque el cliente puede ser manipulado.
\end{cajadefinicion}

\subsection{Tema 3: Codificación de caracteres UTF-8}

Una página que muestra caracteres como \emph{??Hola Mar\'?a??}
en lugar de \emph{¡Hola María!} casi siempre tiene un problema de
codificación. UTF-8 debe aparecer en cuatro lugares: HTML, base de
datos, conexión y cabecera HTTP. Si falla en uno solo, se rompen
las tildes.


\begin{cajadefinicion}{Resumen ejecutivo}
\textbf{HTML}: \texttt{<meta charset=\char34{}UTF-8\char34{}>} dentro de
\texttt{<head>}.

\textbf{PHP}: archivos guardados como UTF-8 sin BOM.

\textbf{BD MySQL}: tabla con \texttt{CHARACTER SET utf8mb4}.

\textbf{HTTP}: cabecera
\texttt{Content-Type: text/html; charset=utf-8}.

Si alguno de los cuatro falla, aparecen los famosos «ñ» como «Ã±».
\end{cajadefinicion}

\section{Preparación para la EP1}

La EP1 cubre las semanas 1--7. Sigue la estructura habitual:
\begin{itemize}
\item Parte A: selección múltiple + V/F (20 pts).
\item Parte B: respuesta corta (20 pts).
\item Parte C: análisis de código HTML/CSS/JS/PHP (30 pts).
\item Parte D: ejercicio práctico integrador (30 pts).
\end{itemize}

\textbf{Estrategia:} resolver los ejercicios resueltos del libro
SIN mirar la solución; comparar al final.

%%% ==== SEMANA 09 v3 ====
% =====================================================================
%   SEMANA 9 - VERSIÓN V3 MEJORADA
%   Seguridad en aplicaciones web; Evaluación Parcial Primer Corte
% =====================================================================

\chapter{Seguridad en aplicaciones web; Evaluación Parcial Primer Corte}
\label{cap:semana9}

\epigraph{«El software seguro no es un producto: es un proceso
continuo. Cada commit es una oportunidad de introducir o eliminar una
vulnerabilidad.»}{--- \textit{Adaptado de Bruce Schneier}}

\section*{Resultado de aprendizaje}
Al concluir la semana 9, el estudiante identifica las diez
vulnerabilidades más críticas de aplicaciones web (OWASP Top 10:2021
\cite{owasp2021top10}), aplica controles defensivos básicos
(validación, escape, parametrización, cabeceras de seguridad), y
rinde la Evaluación Parcial Primer Corte.

\section{Historia del OWASP y la cultura de seguridad web}

La cultura de seguridad web no surgió por generación
espontánea: fue resultado de un trabajo organizado durante dos
décadas, principalmente bajo el paraguas de OWASP. Conocer esa
historia ayuda a entender por qué hoy hablamos en términos de
\emph{Top 10}, \emph{CWE} y \emph{vulnerability disclosure}.


\begin{cajahistoria}{OWASP --- 25 años de evangelización de la seguridad}
La \textbf{Open Web Application Security Project} (OWASP) fue
fundada en 2001 por Mark Curphey, ante la frustración de que la
industria del software seguía cometiendo los mismos errores de
seguridad sin aprender de ellos. OWASP es hoy una fundación sin
fines de lucro, con capítulos en más de 100 países, que mantiene
estándares abiertos, herramientas (ZAP, ModSecurity), guías de
desarrollo seguro y, sobre todo, el famoso \textbf{Top 10}: la
lista de las diez vulnerabilidades más críticas de aplicaciones
web, publicada cada 3--4 años \cite{owasp2021top10}.

El Top 10 ha pasado por revisiones en 2004, 2007, 2010, 2013,
2017 y 2021. La versión 2025 está en consulta pública y mantiene
las categorías principales con ajustes. La adopción del Top 10 se
ha vuelto contractual: muchas licitaciones públicas y privadas
exigen explícitamente cumplir sus controles.
\end{cajahistoria}

\section{OWASP Top 10:2021 \cite{owasp2021top10}}

El OWASP Top 10 es la lista consensuada de las diez
vulnerabilidades más críticas de aplicaciones web. La revisión 2021
es la vigente hasta que se publique la de 2025. La tabla siguiente
las cataloga.


\begin{table}[htbp]
\centering
\caption{OWASP Top 10:2021 \cite{owasp2021top10}.}
\footnotesize
\begin{tabularx}{\textwidth}{llX}
\toprule
\textbf{ID} & \textbf{Categoría} & \textbf{Descripción y defensa}\\
\midrule
A01 & Broken Access Control & Permitir acceso a recursos ajenos. Defensa: autorización en cada endpoint, no en cliente.\\
A02 & Cryptographic Failures & Algoritmos débiles, datos sin cifrar. Defensa: TLS 1.3, AES-256, BCrypt.\\
A03 & Injection & SQL/NoSQL/LDAP injection. Defensa: prepared statements, validación.\\
A04 & Insecure Design & Diseñar sin pensar en amenazas. Defensa: threat modeling.\\
A05 & Security Misconfiguration & Defaults inseguros, debug expuesto. Defensa: hardening checklist.\\
A06 & Vulnerable Components & Librerías obsoletas. Defensa: SCA (Snyk, Dependabot).\\
A07 & Identification/Authentication Failures & Login sin rate limiting, sin MFA. Defensa: BCrypt, MFA, rate limit.\\
A08 & Software/Data Integrity Failures & Deserialización insegura, CI/CD comprometido. Defensa: firmas digitales.\\
A09 & Security Logging/Monitoring Failures & Sin logs ni alertas. Defensa: logs estructurados, SIEM.\\
A10 & Server-Side Request Forgery (SSRF) & Servidor hace peticiones a hosts internos. Defensa: allowlist URLs.\\
\bottomrule
\end{tabularx}
\end{table}

\section{Defensas técnicas básicas}

Conocer las vulnerabilidades no basta: el estudiante debe
implementar las defensas correctas en cada lenguaje del curso. Las
subsecciones siguientes muestran las cinco defensas mínimas que
todo backend debe incorporar.


\subsection{Anti-SQL Injection: prepared statements}

SQL Injection es la vulnerabilidad número 3 del OWASP Top 10
y la causa de algunos de los ataques más devastadores de la
historia. La única defensa real son las \emph{sentencias preparadas}.

\textbf{Ejemplo corto comentado (PHP con PDO) (Listado~\ref{lst:9.1}):}
\begin{lstlisting}[language=PHP,numbers=none,caption={Ejemplo de PHP},label=lst:9.1]
<?php
// MAL: concatenar datos del usuario en SQL (vulnerable a SQLi)
// $sql = "SELECT * FROM users WHERE email = '$email'";

// BIEN: sentencia preparada con parametro nombrado
$sql = "SELECT * FROM users WHERE email = :email";
$stmt = $pdo->prepare($sql);                  // prepara la consulta
$stmt->execute([':email' => $email]);         // pasa el dato aparte
$user = $stmt->fetch();                       // recupera el resultado
?>
\end{lstlisting}


\begin{cajaejemplo}{Ejemplo 9.1 (RESUELTO) --- Defensa anti-SQLi en PHP}


\textbf{Enunciado.} Defender el código PHP contra SQL Injection usando PDO con sentencias preparadas. Se contrasta el código vulnerable con la versión segura.

\textbf{Anti-patrón (NUNCA hacer) (Listado~\ref{lst:9.2}):}
\begin{lstlisting}[style=php,numbers=none,caption={Ejemplo de PHP},label=lst:9.2]
$sql = "SELECT * FROM users WHERE email='" . $_POST['email'] . "'";
$res = $pdo->query($sql);
\end{lstlisting}

\textbf{Patrón correcto (Listado~\ref{lst:9.3}):}
\begin{lstlisting}[style=php,numbers=none,caption={Ejemplo de PHP},label=lst:9.3]
// prepara consulta SQL (defensa anti-SQLi)
// prepara consulta SQL (defensa anti-SQLi)
$stmt = $pdo->prepare('SELECT * FROM users WHERE email = :em');
$stmt->execute([':em' => $_POST['email']]);                 // ejecuta la consulta con parametros enlazados
$user = $stmt->fetch();                                     // recupera la siguiente fila como array
\end{lstlisting}

\textbf{Resultado de auditoría:} con prepared statements, los intentos
de inyección (\texttt{' OR '1'='1}) se interpretan como cadenas
literales, no como código SQL. La defensa es 100\% efectiva si se
aplica consistentemente.
\end{cajaejemplo}

\subsection{Anti-XSS: escape de salida}

XSS (\emph{Cross-Site Scripting}) ocurre cuando datos del
usuario terminan dentro del HTML sin escapar, permitiendo inyectar
JavaScript malicioso. La defensa es \emph{escapar la salida} en
cada inserción dinámica.

\textbf{Ejemplo corto comentado (Listado~\ref{lst:9.4}):}
\begin{lstlisting}[language=PHP,numbers=none,caption={Ejemplo de PHP},label=lst:9.4]
<?php
// Si el usuario introduce: <script>alert('XSS')</script>
$comentario = $_POST['comentario'];

// MAL: imprime directo (ejecuta el script en el navegador)
// echo "<p>$comentario</p>";

// BIEN: escapa los caracteres HTML especiales (< > " ' &)
// escapa caracteres HTML (defensa anti-XSS)
// escapa caracteres HTML (defensa anti-XSS)
echo "<p>" . htmlspecialchars($comentario, ENT_QUOTES, 'UTF-8') . "</p>";
?>
\end{lstlisting}


\begin{cajaejemplo}{Ejemplo 9.2 (RESUELTO) --- Defensa anti-XSS}


\textbf{Enunciado.} Defender una aplicación contra XSS escapando todas las salidas dinámicas con \texttt{htmlspecialchars()} en PHP y equivalentes en otras plataformas.

\textbf{Anti-patrón (Listado~\ref{lst:9.5}):}
\begin{lstlisting}[style=php,numbers=none,caption={Ejemplo de PHP},label=lst:9.5]
echo "Bienvenido, " . $_GET['nombre'];                      // imprime al output del servidor
\end{lstlisting}

Si \texttt{nombre=\char34<script>alert('XSS')</script>\char34}, el script
se ejecuta.

\textbf{Patrón correcto (Listado~\ref{lst:9.6}):}
\begin{lstlisting}[style=php,numbers=none,caption={Ejemplo de PHP},label=lst:9.6]
echo "Bienvenido, " . htmlspecialchars($_GET['nombre'],     // escapa caracteres HTML (defensa anti-XSS)
                                        ENT_QUOTES, 'UTF-8');
\end{lstlisting}

\textbf{Resultado:} ahora aparece literalmente
\texttt{\&lt;script\&gt;alert('XSS')\&lt;/script\&gt;}, inocuo.
\end{cajaejemplo}

\subsection{Cookies de sesión seguras}

Una cookie de sesión sin \texttt{HttpOnly} es robable por XSS;
sin \texttt{Secure} viaja en claro por HTTP; sin \texttt{SameSite}
es vulnerable a CSRF. Estas tres banderas son obligatorias en
producción.

\textbf{Ejemplo corto comentado (PHP) (Listado~\ref{lst:9.7}):}
\begin{lstlisting}[language=PHP,numbers=none,caption={Ejemplo de PHP},label=lst:9.7]
<?php
// Configurar TODOS los flags ANTES de session_start()
session_set_cookie_params([                                 // configura flags de la cookie de sesion
  'lifetime' => 1800,        // 30 minutos
  'path'     => '/',         // valida en todo el sitio
  'secure'   => true,        // SOLO por HTTPS
  'httponly' => true,        // JavaScript NO la puede leer
  'samesite' => 'Strict',    // No se envia desde otros dominios
]);
session_start();             // ahora si: arrancar sesion
?>
\end{lstlisting}


\begin{lstlisting}[style=php,numbers=none,caption={Ejemplo de PHP},label=lst:9.8]
session_set_cookie_params([                                 // configura flags de la cookie de sesion
    'lifetime' => 1800,
    'path'     => '/',
    'secure'   => true,
    'httponly' => true,
    'samesite' => 'Strict',
]);
session_start();                                            // arranca o reanuda la sesion
\end{lstlisting}

\subsection{Hash de contraseñas con BCrypt}

Guardar contraseñas en texto plano es un error de carrera;
guardar su MD5 o SHA1 también lo es desde hace 20 años. La única
forma profesional es usar funciones de hash lentas diseñadas para
contraseñas: BCrypt, Argon2 o scrypt.

\textbf{Ejemplo corto comentado (Listado~\ref{lst:9.9}):}
\begin{lstlisting}[language=PHP,numbers=none,caption={Ejemplo de PHP},label=lst:9.9]
<?php
// Al REGISTRAR: hashear antes de guardar
// hashea contrasena con BCrypt
// hashea contrasena con BCrypt
$pwHash = password_hash($password, PASSWORD_BCRYPT, ['cost' => 12]);
// $pwHash sera algo como: $2y$12$N9qo8uLOickgx2ZMRZoMye...

// Al LOGIN: comparar el password introducido con el hash guardado
if (password_verify($passwordIntroducido, $pwHashEnBD)) {   // verifica contrasena contra hash
    echo "Autenticacion correcta";                          // imprime al output del servidor
} else {
    echo "Credenciales invalidas";                          // imprime al output del servidor
}
?>
\end{lstlisting}


\begin{lstlisting}[style=php,numbers=none,caption={Ejemplo de PHP},label=lst:9.10]
// Al registrar
$hash = password_hash($passwordPlano, PASSWORD_BCRYPT,      // hashea contrasena con BCrypt
                      ['cost' => 12]);

// Al verificar (resistente a timing attacks)
if (password_verify($passwordIngresado, $hashGuardado)) {   // verifica contrasena contra hash
    // login exitoso
}
\end{lstlisting}

\subsection{Cabeceras de seguridad HTTP}

Las cabeceras HTTP \emph{de seguridad} instruyen al
navegador sobre cómo proteger al usuario: bloquear scripts no
confiables, prevenir \emph{clickjacking}, forzar HTTPS. Configurarlas
es trabajo del servidor y debe hacerse en todo proyecto.

\textbf{Ejemplo corto comentado (Listado~\ref{lst:9.11}):}
\begin{lstlisting}[language=PHP,numbers=none,caption={Ejemplo de PHP},label=lst:9.11]
<?php
// CSP: solo carga scripts/styles del propio dominio
header("Content-Security-Policy: default-src 'self'");      // envia cabecera HTTP al cliente
// HSTS: fuerza HTTPS durante 1 ano
// envia cabecera HTTP al cliente
// envia cabecera HTTP al cliente
header("Strict-Transport-Security: max-age=31536000; includeSubDomains");
// Bloquea framing externo (anti-clickjacking)
header("X-Frame-Options: DENY");                            // envia cabecera HTTP al cliente
// Evita sniffing del Content-Type por el navegador
header("X-Content-Type-Options: nosniff");                  // envia cabecera HTTP al cliente
?>
\end{lstlisting}


\begin{lstlisting}[style=php,numbers=none,caption={Ejemplo de PHP},label=lst:9.12]
// envia cabecera HTTP al cliente
// envia cabecera HTTP al cliente
header('Strict-Transport-Security: max-age=31536000; includeSubDomains');
header("Content-Security-Policy: default-src 'self'");      // envia cabecera HTTP al cliente
header('X-Frame-Options: DENY');                            // envia cabecera HTTP al cliente
header('X-Content-Type-Options: nosniff');                  // envia cabecera HTTP al cliente
header('Referrer-Policy: strict-origin-when-cross-origin'); // envia cabecera HTTP al cliente
\end{lstlisting}

\section{Ejercicios}

Los siguientes ejercicios consolidan las defensas de
seguridad. El estudiante debe entregar código que aplique todas
las protecciones vistas.


\begin{cajaejercicio}{Ejercicio 9.1 (RESUELTO) --- Login seguro completo}


\textbf{Enunciado.} Construir un sistema de login seguro completo: validación, hash BCrypt, sesiones con flags, cabeceras de seguridad, defensa anti-CSRF.

Implementar un login PHP con TODAS las defensas: BCrypt, PDO prepared
statements, rate limiting (5 intentos), regeneración de session ID,
flags de cookie. Código de referencia en sem 10 ejemplo 10.5.

\textbf{Resultado esperado:}

\fbox{\begin{minipage}{0.75\textwidth}
\vspace{0.4em}
{\Large\bfseries\color{uteqverde} Iniciar sesión}\\[0.5em]
Correo:\\ \fbox{\parbox{0.85\linewidth}{admin@uteq.edu.ec}}\\
Contraseña:\\ \fbox{\parbox{0.85\linewidth}{$\bullet\bullet\bullet\bullet\bullet\bullet\bullet\bullet$}}\\[0.4em]
\colorbox{uteqverde}{\textcolor{white}{\bfseries\ Ingresar\ }}
\vspace{0.4em}
\end{minipage}}

Tras login exitoso, en DevTools (F12 $\rightarrow$ Application
$\rightarrow$ Cookies) la cookie debe tener: \texttt{HttpOnly=true},
\texttt{Secure=true}, \texttt{SameSite=Strict}.
\end{cajaejercicio}

\begin{cajaejercicio}{Ejercicio 9.2 (PROPUESTO) --- Auditoría de un sitio web}


\textbf{Enunciado.} Realizar una auditoría de seguridad sobre un sitio existente (propio o público) usando OWASP ZAP o herramientas equivalentes.

Tomar un sitio web pequeño que el estudiante haya construido en
semanas previas, y auditarlo manualmente contra los 10 controles
OWASP. Producir un informe con:
\begin{itemize}
\item Control evaluado.
\item Estado (cumple/no cumple/parcial).
\item Evidencia (captura, comando).
\item Recomendación.
\end{itemize}
\end{cajaejercicio}

\section{Evaluación Parcial Primer Corte}

La EP1 se rinde al final de esta semana y cubre las semanas 1--8.

\subsection{Estructura}

El EP1 cubrirá las semanas 1--8 con una estructura conocida
de antemano para que el estudiante pueda preparar específicamente.


\begin{tabularx}{\textwidth}{lXl}
\toprule
\textbf{Parte} & \textbf{Tipo} & \textbf{Pts}\\
\midrule
A & Selección múltiple + V/F & 20\\
B & Respuesta corta (3 líneas máx) & 20\\
C & Análisis de código (errores) & 30\\
D & Ejercicio integrador & 30\\
\bottomrule
\end{tabularx}

\subsection{Temas garantizados}

Los siguientes temas aparecen en TODA edición del EP1 sin
excepción. Dominarlos garantiza un mínimo de 60\% del puntaje.


\begin{itemize}
\item HTML5 semántico y accesibilidad WCAG.
\item CSS3, Flexbox, Grid, mobile-first.
\item JavaScript moderno, DOM, fetch/async.
\item PHP 8 con PDO, validación, sesiones.
\item Spring Boot 3 con JPA básico.
\item ASP.NET Core 8 con EF Core básico.
\item OWASP Top 10 (especialmente A03 SQL Injection y A07
Authentication).
\end{itemize}

\section{Buenas prácticas de seguridad}

La seguridad es transversal: no se ``añade al final'' del
proyecto. La caja siguiente compila quince reglas profesionales que
el estudiante debe incorporar desde la línea uno de código.


\begin{cajabuenas}{Quince reglas profesionales 2026}
\begin{enumerate}
\item HTTPS obligatorio. TLS 1.3 o 1.2 mínimo.
\item Nunca confiar en datos del cliente: revalidar siempre.
\item Prepared statements en TODA consulta SQL.
\item \texttt{htmlspecialchars} (o equivalente) en TODA salida.
\item BCrypt o Argon2 para contraseñas. JAMÁS MD5/SHA1 puros.
\item Cookies de sesión con \texttt{HttpOnly}, \texttt{Secure},
\texttt{SameSite}.
\item Rate limiting en login y endpoints sensibles.
\item Cabeceras de seguridad (HSTS, CSP, XFO).
\item Mantener dependencias actualizadas (Dependabot, Snyk).
\item Logs de eventos de seguridad (login, fallos, accesos
denegados).
\item No exponer detalles en mensajes de error al usuario.
\item Sin secretos en el repositorio (\texttt{.env}, vault).
\item Backups cifrados y restaurables periódicamente.
\item Threat modeling al inicio de cada feature.
\item Pentest externo antes de release mayor.
\end{enumerate}
\end{cajabuenas}


% =====================================================================
%   PARTE 2 AMPLIADA --- SEMANAS 10 A 14
%
%   Para integrar con la Parte 1 Ampliada:
%   - Agregar al preámbulo el contenido de "preambulo_modificaciones.tex"
%   - Eliminar el "\end{document}" de la Parte 1 Ampliada (después de la
%     bibliografía) y pegar este contenido en su lugar.
%   - Actualizar la bibliografía con las entradas indicadas.
%   - Agregar al final un nuevo "\end{document}".
% =====================================================================

% Reseteamos el contador de capítulos para continuar la numeración
% adecuada (las semanas 1-9 son los capítulos 1-9 de la Parte 1).
\setcounter{chapter}{9}

% =====================================================================
%       SEMANA 10: GESTIÓN DE ESTADO Y OBJETOS INTEGRADOS
% =====================================================================
\chapter{Gestión de estado y objetos integrados en aplicaciones web}
\label{cap:semana10}

\epigraph{«HTTP es el protocolo de la amnesia colectiva: cada petición
nace y muere sin memoria. Programar la web es, en buena medida, el arte
de simular una memoria que el protocolo nos negó.»}{--- \textit{Adaptado
de Roy Fielding}}

\section*{Resultado de aprendizaje de la semana}
Al concluir la semana 10, el estudiante diseña e implementa soluciones
web que integran mecanismos robustos de gestión de estado --- cookies,
sesiones de servidor, almacenamiento del cliente y tokens --- y domina
los objetos integrados \texttt{Request}, \texttt{Response},
\texttt{Session}, \texttt{Application} y \texttt{Page} en PHP, Java y
ASP.NET Core, aplicando los flags de seguridad estándar que la
industria exige en 2026.

\section{Introducción: el dilema fundamental de HTTP}

HTTP es un protocolo \emph{stateless}: cada petición se
trata independientemente, sin recordar las anteriores. Esto es una
fortaleza para escalar (cualquier servidor puede atender cualquier
petición), pero un problema para construir aplicaciones que necesitan
recordar al usuario entre clics. La resolución de esta tensión es el
tema central de la semana.


\begin{cajahistoria}{HTTP --- el protocolo amnésico que cambió el mundo}
Cuando Tim Berners-Lee diseñó HTTP en 1989, lo hizo deliberadamente
como un protocolo \emph{stateless} (sin estado): cada petición es
independiente, autocontenida, y el servidor no recuerda nada de la
petición anterior. Esta decisión, tomada en una época de servidores
con poca memoria y conexiones lentas, parecía limitada pero resultó
genial: permitió que la web escalara a miles de millones de usuarios
sin colapsar bajo el peso de la memoria de sesiones.

Pero la web pronto necesitó \emph{algo} que pareciera memoria: ¿cómo
saber qué usuario está autenticado?, ¿cómo recordar los productos en
el carrito?, ¿cómo personalizar la experiencia? La respuesta llegó en
1994 cuando Lou Montulli, ingeniero de Netscape, inventó las
\textbf{cookies HTTP} originalmente como solución para un cliente de
e-commerce. Las cookies fueron una idea simple: el servidor envía un
identificador en la respuesta; el navegador se lo guarda; en cada
petición posterior, el navegador devuelve ese identificador. Con eso,
el servidor puede asociar peticiones del mismo usuario.

Sobre las cookies se construyeron las \textbf{sesiones de servidor}:
el servidor guarda los datos del usuario en su propia memoria
asociados a un \emph{session ID} aleatorio, y solo el ID viaja en la
cookie. Treinta años después, este modelo sigue siendo el sustrato de
la mayoría de aplicaciones web tradicionales del mundo. Pero la
arquitectura moderna --- microservicios, escalado horizontal, SPAs
--- llevó a la emergencia de los \textbf{tokens stateless} (JWT) que
trasladan el estado de vuelta al cliente, permitiendo que el servidor
permanezca verdaderamente sin memoria.

La literatura técnica reciente confirma que la gestión de estado sigue
siendo uno de los aspectos más críticos de la arquitectura web:
\cite{fielding2000rest} en su tesis doctoral fundacional estableció
que \emph{statelessness} es una restricción arquitectónica deseable
para escalabilidad; \cite{kleppmann2017designing} discute en
profundidad cómo el estado distribuido es la fuente de la mayoría de
los bugs sutiles en sistemas modernos; y
\cite{jones2015jwt} (RFC 7519) estandarizó JWT como el mecanismo
dominante para sesiones distribuidas, hoy ampliamente adoptado en
APIs públicas \cite{jones2015jwt}.
\end{cajahistoria}

\subsection{Las cinco estrategias de gestión de estado}

Hay cinco estrategias diferentes para preservar estado entre
peticiones HTTP. Cada una tiene ventajas, desventajas y casos donde
es la mejor opción. La tabla siguiente las contrasta.


\begin{table}[htbp]
\centering
\caption{Estrategias de gestión de estado en aplicaciones web modernas.}
\footnotesize
\begin{tabularx}{\textwidth}{lXXl}
\toprule
\textbf{Mecanismo} & \textbf{Ubicación del estado} & \textbf{Caso de uso ideal} & \textbf{Tamaño}\\
\midrule
Cookie clásica & Cliente (cookie HTTP) & Preferencias simples (idioma, tema) & 4 KB\\
Sesión servidor & Servidor (memoria/Redis) & App tradicional multi-página & MB \\
\texttt{localStorage} & Cliente (navegador) & Datos persistentes locales & 5--10 MB\\
\texttt{sessionStorage} & Cliente (pestaña) & Datos volátiles de pestaña & 5--10 MB\\
JWT & Cliente (cookie/header) & APIs REST stateless, SPAs & 1--4 KB\\
IndexedDB & Cliente (navegador) & Apps offline, gran volumen & GB\\
\bottomrule
\end{tabularx}
\end{table}

\subsection{Los flags de seguridad de cookies en 2026}

Antes de profundizar en cualquier mecanismo, hay que dominar los
\emph{flags} de seguridad que toda cookie debe llevar en producción:

\begin{cajadefinicion}{Los cuatro flags imprescindibles}
\begin{description}[leftmargin=1.5cm]
\item[\texttt{Secure}.] La cookie solo viaja por HTTPS, nunca por HTTP
plano. \emph{Sin este flag, la cookie puede ser interceptada en redes
WiFi públicas.}
\item[\texttt{HttpOnly}.] JavaScript del navegador no puede leer la
cookie a través de \texttt{document.cookie}. \emph{Defensa principal
contra XSS para robar la sesión.}
\item[\texttt{SameSite}.] Controla cuándo el navegador envía la cookie
en peticiones de otros sitios. Tres valores: \texttt{Strict} (nunca
desde otros sitios), \texttt{Lax} (solo en navegación de nivel
superior), \texttt{None} (siempre, requiere \texttt{Secure}).
\emph{Defensa principal contra CSRF.}
\item[\texttt{Max-Age} o \texttt{Expires}.] Tiempo de vida de la
cookie. Sin esto, es \emph{session cookie} (se borra al cerrar el
navegador).
\end{description}
\end{cajadefinicion}

\begin{cajaadvertencia}{Combinación correcta vs incorrecta}
\textbf{Correcto para sesiones de autenticación (Listado~\ref{lst:10.1}):}
\begin{lstlisting}[style=shell,numbers=none,caption={Ejemplo de Shell},label=lst:10.1]
Set-Cookie: SESSID=abc123; Path=/; Secure; HttpOnly; SameSite=Strict; Max-Age=1800
\end{lstlisting}
\textbf{Incorrecto (frecuentemente visto en código de aprendizaje) (Listado~\ref{lst:10.2}):}
\begin{lstlisting}[style=shell,numbers=none,caption={Ejemplo de Shell},label=lst:10.2]
Set-Cookie: SESSID=abc123
\end{lstlisting}
La segunda forma expone la cookie a HTTP plano (puede leerse en redes
inseguras), permite que JavaScript la lea (vulnerabilidad XSS), y se
envía en peticiones cross-site (vulnerabilidad CSRF).
\end{cajaadvertencia}

\section{El objeto Request: lectura segura de la petición}

El objeto \texttt{Request} encapsula toda la información que el
cliente envía al servidor. Su correcta lectura es la primera línea de
defensa de la aplicación: nunca se debe confiar en datos del cliente
sin filtrarlos.

\subsection{Componentes del objeto Request}

El objeto \emph{Request} encapsula todo lo que el cliente
envía: URL, método, cabeceras, cuerpo, cookies. Conocer sus
componentes es prerrequisito para programar cualquier endpoint
profesional.

\textbf{Ejemplo corto comentado (Listado~\ref{lst:10.3}):}
\begin{lstlisting}[language=PHP,numbers=none,caption={Ejemplo de PHP},label=lst:10.3]
<?php
// Componentes principales del Request HTTP en PHP
$_SERVER['REQUEST_METHOD'];   // 'GET', 'POST', 'PUT', 'DELETE'
$_SERVER['REQUEST_URI'];      // '/api/v1/usuarios?id=42'
$_SERVER['HTTP_HOST'];        // 'tienda.uteq.edu.ec'
$_SERVER['HTTP_USER_AGENT'];  // navegador del cliente
$_GET['id'];                  // parametros de la URL
$_POST['email'];              // datos del cuerpo (form-urlencoded)
$_COOKIE['SESSID'];           // cookies enviadas por el cliente
file_get_contents('php://input'); // cuerpo crudo (JSON, XML)
?>
\end{lstlisting}


\begin{itemize}
\item \textbf{Método HTTP}: \texttt{GET}, \texttt{POST}, \texttt{PUT},
\texttt{DELETE}, \texttt{PATCH}.
\item \textbf{URL completa}: ruta, query string, fragmento.
\item \textbf{Encabezados}: \texttt{Content-Type},
\texttt{Authorization}, \texttt{User-Agent}, \texttt{Accept-Language}.
\item \textbf{Cookies}: enviadas por el navegador.
\item \textbf{Cuerpo} (en POST/PUT/PATCH): JSON, form-urlencoded,
multipart, raw.
\item \textbf{Metadatos del cliente}: IP origen, protocolo (HTTP/HTTPS).
\end{itemize}

\begin{cajaejemplo}{Ejemplo 10.1 --- Lectura segura del Request en PHP 8.3}

\textbf{Enunciado.} Mostrar la lectura segura del objeto Request en PHP 8.3 obteniendo método, URL, cabeceras, query string y cuerpo, con validación básica de cada campo.

\begin{lstlisting}[style=php,caption={request-helpers.php},label=lst:10.4]
<?php
declare(strict_types=1);

/**
 * Obtiene la IP real del cliente, considerando proxy reverso.
 * Verifica las cabeceras X-Forwarded-For y X-Real-IP que un proxy
 * Nginx o un balanceador de carga (HAProxy, AWS ELB) inyecta.
 */
function obtenerIpCliente(): string {
    $cabeceras = [
        'HTTP_X_FORWARDED_FOR',  // proxy estandar
        'HTTP_X_REAL_IP',        // Nginx
        'HTTP_CF_CONNECTING_IP', // Cloudflare
        'REMOTE_ADDR'            // conexion directa
    ];
    foreach ($cabeceras as $cabecera) {                     // itera sobre colecion
        if (!empty($_SERVER[$cabecera])) {                  // condicional
            // X-Forwarded-For puede tener varias IPs separadas por coma
            $ip = explode(',', $_SERVER[$cabecera])[0];
            $ip = trim($ip);
            if (filter_var($ip, FILTER_VALIDATE_IP,         // valida/sanea entrada con filtro
                FILTER_FLAG_NO_PRIV_RANGE | FILTER_FLAG_NO_RES_RANGE)) {
                return $ip;                                 // retorna valor
            }
        }
    }
    return '0.0.0.0';                                       // retorna valor
}

/**
 * Lectura segura de un parametro GET.
 * NUNCA acceder directamente a $_GET['param'] sin filtrar.
 */
function leerEntero(string $nombre, int $minimo = 1,
                    int $maximo = PHP_INT_MAX, int $defecto = 1): int {
    $val = filter_input(INPUT_GET, $nombre, FILTER_VALIDATE_INT,
        ['options' => [
            'default'   => $defecto,
            'min_range' => $minimo,
            'max_range' => $maximo
        ]]);
    return $val ?? $defecto;                                // retorna valor
}

function leerCadena(string $nombre, int $maxLen = 100,
                    string $defecto = ''): string {
    $raw = $_REQUEST[$nombre] ?? $defecto;
    $val = trim((string)$raw);
    // condicional
    // condicional
    if (strlen($val) > $maxLen) $val = substr($val, 0, $maxLen);
    return $val;                                            // retorna valor
}

function exigirMetodo(string $metodo): void {
    // condicional
    // condicional
    if ($_SERVER['REQUEST_METHOD'] !== strtoupper($metodo)) {
        http_response_code(405);
        header('Allow: ' . strtoupper($metodo));            // envia cabecera HTTP al cliente
        exit("Metodo no permitido");                        // termina ejecucion
    }
}

// Uso conjunto
exigirMetodo('GET');
$pagina  = leerEntero('page', 1, 1000, 1);
$busca   = leerCadena('q', 80);
$ip      = obtenerIpCliente();

// Registrar acceso para auditoria
error_log("[" . date('c') . "] $ip - GET ?page=$pagina&q=$busca");
\end{lstlisting}

\textbf{Por qué este código es importante:}
\begin{itemize}
\item Centraliza la lectura segura en funciones helper reutilizables.
\item Detecta proxy reverso correctamente (un error común en
producción).
\item Filtra IPs privadas y reservadas con flags de
\texttt{filter\_var}.
\item Aplica límites estrictos a los enteros (rangos).
\item Trunca cadenas demasiado largas para prevenir ataques DoS por
memoria.
\item Verifica el método HTTP esperado.
\end{itemize}
\end{cajaejemplo}

\subsection{Lectura del Request en Java Servlets / Spring Boot}

En Java/Spring Boot, el equivalente al \texttt{\$\_SERVER}
de PHP es la interfaz \texttt{HttpServletRequest} o las anotaciones
de Spring (\texttt{@RequestParam}, \texttt{@RequestHeader},
\texttt{@RequestBody}). El ejemplo siguiente las recorre.


\begin{cajaejemplo}{Ejemplo 10.2 --- Acceso al Request en Spring Boot 3.4}

\textbf{Enunciado.} Mostrar el acceso al objeto Request en Spring Boot 3.4 usando las anotaciones idiomáticas: \texttt{@RequestParam}, \texttt{@RequestHeader}, \texttt{@RequestBody}.

\begin{lstlisting}[style=java,caption={RequestInspectorController.java},label=lst:10.5]
package ec.edu.uteq.appweb.controller;                      // paquete del archivo

import jakarta.servlet.http.HttpServletRequest;             // importacion de clase externa
import org.springframework.web.bind.annotation.*;           // importacion de clase externa
import java.util.*;                                         // importacion de clase externa

@RestController                                             // expone endpoints REST con JSON
@RequestMapping("/api/inspeccion")                          // ruta base del controlador
public class RequestInspectorController {                   // declaracion de clase publica

    @GetMapping                                             // endpoint HTTP GET
    public Map<String, Object> inspeccionar(                // metodo publico
            HttpServletRequest req,
            @RequestParam(defaultValue = "") String q,      // extrae parametro de query string
            @RequestParam(defaultValue = "1") int page,     // extrae parametro de query string
            @RequestHeader(value = "User-Agent",            // extrae cabecera HTTP
                           defaultValue = "?") String userAgent) {

        // Headers selectivos
        var headers = new LinkedHashMap<String, String>();
        var nombres = req.getHeaderNames();
        while (nombres.hasMoreElements()) {                 // bucle while
            var n = nombres.nextElement();
            headers.put(n, req.getHeader(n));
        }

        return Map.of(                                      // devuelve el resultado al llamador
            "metodo",      req.getMethod(),
            "url",         req.getRequestURL().toString(),
            "remoteAddr",  obtenerIp(req),
            "userAgent",   userAgent,
            "queryParams", Map.of("q", q, "page", page),
            "headersSize", headers.size()
        );
    }

    private String obtenerIp(HttpServletRequest r) {        // campo privado
        var fwd = r.getHeader("X-Forwarded-For");
        if (fwd != null && !fwd.isBlank()) {                // condicional
            return fwd.split(",")[0].trim();                // devuelve el resultado al llamador
        }
        return r.getRemoteAddr();                           // devuelve el resultado al llamador
    }
}
\end{lstlisting}
\end{cajaejemplo}

\subsection{Lectura del Request en ASP.NET Core 8}

ASP.NET Core usa \texttt{HttpContext.Request} para acceder a
los datos de la petición, con una API moderna basada en propiedades
fuertemente tipadas. El ejemplo siguiente la ilustra.


\begin{cajaejemplo}{Ejemplo 10.3 --- Endpoint que inspecciona el Request en ASP.NET Core}

\textbf{Enunciado.} Implementar un endpoint en ASP.NET Core que inspeccione el Request: método, ruta, cabeceras, query string. Se usará el objeto \texttt{HttpContext.Request}.

\begin{lstlisting}[style=cs,caption={Controllers/InspeccionController.cs},label=lst:10.6]
using Microsoft.AspNetCore.Mvc;                             // importa namespace

namespace AppwebApi.Controllers;                            // namespace del archivo

[ApiController]                                             // controlador de Web API (validacion automatica)
[Route("api/[controller]")]                                 // ruta base del controlador
public class InspeccionController : ControllerBase          // declaracion de clase publica
{
    [HttpGet]                                               // endpoint HTTP GET
    public IActionResult Inspeccionar(                      // metodo o miembro publico
        [FromQuery] string q = "",                          // mapea query string
        [FromQuery] int page = 1)                           // mapea query string
    {
        // IP real considerando X-Forwarded-For
        var ip = HttpContext.Request.Headers["X-Forwarded-For"]
                    .FirstOrDefault()?.Split(',')[0].Trim()
                 ?? HttpContext.Connection.RemoteIpAddress?.ToString()
                 ?? "0.0.0.0";

        return Ok(new {                                     // HTTP 200 con cuerpo
            metodo      = Request.Method,
            url         = $"{Request.Scheme}://{Request.Host}{Request.Path}",
            remoteAddr  = ip,
            userAgent   = Request.Headers.UserAgent.ToString(),
            queryParams = new { q, page },
            cookieCount = Request.Cookies.Count
        });
    }
}
\end{lstlisting}
\end{cajaejemplo}

\section{El objeto Response: control total de la respuesta HTTP}

El objeto \texttt{Response} permite controlar tres aspectos:
código de estado, cabeceras y cuerpo.

\subsection{Códigos de estado HTTP --- la guía exhaustiva}

Los códigos de estado HTTP son la \emph{lengua común} entre
cliente y servidor para indicar el resultado de una petición.
Confundirlos (devolver 200 cuando hubo un error, por ejemplo) es uno
de los defectos más comunes en APIs amateur.


\begin{table}[htbp]
\centering
\caption{Códigos de estado HTTP de uso obligado en aplicaciones web profesionales.}
\footnotesize
\begin{tabularx}{\textwidth}{llX}
\toprule
\textbf{Código} & \textbf{Nombre} & \textbf{Cuándo usarlo (con ejemplos)}\\
\midrule
\multicolumn{3}{l}{\textbf{2xx --- Éxito}}\\
200 & OK                  & Operación exitosa con cuerpo en respuesta. GET, PUT exitosos.\\
201 & Created             & Recurso creado. POST exitoso. Incluir \texttt{Location}.\\
204 & No Content          & Operación exitosa sin cuerpo. DELETE típicamente.\\
\midrule
\multicolumn{3}{l}{\textbf{3xx --- Redirección}}\\
301 & Moved Permanently   & Cambio definitivo de URL. SEO sigue al nuevo URL.\\
302 & Found               & Redirección temporal. POST $\to$ GET (patrón PRG).\\
304 & Not Modified        & Caché válida (responde a \texttt{If-None-Match}).\\
\midrule
\multicolumn{3}{l}{\textbf{4xx --- Error del cliente}}\\
400 & Bad Request         & JSON mal formado, parámetro faltante crítico.\\
401 & Unauthorized        & No autenticado. Falta JWT o cookie de sesión.\\
403 & Forbidden           & Autenticado pero sin permiso. Rol insuficiente.\\
404 & Not Found           & Recurso inexistente. ID inválido.\\
405 & Method Not Allowed  & Endpoint existe pero no acepta ese verbo. Incluir \texttt{Allow}.\\
409 & Conflict            & Email duplicado, optimistic locking failure.\\
422 & Unprocessable Entity & Validación de negocio falla (campos inválidos).\\
429 & Too Many Requests   & Rate limit excedido. Incluir \texttt{Retry-After}.\\
\midrule
\multicolumn{3}{l}{\textbf{5xx --- Error del servidor}}\\
500 & Internal Server Error & Excepción no controlada. NUNCA mostrar stack trace al cliente.\\
502 & Bad Gateway         & Proxy recibió respuesta inválida del upstream.\\
503 & Service Unavailable & Mantenimiento programado. Incluir \texttt{Retry-After}.\\
504 & Gateway Timeout     & Upstream tardó demasiado.\\
\bottomrule
\end{tabularx}
\end{table}

\subsection{Errores estandarizados con RFC 7807 ProblemDetails}

El RFC 7807 \cite{rfc7807} estandariza un formato JSON para errores
HTTP. Es el estándar de la industria en 2026 y los frameworks modernos
(ASP.NET Core, Spring Boot) lo soportan nativamente (Listado~\ref{lst:10.7}):

\begin{lstlisting}[style=js,numbers=none,caption={Respuesta de error según RFC 7807},label=lst:10.7]
HTTP/1.1 422 Unprocessable Entity
Content-Type: application/problem+json

{
  "type":   "https://uteq.edu.ec/errors/validacion",
  "title":  "Datos invalidos",
  "status": 422,
  "detail": "El campo email no cumple el formato esperado",
  "instance": "/api/usuarios",
  "errores": {
    "email":  ["debe tener formato valido"],
    "edad":   ["debe ser mayor a 0"]
  }
}
\end{lstlisting}

\subsection{Cabeceras de seguridad obligatorias en producción}

En producción, además de las cabeceras de seguridad básicas
vistas en la semana 9, conviene configurar otras que protegen contra
ataques específicos. La tabla y los ejemplos siguientes las catalogan.


\begin{cajabuenas}{Las siete cabeceras de seguridad imprescindibles}
\begin{enumerate}
\item \textbf{Strict-Transport-Security}:
\texttt{max-age=31536000; includeSubDomains; preload}\\
Fuerza HTTPS por un año en todo el dominio.

\item \textbf{Content-Security-Policy}:
\texttt{default-src 'self'; script-src 'self' 'nonce-XYZ'}\\
Mitiga XSS limitando fuentes de scripts y estilos.

\item \textbf{X-Frame-Options}: \texttt{DENY}\\
No permite que el sitio se cargue en \texttt{<iframe>}: anti-clickjacking.

\item \textbf{X-Content-Type-Options}: \texttt{nosniff}\\
Impide que el navegador deduzca el tipo MIME (anti-MIME-sniffing).

\item \textbf{Referrer-Policy}:
\texttt{strict-origin-when-cross-origin}\\
Equilibrio entre privacidad y funcionalidad.

\item \textbf{Permissions-Policy}:
\texttt{geolocation=(), microphone=(), camera=()}\\
Restringe APIs sensibles del navegador.

\item \textbf{Cross-Origin-Opener-Policy}: \texttt{same-origin}\\
Aísla el contexto de navegación de otros orígenes (mitigación
Spectre).
\end{enumerate}
\end{cajabuenas}

\begin{cajaejemplo}{Ejemplo 10.4 --- Cabeceras de seguridad en las tres plataformas}

\textbf{Enunciado.} Aplicar las cabeceras de seguridad obligatorias en producción (\texttt{Content-Security-Policy}, \texttt{X-Frame-Options}, \texttt{Strict-Transport-Security}, etc.) en las tres plataformas del curso.

\textbf{En PHP (Listado~\ref{lst:10.8}):}
\begin{lstlisting}[style=php,numbers=none,caption={Ejemplo de PHP},label=lst:10.8]
<?php
// envia cabecera HTTP al cliente
// envia cabecera HTTP al cliente
header('Strict-Transport-Security: max-age=31536000; includeSubDomains');
header("Content-Security-Policy: default-src 'self'");      // envia cabecera HTTP al cliente
header('X-Frame-Options: DENY');                            // envia cabecera HTTP al cliente
header('X-Content-Type-Options: nosniff');                  // envia cabecera HTTP al cliente
header('Referrer-Policy: strict-origin-when-cross-origin'); // envia cabecera HTTP al cliente
// envia cabecera HTTP al cliente
// envia cabecera HTTP al cliente
header('Permissions-Policy: geolocation=(), microphone=(), camera=()');
\end{lstlisting}

\textbf{En Spring Boot 3 (Listado~\ref{lst:10.9}):}
\begin{lstlisting}[style=java,numbers=none,caption={Ejemplo de Java},label=lst:10.9]
@Configuration                                              // clase de configuracion Spring
public class SeguridadConfig {                              // declaracion de clase publica
    @Bean
    // metodo publico
    // metodo publico
    public SecurityFilterChain seguridad(HttpSecurity http) throws Exception {
        http.headers(h -> h
            .httpStrictTransportSecurity(hsts -> hsts
                .maxAgeInSeconds(31536000).includeSubDomains(true))
            .frameOptions(f -> f.deny())
            .contentSecurityPolicy(csp -> csp
                .policyDirectives("default-src 'self'"))
            .referrerPolicy(rp -> rp.policy(
                ReferrerPolicy.STRICT_ORIGIN_WHEN_CROSS_ORIGIN))
        );
        return http.build();                                // devuelve el resultado al llamador
    }
}
\end{lstlisting}

\textbf{En ASP.NET Core 8 (Listado~\ref{lst:10.10}):}
\begin{lstlisting}[style=cs,numbers=none,caption={Ejemplo de C\#},label=lst:10.10]
app.Use(async (ctx, next) => {
    ctx.Response.Headers["Strict-Transport-Security"]
        = "max-age=31536000; includeSubDomains";
    ctx.Response.Headers["X-Frame-Options"] = "DENY";
    ctx.Response.Headers["X-Content-Type-Options"] = "nosniff";
    ctx.Response.Headers["Referrer-Policy"]
        = "strict-origin-when-cross-origin";
    ctx.Response.Headers["Content-Security-Policy"]
        = "default-src 'self'";
    await next();
});
\end{lstlisting}
\end{cajaejemplo}

\section{El objeto Session: estado entre peticiones}

La \emph{sesión} es el mecanismo más antiguo y aún
ampliamente usado para mantener estado entre peticiones: el servidor
guarda los datos del usuario en memoria (o en Redis) y envía al
cliente solo un identificador.


\subsection{Mecánica detallada de las sesiones de servidor}

Comprender la mecánica detallada de las sesiones de servidor
es clave porque la mayoría de fallos de seguridad ocurren por
desconocimiento de cómo se genera, transporta y persiste el
\emph{Session ID}.


\begin{figure}[htbp]
\centering
\begin{tikzpicture}[
  node distance=1.2cm and 1.5cm,
  every node/.style={font=\scriptsize,align=center},
  caja/.style={rectangle,draw=uteqverde,thick,minimum width=2cm,
               minimum height=0.8cm,fill=uteqverde!10,rounded corners=2pt},
  flecha/.style={->,>=stealth,thick}]

  \node[caja] (cli)  {Navegador};
  \node[caja,right=4cm of cli] (srv) {Servidor};
  \node[cylinder,draw=uteqverde,fill=uteqverde!15,thick,
        shape border rotate=90,aspect=0.3,minimum width=1.5cm,
        minimum height=1.7cm,below=0.5cm of srv] (sto)
        {Session\\Store};

  \draw[flecha] ([yshift=0.4cm]cli.east) --
    node[above]{1. GET /login (sin cookie)} ([yshift=0.4cm]srv.west);
  \draw[flecha] ([yshift=-0.4cm]srv.west) --
    node[below,align=center]{\scriptsize 2. Set-Cookie: SESSID=abc...\\\scriptsize Secure; HttpOnly; SameSite}
    ([yshift=-0.4cm]cli.east);
  \draw[flecha] (srv) -- node[right]{3. Crear} (sto);

\end{tikzpicture}
\caption{Establecimiento de una sesión: el servidor genera un ID
aleatorio, lo asocia con datos en su \emph{Session Store} (memoria,
Redis o BD) y lo envía al cliente como cookie con flags de seguridad.}
\end{figure}

\subsection{Implementación comparada en las tres plataformas}

Cada plataforma implementa sesiones con su propia API, pero
los conceptos son los mismos. La tabla siguiente las contrasta para
PHP, Spring Boot y ASP.NET Core.


\begin{cajacompara}{Comparativa de APIs de sesión}
\centering
\footnotesize
\begin{tabularx}{\textwidth}{lXXX}
\toprule
\textbf{Operación} & \textbf{PHP 8.3} & \textbf{ASP.NET Core 8} & \textbf{Spring Boot 3}\\
\midrule
Iniciar     & \texttt{session\_start()} & \texttt{app.UseSession()} & Automático (Servlet)\\
Guardar     & \texttt{\$\_SESSION['k']=\$v} & \texttt{ctx.Session.SetString('k',\$v)} & \texttt{session.setAttribute('k',v)}\\
Leer        & \texttt{\$\_SESSION['k']} & \texttt{ctx.Session.GetString('k')} & \texttt{session.getAttribute('k')}\\
Borrar      & \texttt{unset(\$\_SESSION['k'])} & \texttt{ctx.Session.Remove('k')} & \texttt{session.removeAttribute('k')}\\
Destruir    & \texttt{session\_destroy()} & \texttt{ctx.Session.Clear()} & \texttt{session.invalidate()}\\
Regenerar ID & \texttt{session\_regenerate\_id(true)} & \texttt{Refresh()} & \texttt{request.changeSessionId()}\\
Cookie defecto & \texttt{PHPSESSID} & \texttt{.AspNetCore.Session} & \texttt{JSESSIONID}\\
Almacenamiento & Archivos /tmp & Memoria/Redis/SQL & Memoria JVM/Redis\\
Timeout defecto & 24 min & 20 min & 30 min\\
\bottomrule
\end{tabularx}
\end{cajacompara}

\subsection{Sesiones en PHP 8.3 con configuración profesional}

PHP gestiona sesiones con la función nativa
\texttt{session\_start()}. El ejemplo siguiente la configura con todos
los flags de seguridad de producción y muestra los anti-patrones
comunes a evitar.


\begin{cajaejemplo}{Ejemplo 10.5 --- Sesión segura completa con login PHP}

\textbf{Enunciado.} Construir una sesión segura completa con login PHP: hash BCrypt, cookies con todos los flags de seguridad, regeneración de ID y logout robusto.

\begin{lstlisting}[style=php,caption={config/session.php (incluir antes de session\_start)},label=lst:10.11]
<?php
declare(strict_types=1);

/**
 * Configuracion de sesion para produccion.
 * DEBE llamarse ANTES de session_start().
 */
function configurarSesionSegura(): void {
    // Solo cookies HTTPS
    ini_set('session.cookie_secure',  '1');
    // JavaScript no puede leer la cookie
    ini_set('session.cookie_httponly', '1');
    // Anti-CSRF
    ini_set('session.cookie_samesite', 'Strict');
    // Forzar regeneracion de ID si llega uno desconocido
    ini_set('session.use_strict_mode', '1');
    // Timeout 30 min de inactividad
    ini_set('session.gc_maxlifetime',  '1800');
    // Cookie expira al cerrar navegador (=0) o tras 30 min (1800)
    session_set_cookie_params([                             // configura flags de la cookie de sesion
        'lifetime' => 1800,
        'path'     => '/',
        'secure'   => true,
        'httponly' => true,
        'samesite' => 'Strict',
    ]);
}
\end{lstlisting}

\begin{lstlisting}[style=php,caption={login.php con todas las prácticas},label=lst:10.12]
<?php
declare(strict_types=1);
require_once __DIR__ . '/config/session.php';               // incluye otro archivo PHP (fatal si falla)
require_once __DIR__ . '/config/db.php';                    // incluye otro archivo PHP (fatal si falla)

configurarSesionSegura();
session_start();                                            // arranca o reanuda la sesion

// Si ya esta autenticado, redirigir
if (!empty($_SESSION['user_id'])) {                         // condicional
    header('Location: /dashboard.php');                     // envia cabecera HTTP al cliente
    exit;
}

if ($_SERVER['REQUEST_METHOD'] === 'POST') {                // condicional
    // Validar token CSRF (anti-CSRF en formularios)
    $tokenEnviado = $_POST['_csrf'] ?? '';
    $tokenSesion  = $_SESSION['_csrf'] ?? '';
    if (!hash_equals($tokenSesion, $tokenEnviado)) {        // condicional
        http_response_code(419);
        exit('Token CSRF invalido');                        // termina ejecucion
    }

    $email = filter_input(INPUT_POST, 'email',
                          FILTER_VALIDATE_EMAIL);
    $pass  = $_POST['password'] ?? '';

    if (!$email || strlen($pass) < 8) {                     // condicional
        http_response_code(422);
        echo 'Datos invalidos';                             // imprime al output del servidor
        exit;
    }

    // Rate limiting en memoria (en produccion: Redis)
    $intentosKey = 'intentos:' . $_SERVER['REMOTE_ADDR'];
    if (($_SESSION[$intentosKey] ?? 0) >= 5) {              // condicional
        http_response_code(429);
        header('Retry-After: 900');  // 15 min
        exit('Demasiados intentos. Vuelva en 15 minutos.'); // termina ejecucion
    }

    $pdo = obtenerPDO();
    $stmt = $pdo->prepare(                                  // prepara consulta SQL (defensa anti-SQLi)
        // hashea contrasena con BCrypt
        // hashea contrasena con BCrypt
        'SELECT id, password_hash, rol, activo, intentos_fallidos
         FROM usuarios WHERE email = :em LIMIT 1');
    $stmt->execute([':em' => $email]);                      // ejecuta la consulta con parametros enlazados
    $u = $stmt->fetch(PDO::FETCH_ASSOC);                    // recupera la siguiente fila como array

    // password_verify usa tiempo constante (resistente a timing attacks)
    // hashea contrasena con BCrypt
    // hashea contrasena con BCrypt
    if (!$u || !$u['activo'] || !password_verify($pass, $u['password_hash'])) {
        $_SESSION[$intentosKey] = ($_SESSION[$intentosKey] ?? 0) + 1;
        http_response_code(401);
        echo 'Credenciales invalidas';                      // imprime al output del servidor
        exit;
    }

    // CRITICAL: regenerar el ID para prevenir session fixation
    session_regenerate_id(true);                            // cambia el ID para evitar session fixation

    $_SESSION['user_id']    = (int)$u['id'];
    $_SESSION['user_rol']   = $u['rol'];
    $_SESSION['login_time'] = time();
    $_SESSION['ip']         = $_SERVER['REMOTE_ADDR'];
    $_SESSION['ua']         = $_SERVER['HTTP_USER_AGENT'] ?? '';
    unset($_SESSION[$intentosKey]);

    // Auditoria
    error_log("[LOGIN] user_id={$u['id']} ip={$_SERVER['REMOTE_ADDR']}");

    header('Location: /dashboard.php');                     // envia cabecera HTTP al cliente
    exit;
}

// Generar token CSRF para mostrar el formulario
if (empty($_SESSION['_csrf'])) {                            // condicional
    $_SESSION['_csrf'] = bin2hex(random_bytes(32));
}
$csrf = $_SESSION['_csrf'];
?>
<!DOCTYPE html>
<html lang="es-EC">
<head><meta charset="UTF-8"><title>Login UTEQ</title></head>
<body>
<form method="POST">
  // escapa caracteres HTML (defensa anti-XSS)
  // escapa caracteres HTML (defensa anti-XSS)
  <input type="hidden" name="_csrf" value="<?= htmlspecialchars($csrf) ?>">
  <label>Email: <input type="email" name="email" required></label>
  <label>Password: <input type="password" name="password" required minlength="8"></label>
  <button type="submit">Ingresar</button>
</form>
</body></html>
\end{lstlisting}

\textbf{Defensas implementadas en este código (todas combinadas):}
\begin{enumerate}
\item Cookies con \texttt{Secure}, \texttt{HttpOnly},
\texttt{SameSite=Strict}.
\item Token CSRF generado por sesión y validado con
\texttt{hash\_equals} (tiempo constante).
\item Validación email con \texttt{FILTER\_VALIDATE\_EMAIL}.
\item Rate limiting (5 intentos por IP).
\item \texttt{password\_verify} (BCrypt, resistente a timing attacks).
\item \texttt{session\_regenerate\_id(true)}: previene session
fixation.
\item Auditoría con \texttt{error\_log}.
\item \texttt{htmlspecialchars} al renderizar (anti-XSS).
\item Verificación de cuenta activa.
\end{enumerate}
\end{cajaejemplo}

\section{Almacenamiento del cliente: localStorage, sessionStorage e IndexedDB}

Además del almacenamiento del servidor (cookies de sesión,
JWT), el navegador ofrece tres APIs nativas para guardar datos del
lado cliente: \texttt{localStorage}, \texttt{sessionStorage} e
\texttt{IndexedDB}. Conocer cuándo usar cada una es esencial para
SPAs modernas.


\subsection{Comparativa de las APIs de almacenamiento del cliente}

La tabla siguiente contrasta las tres APIs de almacenamiento
cliente por capacidad, persistencia, tipo de datos y casos de uso
recomendados.


\begin{table}[htbp]
\centering
\caption{Almacenamiento del cliente en navegadores modernos.}
\footnotesize
\begin{tabularx}{\textwidth}{lXll}
\toprule
\textbf{Mecanismo} & \textbf{Características} & \textbf{Capacidad} & \textbf{API}\\
\midrule
Cookies & Enviadas con cada petición; \emph{HttpOnly} oculta de JS & 4 KB & \texttt{document.cookie}\\
\texttt{localStorage} & Persistente, mismo origen, síncrono, solo strings & 5--10 MB & Síncrona\\
\texttt{sessionStorage} & Solo durante la pestaña, mismo origen & 5--10 MB & Síncrona\\
IndexedDB & Persistente, transaccional, asíncrono, objetos & GB & Asíncrona\\
Cache API & Almacena objetos \texttt{Response} (Service Workers) & GB & Asíncrona\\
\bottomrule
\end{tabularx}
\end{table}

\begin{cajaadvertencia}{La regla de oro --- ¿dónde guardar el JWT?}
\textbf{Nunca} guarde el JWT en \texttt{localStorage} si su aplicación
tiene cualquier riesgo de XSS (lo cual es prácticamente toda
aplicación). El JavaScript puede leer \texttt{localStorage}; un solo
script malicioso inyectado roba todos los tokens.

\textbf{Solución correcta:} cookie \texttt{HttpOnly} + \texttt{Secure}
+ \texttt{SameSite=Strict}. JavaScript no puede leer cookies
\texttt{HttpOnly}; el navegador las envía automáticamente al servidor
en cada petición al mismo origen.
\end{cajaadvertencia}

\subsection{Ejemplos prácticos de almacenamiento}

Los siguientes ejemplos breves muestran las tres APIs en
acción para el caso más frecuente: persistir preferencias de UI.

\textbf{Ejemplo corto comentado (localStorage) (Listado~\ref{lst:10.13}):}
\begin{lstlisting}[style=js,numbers=none,caption={Ejemplo de JavaScript},label=lst:10.13]
// localStorage persiste INDEFINIDAMENTE (hasta que el usuario lo borre)
// Solo guarda STRINGS, asi que objetos hay que serializarlos con JSON
localStorage.setItem('tema', 'oscuro');           // escribe
const tema = localStorage.getItem('tema');         // lee
localStorage.removeItem('tema');                   // borra una clave
localStorage.clear();                              // borra todo

// Guardar un objeto: serializar con JSON.stringify
const prefs = { tema: 'oscuro', idioma: 'es' };             // constante (no reasignable)
// serializa objeto a string JSON
// serializa objeto a string JSON
localStorage.setItem('preferencias', JSON.stringify(prefs));
// Recuperar y parsear
// constante (no reasignable)
// constante (no reasignable)
const guardado = JSON.parse(localStorage.getItem('preferencias') || '{}');
\end{lstlisting}


\begin{cajaejemplo}{Ejemplo 10.6 --- Persistencia de preferencias UI con localStorage}

\textbf{Enunciado.} Persistir preferencias de UI (tema oscuro, idioma) en \texttt{localStorage} desde JavaScript, mostrando cómo serializar objetos con JSON y manejar errores de cuota.

\begin{lstlisting}[style=js,caption={preferencias.js},label=lst:10.14]
"use strict";

const PREF_KEY = "uteq:preferencias:v1";                    // constante (no reasignable)

const Preferencias = {                                      // constante (no reasignable)
  cargar() {
    try {
      const raw = localStorage.getItem(PREF_KEY);           // constante (no reasignable)
      return raw ? JSON.parse(raw) : this.defectos();       // devuelve valor de la funcion
    } catch (e) {
      console.warn("Error leyendo preferencias", e);
      return this.defectos();                               // devuelve valor de la funcion
    }
  },
  guardar(p) {
    try {
      localStorage.setItem(PREF_KEY, JSON.stringify(p));    // serializa objeto a string JSON
    } catch (e) {
      console.error("Error guardando preferencias", e);
    }
  },
  defectos() {
    return { tema: "auto", idioma: "es-EC",                 // devuelve valor de la funcion
             fontSize: 1.0, animaciones: true };
  },
  aplicar(p) {
    document.documentElement.dataset.tema = p.tema;
    document.documentElement.lang = p.idioma;
    document.documentElement.style.fontSize = `${p.fontSize}rem`;
    if (!p.animaciones) {                                   // condicional
      document.documentElement.classList.add("sin-animaciones");
    }
  }
};

// Al cargar
document.addEventListener("DOMContentLoaded", () => {       // registra manejador de evento
  const p = Preferencias.cargar();                          // constante (no reasignable)
  Preferencias.aplicar(p);

  // Boton de cambio de tema
  // busca elemento por id
  // busca elemento por id
  document.getElementById("toggle-tema")?.addEventListener("click", () => {
    p.tema = p.tema === "claro" ? "oscuro" : "claro";
    Preferencias.guardar(p);
    Preferencias.aplicar(p);
  });
});
\end{lstlisting}
\end{cajaejemplo}

\section{Práctica integrada: carrito de compras con sesiones}

La práctica de la semana es construir un carrito de compras
completo con sesiones de servidor, integrando todos los conceptos:
Request, Response, Session, Cookie segura y validación.


\begin{cajaejercicio}{Ejercicio 10.1 --- Carrito de compras stateful en las tres plataformas (Sumativa GA-Práctica)}

\textbf{Objetivo:} implementar un carrito de compras con sesiones de
servidor en PHP 8.3, ASP.NET Core 8 y Spring Boot 3, evidenciando las
diferencias de API y cómo cada plataforma maneja seguridad de cookies.

\textbf{Especificación funcional:}
\begin{itemize}
\item Página de listado con productos hardcodeados (nombre, precio).
\item Botón «Agregar» para cada producto que lo añade al carrito.
\item Página del carrito que muestra los productos agregados y total.
\item Botón «Eliminar» por producto en el carrito.
\item Botón «Vaciar carrito».
\item La cookie de sesión debe tener \texttt{HttpOnly},
\texttt{Secure} y \texttt{SameSite=Strict}.
\end{itemize}

\subsubsection*{IDE recomendado}
\textbf{IntelliJ IDEA Ultimate} con plugins:
\begin{itemize}
\item \emph{PHP} (JetBrains, instalado por defecto en Ultimate).
\item \emph{Spring Boot} (incluido en Ultimate).
\item \emph{.NET Core} (Marketplace) o usar JetBrains Rider para C\#.
\end{itemize}

\subsubsection*{Plataforma 1 --- PHP 8.3 con XAMPP}

\textbf{Paso 1:} verificar XAMPP. Iniciar el panel
(\texttt{xampp-control}). Pulsar \texttt{Start} en Apache. Verificar
en \url{http://localhost} que aparece el dashboard XAMPP.

\textbf{Paso 2:} crear carpeta \texttt{C:\textbackslash{}xampp\textbackslash{}htdocs\textbackslash{}carrito-php}.
En IntelliJ: \texttt{File $\rightarrow$ Open}, seleccionar la carpeta.

\textbf{Paso 3:} crear \texttt{config.php} (Listado~\ref{lst:10.15}):

\begin{lstlisting}[style=php,caption={config.php},label=lst:10.15]
<?php
declare(strict_types=1);

ini_set('session.cookie_httponly', '1');
ini_set('session.cookie_secure',   '1');
ini_set('session.cookie_samesite', 'Strict');
ini_set('session.use_strict_mode', '1');
ini_set('session.gc_maxlifetime',  '1800');

session_name('CARRITO_PHP');
session_start();                                            // arranca o reanuda la sesion

// Inicializar carrito si no existe
if (!isset($_SESSION['carrito'])) {                         // condicional
    $_SESSION['carrito'] = [];
}

// Catalogo hardcodeado
const CATALOGO = [
    1 => ['id'=>1, 'nombre'=>'Cafe galletti 500g', 'precio'=>4.50],
    2 => ['id'=>2, 'nombre'=>'Cacao Pacari 70%',   'precio'=>6.00],
    3 => ['id'=>3, 'nombre'=>'Te verde organico',  'precio'=>3.20],
    4 => ['id'=>4, 'nombre'=>'Panela en bloque',   'precio'=>2.80],
];
\end{lstlisting}

\textbf{Paso 4:} crear \texttt{index.php} (Listado~\ref{lst:10.16}):

\begin{lstlisting}[style=php,caption={index.php},label=lst:10.16]
<?php require_once 'config.php'; ?>
<!DOCTYPE html>
<html lang="es-EC">
<head><meta charset="UTF-8"><title>Tienda UTEQ</title>
<style>
  body{font-family:system-ui;max-width:700px;margin:2rem auto;padding:0 1rem}
  h1{color:#006633} table{width:100%;border-collapse:collapse}
  td,th{padding:.5rem;text-align:left;border-bottom:1px solid #eee}
  .btn{padding:.4rem .8rem;background:#006633;color:white;
       border:0;border-radius:4px;cursor:pointer}
</style>
</head>
<body>
<h1>Productos UTEQ</h1>
<p><a href="carrito.php">Ver carrito (<?= count($_SESSION['carrito']) ?>)</a></p>
<table>
  <thead><tr><th>Producto</th><th>Precio</th><th></th></tr></thead>
  <tbody>
  <?php foreach (CATALOGO as $p): ?>
    <tr>
      <td><?= htmlspecialchars($p['nombre']) ?></td>        // escapa caracteres HTML (defensa anti-XSS)
      <td>$<?= number_format($p['precio'], 2) ?></td>
      <td>
        <form action="agregar.php" method="post" style="margin:0">
          <input type="hidden" name="id" value="<?= $p['id'] ?>">
          <button type="submit" class="btn">Agregar</button>
        </form>
      </td>
    </tr>
  <?php endforeach; ?>
  </tbody>
</table>
</body></html>
\end{lstlisting}

\textbf{Paso 5:} crear \texttt{agregar.php}, \texttt{eliminar.php},
\texttt{vaciar.php}, \texttt{carrito.php} (continuación obvia siguiendo
el patrón).

\textbf{Paso 6:} probar.
\begin{enumerate}
\item Abrir \url{http://localhost/carrito-php/index.php}.
\item Agregar 3 productos diferentes.
\item Ir a Carrito: aparecen los 3.
\item Cerrar la pestaña, abrir nueva: el carrito persiste.
\item Cerrar el navegador completamente, abrirlo de nuevo: el carrito
sigue ahí (porque \texttt{lifetime=1800}).
\end{enumerate}

\textbf{Paso 7:} verificar la cookie. F12 $\rightarrow$ Application
$\rightarrow$ Cookies $\rightarrow$ \url{http://localhost}. Debe
aparecer \texttt{CARRITO\_PHP} con: \texttt{HttpOnly = true},
\texttt{Secure = true}, \texttt{SameSite = Strict}.

\textbf{Resultado esperado:} carrito persistente entre páginas y entre
sesiones del navegador, con cookie hardenizada visible en DevTools.

\subsubsection*{Plataforma 2 --- ASP.NET Core 8}

\textbf{Paso 1:} verificar .NET 8. \texttt{dotnet --version} debe
devolver \texttt{8.0.x}.

\textbf{Paso 2:} crear el proyecto desde la terminal de IntelliJ
(\texttt{Alt+F12}) (Listado~\ref{lst:10.17}):
\begin{lstlisting}[style=shell,numbers=none,caption={Ejemplo de Shell},label=lst:10.17]
dotnet new mvc -n CarritoAsp -o CarritoAsp
cd CarritoAsp                                               # cambia directorio
\end{lstlisting}

\textbf{Paso 3:} configurar sesiones en \texttt{Program.cs} (Listado~\ref{lst:10.18}):

\begin{lstlisting}[style=cs,caption={Program.cs},label=lst:10.18]
using Microsoft.AspNetCore.Builder;                         // importa namespace
using Microsoft.Extensions.DependencyInjection;             // importa namespace

var builder = WebApplication.CreateBuilder(args);           // crea el builder de la app

builder.Services.AddControllersWithViews();                 // configuracion de servicios DI
builder.Services.AddDistributedMemoryCache();               // configuracion de servicios DI
builder.Services.AddSession(o => {                          // configuracion de servicios DI
    o.IdleTimeout = TimeSpan.FromMinutes(30);
    o.Cookie.Name        = "CARRITO_ASP";
    o.Cookie.HttpOnly    = true;
    o.Cookie.IsEssential = true;
    o.Cookie.SecurePolicy = CookieSecurePolicy.Always;
    o.Cookie.SameSite     = SameSiteMode.Strict;
});

var app = builder.Build();                                  // construye la app a partir del builder
app.UseHttpsRedirection();
app.UseStaticFiles();
app.UseRouting();                                           // activa el routing del pipeline
app.UseSession();
app.UseAuthorization();
app.MapControllerRoute(
    name: "default",
    pattern: "{controller=Tienda}/{action=Index}/{id?}");
app.Run();                                                  // arranca la aplicacion
\end{lstlisting}

\textbf{Paso 4:} crear \texttt{Models/Producto.cs} y
\texttt{Controllers/TiendaController.cs} (con acciones
\texttt{Index}, \texttt{Agregar}, \texttt{Carrito},
\texttt{Eliminar}, \texttt{Vaciar} usando
\texttt{HttpContext.Session.SetString} con JSON).

\textbf{Paso 5:} ejecutar con \texttt{dotnet run}. Abrir
\url{https://localhost:5001}.

\textbf{Paso 6:} verificar la cookie \texttt{CARRITO\_ASP} en
DevTools: debe tener \texttt{HttpOnly}, \texttt{Secure},
\texttt{SameSite=Strict}.

\subsubsection*{Plataforma 3 --- Spring Boot 3}

\textbf{Paso 1:} en IntelliJ Ultimate:
\texttt{File $\rightarrow$ New $\rightarrow$ Project $\rightarrow$ Spring Boot}.
Java 21, Spring Boot 3.4, dependencias \texttt{Spring Web} y
\texttt{Thymeleaf}.

\textbf{Paso 2:} configurar \texttt{src/main/resources/application.yml} (Listado~\ref{lst:10.19}):

\begin{lstlisting}[style=yaml,caption={application.yml},label=lst:10.19]
server:                                                     # configuracion del servidor embebido
  port: 8080                                                # puerto donde escucha la aplicacion
  servlet:
    session:
      timeout: 30m
      cookie:
        name:      CARRITO_SPRING
        http-only: true
        secure:    false   # true en HTTPS de produccion
        same-site: strict
\end{lstlisting}

\textbf{Paso 3:} crear \texttt{Producto.java} (record \texttt{Serializable}),
\texttt{TiendaController.java} con \texttt{HttpSession}.

\textbf{Paso 4:} ejecutar y probar en
\url{http://localhost:8080}.

\subsubsection*{Tabla comparativa de la práctica (entrega)}

Cada estudiante debe completar y entregar:

\begin{tabularx}{\textwidth}{lXXX}
\toprule
\textbf{Aspecto} & \textbf{PHP} & \textbf{ASP.NET} & \textbf{Spring}\\
\midrule
Líneas de código del proyecto & & & \\
Tiempo de arranque & & & \\
Cookie tiene HttpOnly & ¿sí/no? & ¿sí/no? & ¿sí/no?\\
Cookie tiene Secure & & & \\
Cookie tiene SameSite=Strict & & & \\
Carrito persiste entre pestañas & & & \\
\bottomrule
\end{tabularx}

\subsubsection*{Reflexión final}

Responder por escrito en máximo 200 palabras: ¿qué plataforma le
resultó más fácil? ¿Cuál considera más segura por defecto? ¿Cuál usaría
en un proyecto real y por qué?

\textbf{Esta es la \emph{Sumativa GA-Práctica en clase} de la semana
10.}
\end{cajaejercicio}

\section{Buenas prácticas profesionales en gestión de estado}

Las técnicas de gestión de estado tienen su propio conjunto
de reglas profesionales. La caja siguiente las reúne en 15 puntos.


\begin{cajabuenas}{Quince reglas profesionales 2026}
\begin{enumerate}
\item Cookies de sesión SIEMPRE con \texttt{Secure},
\texttt{HttpOnly}, \texttt{SameSite=Strict}.
\item Regenerar el ID de sesión al hacer login (anti-fixation).
\item Tiempo de inactividad razonable: 15--30 min para apps
sensibles, 60 min para uso general.
\item Logout debe invalidar el ID de sesión, no solo limpiar variables.
\item Para APIs REST que sirven SPAs: preferir JWT en cookie
\texttt{HttpOnly}.
\item NUNCA guardar JWT en \texttt{localStorage} en sitios con riesgo
de XSS.
\item Validar token CSRF en todo POST/PUT/PATCH/DELETE.
\item En clusters: usar Redis o BD compartida para sesiones, no
memoria local.
\item No almacenar información sensible en la cookie misma; solo el ID.
\item Auditar logins exitosos y fallidos con IP y User-Agent.
\item Implementar rate limiting (máx. 5--10 intentos por IP por minuto).
\item Cookies con \texttt{Domain} restrictivo (no
\texttt{example.com} si solo se necesita \texttt{app.example.com}).
\item Nombres de cookies que NO revelen la tecnología
(\texttt{SESS} en lugar de \texttt{PHPSESSID}).
\item En SPAs Angular/React: usar \texttt{credentials: 'include'} y
configurar CORS con \texttt{Access-Control-Allow-Credentials: true}.
\item Documentar la estrategia de gestión de estado en la ADR-001 del
proyecto.
\end{enumerate}
\end{cajabuenas}

\section{Ejercicio resuelto adicional con resultado visual}

Como cierre de la semana, el siguiente ejercicio adicional
está completamente resuelto e incluye captura del resultado visual.
Se complementa con un ejercicio propuesto para que el estudiante
practique.


\begin{cajaejemplo}{Ejercicio resuelto 10.1 --- Contador de visitas con sesión PHP}
\textbf{Enunciado.} Implementar una página PHP que muestre cuántas
veces el usuario ha visitado el sitio en la sesión actual, con
botones para incrementar manualmente y para reiniciar el contador.
La cookie de sesión debe llevar todos los flags de seguridad de
producción.

\textbf{Código completo (\texttt{contador.php}):}

\begin{lstlisting}[language=PHP,caption={Contador stateful con sesión segura},label=lst:10.20]
<?php
// Configuracion ANTES de session_start()
session_set_cookie_params([                                 // configura flags de la cookie de sesion
  'lifetime' => 1800,
  'path'     => '/',
  'domain'   => '',
  'secure'   => false,   // true en HTTPS de produccion
  'httponly' => true,
  'samesite' => 'Strict'
]);
session_name('VISITAS_SESS');
session_start();                                            // arranca o reanuda la sesion

// Inicializar contador
if (!isset($_SESSION['visitas'])) {                         // condicional
    $_SESSION['visitas'] = 1;
} else {
    // Solo incrementar en GET sin parametros (visita real)
    // condicional
    // condicional
    if ($_SERVER['REQUEST_METHOD'] === 'GET' && empty($_GET)) {
        $_SESSION['visitas']++;
    }
}

// Acciones
if (isset($_GET['accion'])) {                               // condicional
    // condicional
    // condicional
    if ($_GET['accion'] === 'incrementar') $_SESSION['visitas']++;
    // condicional
    // condicional
    if ($_GET['accion'] === 'reiniciar')   $_SESSION['visitas'] = 0;
    header('Location: contador.php');                       // envia cabecera HTTP al cliente
    exit;
}
?>
<!DOCTYPE html>
<html lang="es">
<head><meta charset="UTF-8"><title>Contador de visitas</title></head>
<body>
  <h1>Contador de visitas</h1>
  <p>Visitas en esta sesion: <strong><?= $_SESSION['visitas'] ?></strong></p>
  // escapa caracteres HTML (defensa anti-XSS)
  // escapa caracteres HTML (defensa anti-XSS)
  <p>Session ID: <code><?= htmlspecialchars(session_id()) ?></code></p>
  <a href="?accion=incrementar">Incrementar +1</a> |
  <a href="?accion=reiniciar">Reiniciar</a>
</body>
</html>
\end{lstlisting}

\textbf{Resultado visual en el navegador} (simulación de cómo se
verá la página en la pantalla del estudiante tras varias recargas
e incrementos manuales):

\begin{center}
\fbox{\begin{minipage}{0.85\textwidth}
\vspace{0.2cm}
\hspace{0.3cm}\textbf{\large Contador de visitas}\\[0.4cm]
\hspace{0.3cm}Visitas en esta sesi\'on: \textbf{7}\\[0.2cm]
\hspace{0.3cm}Session ID: \texttt{a8f3d9c2e1b4...}\\[0.3cm]
\hspace{0.3cm}\textcolor{blue}{\underline{Incrementar +1}} \textbar{}
\textcolor{blue}{\underline{Reiniciar}}\\
\vspace{0.2cm}
\end{minipage}}
\end{center}

\textbf{Verificación de los flags de seguridad} con DevTools del
navegador (pestaña \emph{Application $\to$ Cookies}):

\begin{center}
\fbox{\begin{minipage}{0.85\textwidth}
\vspace{0.15cm}\footnotesize
\hspace{0.3cm}\textbf{Cookie:} \texttt{VISITAS\_SESS}\\
\hspace{0.3cm}Value: \texttt{a8f3d9c2e1b4...}\\
\hspace{0.3cm}HttpOnly: $\checkmark$ \quad Secure: $\times$ \quad SameSite: \texttt{Strict}\\
\hspace{0.3cm}Expires: Session\\
\vspace{0.15cm}
\end{minipage}}
\end{center}

\textbf{Discusión.} El contador demuestra que la sesión persiste
entre recargas (el navegador reenvía la cookie automáticamente) y
que el servidor reconoce al mismo usuario. La cookie es invisible a
JavaScript (\texttt{HttpOnly}), bloqueando ataques XSS de robo de
sesión, y solo se envía desde el propio dominio
(\texttt{SameSite=Strict}), bloqueando CSRF. En producción HTTPS,
añadir \texttt{secure=true} bloquea su envío por HTTP plano.
\end{cajaejemplo}

\begin{cajaejercicio}{Ejercicio propuesto 10.2 --- Autenticación con login persistente «Recuérdame»}
\textbf{Enunciado.} Implementar en cualquiera de las tres plataformas
(PHP, Spring Boot o ASP.NET Core) un flujo de login con la opción
\emph{«Recuérdame»} que mantenga al usuario autenticado durante 30
días si la marca, o solo durante la sesión del navegador si no.

\textbf{Requisitos técnicos:}
\begin{enumerate}
\item Formulario de login con campos \texttt{usuario},
\texttt{contraseña} y checkbox \texttt{recordarme}.
\item Si \texttt{recordarme} está marcado: crear una cookie
adicional \texttt{REMEMBER\_TOKEN} con un valor aleatorio de 256 bits
(\texttt{random\_bytes(32)} en PHP, \texttt{SecureRandom} en Java),
\texttt{Max-Age=2592000} (30 días), todos los flags de seguridad.
\item Persistir el hash del token (\texttt{password\_hash()} o
\texttt{BCrypt}) en una tabla \texttt{remember\_tokens(user\_id,
token\_hash, expires\_at)}.
\item En cada nueva sesión, si el usuario no está autenticado pero
existe \texttt{REMEMBER\_TOKEN}, validar contra la tabla y
autenticar automáticamente.
\item Al hacer logout, eliminar la fila correspondiente y borrar la
cookie con \texttt{Max-Age=0}.
\item Implementar rotación del token tras cada uso (defensa contra
robo).
\end{enumerate}

\textbf{Criterios de evaluación:}
\begin{itemize}
\item Token nunca se almacena en plano (solo el hash).
\item Cookies con todos los flags (\texttt{HttpOnly},
\texttt{Secure}, \texttt{SameSite=Strict}, \texttt{Max-Age}).
\item Rotación correcta del token tras autenticación automática.
\item Caso de logout limpia BD y navegador.
\item Documentación con capturas de Postman o navegador mostrando
las cookies generadas.
\end{itemize}

\textbf{Lecturas recomendadas:} \cite{jones2015jwt} para
comparar con JWT; \cite{ortega2022securecoding} para defensas
contra robo de sesión.
\end{cajaejercicio}

% =====================================================================
%       SEMANA 11: ACCESO A DATOS DESDE APLICACIONES WEB
% =====================================================================
\chapter{Acceso a datos desde aplicaciones web: SQL parametrizado y ORM}
\label{cap:semana11}

\epigraph{«Toda aplicación que persiste datos hereda, tarde o
temprano, la disciplina o el caos de cómo accede a ellos. Elegir bien
es proteger el sistema durante años.»}{--- \textit{Adaptado de Martin Fowler}}

\section*{Resultado de aprendizaje de la semana}
Al concluir la semana 11, el estudiante diseña e implementa el acceso a
datos relacionales desde aplicaciones web utilizando SQL parametrizado
y mapeo objeto-relacional (ORM), aplicando los patrones Repository y
Unit of Work, prevenir totalmente la inyección SQL y comprende el
problema N+1 y su solución.

\section{Historia y panorama del acceso a datos}

El acceso a datos es probablemente la decisión arquitectónica
con mayor impacto en el rendimiento y la mantenibilidad de una
aplicación web. La elección entre SQL plano, JPA/Hibernate, Entity
Framework Core o Eloquent no es de detalle: condiciona todo el código
del backend.


\begin{cajahistoria}{Del SQL embebido a los ORM modernos --- 50 años de iteración}
La historia del acceso a datos desde aplicaciones web comienza con SQL
en su forma más cruda: cadenas de texto concatenadas dentro del código
de la aplicación. En 1995--2005 era habitual ver código PHP así:
\texttt{\$sql = \char34 SELECT * FROM users WHERE name = '\char34{} . \$\_POST['name'] . \char34 '\char34;}.
Este patrón --- inocente en apariencia --- destruyó la seguridad de
miles de aplicaciones durante una década por SQL Injection.

La primera respuesta fueron las \textbf{sentencias preparadas}
(\emph{prepared statements}), formalizadas en SQL-92 y ampliamente
disponibles desde principios de los 2000: el SQL se compila una vez
con marcadores (\texttt{?} o \texttt{:nombre}), y los valores viajan
separadamente sin interpretarse como código. Esto cerró la mayor
puerta de entrada a los ataques.

Una segunda evolución fueron los \textbf{Query Builders}: APIs
fluidas para construir SQL programáticamente sin escribirlo a mano
(\texttt{Doctrine DBAL} en PHP, \texttt{jOOQ} en Java). Permiten
componer consultas dinámicas sin riesgo de inyección.

El salto cualitativo vino con los \textbf{ORM} (Object-Relational
Mappers), que abstraen completamente la base de datos detrás de
objetos del lenguaje. Hibernate en Java (Gavin King, 2001), Active
Record en Ruby on Rails (David Heinemeier Hansson, 2004), Doctrine en
PHP (Roman Borschel, 2006), Entity Framework en .NET (Microsoft, 2008)
y Eloquent en Laravel (Taylor Otwell, 2011) representan distintas
filosofías sobre el mismo problema. Hoy un desarrollador web rara vez
escribe SQL a mano; pero comprender lo que hace el ORM por debajo es
indispensable.

La literatura técnica confirma que la elección del mecanismo de
acceso a datos es una decisión arquitectónica de primer orden.
\cite{fowler2002patterns} describe formalmente los patrones que hoy
dominan (Data Mapper, Active Record, Unit of Work, Identity Map);
\cite{ottinger2017hibernate} y \cite{bauer2015java} son las
referencias canónicas para Hibernate; \cite{smith2018entity} para
Entity Framework Core; \cite{otwell2024laravel} para Eloquent. El
trabajo empírico de \cite{salah2017orm} compara performance entre
los principales ORM y documenta que la diferencia de rendimiento
respecto a SQL plano oscila entre 10\% y 30\%, justificándose
ampliamente por la productividad ganada. Estudios más recientes
\cite{ottinger2017hibernate} en revistas indexadas confirman que el problema
N+1 sigue siendo la causa número uno de degradación de performance
en aplicaciones Java EE y Spring Boot en producción.
\end{cajahistoria}

\subsection{Las cuatro estrategias de acceso a datos}

En la industria conviven cuatro estrategias para acceder a
datos desde una aplicación web. La tabla siguiente las contrasta
con sus ventajas, desventajas y cuándo elegir cada una.


\begin{table}[htbp]
\centering
\caption{Estrategias de acceso a datos en aplicaciones web (2026).}
\footnotesize
\begin{tabularx}{\textwidth}{lXX}
\toprule
\textbf{Estrategia} & \textbf{Características} & \textbf{Cuándo usarla}\\
\midrule
Driver nativo + SQL & Máximo control y rendimiento; máxima responsabilidad. PDO (PHP), JDBC (Java), ADO.NET (C\#). & Reportes complejos, consultas analíticas, ETL.\\
Query Builder & SQL programático sin concatenación. Doctrine DBAL, jOOQ, knex.js. & Consultas dinámicas con muchos filtros opcionales.\\
ORM Active Record & La entidad \emph{es} su unidad de persistencia: \texttt{User::find(1)->save()}. Eloquent, Rails AR. & CRUDs estándar; prototipado rápido.\\
ORM Data Mapper & Entidades de dominio agnósticas; mapper externo. Hibernate, Doctrine ORM, EF Core. & Aplicaciones empresariales medianas/grandes con DDD.\\
\bottomrule
\end{tabularx}
\end{table}

\subsection{Tabla evolutiva de los ORM principales}

Los ORM dominantes han evolucionado durante dos décadas
hasta sus versiones actuales. La tabla siguiente recoge las
versiones de referencia 2026 con su licencia y características
notables.


\begin{table}[htbp]
\centering
\caption{Versiones actuales de ORM en 2026.}
\footnotesize
\begin{tabular}{llll}
\toprule
\textbf{ORM} & \textbf{Lenguaje} & \textbf{Versión 2026} & \textbf{Filosofía}\\
\midrule
Hibernate / JPA      & Java       & Hibernate 6.6 / JPA 3.2 & Data Mapper\\
Spring Data JPA      & Java       & 3.4                     & Repository sobre JPA\\
Eloquent             & PHP        & Laravel 11 incluido     & Active Record\\
Doctrine ORM         & PHP        & 3.x                     & Data Mapper\\
Entity Framework Core & C\#       & 8.0                     & Code-First, Data Mapper\\
TypeORM              & TypeScript & 0.3                     & Híbrido AR/DM\\
Prisma               & TypeScript & 5.x                     & Schema-first generador\\
SQLAlchemy           & Python     & 2.x                     & Data Mapper + AR\\
Active Record (Rails) & Ruby      & 8.x                     & Active Record canónico\\
GORM                 & Go         & 1.25                    & Active Record\\
\bottomrule
\end{tabular}
\end{table}

\section{Consultas SQL desde aplicaciones web}

Aunque el ORM oculte el SQL, el desarrollador web profesional
debe conocer cómo se ejecutan las consultas desde aplicaciones,
qué amenazas existen y qué defensas debe aplicar siempre.


\subsection{La amenaza fundamental: SQL Injection}

SQL Injection no es una vulnerabilidad teórica: ha sido la
causa de ataques contra Sony, LinkedIn, Yahoo y miles de
organizaciones. Sigue siendo la \emph{Top 3} del OWASP en 2021. Su
mecánica es engañosamente simple.

\textbf{Ejemplo corto comentado (anatomía de un SQL Injection) (Listado~\ref{lst:11.1}):}
\begin{lstlisting}[language=PHP,numbers=none,caption={Ejemplo de PHP},label=lst:11.1]
<?php
// Codigo vulnerable: concatena directamente la entrada del usuario
$email = $_POST['email'];
$sql = "SELECT * FROM users WHERE email = '$email'";

// El atacante introduce en el campo email:
//   ' OR '1'='1
// La consulta resultante se convierte en:
//   SELECT * FROM users WHERE email = '' OR '1'='1'
// Que devuelve TODOS los usuarios. El atacante entra sin password.
?>
\end{lstlisting}


\begin{cajaadvertencia}{Anti-patrón --- el código que NUNCA debe escribirse}
\begin{lstlisting}[style=php,numbers=none,caption={Ejemplo de PHP},label=lst:11.2]
// VULNERABLE - PHP
$sql = "SELECT * FROM users WHERE email = '" . $_POST['email'] . "'";
$res = $pdo->query($sql);
\end{lstlisting}
\textbf{Ejemplo de ataque trivial:}\\
Si \texttt{\$\_POST['email'] = \char34{}admin' OR '1'='1' -{}-\char34{}}, la consulta
ejecutada será:\\
\texttt{SELECT * FROM users WHERE email = 'admin' OR '1'='1' -{}-'}\\
que devuelve \emph{todos} los usuarios. Variantes más sofisticadas
permiten dropear tablas, exfiltrar contraseñas, leer archivos del
servidor.

\textbf{Estadística reveladora:} en el \emph{Verizon Data Breach
Investigations Report 2025} \cite{verizon2025dbir}, SQL Injection
sigue siendo responsable del 14\% de las brechas de aplicaciones web,
30 años después de su descubrimiento.
\end{cajaadvertencia}

\subsection{La defensa: sentencias preparadas en las tres plataformas}

La defensa universal contra SQLi es la \emph{sentencia
preparada}: el SQL y los datos viajan por canales separados, de
modo que los datos no pueden alterar la estructura del SQL. Las
tres plataformas del curso la implementan.


\subsubsection{PHP con PDO}

PHP dispone de PDO (\emph{PHP Data Objects}), una capa de
abstracción que soporta múltiples motores (MySQL, PostgreSQL, SQLite,
SQL Server) con una API uniforme.


\begin{cajaejemplo}{Ejemplo 11.1 --- PDO con parámetros nominales y posicionales}

\textbf{Enunciado.} Implementar el CRUD de la entidad \texttt{Libro} con Spring Data JPA y H2: definir entidad, repositorio y controlador REST. Se exponen las cuatro operaciones básicas.

\begin{lstlisting}[style=php,caption={UsuarioRepository.php},label=lst:11.3]
<?php
declare(strict_types=1);

class UsuarioRepository                                     // declara clase PHP
{
    public function __construct(private PDO $pdo) {}        // metodo publico

    /** Busqueda con parametros NOMINALES (recomendado). */
    public function buscarPorEmail(string $email): ?array { // metodo publico
        $sql = 'SELECT id, nombre, email, rol, activo
                FROM usuarios
                WHERE email = :em AND activo = 1';
        $stmt = $this->pdo->prepare($sql);                  // prepara consulta SQL (defensa anti-SQLi)
        $stmt->execute([':em' => $email]);                  // ejecuta la consulta con parametros enlazados
        $u = $stmt->fetch(PDO::FETCH_ASSOC);                // recupera la siguiente fila como array
        return $u ?: null;                                  // retorna valor
    }

    /** Busqueda con parametros POSICIONALES. */
    public function buscarPorId(int $id): ?array {          // metodo publico
        $stmt = $this->pdo->prepare(                        // prepara consulta SQL (defensa anti-SQLi)
            'SELECT id, nombre, email FROM usuarios WHERE id = ?');
        $stmt->execute([$id]);                              // ejecuta la consulta con parametros enlazados
        return $stmt->fetch(PDO::FETCH_ASSOC) ?: null;      // recupera la siguiente fila como array
    }

    /** Busqueda con LIKE seguro - nunca concatenar el % */
    // metodo publico
    // metodo publico
    public function buscarPorNombre(string $patron): array {
        $stmt = $this->pdo->prepare(                        // prepara consulta SQL (defensa anti-SQLi)
            'SELECT id, nombre, email FROM usuarios
             WHERE nombre LIKE :p
             ORDER BY nombre LIMIT 50');
        // El % se anade aqui, NO concatenado en el SQL
        $stmt->execute([':p' => '%' . $patron . '%']);      // ejecuta la consulta con parametros enlazados
        return $stmt->fetchAll(PDO::FETCH_ASSOC);           // recupera todas las filas restantes
    }

    /** INSERT con retorno del ID generado. */
    public function crear(string $email, string $nombre,    // metodo publico
                          string $passwordPlano): int {
        // hashea contrasena con BCrypt
        // hashea contrasena con BCrypt
        $hash = password_hash($passwordPlano, PASSWORD_BCRYPT,
                              ['cost' => 12]);
        $stmt = $this->pdo->prepare(                        // prepara consulta SQL (defensa anti-SQLi)
            // hashea contrasena con BCrypt
            // hashea contrasena con BCrypt
            'INSERT INTO usuarios (email, nombre, password_hash, rol)
             VALUES (:em, :nom, :hash, :rol)');
        $stmt->execute([                                    // ejecuta la consulta con parametros enlazados
            ':em'   => $email,
            ':nom'  => $nombre,
            ':hash' => $hash,
            ':rol'  => 'ROLE_USER',
        ]);
        return (int)$this->pdo->lastInsertId();             // retorna valor
    }

    /** Transaccion explicita para operaciones multi-tabla. */
    // metodo publico
    // metodo publico
    public function transferirPuntos(int $de, int $a, int $monto): bool {
        try {
            $this->pdo->beginTransaction();
            $stmt1 = $this->pdo->prepare(                   // prepara consulta SQL (defensa anti-SQLi)
                'UPDATE saldos SET puntos = puntos - ? WHERE user_id = ?');
            $stmt1->execute([$monto, $de]);                 // ejecuta la consulta con parametros enlazados
            $stmt2 = $this->pdo->prepare(                   // prepara consulta SQL (defensa anti-SQLi)
                'UPDATE saldos SET puntos = puntos + ? WHERE user_id = ?');
            $stmt2->execute([$monto, $a]);                  // ejecuta la consulta con parametros enlazados
            $this->pdo->commit();
            return true;                                    // retorna valor
        } catch (\Throwable $e) {
            $this->pdo->rollBack();
            error_log("Error transferir: " . $e->getMessage());
            return false;                                   // retorna valor
        }
    }
}
\end{lstlisting}

\textbf{Configuración inicial obligatoria del PDO (Listado~\ref{lst:11.4}):}
\begin{lstlisting}[style=php,numbers=none,caption={Ejemplo de PHP},label=lst:11.4]
$pdo = new PDO(                                             // conecta a la base de datos via PDO
    'mysql:host=localhost;dbname=appweb;charset=utf8mb4',
    'usuario', 'password',
    [
        // Errores como excepciones (NO codigos de error silenciosos)
        PDO::ATTR_ERRMODE            => PDO::ERRMODE_EXCEPTION,
        // FETCH_ASSOC por defecto (arrays con keys, sin duplicar)
        PDO::ATTR_DEFAULT_FETCH_MODE => PDO::FETCH_ASSOC,
        // Prepared statements REALES (no emulados, mas seguros)
        PDO::ATTR_EMULATE_PREPARES   => false,
    ]
);
\end{lstlisting}
\end{cajaejemplo}

\subsubsection{Java con JDBC y try-with-resources}

Java accede a bases de datos a través de JDBC. Desde Java 7,
la sintaxis \emph{try-with-resources} garantiza que las conexiones,
\texttt{Statement}s y \texttt{ResultSet}s se cierren automáticamente,
incluso si hay excepciones.


\begin{cajaejemplo}{Ejemplo 11.2 --- JDBC moderno con try-with-resources}

\textbf{Enunciado.} Mostrar el patrón profesional de acceso a datos en Java con JDBC moderno: \texttt{try-with-resources} que gestiona automáticamente el cierre de \texttt{Connection}, \texttt{Statement} y \texttt{ResultSet}.

\begin{lstlisting}[style=java,caption={UsuarioJdbcRepository.java},label=lst:11.5]
package ec.edu.uteq.appweb.repository;                      // paquete del archivo

import java.sql.*;                                          // importacion de clase externa
import java.util.*;                                         // importacion de clase externa
import javax.sql.DataSource;                                // importacion de clase externa

public class UsuarioJdbcRepository {                        // declaracion de clase publica

    private final DataSource ds;                            // campo inmutable (mejor practica)

    // metodo publico
    // metodo publico
    public UsuarioJdbcRepository(DataSource ds) { this.ds = ds; }

    public Optional<Usuario> buscarPorEmail(String email) { // metodo publico
        String sql = """
            SELECT id, nombre, email, rol
            FROM usuarios
            WHERE email = ? AND activo = TRUE
            """;
        try (Connection cn = ds.getConnection();
             PreparedStatement ps = cn.prepareStatement(sql)) {
            ps.setString(1, email);
            try (ResultSet rs = ps.executeQuery()) {
                if (rs.next()) {                            // condicional
                    return Optional.of(new Usuario(         // devuelve el resultado al llamador
                        rs.getLong("id"),
                        rs.getString("nombre"),
                        rs.getString("email"),
                        rs.getString("rol")
                    ));
                }
                return Optional.empty();                    // devuelve el resultado al llamador
            }
        } catch (SQLException e) {
            // lanza excepcion explicita
            // lanza excepcion explicita
            throw new RepositoryException("Error buscando por email", e);
        }
    }

    public List<Usuario> buscarPorPatron(String patron) {   // metodo publico
        String sql = """
            SELECT id, nombre, email, rol
            FROM usuarios
            WHERE nombre LIKE ?
            ORDER BY nombre LIMIT 50
            """;
        List<Usuario> resultado = new ArrayList<>();
        try (Connection cn = ds.getConnection();
             PreparedStatement ps = cn.prepareStatement(sql)) {
            ps.setString(1, "%" + patron + "%");
            try (ResultSet rs = ps.executeQuery()) {
                while (rs.next()) {                         // bucle while
                    resultado.add(new Usuario(
                        rs.getLong("id"),
                        rs.getString("nombre"),
                        rs.getString("email"),
                        rs.getString("rol")
                    ));
                }
            }
        } catch (SQLException e) {
            // lanza excepcion explicita
            // lanza excepcion explicita
            throw new RepositoryException("Error buscando por patron", e);
        }
        return resultado;                                   // devuelve el resultado al llamador
    }

    // metodo publico
    // metodo publico
    public long crear(String email, String nombre, String hash) {
        String sql = """
            INSERT INTO usuarios (email, nombre, password_hash, rol)
            VALUES (?, ?, ?, 'ROLE_USER')
            """;
        try (Connection cn = ds.getConnection();
             PreparedStatement ps = cn.prepareStatement(sql,
                                       Statement.RETURN_GENERATED_KEYS)) {
            ps.setString(1, email);
            ps.setString(2, nombre);
            ps.setString(3, hash);
            ps.executeUpdate();
            try (ResultSet keys = ps.getGeneratedKeys()) {
                if (keys.next()) return keys.getLong(1);    // condicional
                // lanza excepcion explicita
                // lanza excepcion explicita
                throw new RepositoryException("No se obtuvo ID generado");
            }
        } catch (SQLException e) {
            // lanza excepcion explicita
            // lanza excepcion explicita
            throw new RepositoryException("Error creando usuario", e);
        }
    }
}
\end{lstlisting}
\end{cajaejemplo}

\subsubsection{C\# con ADO.NET}

C\# accede a bases de datos a través de ADO.NET, la
biblioteca clásica de .NET. Aunque Entity Framework Core es lo
preferido en proyectos nuevos, ADO.NET sigue siendo la base sobre
la que se construye y aparece en código heredado.


\begin{cajaejemplo}{Ejemplo 11.3 --- ADO.NET puro con SqlClient}

\textbf{Enunciado.} Mostrar el acceso a datos en C\# con ADO.NET puro (\texttt{SqlClient}), antes de pasar a Entity Framework Core. Útil para entender lo que el ORM oculta.

\begin{lstlisting}[style=cs,caption={UsuarioRepository.cs},label=lst:11.6]
using Microsoft.Data.SqlClient;                             // importa namespace
using AppwebApi.Models;                                     // importa namespace

public class UsuarioRepository                              // declaracion de clase publica
{
    private readonly string _connStr;                       // metodo o miembro privado

    public UsuarioRepository(IConfiguration cfg) =>         // metodo o miembro publico
        _connStr = cfg.GetConnectionString("Default")!;

    // metodo o miembro publico
    // metodo o miembro publico
    public async Task<Usuario?> BuscarPorEmailAsync(string email)
    {
        const string sql = @"
            SELECT id, nombre, email, rol
            FROM usuarios
            WHERE email = @email AND activo = 1";

        await using var cn = new SqlConnection(_connStr);
        await cn.OpenAsync();

        await using var cmd = new SqlCommand(sql, cn);
        cmd.Parameters.AddWithValue("@email", email);

        await using var rdr = await cmd.ExecuteReaderAsync();
        if (await rdr.ReadAsync()) {                        // condicional
            return new Usuario {                            // devuelve el resultado al llamador
                Id     = rdr.GetInt64(0),
                Nombre = rdr.GetString(1),
                Email  = rdr.GetString(2),
                Rol    = rdr.GetString(3)
            };
        }
        return null;                                        // devuelve el resultado al llamador
    }
}
\end{lstlisting}
\end{cajaejemplo}

\section{El patrón Repository: encapsular el acceso a datos}

El patrón Repository encapsula el acceso a datos detrás de
una interfaz limpia, ocultando los detalles de SQL/JPA al resto de
la aplicación. Esto mejora la testabilidad (se puede mockear) y la
mantenibilidad (cambiar de SQL puro a JPA no afecta a los servicios).


\begin{cajadefinicion}{Patrón Repository según Fowler}
\emph{«Mediates between the domain and data mapping layers using a
collection-like interface for accessing domain objects.»}
\cite{fowler2002patterns}

Un Repository expone una API estilo colección sobre las entidades de
dominio, ocultando los detalles del SGBD. Sus responsabilidades:
\begin{itemize}
\item Encapsular consultas (no exponer SQL al servicio).
\item Devolver objetos de dominio o DTOs, no \texttt{ResultSet}.
\item Centralizar la traducción objeto $\leftrightarrow$ tabla.
\item Facilitar el testing con repositorios fake o mocks.
\end{itemize}
\end{cajadefinicion}

\subsection{Repositorio típico en Spring Data JPA}

Spring Data JPA reduce el boilerplate al mínimo: basta con declarar
una interfaz que extiende \texttt{JpaRepository}; Spring genera la
implementación en tiempo de ejecución.

\begin{cajaejemplo}{Ejemplo 11.4 --- ProductoRepository de Spring Data}

\textbf{Enunciado.} Implementar un repositorio Spring Data JPA para la entidad \texttt{Producto} declarando solo la interfaz; Spring genera la implementación con todas las operaciones CRUD.

\begin{lstlisting}[style=java,caption={ProductoRepository.java},label=lst:11.7]
package ec.edu.uteq.appweb.repository;                      // paquete del archivo

import ec.edu.uteq.appweb.model.Producto;                   // importacion de clase externa
import org.springframework.data.domain.Page;                // importacion de clase externa
import org.springframework.data.domain.Pageable;            // importacion de clase externa
import org.springframework.data.jpa.repository.*;           // importacion de clase externa
import org.springframework.data.repository.query.Param;     // importacion de clase externa
import java.math.BigDecimal;                                // importacion de clase externa
import java.util.List;                                      // importacion de clase externa

public interface ProductoRepository                         // declaracion de interfaz
        extends JpaRepository<Producto, Long> {

    // Consulta DERIVADA del nombre del metodo (sin SQL)
    List<Producto> findByStockGreaterThanOrderByNombreAsc(int min);

    Page<Producto> findByCategoriaIgnoreCase(String cat, Pageable pag);

    // JPQL personalizado con parametros nominales
    @Query("""                                              // consulta JPQL personalizada
        SELECT p FROM Producto p
        WHERE p.precio BETWEEN :min AND :max
          AND p.stock > 0
        ORDER BY p.precio ASC
        """)
    List<Producto> rangoPrecio(@Param("min") BigDecimal min,
                               @Param("max") BigDecimal max);

    // Update masivo
    @Modifying                                              // consulta que modifica datos (no SELECT)
    // consulta JPQL personalizada
    // consulta JPQL personalizada
    @Query("UPDATE Producto p SET p.stock = p.stock - :n WHERE p.id = :id AND p.stock >= :n")
    int decrementarStock(@Param("id") Long id, @Param("n") int n);

    // Native SQL cuando JPQL no alcanza
    @Query(value = """                                      // consulta JPQL personalizada
        SELECT p.*, c.nombre AS cat_nombre
        FROM productos p
        JOIN categorias c ON c.id = p.categoria_id
        WHERE p.stock > 0
        """, nativeQuery = true)
    List<Object[]> reporteStock();
}
\end{lstlisting}
\end{cajaejemplo}

\section{El problema N+1 y su solución}

El problema N+1 es la patología de rendimiento más común en
aplicaciones que usan ORM. Ocurre cuando, al iterar una colección
de N entidades con relaciones, se ejecutan N consultas adicionales
en lugar de una sola con \texttt{JOIN}.


\begin{cajaadvertencia}{Problema N+1 --- el bug silencioso que mata el rendimiento}
Si tiene una entidad \texttt{Producto} con relación
\texttt{@ManyToOne Categoria categoria}, y carga una lista de 100
productos para mostrarlos con su categoría, ocurre lo siguiente sin
optimización:

\begin{enumerate}
\item Una consulta para obtener los 100 productos.
\item Cien consultas adicionales (una por cada producto) para
obtener su categoría asociada.
\end{enumerate}

\textbf{Total: 101 consultas.} Si la página muestra 1000 productos,
son 1001 consultas. La latencia se dispara, la BD colapsa.

\textbf{La solución universal:} \emph{eager loading} explícito, que
hace UN SOLO JOIN en lugar de N+1 consultas.
\end{cajaadvertencia}

\subsection{Solución en cada ORM}

\textbf{Spring Data JPA --- usar JOIN FETCH (Listado~\ref{lst:11.8}):}
\begin{lstlisting}[style=java,numbers=none,caption={Ejemplo de Java},label=lst:11.8]
@Query("""                                                  // consulta JPQL personalizada
    SELECT p FROM Producto p
    JOIN FETCH p.categoria
    WHERE p.stock > 0
    """)
List<Producto> findAllConCategoria();
\end{lstlisting}

\textbf{Eloquent (Laravel) --- usar with (Listado~\ref{lst:11.9}):}
\begin{lstlisting}[style=php,numbers=none,caption={Ejemplo de PHP},label=lst:11.9]
$productos = Producto::with('categoria')
    ->where('stock','>',0)
    ->get();
\end{lstlisting}

\textbf{Entity Framework Core --- usar Include (Listado~\ref{lst:11.10}):}
\begin{lstlisting}[style=cs,numbers=none,caption={Ejemplo de C\#},label=lst:11.10]
var productos = await _ctx.Productos
    .Include(p => p.Categoria)
    .Where(p => p.Stock > 0)
    .ToListAsync();
\end{lstlisting}

\textbf{Doctrine ORM --- usar fetch en DQL (Listado~\ref{lst:11.11}):}
\begin{lstlisting}[style=php,numbers=none,caption={Ejemplo de PHP},label=lst:11.11]
$query = $em->createQuery(
    'SELECT p, c FROM App\Entity\Producto p
     JOIN p.categoria c WHERE p.stock > 0');
$productos = $query->getResult();
\end{lstlisting}

\section{Migraciones de esquema con Flyway y Liquibase}

Las migraciones permiten versionar el esquema de la BD como código.
Cada migración es un script SQL numerado que se aplica una sola vez.
La BD «sabe» qué migraciones ha aplicado mediante una tabla de control
(\texttt{flyway\_schema\_history} o \texttt{databasechangelog}).

\begin{cajaejemplo}{Ejemplo 11.5 --- Estructura Flyway en Spring Boot}

\textbf{Enunciado.} Configurar Flyway en un proyecto Spring Boot para gestionar migraciones de base de datos versionadas y reproducibles entre entornos.

\begin{lstlisting}[style=shell,numbers=none,caption={Ejemplo de Shell},label=lst:11.12]
src/main/resources/db/migration/
  V1__crear_usuarios.sql
  V2__crear_productos.sql
  V3__agregar_indice_email_usuarios.sql
  V4__crear_pedidos.sql
\end{lstlisting}

\begin{lstlisting}[style=sql,caption={V1\_\_crear\_usuarios.sql},label=lst:11.13]
-- crea tabla nueva en la BD
CREATE TABLE usuarios (                                     -- crea tabla nueva en la BD
    id            BIGSERIAL PRIMARY KEY,
    -- columna obligatoria (no admite NULL)
    email         VARCHAR(180) NOT NULL UNIQUE,             -- columna obligatoria (no admite NULL)
    -- columna obligatoria (no admite NULL)
    nombre        VARCHAR(100) NOT NULL,                    -- columna obligatoria (no admite NULL)
    -- columna obligatoria (no admite NULL)
    password_hash VARCHAR(255) NOT NULL,                    -- columna obligatoria (no admite NULL)
    -- columna obligatoria (no admite NULL)
    -- columna obligatoria (no admite NULL)
    rol           VARCHAR(30)  NOT NULL DEFAULT 'ROLE_USER',
    -- columna obligatoria (no admite NULL)
    activo        BOOLEAN      NOT NULL DEFAULT TRUE,       -- columna obligatoria (no admite NULL)
    -- columna obligatoria (no admite NULL)
    creado_en     TIMESTAMPTZ  NOT NULL DEFAULT NOW(),      -- columna obligatoria (no admite NULL)
    actualizado_en TIMESTAMPTZ
);

-- crea indice para acelerar busquedas
CREATE INDEX idx_usuarios_email_activo                      -- crea indice para acelerar busquedas
    ON usuarios(email) WHERE activo = TRUE;
\end{lstlisting}

\textbf{Configuración mínima en \texttt{application.yml}:}
\begin{lstlisting}[style=yaml,numbers=none,caption={Ejemplo de YAML},label=lst:11.14]
spring:                                                     # configuracion de Spring
  datasource:                                               # configuracion de origen de datos (BD)
    url:      jdbc:postgresql://localhost:5432/appweb
    username: appweb                                        # usuario de la BD
    password: ${DB_PASSWORD}                                # contrasena (mejor leer de env var)
  jpa:                                                      # configuracion JPA/Hibernate
    hibernate:                                              # configuracion especifica de Hibernate
      ddl-auto: validate    # NO usar 'update' con Flyway
    properties:
      hibernate.dialect: org.hibernate.dialect.PostgreSQLDialect
  flyway:
    enabled: true
    baseline-on-migrate: true
\end{lstlisting}
\end{cajaejemplo}

\section{Procedimientos almacenados: cuándo aún tiene sentido}

Los procedimientos almacenados (\emph{stored procedures})
fueron durante décadas la forma canónica de encapsular lógica de
datos en el motor de BD. Hoy se usan menos, pero hay casos donde aún
tienen sentido y conviene conocerlos.


\begin{cajacompara}{Procedimientos almacenados vs lógica en aplicación}
\centering\footnotesize
\begin{tabularx}{\textwidth}{lXX}
\toprule
\textbf{Aspecto} & \textbf{Stored Procedures} & \textbf{Lógica en aplicación}\\
\midrule
Rendimiento & Excelente (compilados, cerca del dato) & Bueno con cache\\
Versionado con Git & Difícil (suele estar fuera del repo) & Natural\\
Testing & Difícil (requiere BD viva) & Fácil con mocks\\
Portabilidad SGBD & Acoplado a un motor & Multi-SGBD via ORM\\
Lógica compleja & Mantenimiento difícil & Más limpia\\
Operaciones masivas & Imbatible (sin round-trips) & Ineficiente\\
\bottomrule
\end{tabularx}
\end{cajacompara}

\textbf{Cuándo usar stored procedures en 2026:}
\begin{itemize}
\item Operaciones masivas que tocan miles de filas (ETL, batch).
\item Lógica que requiere consistencia transaccional fina con
múltiples tablas.
\item Reportes complejos con muchos JOINs.
\item Sistemas heredados donde la lógica ya está en SP.
\end{itemize}

\section{Práctica integrada: CRUD ORM vs SQL puro}

La práctica integrada de la semana confronta directamente
las dos estrategias dominantes: JDBC puro vs Spring Data JPA. El
estudiante implementa el mismo CRUD con ambas y produce una tabla
comparativa.


\begin{cajaejercicio}{Ejercicio 11.1 --- Comparación experimental SQL puro vs ORM (Sumativa GA-PE Unidad III)}

\textbf{Objetivo:} implementar el mismo CRUD en dos versiones (SQL
puro y ORM) y comparar líneas de código, rendimiento y mantenibilidad.

\subsubsection*{IDE: IntelliJ IDEA Ultimate}
Plugins necesarios (todos incluidos en Ultimate):
\begin{itemize}
\item \emph{Database Tools and SQL} (cliente BD nativo).
\item \emph{Spring Boot}.
\item \emph{Lombok}.
\end{itemize}

\subsubsection*{Paso 1 --- Preparar PostgreSQL en Docker}

Más simple que instalar PostgreSQL nativo (Listado~\ref{lst:11.15}):
\begin{lstlisting}[style=shell,numbers=none,caption={Ejemplo de Shell},label=lst:11.15]
docker run -d --name pg-appweb \                            # arranca un contenedor
  -e POSTGRES_USER=appweb \
  -e POSTGRES_PASSWORD=appweb \
  -e POSTGRES_DB=appweb \
  -p 5432:5432 postgres:16
\end{lstlisting}

Verificar (Listado~\ref{lst:11.16}):
\begin{lstlisting}[style=shell,numbers=none,caption={Ejemplo de Shell},label=lst:11.16]
# ejecuta comando en contenedor activo
# ejecuta comando en contenedor activo
docker exec -it pg-appweb psql -U appweb -d appweb -c "SELECT version();"
\end{lstlisting}

\subsubsection*{Paso 2 --- Conectar IntelliJ a la BD}

Una vez levantada PostgreSQL, el siguiente paso es configurar
IntelliJ IDEA Ultimate para conectarse a ella desde la propia IDE.
Esto permite ejecutar SQL ad-hoc, ver datos y probar consultas sin
salir del entorno.


\begin{enumerate}
\item Abrir el panel \emph{Database} (lateral derecho).
\item Pulsar \texttt{+} $\rightarrow$ \texttt{Data Source}
$\rightarrow$ \texttt{PostgreSQL}.
\item Host: \texttt{localhost}, Port: \texttt{5432},
Database: \texttt{appweb}, User: \texttt{appweb}, Password:
\texttt{appweb}.
\item Pulsar \texttt{Test Connection}. Si pide descargar el driver,
aceptar.
\item Pulsar \texttt{OK}.
\end{enumerate}

\subsubsection*{Paso 3 --- Crear tabla y datos de prueba}

Click derecho sobre la conexión $\rightarrow$ \texttt{New}
$\rightarrow$ \texttt{Query Console}. Ejecutar (Listado~\ref{lst:11.17}):

\begin{lstlisting}[style=sql,numbers=none,caption={Ejemplo de SQL},label=lst:11.17]
-- crea tabla nueva en la BD
CREATE TABLE productos (                                    -- crea tabla nueva en la BD
    id        BIGSERIAL PRIMARY KEY,
    -- columna obligatoria (no admite NULL)
    nombre    VARCHAR(120) NOT NULL,                        -- columna obligatoria (no admite NULL)
    -- columna obligatoria (no admite NULL)
    precio    NUMERIC(10,2) NOT NULL CHECK (precio >= 0),   -- columna obligatoria (no admite NULL)
    -- columna obligatoria (no admite NULL)
    -- columna obligatoria (no admite NULL)
    stock     INTEGER NOT NULL DEFAULT 0 CHECK (stock >= 0),
    -- columna obligatoria (no admite NULL)
    creado    TIMESTAMPTZ NOT NULL DEFAULT NOW()            -- columna obligatoria (no admite NULL)
);

-- Cargar 1000 productos para benchmark
-- inserta una fila nueva
INSERT INTO productos (nombre, precio, stock)               -- inserta una fila nueva
SELECT
    'Producto ' || i,
    ROUND((random() * 100 + 1)::numeric, 2),
    (random() * 100)::int
-- tabla de origen
FROM generate_series(1, 1000) i;                            -- tabla de origen
\end{lstlisting}

Verificar: \texttt{SELECT COUNT(*) FROM productos;} debe devolver 1000.

\subsubsection*{Paso 4 --- Crear proyecto Spring Boot}

En IntelliJ: \texttt{File $\rightarrow$ New $\rightarrow$ Project
$\rightarrow$ Spring Boot}. Configurar:
\begin{itemize}
\item Language: Java; Build: Maven; JDK: 21.
\item Group: \texttt{ec.edu.uteq}; Artifact: \texttt{benchmark-orm}.
\item Spring Boot 3.4.x.
\item Dependencies: \emph{Spring Web}, \emph{Spring Data JPA},
\emph{PostgreSQL Driver}, \emph{Lombok}, \emph{Spring Boot DevTools}.
\end{itemize}

\subsubsection*{Paso 5 --- Configurar la conexión}

Editar \texttt{src/main/resources/application.yml} (Listado~\ref{lst:11.18}):

\begin{lstlisting}[style=yaml,numbers=none,caption={Ejemplo de YAML},label=lst:11.18]
spring:                                                     # configuracion de Spring
  datasource:                                               # configuracion de origen de datos (BD)
    url:      jdbc:postgresql://localhost:5432/appweb
    username: appweb                                        # usuario de la BD
    password: appweb                                        # contrasena (mejor leer de env var)
  jpa:                                                      # configuracion JPA/Hibernate
    hibernate.ddl-auto: none   # tabla ya existe
    show-sql: true             # ver SQL generado
    properties.hibernate.format_sql: true
\end{lstlisting}

\subsubsection*{Paso 6 --- Versión A: JDBC puro}

La versión A implementa el CRUD usando solo JDBC puro: el
estudiante escribe el SQL a mano, gestiona la conexión, prepara las
sentencias y procesa los resultados. Es laborioso pero ilustrativo. V\'ease el Listado~\ref{lst:11.19}.


\begin{lstlisting}[style=java,caption={JdbcController.java},label=lst:11.19]
package ec.edu.uteq.benchmark.jdbc;                         // paquete del archivo

import org.springframework.jdbc.core.JdbcTemplate;          // importacion de clase externa
import org.springframework.web.bind.annotation.*;           // importacion de clase externa
import java.sql.Timestamp;                                  // importacion de clase externa
import java.util.*;                                         // importacion de clase externa

@RestController                                             // expone endpoints REST con JSON
@RequestMapping("/api/jdbc/productos")                      // ruta base del controlador
public class JdbcController {                               // declaracion de clase publica

    private final JdbcTemplate jdbc;                        // campo inmutable (mejor practica)

    // metodo publico
    // metodo publico
    public JdbcController(JdbcTemplate jdbc) { this.jdbc = jdbc; }

    @GetMapping                                             // endpoint HTTP GET
    public List<Map<String, Object>> listar() {             // metodo publico
        return jdbc.queryForList(                           // devuelve el resultado al llamador
            "SELECT id, nombre, precio, stock, creado FROM productos "
          + "ORDER BY id");
    }

    @GetMapping("/{id}")                                    // endpoint HTTP GET
    // metodo publico
    // metodo publico
    public Map<String, Object> obtener(@PathVariable long id) {
        return jdbc.queryForMap(                            // devuelve el resultado al llamador
            "SELECT id, nombre, precio, stock, creado FROM productos "
          + "WHERE id = ?", id);
    }
}
\end{lstlisting}

\subsubsection*{Paso 7 --- Versión B: Spring Data JPA}

La versión B implementa el mismo CRUD usando Spring Data
JPA: el estudiante define la entidad, declara la interfaz del
repositorio, y JPA genera el SQL automáticamente. El contraste con
la versión A es revelador. V\'ease el Listado~\ref{lst:11.20}.


\begin{lstlisting}[style=java,caption={Producto.java},label=lst:11.20]
package ec.edu.uteq.benchmark.jpa;                          // paquete del archivo

import jakarta.persistence.*;                               // importacion de clase externa
import lombok.*;                                            // importacion de clase externa
import java.math.BigDecimal;                                // importacion de clase externa
import java.time.OffsetDateTime;                            // importacion de clase externa

@Entity @Table(name = "productos")                          // entidad JPA (mapea a tabla)
@Getter @Setter @NoArgsConstructor @AllArgsConstructor
public class Producto {                                     // declaracion de clase publica
    @Id @GeneratedValue(strategy = GenerationType.IDENTITY) // campo clave primaria
    private Long id;                                        // campo privado
    private String nombre;                                  // campo privado
    private BigDecimal precio;                              // campo privado
    private Integer stock;                                  // campo privado
    private OffsetDateTime creado;                          // campo privado
}
\end{lstlisting}

\begin{lstlisting}[style=java,caption={ProductoRepository.java + JpaController.java},label=lst:11.21]
package ec.edu.uteq.benchmark.jpa;                          // paquete del archivo

// importacion de clase externa
// importacion de clase externa
import org.springframework.data.jpa.repository.JpaRepository;
public interface ProductoRepository                         // declaracion de interfaz
    extends JpaRepository<Producto, Long> {}

// ---

import org.springframework.web.bind.annotation.*;           // importacion de clase externa
import java.util.*;                                         // importacion de clase externa

@RestController                                             // expone endpoints REST con JSON
@RequestMapping("/api/jpa/productos")                       // ruta base del controlador
public class JpaController {                                // declaracion de clase publica
    private final ProductoRepository repo;                  // campo inmutable (mejor practica)
    // metodo publico
    // metodo publico
    public JpaController(ProductoRepository r) { this.repo = r; }

    // endpoint HTTP GET
    // endpoint HTTP GET
    @GetMapping public List<Producto> listar() { return repo.findAll(); }

    @GetMapping("/{id}")                                    // endpoint HTTP GET
    public Producto obtener(@PathVariable Long id) {        // metodo publico
        return repo.findById(id).orElseThrow();             // devuelve el resultado al llamador
    }
}
\end{lstlisting}

\subsubsection*{Paso 8 --- Benchmark con Apache Bench}

Ejecutar la app: \texttt{Shift+F10} sobre la clase principal. Servidor
en \url{http://localhost:8080}.

Desde otra terminal, instalar Apache Bench (\texttt{ab}). En Ubuntu:
\texttt{sudo apt install apache2-utils}; en macOS: viene preinstalado.

Ejecutar (Listado~\ref{lst:11.22}):
\begin{lstlisting}[style=shell,numbers=none,caption={Ejemplo de Shell},label=lst:11.22]
# 1000 peticiones, 10 concurrentes, version JDBC
ab -n 1000 -c 10 http://localhost:8080/api/jdbc/productos > jdbc.txt

# Misma prueba para JPA
ab -n 1000 -c 10 http://localhost:8080/api/jpa/productos > jpa.txt

cat jdbc.txt | grep "Requests per second"
cat jpa.txt  | grep "Requests per second"
\end{lstlisting}

\subsubsection*{Paso 9 --- Tabla comparativa para entregar}

Tras implementar las dos versiones, el estudiante debe
producir una tabla comparativa que cuantifique las diferencias.


\begin{tabularx}{\textwidth}{Xll}
\toprule
\textbf{Métrica} & \textbf{Versión JDBC} & \textbf{Versión JPA}\\
\midrule
Líneas de código del controller & & \\
Líneas de código del repository & & \\
Requests/segundo (1000 req, 10 conc) & & \\
Tiempo medio por request (ms) & & \\
SQL generado (consultas en log) & 1 fija & varía con caché L1/L2 \\
Vulnerabilidad SQL Injection & Eliminada (\texttt{?}) & Eliminada (JPA) \\
Mantenibilidad (1--5) & & \\
Productividad para CRUD nuevo (1--5) & & \\
\bottomrule
\end{tabularx}

\subsubsection*{Reflexión final escrita (entregable)}

Responder en máximo 300 palabras:
\begin{enumerate}
\item ¿Cuál fue más rápido en el benchmark?
\item ¿Qué versión escribiría más rápido un desarrollador nuevo?
\item ¿En qué casos el ORM es perjudicial (ineficiente)?
\item ¿Por qué muchas empresas mantienen \emph{ambos} estilos en el
mismo proyecto?
\end{enumerate}

\textbf{Esta es la \emph{Sumativa GA-PE Unidad III} de la semana 11.}
\end{cajaejercicio}

\section{Buenas prácticas profesionales en acceso a datos}

La caja siguiente compila las reglas profesionales del
acceso a datos: principios que separan el código mantenible del
código que se vuelve ingobernable a los seis meses.


\begin{cajabuenas}{Quince reglas profesionales 2026}
\begin{enumerate}
\item NUNCA concatenar datos del usuario en SQL: usar siempre
parámetros.
\item Configurar el driver con \texttt{ATTR\_EMULATE\_PREPARES = false}
en PHP.
\item DTOs separados de entidades: nunca exponer entidades JPA en
controllers.
\item Consultas paginadas siempre (\texttt{LIMIT} + \texttt{OFFSET})
para listados.
\item Resolver el problema N+1 con \emph{eager loading} explícito.
\item Migraciones con Flyway o Liquibase: jamás
\texttt{ddl-auto=update} en producción.
\item Índices en columnas de búsqueda frecuente y en foreign keys.
\item Conexiones desde un \emph{pool} (HikariCP en Spring, pgBouncer
externo).
\item Transacciones explícitas para operaciones multi-tabla.
\item Soft delete (\texttt{activo=false}) sobre delete físico cuando
hay relaciones.
\item Auditar cambios sensibles (\texttt{creado\_en},
\texttt{actualizado\_en}, \texttt{actualizado\_por}).
\item Cuenta BD con permisos mínimos: la app no necesita
\texttt{DROP TABLE}.
\item Encriptar la conexión con SSL/TLS en producción.
\item Tests con BD en memoria (H2) o testcontainers para integración
real.
\item Documentar el modelo entidad-relación en el repositorio
(ER en draw.io).
\end{enumerate}
\end{cajabuenas}

\section{Ejercicio resuelto adicional con resultado visual}

Como cierre de la semana 11, el siguiente ejercicio adicional
está completamente resuelto y se complementa con un propuesto.


\begin{cajaejemplo}{Ejercicio resuelto 11.1 --- CRUD de libros con Spring Data JPA y H2}
\textbf{Enunciado.} Implementar un CRUD completo de la entidad
\texttt{Libro} (id, titulo, autor, anio) con Spring Data JPA y H2
embebido, exponer los endpoints REST y verificar visualmente las
operaciones desde Postman.

\textbf{Entidad (\texttt{Libro.java}):}
\begin{lstlisting}[language=Java,caption={Entidad JPA con validaciones},label=lst:11.23]
@Entity                                                     // entidad JPA (mapea a tabla)
@Table(name = "libros")                                     // nombre de la tabla en BD
public class Libro {                                        // declaracion de clase publica
  @Id @GeneratedValue(strategy = GenerationType.IDENTITY)   // campo clave primaria
  private Long id;                                          // campo privado
  @NotBlank @Size(max = 200)                                // validacion: no nulo y no vacio
  private String titulo;                                    // campo privado
  @NotBlank @Size(max = 100)                                // validacion: no nulo y no vacio
  private String autor;                                     // campo privado
  @Min(1900) @Max(2100)                                     // validacion: valor minimo
  private int anio;                                         // campo privado
  // getters/setters omitidos por brevedad
}
\end{lstlisting}

\textbf{Repositorio (\texttt{LibroRepository.java}):}
\begin{lstlisting}[language=Java,caption={Repositorio Spring Data},label=lst:11.24]
// declaracion de interfaz
// declaracion de interfaz
public interface LibroRepository extends JpaRepository<Libro, Long> {
  List<Libro> findByAutor(String autor);
  @Query("SELECT l FROM Libro l WHERE l.anio >= :desde")    // consulta JPQL personalizada
  List<Libro> publicadosDesde(@Param("desde") int desde);
}
\end{lstlisting}

\textbf{Controlador REST (\texttt{LibroController.java}):}
\begin{lstlisting}[language=Java,caption={Endpoints REST},label=lst:11.25]
@RestController @RequestMapping("/api/v1/libros")           // expone endpoints REST con JSON
public class LibroController {                              // declaracion de clase publica
  private final LibroRepository repo;                       // campo inmutable (mejor practica)
  // metodo publico
  // metodo publico
  public LibroController(LibroRepository r) { this.repo = r; }

  // endpoint HTTP GET
  // endpoint HTTP GET
  @GetMapping public List<Libro> listar() { return repo.findAll(); }
  // endpoint HTTP GET
  // endpoint HTTP GET
  @GetMapping("/{id}") public ResponseEntity<Libro> obtener(@PathVariable Long id) {
    return repo.findById(id).map(ResponseEntity::ok)        // devuelve el resultado al llamador
      .orElse(ResponseEntity.notFound().build());
  }
  // endpoint HTTP POST (crear)
  // endpoint HTTP POST (crear)
  @PostMapping public ResponseEntity<Libro> crear(@Valid @RequestBody Libro l) {
    Libro guardado = repo.save(l);
    return ResponseEntity.status(201).body(guardado);       // devuelve el resultado al llamador
  }
  // endpoint HTTP PUT (actualizar completo)
  // endpoint HTTP PUT (actualizar completo)
  @PutMapping("/{id}") public ResponseEntity<Libro> actualizar(
      @PathVariable Long id, @Valid @RequestBody Libro l) { // extrae valor de la URL
    return repo.findById(id).map(existente -> {             // devuelve el resultado al llamador
      existente.setTitulo(l.getTitulo());
      existente.setAutor(l.getAutor());
      existente.setAnio(l.getAnio());
      return ResponseEntity.ok(repo.save(existente));       // devuelve el resultado al llamador
    }).orElse(ResponseEntity.notFound().build());
  }
  // endpoint HTTP DELETE (eliminar)
  // endpoint HTTP DELETE (eliminar)
  @DeleteMapping("/{id}") public ResponseEntity<Void> eliminar(@PathVariable Long id) {
    return repo.findById(id).map(l -> {                     // devuelve el resultado al llamador
      repo.delete(l);
      return ResponseEntity.noContent().<Void>build();      // devuelve el resultado al llamador
    }).orElse(ResponseEntity.notFound().build());
  }
}
\end{lstlisting}

\textbf{Resultado visual} de \texttt{GET /api/v1/libros} en Postman
tras insertar tres libros:

\begin{center}
\fbox{\begin{minipage}{0.9\textwidth}
\vspace{0.2cm}\footnotesize
\hspace{0.3cm}\textbf{GET} \texttt{http://localhost:8080/api/v1/libros} $\quad$ \textbf{200 OK}\\[0.3cm]
\hspace{0.3cm}\texttt{[}\\
\hspace{0.6cm}\texttt{\{ \char34{}id\char34{}: 1, \char34{}titulo\char34{}: \char34{}Clean Code\char34{}, \char34{}autor\char34{}: \char34{}Robert C. Martin\char34{}, \char34{}anio\char34{}: 2008 \},}\\
\hspace{0.6cm}\texttt{\{ \char34{}id\char34{}: 2, \char34{}titulo\char34{}: \char34{}Designing Data-Intensive Applications\char34{},}\\
\hspace{1.0cm}\texttt{\char34{}autor\char34{}: \char34{}Martin Kleppmann\char34{}, \char34{}anio\char34{}: 2017 \},}\\
\hspace{0.6cm}\texttt{\{ \char34{}id\char34{}: 3, \char34{}titulo\char34{}: \char34{}Software Architecture in Practice\char34{},}\\
\hspace{1.0cm}\texttt{\char34{}autor\char34{}: \char34{}Len Bass\char34{}, \char34{}anio\char34{}: 2021 \}}\\
\hspace{0.3cm}\texttt{]}\\
\vspace{0.1cm}
\end{minipage}}
\end{center}

\textbf{Vista de la tabla H2} (consola H2 en \url{http://localhost:8080/h2-console}):

\begin{center}
\fbox{\begin{minipage}{0.9\textwidth}
\vspace{0.15cm}\footnotesize\centering
\begin{tabular}{|c|l|l|c|}
\hline
\textbf{ID} & \textbf{TITULO} & \textbf{AUTOR} & \textbf{ANIO} \\
\hline
1 & Clean Code & Robert C. Martin & 2008 \\
2 & Designing Data-Intensive Applications & Martin Kleppmann & 2017 \\
3 & Software Architecture in Practice & Len Bass & 2021 \\
\hline
\end{tabular}\\[0.1cm]
\end{minipage}}
\end{center}

\textbf{Discusión.} JPA generó automáticamente el SQL apropiado
(\texttt{INSERT}, \texttt{SELECT}, \texttt{UPDATE}, \texttt{DELETE})
y vinculó los parámetros por nombre, eliminando por completo la
posibilidad de SQL Injection. Las validaciones \texttt{@NotBlank},
\texttt{@Size} y \texttt{@Min}/\texttt{@Max} devuelven HTTP 400 con
detalle si el cliente envía datos inválidos.
\end{cajaejemplo}

\begin{cajaejercicio}{Ejercicio propuesto 11.2 --- Catálogo de cursos con relaciones y paginación}
\textbf{Enunciado.} Implementar un catálogo de cursos universitarios
con relación \emph{«un curso tiene muchos estudiantes»} y
\emph{«un estudiante puede estar en muchos cursos»} (\textbf{relación
muchos-a-muchos}). Resolver el problema N+1 mediante \texttt{JOIN
FETCH} y exponer endpoints REST con paginación, búsqueda y
ordenamiento.

\textbf{Entidades:}
\begin{itemize}
\item \texttt{Curso(id, codigo, nombre, creditos)}.
\item \texttt{Estudiante(id, cedula, nombre, email)}.
\item Tabla intermedia \texttt{inscripciones(curso\_id, estudiante\_id, fecha)}.
\end{itemize}

\textbf{Requisitos técnicos:}
\begin{enumerate}
\item Usar Spring Boot 3.4 + PostgreSQL 16 (no H2).
\item Mapear la relación muchos-a-muchos con \texttt{@ManyToMany}.
\item Endpoint \texttt{GET /api/v1/cursos?page=0\&size=10\&sort=nombre,asc}
debe usar \texttt{Pageable}.
\item Endpoint \texttt{GET /api/v1/cursos/\{id\}/estudiantes} debe
usar \texttt{JOIN FETCH} para evitar N+1.
\item Mostrar evidencia con Hibernate logs (\texttt{spring.jpa.show-sql=true})
de que se ejecuta UNA consulta en lugar de N+1.
\item Documentación OpenAPI con springdoc, accesible en
\url{/swagger-ui.html}.
\item Migraciones de BD con Flyway o Liquibase.
\item Cobertura JaCoCo $\geq 70\%$ con tests de integración usando
testcontainers.
\end{enumerate}

\textbf{Criterios de evaluación:}
\begin{itemize}
\item No hay SQL concatenado en ninguna parte del código.
\item El log muestra una sola consulta para listar cursos con sus
estudiantes (no N+1).
\item Validaciones \texttt{@Valid} correctas en todos los endpoints.
\item Documentación Swagger funcional.
\end{itemize}

\textbf{Lecturas recomendadas:} \cite{ottinger2017hibernate} para
optimización de queries; \cite{ottinger2017hibernate} para análisis del
problema N+1 en producción.
\end{cajaejercicio}

% =====================================================================
%       SEMANA 12: PATRONES ARQUITECTÓNICOS PARA APLICACIONES WEB
% =====================================================================
\chapter{Patrones de arquitectura para aplicaciones web}
\label{cap:semana12}

\epigraph{«La arquitectura es el conjunto de decisiones que son
difíciles de cambiar. Un buen arquitecto, por tanto, posterga las
decisiones innecesarias y razona profundamente las inevitables.»}
{--- \textit{Adaptado de Robert C.~Martin}}

\section*{Resultado de aprendizaje de la semana}
Al concluir la semana 12, el estudiante diseña aplicaciones web bajo
patrones arquitectónicos escalables (multicapa, hexagonal, SOA,
microservicios, EDA, P2P), documenta sus decisiones con el modelo C4
\cite{brown2023c4} y los Architectural Decision Records (ADR), y
evalúa los \emph{trade-offs} entre atributos de calidad ISO/IEC 25010
\cite{iso25010}.

\section{Principios fundacionales de arquitectura de software}

Antes de hablar de patrones concretos (capas, hexagonal,
microservicios), conviene fijar los principios fundacionales que los
gobiernan todos: separación de intereses, alta cohesión, bajo
acoplamiento, dependencias unidireccionales.


\begin{cajadefinicion}{¿Qué es la arquitectura de software?}
Según Bass, Clements y Kazman \cite{bass2021software}, la arquitectura
de un sistema computacional es \emph{el conjunto de estructuras
necesarias para razonar sobre el sistema, comprende los elementos del
software, las relaciones entre ellos y las propiedades de ambos}.

Tres ideas clave en esta definición:
\begin{itemize}
\item \textbf{Estructuras múltiples}: no hay UNA arquitectura, sino
varias vistas (estática, dinámica, despliegue).
\item \textbf{Razonamiento}: la arquitectura existe para que un
equipo pueda pensar el sistema sin perderse en detalles.
\item \textbf{Decisiones difíciles de cambiar}: lo arquitectónico es
lo que costaría reescribir si se cambiara.
\end{itemize}
\end{cajadefinicion}

\subsection{Los principios SOLID aplicados a la arquitectura}

Los principios SOLID, formulados originalmente por Robert C.
Martin para diseño orientado a objetos, se generalizan a la
arquitectura. La caja siguiente los traduce al contexto
arquitectónico.


\begin{table}[htbp]
\centering
\caption{SOLID, sus autores y significado en arquitectura.}
\footnotesize
\begin{tabularx}{\textwidth}{lXX}
\toprule
\textbf{Principio} & \textbf{Significado} & \textbf{Aplicación arquitectónica}\\
\midrule
S --- Single Responsibility & Una clase, una razón para cambiar & Cada módulo/servicio tiene un foco\\
O --- Open/Closed & Abierto a extensión, cerrado a modificación & Plugins, hooks, eventos\\
L --- Liskov Substitution & Subtipos sustituibles por sus tipos base & Polimorfismo correcto\\
I --- Interface Segregation & Interfaces pequeñas y enfocadas & APIs minimal por módulo\\
D --- Dependency Inversion & Depender de abstracciones, no concreciones & Inyección de dependencias\\
\bottomrule
\end{tabularx}
\end{table}

\subsection{Atributos de calidad ISO/IEC 25010}

El estándar ISO/IEC 25010 define las características de
calidad de un sistema software. Conocerlas permite razonar sobre
arquitectura de forma estructurada. La tabla siguiente las cataloga.


\begin{figure}[htbp]
\centering
\begin{tikzpicture}[
  every node/.style={font=\scriptsize,align=center},
  centro/.style={circle,draw=uteqverde,thick,fill=uteqverde!20,
                 minimum size=2.2cm,font=\small\bfseries},
  atrib/.style={rectangle,draw=uteqverde!80,fill=uteqverde!8,thick,
                rounded corners=4pt,minimum width=2cm,minimum height=0.7cm}]

  \node[centro] (q) {Calidad\\ISO 25010};

  \node[atrib,above=2cm of q] (func) {Idoneidad\\funcional};
  \node[atrib,right=2cm of func] (perf) {Eficiencia de\\desempeño};
  \node[atrib,right=of q] (comp) {Compatibilidad};
  \node[atrib,below right=2cm and 1cm of q] (usab) {Usabilidad};
  \node[atrib,below=2cm of q] (fiab) {Fiabilidad};
  \node[atrib,below left=2cm and 1cm of q] (segu) {Seguridad};
  \node[atrib,left=of q] (mant) {Mantenibilidad};
  \node[atrib,left=2cm of func] (port) {Portabilidad};

  \foreach \n in {func,perf,comp,usab,fiab,segu,mant,port}
    \draw[uteqverde!50,thick] (q) -- (\n);
\end{tikzpicture}
\caption{Los ocho atributos de calidad de la norma ISO/IEC 25010. Toda
decisión arquitectónica supone un \emph{trade-off} entre algunos de
estos atributos.}
\end{figure}

\section{Patrón 1: Arquitectura por capas (3-capas y 4-capas)}

La arquitectura por capas es el patrón más extendido en
aplicaciones empresariales: la aplicación se divide en niveles
horizontales con responsabilidades específicas y dependencias hacia
abajo.


\subsection{Estructura clásica}

La estructura clásica de tres capas (presentación, lógica,
datos) se ha extendido a cuatro al separar la lógica de aplicación
del dominio. La figura siguiente la ilustra.


\begin{figure}[htbp]
\centering
\begin{tikzpicture}[
  every node/.style={font=\footnotesize,align=center},
  capa/.style={rectangle,draw=uteqverde,thick,minimum width=10cm,
               minimum height=1.1cm,rounded corners=2pt}]
  \node[capa,fill=uteqverde!15]
    (pres) {\textbf{Presentación}\\
            Razor / Thymeleaf / Angular --- Controllers --- DTOs};
  \node[capa,fill=uteqverde!25,below=0.3cm of pres]
    (neg)  {\textbf{Negocio / Servicios}\\
            Casos de uso, reglas de dominio, validación};
  \node[capa,fill=uteqverde!35,below=0.3cm of neg]
    (dat)  {\textbf{Acceso a Datos / Repositorios}\\
            JPA / Doctrine / EF Core};
  \node[capa,fill=uteqverde!50,below=0.3cm of dat]
    (db)   {\textbf{Base de Datos}\\
            PostgreSQL / MySQL / SQL Server};

  \draw[->,>=stealth,thick] ([xshift=-3cm]pres.south)
       -- ([xshift=-3cm]neg.north);
  \draw[->,>=stealth,thick] ([xshift=-3cm]neg.south)
       -- ([xshift=-3cm]dat.north);
  \draw[->,>=stealth,thick] ([xshift=-3cm]dat.south)
       -- ([xshift=-3cm]db.north);
\end{tikzpicture}
\caption{Arquitectura clásica de 4 capas para aplicaciones web. La
dependencia siempre es descendente: la capa superior depende de la
inferior, nunca al revés.}
\end{figure}

\textbf{Regla cardinal:} una capa puede invocar a la inmediatamente
inferior; nunca al revés ni saltarse capas (no permitir que un
controller acceda directamente al SGBD).

\subsection{Implementación canónica en Spring Boot}

Spring Boot facilita la implementación de la arquitectura por
capas mediante sus anotaciones estereotipo
(\texttt{@Controller}, \texttt{@Service}, \texttt{@Repository}). El
ejemplo siguiente muestra una implementación canónica.


\begin{cajaejemplo}{Ejemplo 12.1 --- Estructura por capas en un proyecto Spring Boot}

\textbf{Enunciado.} Modelar arquitectónicamente con el modelo C4 (niveles 1 y 2) una tienda en línea sencilla. Producir el diagrama de Contexto y el de Contenedores.

\begin{lstlisting}[style=shell,numbers=none,caption={Ejemplo de Shell},label=lst:12.1]
src/main/java/ec/edu/uteq/appweb/
  controller/         <- Capa de Presentacion
    ProductoController.java
    ExceptionHandler.java
  service/            <- Capa de Negocio
    ProductoService.java
    impl/ProductoServiceImpl.java
  repository/         <- Capa de Acceso a Datos
    ProductoRepository.java
  model/              <- Entidades de dominio
    Producto.java
  dto/                <- Objetos de transferencia
    ProductoDto.java
    ProductoRequest.java
  exception/
    ResourceNotFoundException.java
    ValidationException.java
  config/
    SecurityConfig.java
    WebConfig.java
\end{lstlisting}
\end{cajaejemplo}

\subsection{Variante: Arquitectura Hexagonal (Ports and Adapters)}

La arquitectura hexagonal (\emph{Ports and Adapters}) de
Alistair Cockburn es una variante refinada de la arquitectura por
capas que coloca el dominio en el centro y vuelve invertibles las
dependencias.


\begin{figure}[htbp]
\centering
\begin{tikzpicture}[
  every node/.style={font=\scriptsize,align=center}]

  \node[regular polygon,regular polygon sides=6,
        draw=uteqverde,thick,fill=uteqverde!10,
        minimum width=4cm] (hex) {};
  \node[font=\small\bfseries,align=center] at (hex) {Dominio\\(Casos de uso)\\Puro Java};

  \node[rectangle,draw=uteqdorado,fill=uteqdorado!15,thick,
        rounded corners,minimum width=2cm,minimum height=0.8cm,
        above left=1cm and -0.5cm of hex] (rest) {REST\\Controller};
  \node[rectangle,draw=uteqdorado,fill=uteqdorado!15,thick,
        rounded corners,minimum width=2cm,minimum height=0.8cm,
        above right=1cm and -0.5cm of hex] (cli) {Cliente\\HTTP};
  \node[rectangle,draw=uteqdorado,fill=uteqdorado!15,thick,
        rounded corners,minimum width=2cm,minimum height=0.8cm,
        below left=1cm and -0.5cm of hex] (jpa) {JPA\\Repository};
  \node[rectangle,draw=uteqdorado,fill=uteqdorado!15,thick,
        rounded corners,minimum width=2cm,minimum height=0.8cm,
        below right=1cm and -0.5cm of hex] (cache) {Redis\\Cache};

  \draw[<->,thick] (rest)  -- (hex);
  \draw[<->,thick] (cli)   -- (hex);
  \draw[<->,thick] (jpa)   -- (hex);
  \draw[<->,thick] (cache) -- (hex);
\end{tikzpicture}
\caption{Arquitectura hexagonal: el dominio (centro) define
\emph{ports} (interfaces) y los \emph{adapters} (implementaciones)
externos lo conectan con tecnologías concretas (REST, JPA, Redis).
Permite cambiar el adaptador de persistencia sin tocar el dominio.}
\end{figure}

\section{Patrón 2: Arquitectura Orientada a Servicios (SOA)}

SOA emergió en los 2000 como enfoque para integrar sistemas
heterogéneos mediante servicios autónomos comunicados por SOAP/XML
sobre un \emph{Enterprise Service Bus} (ESB). Características:

\begin{itemize}
\item Servicios reutilizables expuestos vía contratos formales (WSDL).
\item Comunicación asíncrona vía ESB (TIBCO, MuleSoft, IBM IIB).
\item Gobernanza centralizada (registro de servicios, políticas
unificadas).
\item Datos a menudo compartidos entre servicios.
\end{itemize}

\textbf{Cuándo SOA sigue siendo relevante:} integraciones con
sistemas heredados (banca, gobierno), entornos donde el WSDL es
contractualmente exigido, dominios con altos requisitos de
interoperabilidad B2B.

\section{Patrón 3: Microservicios}

Microservicios es el patrón arquitectónico que más atención
ha recibido en la última década. Implica descomponer el sistema en
servicios pequeños, autónomos, desplegables y escalables
independientemente.


\begin{cajadefinicion}{Microservicios según Newman}
\emph{«Pequeños servicios autónomos que trabajan juntos, modelados
alrededor de un dominio de negocio.»} \cite{newman2021building}

\textbf{Características distintivas:}
\begin{itemize}
\item Servicios pequeños con responsabilidad acotada.
\item Despliegue independiente (cada servicio tiene su pipeline CI/CD).
\item Comunicación vía HTTP/JSON o eventos asíncronos.
\item Datos descentralizados (una BD por servicio).
\item Equipos autónomos (cada equipo es dueño de uno o varios servicios).
\item Heterogeneidad tecnológica permitida.
\end{itemize}
\end{cajadefinicion}

\subsection{Comparativa SOA vs microservicios}

SOA (\emph{Service-Oriented Architecture}) y microservicios
comparten la idea de modularidad por servicios, pero difieren en
escala, autonomía y tecnologías. La tabla siguiente los contrasta.


\begin{cajacompara}{SOA clásico vs microservicios modernos}
\centering\footnotesize
\begin{tabularx}{\textwidth}{lXX}
\toprule
\textbf{Criterio} & \textbf{SOA clásico} & \textbf{Microservicios}\\
\midrule
Tamaño del servicio & Mediano-grande & Pequeño y enfocado\\
Comunicación & SOAP/XML, ESB centralizado & REST/JSON, gRPC, eventos\\
Datos & Compartidos vía ESB & Una BD por servicio\\
Despliegue & Acoplado al ESB & Independiente por servicio\\
Equipos & Centralizados & Autónomos por servicio\\
Lenguajes & Generalmente Java/.NET & Heterogéneo\\
Coreografía/Orquestación & ESB orquesta & Coreografía con eventos\\
Madurez del enfoque & 2000s & 2015+\\
Complejidad operacional & Alta (ESB) & Alta (distribuida)\\
\bottomrule
\end{tabularx}
\end{cajacompara}

\subsection{Cuándo NO usar microservicios}

Microservicios no es la solución para todo problema. Hay
escenarios donde adoptarlos prematuramente genera más complejidad
que valor. La caja siguiente enumera los casos donde un monolito
modular es la mejor opción.


\begin{cajaadvertencia}{El error costoso de los microservicios prematuros}
Sam Newman y Martin Fowler son enfáticos: los microservicios son
una \textbf{arquitectura distribuida}, y los sistemas distribuidos
son \emph{inherentemente complejos}. Los problemas que un monolito
resuelve trivialmente (transacciones ACID, debugging,
deployment) se vuelven retos no triviales en microservicios.

\textbf{No use microservicios si:}
\begin{itemize}
\item Su equipo es pequeño (menos de 8 personas).
\item No tiene CI/CD maduro y observabilidad robusta.
\item No domina contenedores y orquestación (Kubernetes).
\item No ha experimentado los dolores reales de un monolito grande.
\end{itemize}

\textbf{La recomendación moderna} (Fowler 2024) es \emph{«monolith
first»}: empezar con un monolito modular bien diseñado, y migrar
servicios a microservicios solo cuando se justifique con datos.
\end{cajaadvertencia}

\subsection{Patrones esenciales de microservicios}

Los microservicios traen consigo problemas distribuidos
(latencia, fallos parciales, consistencia eventual) que se mitigan
con patrones específicos. La lista siguiente recoge los esenciales.


\begin{table}[htbp]
\centering
\caption{Patrones imprescindibles en arquitectura de microservicios.}
\footnotesize
\begin{tabularx}{\textwidth}{lX}
\toprule
\textbf{Patrón} & \textbf{Propósito}\\
\midrule
API Gateway & Punto único de entrada; auth, rate-limiting, routing.\\
Service Discovery & Registro y descubrimiento dinámico (Eureka, Consul).\\
Circuit Breaker & Cortar llamadas a servicios fallidos (Resilience4j).\\
Saga & Transacciones distribuidas vía secuencia de eventos.\\
CQRS & Separar comandos (escritura) de queries (lectura).\\
Event Sourcing & El estado se deriva de un log de eventos inmutable.\\
Sidecar & Funciones transversales en contenedor adyacente.\\
Bulkhead & Aislamiento de recursos para evitar cascada de fallos.\\
\bottomrule
\end{tabularx}
\end{table}

\section{Patrón 4: Arquitectura dirigida por eventos (EDA)}

En EDA los componentes se comunican mediante \emph{eventos}
(mensajes) en lugar de invocaciones directas. Tres roles:

\begin{itemize}
\item \textbf{Producer}: publica eventos cuando algo relevante ocurre.
\item \textbf{Broker}: encamina eventos (RabbitMQ, Apache Kafka, AWS
SNS/SQS).
\item \textbf{Consumer}: se suscribe y procesa eventos
asíncronamente.
\end{itemize}

\begin{figure}[htbp]
\centering
\begin{tikzpicture}[
  every node/.style={font=\scriptsize,align=center},
  caja/.style={rectangle,draw=uteqverde,thick,minimum width=2cm,
               minimum height=0.8cm,fill=uteqverde!10,rounded corners=2pt},
  flecha/.style={->,>=stealth,thick}]

  \node[caja] (p1) {Servicio\\Pedidos};
  \node[caja,below=of p1] (p2) {Servicio\\Pagos};

  \node[regular polygon,regular polygon sides=8,
        draw=uteqdorado,thick,fill=uteqdorado!15,
        minimum width=2.5cm,right=2cm of p1,yshift=-0.7cm] (br) {};
  \node[font=\bfseries\scriptsize] at (br) {Event\\Broker\\(Kafka)};

  \node[caja,right=2cm of br,yshift=0.8cm] (c1) {Servicio\\Email};
  \node[caja,right=2cm of br] (c2) {Servicio\\Inventario};
  \node[caja,right=2cm of br,yshift=-0.8cm] (c3) {Servicio\\Analitica};

  \draw[flecha] (p1) -- node[above]{publica} (br);
  \draw[flecha] (p2) -- node[below]{publica} (br);
  \draw[flecha] (br) -- node[above]{} (c1);
  \draw[flecha] (br) -- node[above]{} (c2);
  \draw[flecha] (br) -- node[above]{} (c3);
\end{tikzpicture}
\caption{Flujo en una arquitectura dirigida por eventos: los servicios
publican eventos sin saber quién los consumirá; los consumers se
suscriben a los eventos que les interesan.}
\end{figure}

\textbf{Ventajas:} desacoplamiento temporal y espacial,
escalabilidad horizontal trivial, resiliencia ante caídas parciales.

\textbf{Costos:} complejidad operacional, debugging distribuido
difícil, consistencia eventual (no ACID), \emph{message ordering}
problemático.

\section{Patrón 5: Peer-to-peer (P2P)}

Cada nodo es simultáneamente cliente y servidor. Histórica (Napster,
Gnutella, BitTorrent), actualmente vigente en \emph{blockchain},
\emph{IPFS}, dApps y comunicaciones cifradas (Briar, Jami).

\textbf{Para aplicaciones web empresariales tradicionales, el modelo
P2P rara vez es la respuesta.} Lo mencionamos porque el sílabo lo
incluye y porque puede aparecer si el PFC tiene componentes
descentralizados.

\section{Documentación arquitectónica con el modelo C4 de Brown}

El modelo C4 \cite{brown2023c4} estructura la documentación
arquitectónica en cuatro niveles de abstracción crecientes en detalle.

\subsection{Los cuatro niveles del modelo C4}

El modelo C4 de Simon Brown propone documentar la
arquitectura en cuatro niveles de zoom: Contexto, Contenedores,
Componentes y Código. Cada nivel responde preguntas distintas. La
descripción siguiente los detalla.


\begin{figure}[htbp]
\centering
\begin{tikzpicture}[
  every node/.style={font=\scriptsize,align=center},
  nivel/.style={rectangle,draw=uteqverde,thick,fill=uteqverde!10,
                rounded corners=4pt,minimum width=4cm,minimum height=1.5cm}]

  \node[nivel] (n1) {\textbf{Nivel 1: Contexto}\\
                     Cómo el sistema encaja\\con usuarios y sistemas externos};
  \node[nivel,right=0.5cm of n1] (n2)
       {\textbf{Nivel 2: Contenedores}\\
        Apps web, APIs, BDs,\\colas de mensajes};
  \node[nivel,below=0.6cm of n1] (n3)
       {\textbf{Nivel 3: Componentes}\\
        Descomposición interna\\de cada contenedor};
  \node[nivel,right=0.5cm of n3] (n4)
       {\textbf{Nivel 4: Código}\\
        Diagrama de clases UML,\\opcional};

  \draw[->,>=stealth,thick,uteqverde!70!black] (n1) -- (n2);
  \draw[->,>=stealth,thick,uteqverde!70!black] (n2) -- (n3);
  \draw[->,>=stealth,thick,uteqverde!70!black] (n3) -- (n4);
\end{tikzpicture}
\caption{Los cuatro niveles del modelo C4. Cada nivel hace zoom sobre
un elemento del nivel superior. La mayoría de equipos solo necesita
los niveles 1, 2 y 3; el nivel 4 se genera automáticamente desde el
código.}
\end{figure}

\subsection{Ejemplo: diagrama de contexto del PFC}

El diagrama de contexto del PFC muestra el sistema como una
caja única y sus interacciones con usuarios y sistemas externos. Es
el primer artefacto que debe acompañar todo PFC.


\begin{figure}[htbp]
\centering
\begin{tikzpicture}[node distance=1.5cm and 2cm]
  \cfourPerson{est}{(0,0)}{Estudiante}
  \cfourPerson{adm}{(0,-2.5)}{Administrador}
  \cfourSystem{sis}{(5,-1.25)}{Sistema PFC\\(API + SPA Angular)}
  \cfourSystem{smtp}{(10,0)}{Servicio SMTP\\externo (SendGrid)}
  \cfourSystem{idp}{(10,-2.5)}{Identidad UTEQ\\(LDAP)}

  \draw[->,thick,>=stealth] (est) -- node[above,font=\tiny]{HTTPS} (sis);
  \draw[->,thick,>=stealth] (adm) -- node[below,font=\tiny]{HTTPS} (sis);
  \draw[->,thick,>=stealth] (sis) -- node[above,font=\tiny]{SMTP} (smtp);
  \draw[->,thick,>=stealth] (sis) -- node[above,font=\tiny]{LDAP} (idp);
\end{tikzpicture}
\caption{Nivel 1 (Contexto) del modelo C4 para el PFC: usuarios y
sistemas externos.}
\end{figure}

\subsection{Ejemplo: diagrama de contenedores del PFC}

El diagrama de contenedores zoom-in sobre la caja del nivel
1 y revela las piezas tecnológicas mayores: frontend SPA, backend
API, base de datos, cache, servicios externos. La figura siguiente
lo ilustra para el PFC.


\begin{figure}[htbp]
\centering
\begin{tikzpicture}[node distance=1.5cm and 1.8cm]
  \cfourPerson{usr}{(0,0)}{Usuario\\(Estudiante)}

  \cfourContainer{spa}{(4,0)}{Frontend\\Angular 19\\(SPA)}
  \cfourContainer{api}{(8,0)}{API REST\\Spring Boot 3.4\\(Tomcat 10.1)}
  \cfourDb{pg}{(12,0)}{PostgreSQL\\16}
  \cfourDb{rd}{(12,-2.5)}{Redis 7.4\\(cache)}

  \draw[->,thick,>=stealth] (usr.east) -- node[above,font=\tiny]{HTTPS} (spa.west);
  \draw[->,thick,>=stealth] (spa.east) -- node[above,font=\tiny]{JSON} (api.west);
  \draw[->,thick,>=stealth] (api.east) -- node[above,font=\tiny]{JDBC/SSL} (pg.west);
  \draw[->,thick,>=stealth] (api) -- node[right,font=\tiny]{TCP} (rd);
\end{tikzpicture}
\caption{Nivel 2 (Contenedores) del PFC: las unidades desplegables
y la tecnología de cada una.}
\end{figure}

\section{Architectural Decision Records (ADR)}

Un ADR documenta una decisión arquitectónica significativa.
Formato canónico (Michael Nygard, 2011) consta de cuatro secciones:
\textbf{Contexto}, \textbf{Decisión}, \textbf{Estado},
\textbf{Consecuencias}.

\begin{cajaadr}{ADR-001 --- Adoptar JWT en lugar de sesiones de servidor}
\textbf{Estado:} Aceptada el 2026-05-04.

\textbf{Contexto:} El frontend del PFC es una SPA Angular y el
backend es una API REST que requiere escalar horizontalmente con
Docker Compose en el laboratorio y, eventualmente, en cloud. Las
sesiones en memoria del servidor obligarían a \emph{sticky sessions}
en el balanceador o a un \emph{session store} compartido (Redis,
JDBC), lo que añade complejidad operacional.

\textbf{Decisión:} Adoptar autenticación con JSON Web Token (JWT)
firmado HS256 con secreto de 256 bits, transportado en cookie
\texttt{HttpOnly} \texttt{Secure} \texttt{SameSite=Strict}, con
\emph{refresh token} de 7 días y \emph{blacklist} en Redis para
revocación inmediata por JTI.

\textbf{Consecuencias:}
\begin{itemize}
\item[(+)] Backend stateless: cualquier instancia puede validar.
\item[(+)] Compatible con CORS sin sticky sessions.
\item[(+)] Incluye claims (rol, email) sin consulta extra a BD.
\item[(--)] Revocación requiere blacklist (mayor complejidad).
\item[(--)] El JWT es más grande que un session ID (\textasciitilde 350 bytes vs 32).
\item[(--)] Si se compromete el secreto, todos los tokens son
inválidos.
\end{itemize}

\textbf{Alternativas consideradas:}
\begin{itemize}
\item Sesiones server-side con Spring Session + Redis: descartada
por complejidad operacional desproporcionada.
\item OAuth 2.0 con servidor de autorización dedicado (Keycloak):
descartada por excesiva infraestructura para un PFC.
\end{itemize}
\end{cajaadr}

\section{Práctica integrada: documentar el PFC con C4 y ADRs}

La práctica de la semana es documentar el PFC del equipo con
diagramas C4 (niveles 1 y 2) y al menos dos ADRs justificando las
decisiones arquitectónicas mayores.


\begin{cajaejercicio}{Ejercicio 12.1 --- Entrega 2 del PFC --- Arquitectura Completa (Sumativa)}


\textbf{Enunciado.} Producir la Entrega 2 completa del PFC del equipo: arquitectura documentada con C4, API REST funcional con OpenAPI y frontend Angular consumiéndola. Sumativa principal del segundo corte.

\textbf{Esta es la \emph{Sumativa TA --- PFC Entrega 2: Arquitectura
Completa + API REST Documentada + Frontend Angular Funcional}.}

\subsubsection*{IDE y herramientas}

Antes de los pasos, conviene confirmar las herramientas que
el estudiante debe tener instaladas para la práctica.


\begin{itemize}
\item \textbf{IntelliJ IDEA Ultimate} (con plugins Spring Boot,
Database, Docker).
\item \textbf{Structurizr DSL} en línea (\url{https://www.structurizr.com/dsl})
o \textbf{draw.io} integrado en IntelliJ
(plugin \emph{Diagrams.net Integration}).
\item \textbf{Postman} para probar endpoints.
\end{itemize}

\subsubsection*{Paso 1 --- Generar diagrama C4 Nivel 1 (Contexto)}

En \url{https://www.structurizr.com/dsl}, crear un workspace y pegar (Listado~\ref{lst:12.2}):

\begin{lstlisting}[style=shell,numbers=none,caption={Ejemplo de Shell},label=lst:12.2]
workspace "PFC UTEQ" {
  model {
    estudiante = person "Estudiante"
    admin      = person "Administrador"
    sistema    = softwareSystem "Sistema PFC" {
      tags "App Principal"
    }
    smtp = softwareSystem "Servicio SMTP" {
      tags "Externo"
    }
    estudiante -> sistema "Usa via HTTPS"
    admin      -> sistema "Administra"
    sistema    -> smtp    "Envia notificaciones"
  }
  views {
    systemContext sistema "Contexto" {
      include *
      autoLayout
    }
  }
}
\end{lstlisting}

Pulsar \texttt{Render}: aparece el diagrama. Exportar como PNG.

\subsubsection*{Paso 2 --- Generar diagrama C4 Nivel 2 (Contenedores)}

Extender el modelo (Listado~\ref{lst:12.3}):

\begin{lstlisting}[style=shell,numbers=none,caption={Ejemplo de Shell},label=lst:12.3]
sistema = softwareSystem "Sistema PFC" {
  spa     = container "SPA Angular"     "Angular 19" "TypeScript"
  api     = container "API REST"        "Spring Boot 3.4" "Java"
  db      = container "Base de datos"   "PostgreSQL 16"  "BD" {
    tags "Database"
  }
  cache   = container "Cache"           "Redis 7.4"      "BD" {
    tags "Database"
  }
  spa -> api "JSON via HTTPS"
  api -> db  "JDBC/SSL"
  api -> cache "TCP"
}
\end{lstlisting}

\subsubsection*{Paso 3 --- Crear los 5 ADRs en el repositorio Git}

En el repositorio del PFC, crear directorio \texttt{docs/adr/}.
Crear estos archivos en formato Markdown:

\begin{itemize}
\item \texttt{ADR-001-JWT-vs-sesiones.md}
\item \texttt{ADR-002-PostgreSQL-vs-MySQL.md}
\item \texttt{ADR-003-Spring-Boot-vs-Quarkus.md}
\item \texttt{ADR-004-Redis-cache-aside.md}
\item \texttt{ADR-005-Docker-Compose-deployment.md}
\end{itemize}

Cada ADR sigue la plantilla de Nygard mostrada en este capítulo.

\subsubsection*{Paso 4 --- Tabla de atributos de calidad ISO 25010}

Documentar en el informe técnico esta tabla:

\begin{tabularx}{\textwidth}{llXX}
\toprule
\textbf{Atributo} & \textbf{Prioridad} & \textbf{Escenario} & \textbf{Estrategia}\\
\midrule
Seguridad & Crítica & SQL Injection en endpoints & JPA + Bean Validation\\
Eficiencia & Alta & Listado de productos < 200ms & Cache Redis con TTL 60s\\
Usabilidad & Alta & Login intuitivo & Angular con feedback visual\\
Mantenibilidad & Alta & Nuevo dev productivo en 1 día & C4 + ADRs + README\\
Fiabilidad & Media & Disponibilidad 99\% & Health checks + Docker restart\\
Portabilidad & Media & Mover de Windows a Linux & Docker Compose\\
\bottomrule
\end{tabularx}

\subsubsection*{Paso 5 --- API REST documentada con Swagger UI}

Agregar al \texttt{pom.xml} (Listado~\ref{lst:12.4}):

\begin{lstlisting}[style=xml,numbers=none,caption={Ejemplo de XML},label=lst:12.4]
<dependency>
  <groupId>org.springdoc</groupId>
  <artifactId>springdoc-openapi-starter-webmvc-ui</artifactId>
  <version>2.6.0</version>
</dependency>
\end{lstlisting}

Reiniciar la aplicación. Abrir
\url{http://localhost:8080/swagger-ui.html}: aparece toda la API
documentada y probable interactivamente.

\subsubsection*{Paso 6 --- Frontend Angular conectado}

Crear el frontend con Angular CLI (Listado~\ref{lst:12.5}):
\begin{lstlisting}[style=shell,numbers=none,caption={Ejemplo de Shell},label=lst:12.5]
ng new frontend --standalone --routing --style=scss
cd frontend                                                 # cambia directorio
ng generate service auth
ng generate component login --standalone
\end{lstlisting}

Implementar el flujo de login que llame al endpoint
\texttt{POST /api/auth/login} y guarde el JWT en cookie HttpOnly
(que el backend establece).

\subsubsection*{Paso 7 --- Empaquetar todo con Docker Compose}

Crear \texttt{docker-compose.yml} en la raíz (Listado~\ref{lst:12.6}):

\begin{lstlisting}[style=yaml,numbers=none,caption={Ejemplo de YAML},label=lst:12.6]
services:
  postgres:
    image: postgres:16-alpine
    environment:
      POSTGRES_USER: appweb
      POSTGRES_PASSWORD: appweb
      POSTGRES_DB:   appweb
    ports: ["5432:5432"]
    volumes: ["pgdata:/var/lib/postgresql/data"]
    healthcheck:
      test: ["CMD","pg_isready","-U","appweb"]
      interval: 10s

  redis:                                                    # configuracion del cliente Redis
    image: redis:7-alpine                                   # configuracion del cliente Redis
    ports: ["6379:6379"]

  api:
    build: ./backend
    depends_on:
      postgres: { condition: service_healthy }
      redis:    { condition: service_started }              # configuracion del cliente Redis
    environment:
      SPRING_DATASOURCE_URL: jdbc:postgresql://postgres:5432/appweb
      SPRING_REDIS_HOST: redis
    ports: ["8080:8080"]

  web:
    build: ./frontend
    depends_on: [api]
    ports: ["80:80"]

volumes:
  pgdata:
\end{lstlisting}

\subsubsection*{Paso 8 --- Verificar el despliegue completo}

Como verificación final, el estudiante debe levantar el
sistema completo y validar que funciona end-to-end antes de
considerar terminada la entrega. V\'ease el Listado~\ref{lst:12.7}.


\begin{lstlisting}[style=shell,numbers=none,caption={Ejemplo de Shell},label=lst:12.7]
docker compose up -d                                        # levanta los servicios del compose
docker compose ps    # todos healthy
curl http://localhost:8080/actuator/health
# Abrir http://localhost en el navegador: SPA cargada
\end{lstlisting}

\subsubsection*{Entregables}

La entrega de la práctica de la semana 12 debe incluir los
siguientes artefactos consolidados.


\begin{enumerate}
\item Repositorio Git con tag \texttt{v0.5.0}.
\item Documento PDF (15--25 páginas) con: portada UTEQ, diagrama C4
nivel 1 y 2 (PNG), tabla atributos calidad, los 5 ADRs, capturas de
Swagger UI funcional, capturas de la SPA Angular en login,
\texttt{docker-compose.yml} comentado.
\item Demostración en clase: \texttt{docker compose up} arranca todo
en menos de 2 minutos.
\end{enumerate}

\textbf{Esta es la \emph{Sumativa de la semana 12}: PFC Entrega 2.}
\end{cajaejercicio}

\section{Buenas prácticas profesionales en arquitectura}

La caja siguiente compila quince reglas profesionales de
arquitectura que el estudiante debe respetar al diseñar cualquier
sistema.


\begin{cajabuenas}{Doce reglas profesionales 2026}
\begin{enumerate}
\item \emph{Monolith first}: empezar con monolito modular bien
diseñado.
\item Documentar TODA decisión significativa con un ADR.
\item Diagramas C4 nivel 1 y 2 obligatorios; nivel 3 si el sistema
es grande.
\item Atributos de calidad explícitos al inicio del proyecto.
\item Versionar la documentación arquitectónica como código (Git).
\item Revisar la arquitectura cada trimestre con todo el equipo.
\item Evitar acoplamiento entre servicios vía base de datos
compartida.
\item Cada servicio dueño de su esquema (en microservicios).
\item Comunicación asíncrona cuando se pueda; síncrona solo cuando
sea inevitable.
\item Health checks y observabilidad desde el día 1.
\item Tests de arquitectura con ArchUnit (Java) o NetArchTest (.NET).
\item Pensar en \emph{evolutionary architecture}: cambiar es la
norma, no la excepción.
\end{enumerate}
\end{cajabuenas}

\section{Ejercicio resuelto adicional con resultado visual}

Como cierre, el siguiente ejercicio resuelto adicional
incluye un caso completo con diagramas y ADR; lo acompaña un
propuesto.


\begin{cajaejemplo}{Ejercicio resuelto 12.1 --- Diagrama C4 nivel 1 y 2 de una tienda en línea}
\textbf{Enunciado.} Modelar arquitectónicamente con el modelo C4 una
tienda en línea sencilla que tiene: SPA Angular para clientes,
backend Spring Boot REST, PostgreSQL, Redis para caché de catálogo y
servicio externo de pagos (Stripe). Producir los diagramas de nivel 1
(Sistema-Contexto) y nivel 2 (Contenedores).

\textbf{Nivel 1 --- Diagrama de Contexto del Sistema}
(representación textual estilo ASCII; en producción se renderiza
con PlantUML o Structurizr):

\begin{center}
\fbox{\begin{minipage}{0.9\textwidth}
\vspace{0.2cm}\centering\footnotesize
\begin{tabular}{c}
\fbox{\textbf{Cliente}}\\[0.1cm]
$\downarrow$\\[0.1cm]
\fbox{\textbf{Sistema TiendaUTEQ}}\\[0.1cm]
$\downarrow$ \\[0.1cm]
\fbox{\textbf{Stripe API}}\\
\end{tabular}\\[0.2cm]
\textit{El cliente interactúa con TiendaUTEQ; el sistema consume Stripe para pagos.}
\vspace{0.1cm}
\end{minipage}}
\end{center}

\textbf{Nivel 2 --- Diagrama de Contenedores:}

\begin{figure}[H]
\centering
\begin{tikzpicture}[
  every node/.style={font=\footnotesize,align=center},
  persona/.style={circle,draw=uteqverde,thick,fill=uteqverde!15,minimum size=1.4cm},
  contenedor/.style={rectangle,draw=uteqverde,thick,fill=uteqverde!10,
                     rounded corners=4pt,minimum width=2.8cm,minimum height=1.2cm},
  externo/.style={rectangle,draw=uteqdorado,thick,fill=uteqdorado!15,
                  dashed,minimum width=2.8cm,minimum height=1.2cm}]
  \node[persona] (u) {Cliente};
  \node[contenedor,right=1cm of u] (spa) {SPA Angular\\(Browser)};
  \node[contenedor,right=1cm of spa] (api) {API REST\\(Spring Boot)};
  \node[contenedor,below=0.6cm of api] (db) {BD\\(PostgreSQL)};
  \node[contenedor,left=0.6cm of db] (redis) {Cache\\(Redis)};
  \node[externo,right=1cm of api] (stripe) {Stripe};

  \draw[->,thick] (u) -- node[above,font=\tiny]{HTTPS} (spa);
  \draw[->,thick] (spa) -- node[above,font=\tiny]{JSON/HTTPS} (api);
  \draw[->,thick] (api) -- node[right,font=\tiny]{JDBC} (db);
  \draw[->,thick] (api) -- node[above,font=\tiny]{TCP} (redis);
  \draw[->,thick] (api) -- node[above,font=\tiny]{HTTPS} (stripe);
\end{tikzpicture}
\caption{Diagrama C4 nivel 2 (contenedores) de TiendaUTEQ.}
\end{figure}

\textbf{ADR-001 correspondiente:}

\begin{cajaadr}{ADR-001 --- Adopción de arquitectura por capas con caché Redis}
\textbf{Estado:} Aceptado el 5 de mayo de 2026.

\textbf{Contexto.} TiendaUTEQ requiere servir hasta 500
usuarios concurrentes con tiempos de respuesta $<$ 300 ms en el
listado de productos. La BD PostgreSQL muestra latencias de 50--80 ms
en consultas frecuentes.

\textbf{Decisión.} Adoptar arquitectura por capas (presentación,
aplicación, dominio, infraestructura) con Spring Boot y agregar caché
Redis al servicio de catálogo con TTL de 60 segundos.

\textbf{Consecuencias positivas:}
\begin{itemize}
\item Tiempos de respuesta $<$ 50 ms tras el primer hit.
\item Menor carga sobre PostgreSQL.
\end{itemize}

\textbf{Consecuencias negativas:}
\begin{itemize}
\item Complejidad operativa adicional (un contenedor más en Docker).
\item Posibilidad de inconsistencia transitoria.
\end{itemize}

\textbf{Alternativas consideradas:}
\begin{itemize}
\item \emph{Caffeine in-process}: descartado por no compartir caché
entre instancias.
\item \emph{Memcached}: descartado por inferior soporte de
estructuras de datos respecto a Redis.
\end{itemize}
\end{cajaadr}

\textbf{Discusión.} El estudiante debe notar que el diagrama no
muestra clases ni código: el nivel arquitectónico se ocupa de las
piezas grandes y las decisiones difíciles de revertir
\cite{bass2021software}. El ADR documenta la decisión para futuros
mantenedores que se pregunten \emph{«¿por qué Redis?»} en lugar de
adivinar.
\end{cajaejemplo}

\begin{cajaejercicio}{Ejercicio propuesto 12.2 --- Migrar un monolito a arquitectura hexagonal}
\textbf{Enunciado.} Tomar el PFC del equipo (o el monolito tienda en
línea anterior) e identificar los \emph{puertos} y \emph{adaptadores}
para refactorizarlo a arquitectura hexagonal (\emph{Ports and
Adapters} de Alistair Cockburn).

\textbf{Procedimiento:}
\begin{enumerate}
\item Identificar el \textbf{dominio puro}: entidades y servicios
de negocio sin dependencia de Spring, JPA, REST, etc.
\item Definir \textbf{puertos primarios} (interfaces que el dominio
expone para que lo invoquen): casos de uso.
\item Definir \textbf{puertos secundarios} (interfaces que el dominio
necesita): repositorios, servicios externos, notificación.
\item Crear \textbf{adaptadores primarios}: controladores REST,
listeners de eventos, CLI.
\item Crear \textbf{adaptadores secundarios}: implementaciones JPA
de repositorios, clientes HTTP de servicios externos, productores de
eventos.
\item Reorganizar el proyecto en Maven multi-módulo:
\texttt{dominio}, \texttt{aplicacion}, \texttt{adaptadores-rest},
\texttt{adaptadores-persistencia}.
\item Producir test del dominio que NO requieran arrancar Spring
(unit tests puros, $<$ 100 ms cada uno).
\end{enumerate}

\textbf{Criterios de evaluación:}
\begin{itemize}
\item El módulo \texttt{dominio} no tiene ninguna dependencia a
Spring, JPA ni infraestructura.
\item Existen tests del dominio que se ejecutan sin Spring
(\texttt{@SpringBootTest} no aparece).
\item Hay al menos dos adaptadores secundarios alternativos
(ej. \texttt{InMemoryUsuarioRepository} para tests +
\texttt{JpaUsuarioRepository} para producción).
\item Documentación con diagrama de la arquitectura y ADR-002 con la
decisión.
\end{itemize}

\textbf{Lecturas recomendadas:} \cite{ford2017building} para
\emph{evolutionary architecture}; \cite{newman2021building} sobre
descomposición incremental.
\end{cajaejercicio}

% =====================================================================
%       SEMANA 13: DISEÑO DE ARQUITECTURAS WEB ESCALABLES
% =====================================================================
\chapter{Diseño de arquitecturas web escalables}
\label{cap:semana13}

\epigraph{«Hay solo dos cosas difíciles en informática: invalidar la
caché y nombrar las cosas.»}{--- \textit{Phil Karlton}}

\section*{Resultado de aprendizaje de la semana}
Al concluir la semana 13, el estudiante diseña arquitecturas web
escalables aplicando estrategias de balanceo de carga, patrones de
caché distribuida con Redis, técnicas de evaluación arquitectónica
(ATAM) y comprende los \emph{trade-offs} entre escalabilidad,
consistencia y disponibilidad (teorema CAP).

\section{Escalabilidad en aplicaciones web: vertical vs horizontal}

Cuando una aplicación gana usuarios, llega el momento de
escalar. Hay dos estrategias fundamentales, vertical y horizontal,
con implicaciones técnicas y económicas muy diferentes.


\subsection{Definiciones formales}

La distinción entre escalado vertical y horizontal es la
base del razonamiento sobre arquitectura escalable. Definiciones
formales y casos típicos a continuación.


\begin{cajadefinicion}{Las dos dimensiones de la escalabilidad}
\textbf{Escalabilidad vertical} (\emph{scale up}): aumentar la
capacidad del mismo servidor con más CPU, más RAM, más disco. Se
detiene en algún momento (Ley de Moore agotada, costos exponenciales).

\textbf{Escalabilidad horizontal} (\emph{scale out}): añadir más
servidores idénticos detrás de un balanceador de carga. Crece
linealmente hasta el límite del componente menos escalable
(típicamente la BD).
\end{cajadefinicion}

\begin{table}[htbp]
\centering
\caption{Comparativa: escalado vertical vs horizontal.}
\footnotesize
\begin{tabularx}{\textwidth}{lXX}
\toprule
\textbf{Aspecto} & \textbf{Vertical} & \textbf{Horizontal}\\
\midrule
Costo inicial & Bajo (un servidor mayor) & Alto (infraestructura, balanceador)\\
Costo a gran escala & Exponencial & Lineal\\
Complejidad arquitectónica & Mínima & Alta (sesiones, consistencia)\\
Tolerancia a fallos & Baja (SPOF) & Alta (réplicas)\\
Aplicabilidad a BD & Limitada & Sharding requerido\\
Tiempo para añadir capacidad & Minutos (en cloud) & Segundos (autoscaling)\\
Lím. teórico & Capacidad máxima del hw & Prácticamente ilimitado\\
\bottomrule
\end{tabularx}
\end{table}

\subsection{El teorema CAP de Brewer}

El teorema CAP de Eric Brewer ---imposibilidad de tener
consistencia, disponibilidad y tolerancia a particiones
simultáneamente--- es uno de los resultados más citados en
arquitectura distribuida. Su interpretación práctica es más matizada
de lo que parece.


\begin{cajadefinicion}{Teorema CAP --- Eric Brewer 2000}
En un sistema distribuido, ante un \textbf{P}article \textbf{P}artition
(falla de red), un sistema solo puede garantizar
\emph{a lo sumo} dos de las tres propiedades siguientes:

\begin{description}[leftmargin=1.5cm]
\item[Consistency (C).] Toda lectura recibe el dato más reciente.
\item[Availability (A).] Toda petición recibe respuesta (no errores).
\item[Partition tolerance (P).] El sistema sigue funcionando ante
fallos de red entre nodos.
\end{description}

Como las particiones de red en sistemas distribuidos son inevitables,
en la práctica se elige entre \textbf{CP} (consistencia sobre
disponibilidad: PostgreSQL replicado) o \textbf{AP} (disponibilidad
sobre consistencia, con consistencia eventual: Cassandra, DynamoDB).
\end{cajadefinicion}

\section{Estrategias de balanceo de carga}

Un balanceador de carga distribuye las peticiones entrantes
entre múltiples instancias del backend. Es la pieza fundamental para
escalar horizontalmente.


\begin{figure}[htbp]
\centering
\begin{tikzpicture}[
  every node/.style={font=\scriptsize,align=center},
  caja/.style={rectangle,draw=uteqverde,thick,minimum width=1.8cm,
               minimum height=0.8cm,fill=uteqverde!10,rounded corners=2pt}]

  \node[caja] (cli) {Cliente};
  \node[caja,right=1.5cm of cli] (lb) {Load\\Balancer};

  \node[caja,right=1.5cm of lb,yshift=1.5cm] (a1) {App\\Server 1};
  \node[caja,right=1.5cm of lb] (a2) {App\\Server 2};
  \node[caja,right=1.5cm of lb,yshift=-1.5cm] (a3) {App\\Server 3};

  \node[cylinder,draw=uteqverde,fill=uteqverde!15,thick,
        shape border rotate=90,aspect=0.3,minimum width=1.5cm,
        minimum height=1.5cm,right=1.5cm of a2] (db) {BD};

  \draw[->,thick,>=stealth] (cli) -- (lb);
  \draw[->,thick,>=stealth] (lb) -- (a1);
  \draw[->,thick,>=stealth] (lb) -- (a2);
  \draw[->,thick,>=stealth] (lb) -- (a3);
  \draw[->,thick,>=stealth] (a1.east) -- (db.west);
  \draw[->,thick,>=stealth] (a2) -- (db);
  \draw[->,thick,>=stealth] (a3.east) -- (db.west);
\end{tikzpicture}
\caption{Topología clásica con balanceador de carga: el LB recibe las
peticiones del cliente y las distribuye entre N servidores de
aplicación que comparten la BD.}
\end{figure}

\subsection{Algoritmos de balanceo}

Distintos algoritmos de balanceo tienen comportamientos
diferentes ante distintas cargas. La tabla siguiente los cataloga.


\begin{table}[htbp]
\centering
\caption{Algoritmos de balanceo de carga.}
\footnotesize
\begin{tabularx}{\textwidth}{lXX}
\toprule
\textbf{Algoritmo} & \textbf{Descripción} & \textbf{Mejor uso}\\
\midrule
Round-robin & Distribución uniforme rotativa & Servidores homogéneos\\
Least connections & Al servidor con menos conexiones activas & Carga heterogénea\\
Least response time & Al más rápido & APIs con latencia variable\\
IP hash & Hash IP $\to$ servidor (sticky) & Sesiones en memoria\\
Weighted round-robin & RR con pesos por servidor & Servidores asimétricos\\
Random & Aleatorio uniforme & Pruebas, casos especiales\\
\bottomrule
\end{tabularx}
\end{table}

\subsection{Implementaciones populares}

\textbf{Nginx} es el balanceador más usado para aplicaciones web:

\begin{cajaejemplo}{Ejemplo 13.1 --- Nginx como load balancer}

\textbf{Enunciado.} Implementar el patrón cache-aside con Spring Boot 3.4 y Redis 7 para un servicio de catálogo. Medir la latencia antes y después con \texttt{k6}.

\begin{lstlisting}[style=shell,caption={nginx.conf},label=lst:13.1]
upstream app_pool {
    least_conn;
    server 10.0.1.10:8080 max_fails=3 fail_timeout=30s;
    server 10.0.1.11:8080 max_fails=3 fail_timeout=30s;
    server 10.0.1.12:8080 max_fails=3 fail_timeout=30s;
    keepalive 32;
}

server {
    listen 443 ssl http2;
    server_name app.uteq.edu.ec;

    ssl_certificate     /etc/ssl/certs/uteq.crt;
    ssl_certificate_key /etc/ssl/private/uteq.key;
    ssl_protocols       TLSv1.3 TLSv1.2;

    # Cabeceras de seguridad
    add_header Strict-Transport-Security "max-age=31536000" always;
    add_header X-Frame-Options "DENY" always;

    # Comprimir respuestas
    gzip on;
    gzip_types application/json text/css application/javascript;

    location / {
        proxy_pass http://app_pool;
        proxy_set_header Host $host;
        proxy_set_header X-Real-IP $remote_addr;
        proxy_set_header X-Forwarded-For $proxy_add_x_forwarded_for;
        proxy_set_header X-Forwarded-Proto $scheme;
        proxy_connect_timeout 5s;
        proxy_read_timeout    30s;
    }

    # Health check endpoint sin pasar al backend
    location = /lb-health {
        access_log off;
        return 200 'healthy';
    }
}
\end{lstlisting}
\end{cajaejemplo}

\section{Patrones de caché en aplicaciones web}

Una caché bien colocada puede reducir la latencia y la
carga sobre la base de datos en uno o dos órdenes de magnitud. Pero
las cachés son ``armas de doble filo'': mal usadas, producen
inconsistencias que erosionan la confianza del usuario.


\subsection{Los cinco patrones canónicos de caché}

La literatura técnica describe cinco patrones canónicos de
caché, cada uno con su perfil de coherencia y rendimiento. La tabla
siguiente los contrasta.


\begin{table}[htbp]
\centering
\caption{Patrones de caché y sus características.}
\footnotesize
\begin{tabularx}{\textwidth}{lXX}
\toprule
\textbf{Patrón} & \textbf{Mecánica} & \textbf{Cuándo usarlo}\\
\midrule
Cache-aside (lazy loading) & App lee caché; en miss, lee BD y guarda en caché & General; lo más común\\
Read-through & Caché lee BD automáticamente en miss & Cuando se quiere transparencia\\
Write-through & Cada write va a caché y BD sincrónicamente & Consistencia estricta\\
Write-behind & Write a caché; flush a BD asíncrono & Alto throughput de escritura\\
Refresh-ahead & Refresca proactivamente antes de TTL & Datos predecibles, alta latencia\\
\bottomrule
\end{tabularx}
\end{table}

\subsection{Cache-aside con Redis: el patrón dominante}

Cache-aside es el patrón dominante en aplicaciones web
modernas con Redis o Memcached. Su mecánica es simple pero requiere
disciplina para no introducir inconsistencias.


\begin{figure}[htbp]
\centering
\begin{tikzpicture}[
  every node/.style={font=\scriptsize,align=center},
  caja/.style={rectangle,draw=uteqverde,thick,minimum width=2cm,
               minimum height=0.9cm,fill=uteqverde!10,rounded corners=2pt},
  flecha/.style={->,>=stealth,thick}]

  \node[caja] (app) {Aplicación};
  \node[cylinder,draw=red!70,fill=red!10,thick,
        shape border rotate=90,aspect=0.3,minimum width=1.5cm,
        minimum height=1.5cm,right=2.5cm of app] (re) {Redis};
  \node[cylinder,draw=blue!70,fill=blue!10,thick,
        shape border rotate=90,aspect=0.3,minimum width=1.5cm,
        minimum height=1.5cm,right=2.5cm of re] (db) {Postgres};

  \draw[flecha] ([yshift=0.4cm]app.east) --
    node[above,font=\tiny]{1. GET key} ([yshift=0.4cm]re.west);
  \draw[flecha] ([yshift=-0.4cm]re.west) --
    node[below,font=\tiny]{2a. HIT $\to$ value}
    ([yshift=-0.4cm]app.east);
  \draw[flecha] (re.east) --
    node[above,font=\tiny]{2b. MISS} (db.west);
  \draw[flecha,uteqdorado!80!black] ([yshift=-0.3cm]db.west) --
    node[below,font=\tiny]{3. row} ([yshift=-0.3cm]re.east);
  \draw[flecha,uteqdorado!80!black] (re) --
    node[above,font=\tiny,yshift=0.6cm]{4. SET key TTL}
    ([yshift=0cm]re.west);
\end{tikzpicture}
\caption{Patrón cache-aside: la aplicación consulta Redis primero;
si hay HIT, devuelve. Si hay MISS, consulta PostgreSQL, guarda en
Redis con TTL y devuelve.}
\end{figure}

\begin{cajaejemplo}{Ejemplo 13.2 --- Cache-aside con Spring Boot --- Redis y PostgreSQL}


\textbf{Enunciado.} Implementar el patrón cache-aside completo en una aplicación Spring Boot con Redis para cache, PostgreSQL como origen de datos y endpoints REST que permitan medir la diferencia de latencia.

\textbf{Configuración \texttt{application.yml}:}
\begin{lstlisting}[style=yaml,numbers=none,caption={Ejemplo de YAML},label=lst:13.2]
spring:                                                     # configuracion de Spring
  data:
    redis:                                                  # configuracion del cliente Redis
      host: localhost                                       # host del servicio
      port: 6379                                            # puerto donde escucha la aplicacion
      timeout: 2s
      lettuce:
        pool:
          max-active: 16
          max-idle:   8
\end{lstlisting}

\textbf{Servicio con cache-aside (Listado~\ref{lst:13.3}):}
\begin{lstlisting}[style=java,caption={ProductoService.java},label=lst:13.3]
package ec.edu.uteq.appweb.service;                         // paquete del archivo

import com.fasterxml.jackson.databind.ObjectMapper;         // importacion de clase externa
import ec.edu.uteq.appweb.model.Producto;                   // importacion de clase externa
import ec.edu.uteq.appweb.repository.ProductoRepository;    // importacion de clase externa
import lombok.RequiredArgsConstructor;                      // importacion de clase externa
import lombok.extern.slf4j.Slf4j;                           // importacion de clase externa
// importacion de clase externa
// importacion de clase externa
import org.springframework.data.redis.core.StringRedisTemplate;
import org.springframework.stereotype.Service;              // importacion de clase externa
import java.time.Duration;                                  // importacion de clase externa
import java.util.NoSuchElementException;                    // importacion de clase externa

@Service                                                    // capa de servicio (logica de negocio)
@Slf4j
@RequiredArgsConstructor
public class ProductoService {                              // declaracion de clase publica

    private final ProductoRepository repo;                  // campo inmutable (mejor practica)
    private final StringRedisTemplate redis;                // campo inmutable (mejor practica)
    private final ObjectMapper json;                        // campo inmutable (mejor practica)

    // campo privado
    // campo privado
    private static final Duration TTL = Duration.ofMinutes(10);

    public Producto obtener(Long id) throws Exception {     // metodo publico
        String key = "producto:" + id;

        // 1. Intentar cache
        String cached = redis.opsForValue().get(key);
        if (cached != null) {                               // condicional
            log.debug("Cache HIT: {}", key);
            return json.readValue(cached, Producto.class);  // devuelve el resultado al llamador
        }

        // 2. Cache miss: ir a BD
        log.debug("Cache MISS: {}", key);
        Producto p = repo.findById(id).orElseThrow(
            () -> new NoSuchElementException("Producto " + id));

        // 3. Guardar en cache para proximas peticiones
        redis.opsForValue().set(key, json.writeValueAsString(p), TTL);
        return p;                                           // devuelve el resultado al llamador
    }

    /**
     * Al actualizar, INVALIDAR la cache (no actualizarla).
     * Razon: si la operacion falla luego de actualizar la cache,
     * habria inconsistencia. Mas seguro borrar y dejar que la
     * proxima lectura repueble.
     */
    public Producto actualizar(Long id, Producto datos) {   // metodo publico
        Producto guardado = repo.save(datos);
        redis.delete("producto:" + id);
        log.debug("Cache invalidada: producto:{}", id);
        return guardado;                                    // devuelve el resultado al llamador
    }
}
\end{lstlisting}
\end{cajaejemplo}

\subsection{Estrategias de invalidación}

Mantener la caché coherente con la base de datos es el
desafío central: ¿cuándo se invalida la entrada? Cuatro estrategias
canónicas se describen a continuación.


\begin{cajabuenas}{Cuatro estrategias de invalidación de caché}
\begin{enumerate}
\item \textbf{TTL (Time To Live)}: la entrada expira después de
N segundos. Simple, eficaz, fuerza cierta inconsistencia temporal.
\item \textbf{Invalidación por evento}: al modificar la BD se publica
un evento que limpia la caché. Más consistente, más complejo.
\item \textbf{Cache stampede protection}: cuando muchos clientes
piden el mismo dato simultáneamente y la cache expira, todos
golpean la BD. Solución: mutex distribuido (\texttt{SETNX} en Redis).
\item \textbf{Política de evicción}: LRU (Least Recently Used), LFU
(Least Frequently Used), FIFO. Redis usa LRU por defecto cuando se
llena la memoria.
\end{enumerate}
\end{cajabuenas}

\section{Patrones avanzados: CDN, Edge caching y read replicas}

Más allá de la caché de aplicación con Redis, hay otras
capas de caché que conviene conocer para escalar correctamente:
CDN, \emph{edge caching} y réplicas de lectura.


\subsection{CDN para contenido estático}

Un CDN (Content Delivery Network) replica el contenido estático
(HTML, CSS, JS, imágenes, videos) en servidores edge geográficamente
distribuidos. Cloudflare, AWS CloudFront, Fastly y Bunny son
proveedores frecuentes en 2026.

\subsection{Réplicas de lectura para BD}

Cuando el SGBD se vuelve cuello de botella en lecturas:

\begin{itemize}
\item \textbf{Master} acepta escrituras y replica asíncronamente.
\item \textbf{Réplicas read-only} sirven lecturas (90\% del tráfico
típico).
\item \textbf{Sharding} para escalar escrituras: particionar por clave.
\end{itemize}

\section{Métodos de evaluación arquitectónica}

Tomar decisiones arquitectónicas exige criterios objetivos.
Métodos formales como ATAM y TARA permiten evaluar arquitecturas
contra los \emph{atributos de calidad} que el sistema debe satisfacer.


\subsection{ATAM (Architecture Tradeoff Analysis Method)}

\textbf{ATAM}, desarrollado por el Software Engineering Institute de
Carnegie Mellon, es el método de evaluación arquitectónica más
adoptado en la industria. Se desarrolla en 9 pasos en sesiones
grupales con stakeholders:

\begin{enumerate}
\item Presentación de ATAM al equipo.
\item Presentación del negocio.
\item Presentación de la arquitectura.
\item Identificación de approaches arquitectónicos.
\item Generación del árbol de utilidad de calidad.
\item Análisis de approaches arquitectónicos.
\item \emph{Brainstorming} de escenarios de calidad.
\item Análisis de approaches arquitectónicos (segunda iteración).
\item Presentación de resultados.
\end{enumerate}

\subsection{Lightweight ATAM para equipos pequeños}

Para PFCs y proyectos universitarios, una versión ligera (3 horas en
lugar de 3 días):

\begin{enumerate}
\item Listar 3--5 atributos de calidad prioritarios.
\item Para cada uno, escribir 2--3 escenarios concretos de uso.
\item Evaluar la arquitectura propuesta contra cada escenario:
ROJO (no soportado), AMARILLO (parcial), VERDE (soportado).
\item Identificar puntos sensibles (afectan varios atributos) y
puntos de \emph{tradeoff}.
\item Documentar conclusiones en una página.
\end{enumerate}

\section{Práctica integrada: implementar cache-aside con benchmark}

La práctica de la semana implementa el patrón cache-aside
con Spring Boot y Redis, mide la latencia antes y después con k6, y
produce una tabla comparativa.


\begin{cajaejercicio}{Ejercicio 13.1 --- Cache Redis con benchmark medible (Sumativa)}

\textbf{Objetivo:} demostrar empíricamente el impacto del cache Redis
sobre el tiempo de respuesta.

\subsubsection*{IDE: IntelliJ IDEA Ultimate}

La práctica usa IntelliJ IDEA Ultimate para el código y
Docker Compose para los servicios. Configuración recomendada a
continuación.


\subsubsection*{Paso 1 --- Levantar PostgreSQL y Redis con Docker}

El primer paso es levantar los dos servicios de
infraestructura con un único archivo Docker Compose. V\'ease el Listado~\ref{lst:13.4}.


\begin{lstlisting}[style=shell,numbers=none,caption={Ejemplo de Shell},label=lst:13.4]
docker run -d --name pg-bench  -e POSTGRES_USER=appweb \    # arranca un contenedor
  -e POSTGRES_PASSWORD=appweb -e POSTGRES_DB=appweb \
  -p 5432:5432 postgres:16-alpine
# arranca un contenedor
# arranca un contenedor
docker run -d --name redis-bench -p 6379:6379 redis:7-alpine
\end{lstlisting}

Verificar (Listado~\ref{lst:13.5}):
\begin{lstlisting}[style=shell,numbers=none,caption={Ejemplo de Shell},label=lst:13.5]
docker exec -it redis-bench redis-cli ping     # PONG
# ejecuta comando en contenedor activo
docker exec -it pg-bench psql -U appweb -c "SELECT 1;"      # ejecuta comando en contenedor activo
\end{lstlisting}

\subsubsection*{Paso 2 --- Crear proyecto Spring Boot con Redis}

\texttt{File $\rightarrow$ New Project $\rightarrow$ Spring Boot}.
Dependencias: \emph{Spring Web, Spring Data JPA, PostgreSQL, Spring
Data Redis, Lombok, DevTools}.

\subsubsection*{Paso 3 --- Tabla y datos de prueba (10 mil productos)}

Para que el efecto del caché sea visible, la tabla debe
tener un volumen suficiente: al menos 10 mil filas con datos
realistas. El script SQL siguiente los genera. V\'ease el Listado~\ref{lst:13.6}.


\begin{lstlisting}[style=sql,numbers=none,caption={Ejemplo de SQL},label=lst:13.6]
-- crea tabla nueva en la BD
CREATE TABLE productos_bench (                              -- crea tabla nueva en la BD
  id      BIGSERIAL PRIMARY KEY,
  -- columna obligatoria (no admite NULL)
  nombre  VARCHAR(120) NOT NULL,                            -- columna obligatoria (no admite NULL)
  -- columna obligatoria (no admite NULL)
  precio  NUMERIC(10,2) NOT NULL,                           -- columna obligatoria (no admite NULL)
  -- columna obligatoria (no admite NULL)
  stock   INTEGER       NOT NULL                            -- columna obligatoria (no admite NULL)
);
-- inserta una fila nueva
INSERT INTO productos_bench (nombre, precio, stock)         -- inserta una fila nueva
-- consulta datos
SELECT 'P-' || i,                                           -- consulta datos
       ROUND((random()*100+1)::numeric, 2),
       (random()*100)::int
-- tabla de origen
FROM generate_series(1, 10000) i;                           -- tabla de origen
\end{lstlisting}

\subsubsection*{Paso 4 --- Endpoints SIN y CON cache}

Crear dos endpoints idénticos en lógica pero distintos en si usan
caché. El \texttt{/sin-cache/\{id\}} consulta siempre la BD; el
\texttt{/con-cache/\{id\}} usa el patrón cache-aside.

\subsubsection*{Paso 5 --- Benchmark con \texttt{wrk}}

Instalar \texttt{wrk} (\url{https://github.com/wg/wrk}). Ejecutar (Listado~\ref{lst:13.7}):

\begin{lstlisting}[style=shell,numbers=none,caption={Ejemplo de Shell},label=lst:13.7]
# Sin cache: 4 hilos, 100 conexiones, 30 segundos
wrk -t4 -c100 -d30s http://localhost:8080/api/sin-cache/1234

# Con cache (mismo escenario)
wrk -t4 -c100 -d30s http://localhost:8080/api/con-cache/1234
\end{lstlisting}

\subsubsection*{Paso 6 --- Tabla comparativa de resultados}

Tras ejecutar las pruebas de carga con y sin caché, el
estudiante debe producir una tabla comparativa que cuantifique las
diferencias.


\begin{tabularx}{\textwidth}{Xll}
\toprule
\textbf{Métrica} & \textbf{Sin cache} & \textbf{Con cache}\\
\midrule
Requests/segundo  & \rule{1cm}{0.4pt} & \rule{1cm}{0.4pt} \\
Latencia media (ms) & \rule{1cm}{0.4pt} & \rule{1cm}{0.4pt} \\
Latencia p99 (ms)  & \rule{1cm}{0.4pt} & \rule{1cm}{0.4pt} \\
CPU del proceso PG & \rule{1cm}{0.4pt}\% & \rule{1cm}{0.4pt}\% \\
\bottomrule
\end{tabularx}

\textbf{Resultado esperado:} la versión con caché debe mostrar 5x--20x
más requests/segundo y latencia significativamente menor. La CPU de
PostgreSQL debe caer drásticamente.

\subsubsection*{Paso 7 --- Documentar como ADR-004}

Escribir el ADR-004 del PFC con la decisión de adoptar Redis cache
para los listados frecuentes, citando los datos del benchmark como
justificación cuantitativa.
\end{cajaejercicio}

\section{Sumativa de la semana 13: Exposición Frameworks Backend}

La sumativa de la semana 13 es una exposición grupal de un
framework backend a elección. Cada equipo investiga, prepara
material y expone al resto de la clase.


\begin{cajaejercicio}{Sumativa \#12 --- Exposición grupal sobre framework backend asignado}


\textbf{Enunciado.} Investigar a fondo un framework backend asignado por el docente (Laravel, Django, Express, Quarkus o similar), preparar una exposición grupal de 30 minutos y un demo funcional que ilustre los conceptos clave del framework.

\textbf{Modalidad:} grupal (mismo equipo del PFC). Duración: 15--20
min de exposición + 5 min de preguntas. Cada equipo tiene un
framework distinto asignado por el docente.

\textbf{Frameworks disponibles para asignación:}
\begin{itemize}
\item Laravel (PHP), Symfony (PHP)
\item Django (Python), FastAPI (Python)
\item Express.js (Node.js), NestJS (Node.js)
\item Spring Boot (Java), Quarkus (Java)
\item ASP.NET Core (C\#)
\item Ruby on Rails (Ruby)
\item Gin (Go), Echo (Go)
\end{itemize}

\textbf{Estructura obligatoria de la exposición} (6 bloques):

\begin{enumerate}
\item \textbf{Historia y características} (3 min): origen, autor,
licencia, paradigma (MVC, API-first, opinionado o no).
\item \textbf{Arquitectura} (3 min): diagrama de arquitectura del
framework, ciclo de vida de petición.
\item \textbf{Demo en vivo} (5 min): instalación, primer proyecto,
endpoint GET y POST que persistan en BD.
\item \textbf{ORM/driver de BD} (2 min): definición de modelos,
migraciones, prevención SQL Injection.
\item \textbf{Ventajas, desventajas, casos de uso} (2 min):
fortalezas técnicas medibles, debilidades honestas, 2 casos donde
brilla y 1 donde no.
\item \textbf{Comparativa y conclusión} (2 min): tabla comparativa
con otro framework + recomendación final.
\end{enumerate}

\textbf{Entregables en LMS antes de la exposición:}
\begin{itemize}
\item Presentación PDF (mínimo 12 diapositivas, máximo 25).
\item Código fuente del demo en repositorio público (GitHub/GitLab).
\item Lista de referencias bibliográficas (mínimo 5, incluyendo la
documentación oficial y al menos un artículo académico).
\end{itemize}

\textbf{Esta es la \emph{Sumativa GA-Exposición} del segundo corte.}
\end{cajaejercicio}

\section{Buenas prácticas profesionales en escalabilidad}

La caja siguiente compila reglas profesionales de
escalabilidad que el estudiante debe respetar al diseñar y operar
cualquier sistema en producción.


\begin{cajabuenas}{Quince reglas profesionales 2026}
\begin{enumerate}
\item Diseñar stateless desde el inicio para permitir escalado
horizontal trivial.
\item Sesiones distribuidas (Redis, JDBC) si necesita servidor con
estado.
\item TTL en TODO dato cacheado: nunca cache eterna.
\item Cache-aside como patrón por defecto; otros solo si hay razón.
\item Invalidar por evento, no por comparación temporal compleja.
\item Connection pool dimensionado correctamente (regla: 2-4x cores).
\item Read replicas para BDs con alto ratio lectura/escritura.
\item Métricas reales (Prometheus, New Relic) antes de optimizar:
\emph{«premature optimization is the root of all evil»} (Knuth).
\item Health checks que verifiquen dependencias (BD, Redis).
\item Auto-scaling con métricas relevantes (CPU, RPS, latencia p95).
\item CDN para todo contenido estático.
\item Compresión Gzip/Brotli en respuestas (reduce 70-90\%).
\item HTTP/2 o HTTP/3 en producción (multiplexing).
\item Circuit breaker en llamadas a servicios externos.
\item Documentar las decisiones de escalabilidad en ADRs.
\end{enumerate}
\end{cajabuenas}

\section{Contraste bibliográfico de los conceptos clave}

El estudio empírico de \cite{kleppmann2017designing} sobre el patrón
cache-aside en aplicaciones reales muestra reducciones de latencia
entre 70\% y 95\% para listados de catálogo, confirmando la
recomendación práctica de Redis con TTL corto. Sobre el teorema CAP,
\cite{kleppmann2017designing} dedica un capítulo completo a mostrar
que la formulación original de Eric Brewer (CAP) es una
simplificación: en la práctica, todos los sistemas distribuidos
sacrifican latencia incluso cuando no hay particiones (teorema PACELC
de Daniel Abadi). \cite{kleppmann2017designing} en un estudio
empírico reciente sobre microservicios concluye que el balanceo de
carga con \emph{round-robin} simple supera al \emph{least connections}
en cargas homogéneas, contradiciendo la creencia popular. Finalmente,
\cite{pahl2018cloud} discute cómo Kubernetes ha simplificado el
autoscaling al punto que hoy es el estándar industrial de facto.

\section{Ejercicio resuelto adicional con resultado visual}

Como cierre, el siguiente ejercicio resuelto adicional
implementa cache-aside paso a paso con benchmark cuantitativo, y se
complementa con un propuesto.


\begin{cajaejemplo}{Ejercicio resuelto 13.1 --- Implementar cache-aside con Spring Boot y Redis}
\textbf{Enunciado.} Implementar el patrón cache-aside para un servicio
de catálogo de productos en Spring Boot 3.4 con Redis 7, y medir la
diferencia de latencia entre el primer acceso (cache fría) y
accesos sucesivos (cache caliente).

\textbf{Configuración Redis (\texttt{application.yml}):}
\begin{lstlisting}[language=yaml,caption={Configuración Redis Spring Boot},label=lst:13.8]
spring:                                                     # configuracion de Spring
  data:
    redis:                                                  # configuracion del cliente Redis
      host: localhost                                       # host del servicio
      port: 6379                                            # puerto donde escucha la aplicacion
      timeout: 2000ms
  cache:
    type: redis
    redis:                                                  # configuracion del cliente Redis
      time-to-live: 60000   # 60 segundos
\end{lstlisting}

\textbf{Servicio con caché (Listado~\ref{lst:13.9}):}
\begin{lstlisting}[language=Java,caption={Servicio con anotaciones de caché},label=lst:13.9]
@Service                                                    // capa de servicio (logica de negocio)
public class CatalogoService {                              // declaracion de clase publica
  private final ProductoRepository repo;                    // campo inmutable (mejor practica)
  // metodo publico
  // metodo publico
  public CatalogoService(ProductoRepository r) { this.repo = r; }

  @Cacheable(value = "productos", key = "#id")              // guarda el resultado en cache (cache-aside)
  public Producto obtener(Long id) {                        // metodo publico
    try { Thread.sleep(80); } catch (Exception e) { } // simula BD lenta
    return repo.findById(id).orElseThrow();                 // devuelve el resultado al llamador
  }

  @CacheEvict(value = "productos", key = "#p.id")           // invalida la entrada de cache al modificar
  public Producto actualizar(Producto p) {                  // metodo publico
    return repo.save(p);                                    // devuelve el resultado al llamador
  }

  @CacheEvict(value = "productos", allEntries = true)       // invalida la entrada de cache al modificar
  public void invalidarTodo() { }                           // metodo publico
}
\end{lstlisting}

\textbf{Test con prueba de carga k6 (Listado~\ref{lst:13.10}):}
\begin{lstlisting}[language=JavaScript,caption={Script k6 que evidencia el efecto del caché},label=lst:13.10]
import http from 'k6/http';                                 // importacion ES Modules
import { check } from 'k6';                                 // importacion ES Modules
export const options = { vus: 1, iterations: 20 };          // exporta del modulo
export default function () {                                // exporta del modulo
  const r = http.get('http://localhost:8080/api/v1/productos/42');
  check(r, { 'status 200': (x) => x.status === 200 });
}
\end{lstlisting}

\textbf{Resultado visual} (consola tras \texttt{k6 run script.js}):

\begin{center}
\fbox{\begin{minipage}{0.95\textwidth}
\vspace{0.2cm}\footnotesize\ttfamily
\hspace{0.3cm}execution: local\\
\hspace{0.3cm}iteraciones: 20 / 20\\
\hspace{0.3cm}vus\_max......: 1\\[0.2cm]
\hspace{0.3cm}http\_req\_duration..............: avg=8.2ms \quad min=2ms \quad med=4ms\\
\hspace{0.6cm}p(90)=12ms \quad p(95)=85ms \quad max=92ms\\[0.2cm]
\hspace{0.3cm}\textbf{Iteración 1 (cache fría):} 92 ms\\
\hspace{0.3cm}\textbf{Iteraciones 2-20 (cache caliente):} 2-4 ms\\
\hspace{0.3cm}\textbf{Mejora: $\approx$ 95\%}\\
\vspace{0.1cm}
\end{minipage}}
\end{center}

\textbf{Inspección de Redis} (\texttt{redis-cli}):

\begin{center}
\fbox{\begin{minipage}{0.95\textwidth}
\vspace{0.15cm}\footnotesize\ttfamily
\hspace{0.3cm}127.0.0.1:6379{$>$} KEYS productos::*\\
\hspace{0.3cm}1) \char34{}productos::42\char34{}\\[0.2cm]
\hspace{0.3cm}127.0.0.1:6379{$>$} TTL productos::42\\
\hspace{0.3cm}(integer) 54\\[0.2cm]
\hspace{0.3cm}127.0.0.1:6379{$>$} GET productos::42\\
\hspace{0.3cm}\char34{}$\backslash$xac$\backslash$xed... (datos serializados Java)\char34{}\\
\vspace{0.1cm}
\end{minipage}}
\end{center}

\textbf{Discusión.} La primera petición tarda $\sim$92 ms porque
recorre todo el flujo: controlador $\to$ servicio $\to$ JPA $\to$
PostgreSQL. Las siguientes 19 tardan 2-4 ms porque Spring intercepta
\texttt{@Cacheable}, consulta Redis, encuentra el valor y lo devuelve
sin tocar la BD. El TTL de 60 s expirará el valor automáticamente; si
durante esos 60 s alguien llama \texttt{actualizar()},
\texttt{@CacheEvict} lo invalidará explícitamente
\cite{kleppmann2017designing}.
\end{cajaejemplo}

\begin{cajaejercicio}{Ejercicio propuesto 13.2 --- Balanceo de carga con Nginx y 3 instancias}
\textbf{Enunciado.} Desplegar 3 instancias del backend Spring Boot
detrás de un balanceador Nginx con Docker Compose. Configurar
\emph{sticky sessions} con cookies y comparar contra
\emph{round-robin} puro. Medir el impacto en latencia y verificar
que las sesiones se mantienen consistentes.

\textbf{Requisitos:}
\begin{enumerate}
\item \texttt{docker-compose.yml} que levante: 1 PostgreSQL,
1 Redis, 3 instancias Spring Boot (app1, app2, app3), 1 Nginx.
\item Nginx con \texttt{upstream} de 3 backends, política
\texttt{ip\_hash} en una variante y \texttt{round\_robin} en otra.
\item El backend Spring Boot debe imprimir su hostname en cada
respuesta para identificar qué instancia atendió.
\item Sesiones persistidas en Redis con Spring Session
(\texttt{spring-session-data-redis}).
\item Pruebas de carga con k6: 100 VUs durante 1 minuto.
\item Reporte comparativo: latencia p50, p95, p99 entre las dos
políticas.
\end{enumerate}

\textbf{Criterios de evaluación:}
\begin{itemize}
\item Las 3 instancias atienden tráfico (verificable por hostname).
\item Las sesiones son consistentes (un usuario logueado sigue
logueado entre peticiones a cualquier instancia, gracias a Redis).
\item Resultados k6 con tabla comparativa de latencias.
\item Documentación con diagrama de la arquitectura y ADR-003 que
justifique la política de balanceo elegida.
\end{itemize}

\textbf{Lecturas recomendadas:}
\cite{kleppmann2017designing} para benchmarks de balanceo;
\cite{kleppmann2017designing} cap. 5 para consistencia y replicación.
\end{cajaejercicio}

% =====================================================================
%       SEMANA 14: PATRÓN MVC EN APLICACIONES WEB
% =====================================================================
\chapter{El patrón Modelo-Vista-Controlador en aplicaciones web modernas}
\label{cap:semana14}

\epigraph{«MVC no es un framework: es una forma de pensar el software
que separa lo que el usuario ve, lo que el sistema sabe, y la lógica
que conecta ambos mundos.»}{--- \textit{Adaptado de Trygve Reenskaug}}

\section*{Resultado de aprendizaje de la semana}
Al concluir la semana 14, el estudiante diseña aplicaciones web bajo
el patrón Modelo-Vista-Controlador, comprende la trayectoria histórica
desde el MVC original de Smalltalk-80 hasta sus variantes modernas
(server-side MVC vs API REST + SPA), y construye una aplicación
funcional con un \emph{framework} MVC contemporáneo, depurándola paso
a paso en IntelliJ IDEA Ultimate.

\section{Historia del patrón MVC}

MVC es uno de los patrones de diseño más citados en la
historia del software. Comprender su historia ---desde Smalltalk-76
hasta los frameworks modernos--- ayuda a entender por qué hay tantas
variantes y por qué sigue siendo central en 2026.


\begin{cajahistoria}{Reenskaug --- Smalltalk-80 y el nacimiento de un patrón fundamental}
El patrón MVC fue concebido por el ingeniero noruego \textbf{Trygve
Reenskaug} durante una estancia sabática en el Xerox PARC en
diciembre de 1978 mientras trabajaba en la primera versión de
Smalltalk. Reenskaug describió MVC como una solución al problema
fundamental de los sistemas interactivos: cómo permitir que múltiples
vistas reflejen el mismo modelo de datos sin acoplarlas entre sí.

En su nota original \emph{«Models-Views-Controllers»} (10 de
diciembre de 1979), Reenskaug definía:

\begin{itemize}
\item \textbf{Modelo}: la representación del conocimiento sobre el
problema; lo que la aplicación realmente \emph{sabe}.
\item \textbf{Vista}: una representación visual del modelo, no el
modelo mismo.
\item \textbf{Controlador}: el enlace entre el usuario y el sistema,
que interpreta las acciones del usuario y las traduce en cambios al
modelo.
\end{itemize}

MVC fue implementado plenamente en Smalltalk-80 (1980) y durante
quince años permaneció en relativa oscuridad académica. Su
revitalización vino con dos hitos:

\begin{itemize}
\item \textbf{NeXTSTEP / Cocoa} (1989+): Apple adoptó MVC como
patrón estructural de la programación nativa de Mac y, posteriormente,
de iOS.
\item \textbf{Ruby on Rails} (David Heinemeier Hansson, 2004):
implementó una variante \emph{server-side} de MVC para aplicaciones
web. Su éxito masivo inspiró Django (Python, 2005), Symfony
(PHP, 2005), CodeIgniter (PHP, 2006), CakePHP (PHP, 2005), Laravel
(PHP, 2011), ASP.NET MVC (Microsoft, 2009) y Spring MVC (2004).
\end{itemize}

A partir de 2014 emergieron las arquitecturas \textbf{API REST + SPA}
que desplazaron parcialmente el MVC server-side: el servidor expone
JSON, y un cliente Angular/React/Vue se encarga de la vista. Sin
embargo, el patrón MVC sigue siendo fundamental: incluso las SPAs
implementan internamente versiones del patrón (MVVM, Flux, Redux,
Signal-based architectures) que son herederas conceptuales del MVC
original.
\end{cajahistoria}

\subsection{Tabla evolutiva del MVC en frameworks web}

Distintos frameworks implementan MVC con variantes propias.
La tabla siguiente recorre los más influyentes en orden cronológico.


\begin{table}[htbp]
\centering
\caption{Hitos en la historia del MVC en frameworks web.}
\footnotesize
\begin{tabularx}{\textwidth}{lllX}
\toprule
\textbf{Año} & \textbf{Framework} & \textbf{Lenguaje} & \textbf{Aporte}\\
\midrule
1979 & ---           & Smalltalk & Reenskaug formaliza el patrón.\\
1980 & Smalltalk-80  & Smalltalk & Primera implementación completa.\\
1989 & NeXTSTEP      & Objective-C & MVC en aplicaciones nativas de escritorio.\\
1996 & Apache Struts (precursor) & Java & MVC web temprano (Java Servlets).\\
2002 & Spring        & Java & DI + MVC posteriormente (Spring MVC, 2004).\\
2004 & Ruby on Rails & Ruby & Convención sobre configuración; MVC mainstream.\\
2005 & Django        & Python & Modelo MTV (Model-Template-View, variante MVC).\\
2005 & Symfony       & PHP   & MVC con bundles y dependencias.\\
2009 & ASP.NET MVC   & C\#   & Microsoft adopta MVC web.\\
2011 & Laravel       & PHP   & MVC con Eloquent ORM y Blade templates.\\
2014 & Spring Boot   & Java  & Convención sobre Spring MVC, simplificación radical.\\
2014 & Angular 2     & TypeScript & Frontend MVC (versión MVVM).\\
2024 & Laravel 11    & PHP   & Estructura simplificada, Volt, Folio.\\
2024 & Spring Boot 3.4 & Java & Soporte HTTP virtual threads en Tomcat.\\
2024 & ASP.NET Core 8 & C\#  & Razor Pages + MVC + Web API unificados.\\
\bottomrule
\end{tabularx}
\end{table}

\section{Anatomía del patrón MVC en la web}

MVC en aplicaciones web tiene una anatomía específica: el
modelo no es solo una clase, sino una capa entera; la vista no es
una pantalla, sino una plantilla; el controlador no es un widget,
sino un endpoint HTTP.


\subsection{Adaptación del MVC original al contexto HTTP}

El MVC de Reenskaug fue concebido para aplicaciones de escritorio
con observadores reactivos: la vista se suscribía al modelo, que la
notificaba en cada cambio. En la web esto no es posible directamente
(HTTP es \emph{request/response}), por lo que el MVC web introduce
adaptaciones:

\begin{itemize}
\item No hay observadores reactivos: la vista se renderiza en cada
petición.
\item Hay un \textbf{front controller} (DispatcherServlet en Spring,
\texttt{index.php} en Laravel, \texttt{Routes} en ASP.NET Core) que
recibe TODAS las peticiones y las enruta al controller adecuado.
\item Las URLs se mapean a métodos de controller con sistemas de
\emph{routing} declarativos.
\item Las vistas se generan típicamente en el servidor con
\emph{template engines} (Thymeleaf, Blade, Razor, Twig).
\end{itemize}

\subsection{Las tres capas en detalle}

Cada una de las tres capas de MVC tiene responsabilidades
concretas. La descripción siguiente las detalla.


\subsubsection{El Modelo}

El Modelo encapsula:
\begin{itemize}
\item \textbf{Entidades de dominio}: objetos persistentes con
identidad propia (Producto, Pedido, Usuario).
\item \textbf{Servicios de aplicación}: orquestan casos de uso
(crear pedido, calcular descuento, enviar notificación).
\item \textbf{Repositorios}: encapsulan el acceso a la persistencia.
\item \textbf{Reglas de negocio}: invariantes del dominio
(un usuario menor de edad no puede comprar bebidas alcohólicas).
\item \textbf{DTOs} (Data Transfer Objects): contenedores planos
para transferir datos entre capas o servicios.
\end{itemize}

\textbf{Anti-patrón clásico:} el «Anemic Domain Model» (Fowler) en
el que las entidades son meros contenedores sin lógica y toda la
inteligencia está en servicios. En MVC profesional las entidades
tienen métodos significativos.

\subsubsection{La Vista}

La Vista es la capa de presentación. En aplicaciones web tradicionales
suele ser una plantilla HTML con marcadores de sustitución:

\begin{itemize}
\item \textbf{Thymeleaf} (Java/Spring): plantillas HTML que son HTML
válido incluso sin procesar (preview-friendly).
\item \textbf{Blade} (PHP/Laravel): \texttt{@if}, \texttt{@foreach},
\texttt{\{\{ \$variable \}\}}.
\item \textbf{Razor} (C\#/ASP.NET): mezcla \texttt{@} para C\# inline.
\item \textbf{Twig} (PHP/Symfony): sintaxis estilo Django,
\texttt{\{\% if \%\} \{\{ var \}\}}.
\item \textbf{Mustache, Handlebars} (cualquier lenguaje): \emph{logic-less}.
\end{itemize}

En arquitecturas API REST + SPA, la Vista se desplaza al cliente
(Angular, React, Vue), y el servidor solo expone JSON.

\subsubsection{El Controlador}

El Controller es la capa que recibe peticiones HTTP, invoca al Modelo
y selecciona la Vista (o el formato de respuesta JSON). Sus
responsabilidades:

\begin{itemize}
\item \emph{Routing}: mapear URLs a métodos.
\item \emph{Binding}: convertir parámetros HTTP en objetos del modelo.
\item \emph{Validation}: validar entrada (a menudo delegada).
\item \emph{Orquestación}: llamar a servicios del modelo.
\item \emph{Selección de respuesta}: vista a renderizar o DTO a
serializar.
\item \emph{Manejo de errores}: convertir excepciones en respuestas
HTTP apropiadas.
\end{itemize}

\textbf{Regla de oro:} el controller debe ser \emph{delgado} (\emph{thin
controller}). Toda lógica de negocio significativa va en el modelo
(servicios). Un controller con 200 líneas es señal de mal diseño.

\section{Ciclo de vida de una petición MVC en Spring Boot}

Una petición HTTP atraviesa una secuencia precisa de
componentes dentro de Spring Boot. Conocer ese flujo permite
depurar problemas y optimizar puntos críticos.


\begin{figure}[htbp]
\centering
\begin{tikzpicture}[
  every node/.style={font=\scriptsize,align=center},
  caja/.style={rectangle,draw=uteqverde,thick,minimum width=2.4cm,
               minimum height=0.8cm,fill=uteqverde!10,rounded corners=2pt},
  flecha/.style={->,>=stealth,thick}]

  \node[caja] (cli) {1. Cliente\\(Navegador)};
  \node[caja,below=0.5cm of cli] (ds) {2. Dispatcher\\Servlet};
  \node[caja,below=0.5cm of ds] (hm) {3. Handler\\Mapping};
  \node[caja,below=0.5cm of hm] (ctr) {4. Controller\\@RequestMapping};
  \node[caja,right=1.5cm of ctr] (svc) {5. Service\\(reglas negocio)};
  \node[caja,right=1.5cm of svc] (rp) {6. Repository\\(JPA)};
  \node[cylinder,draw=uteqverde,fill=uteqverde!15,thick,
        shape border rotate=90,aspect=0.3,minimum width=1.2cm,
        minimum height=1.2cm,below=0.5cm of rp] (db) {BD};
  \node[caja,left=1.5cm of ctr] (vr) {8. View\\Resolver};
  \node[caja,left=1.5cm of vr] (vw) {9. Vista\\(Thymeleaf)};

  \draw[flecha] (cli) -- node[right]{HTTP Request} (ds);
  \draw[flecha] (ds) -- (hm);
  \draw[flecha] (hm) -- (ctr);
  \draw[flecha] (ctr) -- (svc);
  \draw[flecha] (svc) -- (rp);
  \draw[flecha] (rp) -- (db);
  \draw[flecha,dashed] (db) -- node[right,font=\tiny]{rows} (rp);
  \draw[flecha,dashed] (rp) -- (svc);
  \draw[flecha,dashed] (svc) -- (ctr);
  \draw[flecha] (ctr) -- node[above,font=\tiny]{ModelAndView} (vr);
  \draw[flecha] (vr) -- (vw);
  \draw[flecha,dashed] (vw.north) to[bend left=60]
        node[above,font=\tiny]{HTML} (cli.west);

\end{tikzpicture}
\caption{Ciclo de vida completo de una petición HTTP en Spring MVC. La
ida (línea sólida): Dispatcher $\to$ HandlerMapping $\to$ Controller
$\to$ Service $\to$ Repository $\to$ BD. El regreso (línea punteada):
los datos suben hasta el Controller, que devuelve un
\texttt{ModelAndView}; el ViewResolver elige la plantilla y el HTML
final llega al cliente.}
\end{figure}

\section{Server-side MVC vs SPA + API REST}

En 2026 hay dos arquitecturas dominantes para aplicaciones
web: MVC \emph{server-side} clásico (Thymeleaf/Razor/Blade) y SPA +
API REST (Angular/React + Spring/ASP.NET). La tabla siguiente las
contrasta con sus trade-offs.


\begin{cajacompara}{MVC server-side tradicional vs API REST + SPA Angular/React}
\centering\footnotesize
\begin{tabularx}{\textwidth}{lXX}
\toprule
\textbf{Criterio} & \textbf{MVC server-side} & \textbf{API REST + SPA}\\
\midrule
Ubicación de la vista & Servidor (Thymeleaf, Blade, Razor) & Cliente (Angular, React, Vue)\\
Renderizado inicial & Rápido (HTML completo del primer hit) & Lento (descarga JS, luego renderiza)\\
Tráfico tras la primera carga & Cada navegación = HTML completo & JSON pequeño\\
SEO & Excelente nativo & Requiere SSR/SSG\\
Complejidad del frontend & Baja (HTML + algo de JS) & Alta (framework completo)\\
Reactividad & Limitada (post-back o JS extra) & Nativa\\
Equipo & Un equipo full-stack & Frontend y backend separados\\
Mejor para & CMS, intranet, sitios contenido & SaaS, dashboards, apps tipo Gmail\\
\bottomrule
\end{tabularx}
\end{cajacompara}

\subsection{¿Cuál usar en el PFC?}

El PFC de esta asignatura adopta la combinación \textbf{API REST
(Spring Boot) + SPA Angular}. Es la apuesta arquitectónica moderna y
es la que el sílabo prescribe. Sin embargo, comprender ambos
enfoques es esencial: muchos sistemas reales coexisten (un panel
admin con Razor + una API para móvil + un portal público con SSR).

\section{Implementación canónica: MVC server-side completo}

El ejemplo siguiente implementa un CRUD MVC server-side
completo con Spring Boot 3 y Thymeleaf: el estudiante lo construye
paso a paso y obtiene una aplicación funcional ejecutable.


\begin{cajaejemplo}{Ejemplo 14.1 --- CRUD MVC clásico en Spring Boot 3 + Thymeleaf}


\textbf{Enunciado.} Construir una aplicación MVC server-side de gestión de tareas (to-do) con Spring Boot 3.4, Thymeleaf y H2. Las cuatro operaciones CRUD funcionarán end-to-end.

\textbf{Estructura del proyecto (Listado~\ref{lst:14.1}):}
\begin{lstlisting}[style=shell,numbers=none,caption={Ejemplo de Shell},label=lst:14.1]
src/main/java/ec/edu/uteq/mvc/
  controller/ProductoController.java   <- Controller
  service/ProductoService.java         <- Servicio (modelo)
  repository/ProductoRepository.java   <- Repositorio (modelo)
  model/Producto.java                  <- Entidad de dominio
  dto/ProductoForm.java                <- DTO con validacion
src/main/resources/templates/productos/ <- Vistas
  lista.html
  formulario.html
\end{lstlisting}

\textbf{Entidad \texttt{Producto}:}
\begin{lstlisting}[style=java,caption={Producto.java},label=lst:14.2]
package ec.edu.uteq.mvc.model;                              // paquete del archivo

import jakarta.persistence.*;                               // importacion de clase externa
import lombok.*;                                            // importacion de clase externa
import java.math.BigDecimal;                                // importacion de clase externa

@Entity @Table(name = "productos")                          // entidad JPA (mapea a tabla)
@Getter @Setter @NoArgsConstructor @AllArgsConstructor
public class Producto {                                     // declaracion de clase publica
    @Id @GeneratedValue(strategy = GenerationType.IDENTITY) // campo clave primaria
    private Long id;                                        // campo privado
    private String nombre;                                  // campo privado
    private BigDecimal precio;                              // campo privado
    private Integer stock;                                  // campo privado

    /** Logica de dominio: el modelo NO es anemico. */
    public boolean estaDisponible() {                       // metodo publico
        return stock != null && stock > 0;                  // devuelve el resultado al llamador
    }

    public BigDecimal valorInventario() {                   // metodo publico
        // condicional
        // condicional
        if (precio == null || stock == null) return BigDecimal.ZERO;
        return precio.multiply(BigDecimal.valueOf(stock));  // devuelve el resultado al llamador
    }
}
\end{lstlisting}

\textbf{DTO con validación (Listado~\ref{lst:14.3}):}
\begin{lstlisting}[style=java,caption={ProductoForm.java},label=lst:14.3]
package ec.edu.uteq.mvc.dto;                                // paquete del archivo

import jakarta.validation.constraints.*;                    // importacion de clase externa
import lombok.*;                                            // importacion de clase externa
import java.math.BigDecimal;                                // importacion de clase externa

@Getter @Setter @NoArgsConstructor @AllArgsConstructor
public class ProductoForm {                                 // declaracion de clase publica
    private Long id;                                        // campo privado

    @NotBlank(message = "El nombre es obligatorio")         // validacion: no nulo y no vacio
    @Size(min = 3, max = 120)                               // validacion: longitud min/max
    private String nombre;                                  // campo privado

    @NotNull @DecimalMin("0.00") @DecimalMax("99999.99")    // validacion: no nulo
    private BigDecimal precio;                              // campo privado

    @NotNull @Min(0)                                        // validacion: no nulo
    private Integer stock;                                  // campo privado
}
\end{lstlisting}

\textbf{Repository (Spring Data) (Listado~\ref{lst:14.4}):}
\begin{lstlisting}[style=java,caption={ProductoRepository.java},label=lst:14.4]
package ec.edu.uteq.mvc.repository;                         // paquete del archivo

import ec.edu.uteq.mvc.model.Producto;                      // importacion de clase externa
// importacion de clase externa
// importacion de clase externa
import org.springframework.data.jpa.repository.JpaRepository;

public interface ProductoRepository                         // declaracion de interfaz
    extends JpaRepository<Producto, Long> {}
\end{lstlisting}

\textbf{Service (Modelo) (Listado~\ref{lst:14.5}):}
\begin{lstlisting}[style=java,caption={ProductoService.java},label=lst:14.5]
package ec.edu.uteq.mvc.service;                            // paquete del archivo

import ec.edu.uteq.mvc.model.Producto;                      // importacion de clase externa
import ec.edu.uteq.mvc.repository.ProductoRepository;       // importacion de clase externa
import lombok.RequiredArgsConstructor;                      // importacion de clase externa
import org.springframework.stereotype.Service;              // importacion de clase externa
// importacion de clase externa
// importacion de clase externa
import org.springframework.transaction.annotation.Transactional;
import java.util.*;                                         // importacion de clase externa

@Service                                                    // capa de servicio (logica de negocio)
@RequiredArgsConstructor
public class ProductoService {                              // declaracion de clase publica

    private final ProductoRepository repo;                  // campo inmutable (mejor practica)

    @Transactional(readOnly = true)                         // envuelve el metodo en transaccion
    // metodo publico
    // metodo publico
    public List<Producto> listar() { return repo.findAll(); }

    @Transactional(readOnly = true)                         // envuelve el metodo en transaccion
    public Producto buscar(Long id) {                       // metodo publico
        return repo.findById(id).orElseThrow(               // devuelve el resultado al llamador
            () -> new NoSuchElementException("Producto " + id));
    }

    @Transactional                                          // envuelve el metodo en transaccion
    // metodo publico
    // metodo publico
    public Producto guardar(Producto p) { return repo.save(p); }

    @Transactional                                          // envuelve el metodo en transaccion
    public void eliminar(Long id) { repo.deleteById(id); }  // metodo publico
}
\end{lstlisting}

\textbf{Controller (Listado~\ref{lst:14.6}):}
\begin{lstlisting}[style=java,caption={ProductoController.java},label=lst:14.6]
package ec.edu.uteq.mvc.controller;                         // paquete del archivo

import ec.edu.uteq.mvc.dto.ProductoForm;                    // importacion de clase externa
import ec.edu.uteq.mvc.model.Producto;                      // importacion de clase externa
import ec.edu.uteq.mvc.service.ProductoService;             // importacion de clase externa
import jakarta.validation.Valid;                            // importacion de clase externa
import lombok.RequiredArgsConstructor;                      // importacion de clase externa
import org.springframework.beans.BeanUtils;                 // importacion de clase externa
import org.springframework.stereotype.Controller;           // importacion de clase externa
import org.springframework.ui.Model;                        // importacion de clase externa
import org.springframework.validation.BindingResult;        // importacion de clase externa
import org.springframework.web.bind.annotation.*;           // importacion de clase externa
// importacion de clase externa
// importacion de clase externa
import org.springframework.web.servlet.mvc.support.RedirectAttributes;

@Controller                                                 // controlador MVC server-side (devuelve vistas)
@RequestMapping("/productos")                               // ruta base del controlador
@RequiredArgsConstructor
public class ProductoController {                           // declaracion de clase publica

    private final ProductoService service;                  // campo inmutable (mejor practica)

    /** GET /productos --> lista */
    @GetMapping                                             // endpoint HTTP GET
    public String listar(Model model) {                     // metodo publico
        model.addAttribute("productos", service.listar());
        return "productos/lista";   // -> templates/productos/lista.html
    }

    /** GET /productos/nuevo --> formulario vacio */
    @GetMapping("/nuevo")                                   // endpoint HTTP GET
    public String mostrarFormulario(Model model) {          // metodo publico
        model.addAttribute("form", new ProductoForm());
        return "productos/formulario";                      // devuelve el resultado al llamador
    }

    /** POST /productos --> guardar */
    @PostMapping                                            // endpoint HTTP POST (crear)
    // metodo publico
    // metodo publico
    public String guardar(@Valid @ModelAttribute("form") ProductoForm form,
                          BindingResult br,
                          RedirectAttributes ra) {
        if (br.hasErrors()) {                               // condicional
            // re-renderizar el formulario con los errores
            return "productos/formulario";                  // devuelve el resultado al llamador
        }
        Producto p = new Producto();
        BeanUtils.copyProperties(form, p);
        service.guardar(p);
        ra.addFlashAttribute("mensaje", "Producto guardado");
        return "redirect:/productos";   // patron PRG anti F5
    }

    /** GET /productos/{id}/editar */
    @GetMapping("/{id}/editar")                             // endpoint HTTP GET
    // metodo publico
    // metodo publico
    public String editar(@PathVariable Long id, Model model) {
        Producto p = service.buscar(id);
        ProductoForm form = new ProductoForm();
        BeanUtils.copyProperties(p, form);
        model.addAttribute("form", form);
        return "productos/formulario";                      // devuelve el resultado al llamador
    }

    /** POST /productos/{id}/eliminar */
    @PostMapping("/{id}/eliminar")                          // endpoint HTTP POST (crear)
    // metodo publico
    // metodo publico
    public String eliminar(@PathVariable Long id, RedirectAttributes ra) {
        service.eliminar(id);
        ra.addFlashAttribute("mensaje", "Producto eliminado");
        return "redirect:/productos";                       // devuelve el resultado al llamador
    }
}
\end{lstlisting}

\textbf{Vista Thymeleaf de listado (Listado~\ref{lst:14.7}):}
\begin{lstlisting}[style=html,caption={templates/productos/lista.html},label=lst:14.7]
<!DOCTYPE html>
<html xmlns:th="http://www.thymeleaf.org" lang="es-EC">
<head>
  <meta charset="UTF-8">
  <title>Productos UTEQ</title>
  <link rel="stylesheet" href="/css/app.css">
</head>
<body>
<header>
  <h1>Catalogo de productos</h1>
  <a th:href="@{/productos/nuevo}" class="btn">Nuevo producto</a>
</header>

<div th:if="${mensaje}" class="alerta-exito" th:text="${mensaje}"></div>

<table>
  <thead>
    <tr><th>ID</th><th>Nombre</th><th>Precio</th>
        <th>Stock</th><th>Disponible</th><th>Acciones</th></tr>
  </thead>
  <tbody>
    <tr th:each="p : ${productos}">
      <td th:text="${p.id}"></td>
      <td th:text="${p.nombre}"></td>
      <td th:text="${'$' + #numbers.formatDecimal(p.precio,1,2)}"></td>
      <td th:text="${p.stock}"></td>
      <td th:text="${p.estaDisponible() ? 'Si' : 'No'}"></td>
      <td>
        <a th:href="@{/productos/{id}/editar(id=${p.id})}">Editar</a>
        <form th:action="@{/productos/{id}/eliminar(id=${p.id})}"
              method="post" style="display:inline">
          <button type="submit"
                  onclick="return confirm('Confirmar eliminacion?')">
            Eliminar
          </button>
        </form>
      </td>
    </tr>
  </tbody>
</table>
</body>
</html>
\end{lstlisting}

\textbf{Vista del formulario (Listado~\ref{lst:14.8}):}
\begin{lstlisting}[style=html,caption={templates/productos/formulario.html},label=lst:14.8]
<!DOCTYPE html>
<html xmlns:th="http://www.thymeleaf.org" lang="es-EC">
<head><meta charset="UTF-8"><title>Producto</title></head>
<body>
<h1 th:text="${form.id == null ? 'Nuevo producto' : 'Editar producto'}">
</h1>

<form th:action="@{/productos}" th:object="${form}" method="post">
  <input type="hidden" th:field="*{id}">

  <div>
    <label>Nombre</label>
    <input type="text" th:field="*{nombre}" required>
    <span th:errors="*{nombre}" class="error"></span>
  </div>

  <div>
    <label>Precio</label>
    <input type="number" step="0.01" th:field="*{precio}" required>
    <span th:errors="*{precio}" class="error"></span>
  </div>

  <div>
    <label>Stock</label>
    <input type="number" th:field="*{stock}" min="0" required>
    <span th:errors="*{stock}" class="error"></span>
  </div>

  <button type="submit">Guardar</button>
  <a th:href="@{/productos}">Cancelar</a>
</form>
</body>
</html>
\end{lstlisting}

\textbf{Notas pedagógicas sobre este código:}
\begin{itemize}
\item \emph{Patrón PRG (Post-Redirect-Get):} tras un POST exitoso se
hace \texttt{redirect:/productos}, no se renderiza directamente la
vista. Esto evita el bug clásico de doble envío al pulsar F5.
\item \emph{Flash attributes:} \texttt{RedirectAttributes} permite
pasar mensajes entre el redirect y la vista sin contaminar la URL.
\item \emph{ProductoForm} en lugar de \texttt{Producto} en el
controller: separa la representación de la entidad y la del
formulario; permite validaciones específicas de UI.
\item \emph{El modelo no es anémico:} \texttt{estaDisponible()} es
una regla del dominio.
\item \emph{Thymeleaf-natural:} las plantillas son HTML válido que se
puede abrir en navegador sin servidor (modo preview).
\end{itemize}
\end{cajaejemplo}

\section{Variantes del MVC}

MVC puro no es la única opción: en la historia han surgido
variantes con ventajas específicas (MVP, MVVM, MVI). La tabla
siguiente las catalogará y dirá cuándo usar cada una.


\subsection{MVP (Model-View-Presenter)}

En MVP el Presenter reemplaza al Controller pero acopla más
fuertemente con la Vista. La Vista delega completamente la lógica de
presentación al Presenter. Común en aplicaciones Android pre-MVVM.

\subsection{MVVM (Model-View-ViewModel)}

En MVVM, el ViewModel es una abstracción de la vista que mantiene su
estado. La vista se enlaza al ViewModel por \emph{data binding}
bidireccional. Angular y Vue son herederos parciales de este enfoque.

\subsection{Flux y Redux}

Patrones unidireccionales para SPAs: Action $\to$ Dispatcher $\to$
Store $\to$ View. Resuelven el problema de complejidad del estado en
SPAs grandes (Facebook lo creó para resolver bugs en Messenger).

\subsection{Signals (Angular 17+, SolidJS)}

Modelo reactivo basado en \emph{signals}: variables observables que
notifican automáticamente cambios. Una evolución moderna del MVC
reactivo original de Reenskaug, ahora aplicada al frontend.

\section{Práctica integrada: CRUD MVC con debugging}

La práctica integrada de la semana 14 implementa el CRUD MVC
trazando explícitamente el flujo de cada petición, para que el
estudiante visualice cómo trabaja el framework por debajo.


\begin{cajaejercicio}{Ejercicio 14.1 --- CRUD MVC completo con trazado de la petición en IntelliJ}

\textbf{Objetivo:} construir un CRUD MVC server-side y aprender a
\emph{depurar} cada fase del ciclo de vida de la petición usando los
breakpoints y el inspector de IntelliJ IDEA Ultimate.

\subsubsection*{IDE: IntelliJ IDEA Ultimate}

IntelliJ IDEA Ultimate provee herramientas específicas para
trazar el flujo MVC (debugger con puntos de quiebre por capas,
consola Spring, herramienta HTTP Client). Esta práctica las usa
todas.


\subsubsection*{Paso 1 --- Crear el proyecto}

\texttt{File $\rightarrow$ New $\rightarrow$ Project $\rightarrow$
Spring Boot}. Configurar:
\begin{itemize}
\item Type: Maven, JDK: 21, Spring Boot 3.4.x.
\item Group: \texttt{ec.edu.uteq}; Artifact: \texttt{mvc-debug}.
\item Dependencies: \emph{Spring Web}, \emph{Thymeleaf}, \emph{Spring
Data JPA}, \emph{H2 Database}, \emph{Lombok}, \emph{Validation},
\emph{DevTools}.
\end{itemize}

\subsubsection*{Paso 2 --- Implementar el CRUD del Ejemplo 14.1}

Copiar la estructura completa del ejemplo 14.1 (entidad, DTO,
service, controller, dos vistas Thymeleaf).

Agregar en \texttt{application.yml} (Listado~\ref{lst:14.9}):
\begin{lstlisting}[style=yaml,numbers=none,caption={Ejemplo de YAML},label=lst:14.9]
spring:                                                     # configuracion de Spring
  datasource:                                               # configuracion de origen de datos (BD)
    url: jdbc:h2:mem:appweb;DB_CLOSE_DELAY=-1               # URL JDBC de conexion a la BD
    driver-class-name: org.h2.Driver                        # clase del driver JDBC
    username: sa                                            # usuario de la BD
    password:                                               # contrasena (mejor leer de env var)
  jpa:                                                      # configuracion JPA/Hibernate
    hibernate.ddl-auto: create-drop                         # estrategia de generacion de schema
    show-sql: true                                          # imprime SQL en consola (solo desarrollo)
    properties.hibernate.format_sql: true
  h2.console:
    enabled: true
    path:    /h2                                            # ruta donde se expone
\end{lstlisting}

\subsubsection*{Paso 3 --- Inicializar datos de prueba}

Crear \texttt{DatosIniciales.java} (Listado~\ref{lst:14.10}):
\begin{lstlisting}[style=java,numbers=none,caption={Ejemplo de Java},label=lst:14.10]
@Component                                                  // bean Spring generico
public class DatosIniciales implements CommandLineRunner {  // declaracion de clase publica
    private final ProductoRepository repo;                  // campo inmutable (mejor practica)
    // metodo publico
    // metodo publico
    public DatosIniciales(ProductoRepository r) { this.repo = r; }
    @Override
    public void run(String... args) {                       // metodo publico
        repo.save(new Producto(null, "Cafe galletti", new BigDecimal("4.50"), 30));
        repo.save(new Producto(null, "Cacao Pacari",  new BigDecimal("6.00"),  5));
        repo.save(new Producto(null, "Te organico",   new BigDecimal("3.20"), 12));
    }
}
\end{lstlisting}

\subsubsection*{Paso 4 --- Ejecutar la aplicación}

Pulsar el botón verde de \emph{play} junto a la clase principal.
Esperar el log \texttt{Started McvDebugApplication}.

Abrir \url{http://localhost:8080/productos}: aparece el listado con
los tres productos iniciales.

\subsubsection*{Paso 5 --- Colocar 9 breakpoints estratégicos}

Para entender el ciclo de vida, colocar breakpoints (clic en el
margen izquierdo del editor) en:

\begin{enumerate}
\item \texttt{ProductoController.listar()} --- línea con
\texttt{model.addAttribute}.
\item \texttt{ProductoService.listar()} --- línea con
\texttt{return repo.findAll();}.
\item \texttt{ProductoController.mostrarFormulario()}.
\item \texttt{ProductoController.guardar()} --- línea con
\texttt{br.hasErrors()}.
\item \texttt{ProductoController.guardar()} --- línea con
\texttt{service.guardar(p);}.
\item \texttt{ProductoService.guardar(Producto)} --- la línea con
\texttt{repo.save(p);}.
\item \texttt{ProductoController.editar()}.
\item \texttt{ProductoController.eliminar()}.
\item \texttt{ProductoForm} --- la línea de la anotación
\texttt{@NotBlank} (para ver el orden de validación).
\end{enumerate}

\subsubsection*{Paso 6 --- Reiniciar en modo Debug}

Detener la app. Pulsar el ícono \emph{bug} (\texttt{Shift+F9}) para
arrancar en modo debug. La app está lista para depurar.

\subsubsection*{Paso 7 --- Trazar una petición GET (listar)}

Con la aplicación corriendo, el estudiante debe trazar una
petición \texttt{GET} completa observando cómo se mueve por cada
capa.


\begin{enumerate}
\item En el navegador, recargar
\url{http://localhost:8080/productos}.
\item IntelliJ se detiene en breakpoint 1 (\texttt{listar()}).
\item Observar el panel \emph{Variables}: aparece \texttt{model}
(vacío) y los parámetros del método.
\item Pulsar \texttt{F8} (Step Over): se ejecuta
\texttt{model.addAttribute(\dots, service.listar())}.
\item Pulsar \texttt{F7} (Step Into) en \texttt{service.listar()}:
salta al breakpoint 2.
\item En \emph{Variables} ver que \texttt{repo} es un proxy de Spring.
\item Pulsar \texttt{F8}: se ejecuta \texttt{repo.findAll()};
en la consola aparece el SQL: \texttt{select * from productos}.
\item Pulsar \texttt{F9} (Resume): la petición continúa, Thymeleaf
renderiza la vista, el navegador recibe el HTML.
\end{enumerate}

\subsubsection*{Paso 8 --- Trazar una petición POST con error de validación}

El segundo trazado es una petición \texttt{POST} con datos
inválidos: el objetivo es observar cómo \texttt{BindingResult}
captura los errores y la vista los muestra.


\begin{enumerate}
\item Ir a \url{http://localhost:8080/productos/nuevo}.
\item Llenar solo el campo «Stock» (dejar nombre y precio vacíos).
\item Pulsar Guardar.
\item IntelliJ se detiene en breakpoint 4 (\texttt{br.hasErrors()}).
\item Inspeccionar \texttt{br}: el panel \emph{Variables} muestra
una lista de errores (\texttt{nombre: NotBlank},
\texttt{precio: NotNull}).
\item Pulsar \texttt{F8}: como hay errores, retorna
\texttt{\char34{}productos/formulario\char34{}} con los errores marcados en la vista.
\end{enumerate}

\subsubsection*{Paso 9 --- Trazar un POST exitoso}

El tercer trazado es un POST exitoso: observa cómo se
guarda en BD y cómo se aplica el patrón POST-Redirect-GET para
evitar reenvíos del formulario.


\begin{enumerate}
\item Recargar el formulario; ahora llenar todos los campos
correctamente.
\item Pulsar Guardar.
\item IntelliJ se detiene en breakpoint 4 (\texttt{br.hasErrors()}).
\item \texttt{br.hasErrors()} es \texttt{false}; pulsar \texttt{F8}.
\item Salta al breakpoint 5: a punto de ejecutar
\texttt{service.guardar(p)}.
\item Pulsar \texttt{F7} (Step Into): salta al breakpoint 6
(\texttt{repo.save(p)}).
\item En \emph{Variables}, observar que \texttt{p.id} es \texttt{null}
ANTES del save.
\item Pulsar \texttt{F8}: \texttt{p.id} ahora tiene un valor (la BD
asignó la clave primaria).
\item En la consola, ver el \texttt{INSERT INTO productos} generado.
\item Pulsar \texttt{F9}: la respuesta es un redirect a
\texttt{/productos}, el navegador hace un GET nuevo, listamos
todos los productos (vuelve a parar en breakpoint 1).
\end{enumerate}

\subsubsection*{Paso 10 --- Inspeccionar la BD H2 en vivo}

La consola H2 permite inspeccionar la BD en tiempo real
durante la depuración. Esto es invaluable para verificar que los
datos llegaron a persistirse correctamente.


\begin{enumerate}
\item Sin detener la app, abrir \url{http://localhost:8080/h2} en otra
pestaña.
\item En \emph{JDBC URL} introducir: \texttt{jdbc:h2:mem:appweb}.
User: \texttt{sa}, sin password.
\item Pulsar Connect.
\item Ejecutar \texttt{SELECT * FROM PRODUCTOS;}: aparecen los
productos iniciales más los que se hayan creado.
\end{enumerate}

\subsubsection*{Resultado esperado}

El estudiante puede:
\begin{itemize}
\item Articular el flujo: \texttt{Controller $\to$ Service $\to$
Repository $\to$ BD $\to$ Service $\to$ Controller $\to$ Vista}.
\item Identificar dónde ocurre la validación
(\texttt{@Valid + BindingResult}).
\item Explicar el patrón PRG y por qué se usa.
\item Distinguir Modelo (Producto/Service/Repository) de Vista
(Thymeleaf) de Controller (handler HTTP).
\item Ver el SQL generado por JPA en tiempo real en la consola.
\end{itemize}

\subsubsection*{Reflexión final escrita}

Responder en máximo 250 palabras: ¿Qué pasaría si el controller
llamara directamente al \texttt{ProductoRepository} sin pasar por el
\texttt{ProductoService}? ¿Qué se rompería? ¿Qué clases serían más
difíciles de testear?

\textbf{Esta práctica es la \emph{Sumativa GA-Práctica en clase} de
la semana 14.}
\end{cajaejercicio}

\section{Buenas prácticas profesionales en MVC}

La caja siguiente recoge las reglas profesionales para
escribir MVC mantenible. Aplican tanto a Spring Boot como a
ASP.NET Core y Laravel.


\begin{cajabuenas}{Quince reglas profesionales 2026}
\begin{enumerate}
\item Controllers \emph{thin}: máximo 50 líneas por método de
controller, idealmente 10-20.
\item Toda la lógica de negocio en servicios, no en controllers ni
en vistas.
\item DTOs separados de entidades JPA: nunca exponer entidades
directamente al cliente.
\item Validación con Bean Validation (\texttt{@NotBlank},
\texttt{@Size}, \texttt{@Email}) sobre validación manual.
\item Patrón PRG (Post-Redirect-Get) para evitar reenvíos al pulsar
F5.
\item Flash attributes para mensajes de éxito/error en redirects.
\item Manejo global de excepciones con
\texttt{@ControllerAdvice}/\texttt{@RestControllerAdvice}.
\item Vistas (templates) sin lógica compleja: solo iteración,
condicionales simples y formato.
\item URLs RESTful: \texttt{/productos} para listar, \texttt{/productos/123}
para uno específico, métodos HTTP semánticos.
\item Repositorios definidos como interfaces (Spring Data) para
facilitar tests.
\item Servicios anotados \texttt{@Transactional} en los métodos que
mutan datos.
\item Logging estructurado en cada capa: controller (acceso),
servicio (negocio), repositorio (datos).
\item Tests unitarios de servicio con mock de repositorio; tests de
integración con MockMvc para controllers.
\item Documentación de la API REST con OpenAPI/Swagger.
\item Internacionalización (i18n) desde el inicio: archivos
\texttt{messages\_es.properties}, \texttt{messages\_en.properties}.
\end{enumerate}
\end{cajabuenas}

\section{Contraste bibliográfico de los conceptos clave}

La separación de responsabilidades que propone MVC se ha discutido
profundamente en la literatura. \cite{galloway2011mvc} es la
referencia clásica para entender MVC del lado servidor en
\texttt{ASP.NET}. \cite{walls2022spring} y \cite{spilca2024spring}
muestran cómo Spring MVC implementa el patrón con
\texttt{DispatcherServlet}, \texttt{HandlerMapping} y
\texttt{ViewResolver}. En el ámbito del frontend, las arquitecturas
modernas (Angular, React) han evolucionado hacia variantes
\emph{component-based} que rompen la estricta tripartición MVC pero
mantienen el principio de separación
\cite{richards2020fundamentals}. \cite{richardson2007restful} documenta
cómo el patrón \emph{API REST + SPA} es hoy el reemplazo dominante
del MVC server-side tradicional, con tres ventajas: equipos
desacoplados (frontend/backend), reutilización del backend para
varios clientes (web, móvil, IoT) y mejor experiencia interactiva.

\section{Ejercicio resuelto adicional con resultado visual}

Como cierre, el siguiente ejercicio resuelto adicional
implementa un CRUD completo con Thymeleaf y se complementa con un
propuesto que lo migra a SPA Angular.


\begin{cajaejemplo}{Ejercicio resuelto 14.1 --- App MVC de tareas con Spring Boot y Thymeleaf}
\textbf{Enunciado.} Implementar una aplicación MVC server-side de
gestión de tareas (\emph{to-do list}) con Spring Boot 3.4, Thymeleaf
y H2. Permitir listar, crear, completar y eliminar tareas.

\textbf{Modelo (\texttt{Tarea.java}):}
\begin{lstlisting}[language=Java,caption={Ejemplo de Java},label=lst:14.11]
@Entity @Table(name="tareas")                               // entidad JPA (mapea a tabla)
public class Tarea {                                        // declaracion de clase publica
  @Id @GeneratedValue(strategy=GenerationType.IDENTITY)     // campo clave primaria
  private Long id;                                          // campo privado
  @NotBlank @Size(max=120) private String titulo;           // validacion: no nulo y no vacio
  private boolean completada;                               // campo privado
  // getters/setters
}
\end{lstlisting}

\textbf{Controlador (\texttt{TareaController.java}):}
\begin{lstlisting}[language=Java,caption={Ejemplo de Java},label=lst:14.12]
@Controller @RequestMapping("/tareas")                      // controlador MVC server-side (devuelve vistas)
public class TareaController {                              // declaracion de clase publica
  private final TareaRepository repo;                       // campo inmutable (mejor practica)
  // metodo publico
  // metodo publico
  public TareaController(TareaRepository r) { this.repo = r; }

  @GetMapping                                               // endpoint HTTP GET
  public String listar(Model m) {                           // metodo publico
    m.addAttribute("tareas", repo.findAll());
    m.addAttribute("nueva", new Tarea());
    return "tareas";                                        // devuelve el resultado al llamador
  }

  @PostMapping                                              // endpoint HTTP POST (crear)
  // metodo publico
  // metodo publico
  public String crear(@Valid @ModelAttribute("nueva") Tarea t,
                      BindingResult br) {
    if (br.hasErrors()) return "tareas";                    // condicional
    repo.save(t);
    return "redirect:/tareas";                              // devuelve el resultado al llamador
  }

  @PostMapping("/{id}/completar")                           // endpoint HTTP POST (crear)
  public String completar(@PathVariable Long id) {          // metodo publico
    repo.findById(id).ifPresent(t -> {
      t.setCompletada(!t.isCompletada());
      repo.save(t);
    });
    return "redirect:/tareas";                              // devuelve el resultado al llamador
  }

  @PostMapping("/{id}/eliminar")                            // endpoint HTTP POST (crear)
  public String eliminar(@PathVariable Long id) {           // metodo publico
    repo.deleteById(id);
    return "redirect:/tareas";                              // devuelve el resultado al llamador
  }
}
\end{lstlisting}

\textbf{Vista (\texttt{templates/tareas.html} --- Thymeleaf):}
\begin{lstlisting}[language=HTML5,caption={Ejemplo de HTML},label=lst:14.13]
<!DOCTYPE html>
<html xmlns:th="http://www.thymeleaf.org" lang="es">
<head><meta charset="UTF-8"><title>Mis tareas</title></head>
<body>
  <h1>Mis tareas</h1>
  <form th:action="@{/tareas}" th:object="${nueva}" method="post">
    <input type="text" th:field="*{titulo}" placeholder="Nueva tarea">
    <button type="submit">Agregar</button>
    <span th:errors="*{titulo}" style="color:red"></span>
  </form>
  <ul>
    <li th:each="t : ${tareas}">
      <span th:text="${t.titulo}"
            th:style="${t.completada} ? 'text-decoration:line-through' : ''"></span>
      <form th:action="@{|/tareas/${t.id}/completar|}" method="post" style="display:inline">
        <button type="submit" th:text="${t.completada} ? 'Reabrir' : 'Completar'"></button>
      </form>
      <form th:action="@{|/tareas/${t.id}/eliminar|}" method="post" style="display:inline">
        <button type="submit">Eliminar</button>
      </form>
    </li>
  </ul>
</body>
</html>
\end{lstlisting}

\textbf{Resultado visual} en el navegador:

\begin{center}
\fbox{\begin{minipage}{0.85\textwidth}
\vspace{0.2cm}
\hspace{0.3cm}\textbf{\Large Mis tareas}\\[0.4cm]
\hspace{0.3cm}\fbox{\texttt{Nueva tarea\ldots\ldots\ldots\ldots\ldots}} \fbox{Agregar}\\[0.5cm]
\hspace{0.3cm}$\bullet$ \emph{Leer cap.\ 14 del libro} \fbox{\small Reabrir} \fbox{\small Eliminar}\\[0.2cm]
\hspace{0.3cm}$\bullet$ Estudiar Spring MVC \fbox{\small Completar} \fbox{\small Eliminar}\\[0.2cm]
\hspace{0.3cm}$\bullet$ Hacer el ejercicio propuesto 14.2 \fbox{\small Completar} \fbox{\small Eliminar}\\[0.2cm]
\vspace{0.1cm}
\end{minipage}}
\end{center}

\textbf{Discusión.} El flujo MVC completo se ve en una sola pantalla:
el navegador hace \texttt{GET /tareas} $\to$
\texttt{DispatcherServlet} despacha a \texttt{TareaController.listar()}
$\to$ el controlador consulta el repositorio (modelo) y prepara los
datos en el \texttt{Model} $\to$ Thymeleaf renderiza la vista con esos
datos $\to$ el navegador recibe HTML. Al hacer clic en
\emph{Completar}, ocurre el patrón \emph{POST-Redirect-GET}
\cite{walls2022spring}: \texttt{POST} ejecuta la acción y el
servidor responde \texttt{302 Found} con \texttt{Location: /tareas},
forzando un nuevo \texttt{GET}. Esto evita el reenvío accidental del
formulario al recargar.
\end{cajaejemplo}

\begin{cajaejercicio}{Ejercicio propuesto 14.2 --- Migrar la app de tareas a SPA Angular + REST}
\textbf{Enunciado.} Tomar la aplicación MVC server-side del ejercicio
14.1 y refactorizarla a una arquitectura de dos niveles: backend Spring
Boot expone una API REST y un frontend Angular 19 SPA consume la API.
Comparar las dos arquitecturas en términos de líneas de código,
acoplamiento, experiencia de usuario y testabilidad.

\textbf{Requisitos:}
\begin{enumerate}
\item Convertir \texttt{TareaController} de \texttt{@Controller} a
\texttt{@RestController}, eliminando Thymeleaf.
\item Endpoints: \texttt{GET /api/v1/tareas},
\texttt{POST /api/v1/tareas}, \texttt{PATCH /api/v1/tareas/\{id\}/completar},
\texttt{DELETE /api/v1/tareas/\{id\}}.
\item Frontend Angular con un componente \texttt{TareasComponent}
que consuma la API mediante \texttt{HttpClient}.
\item Manejar CORS correctamente con
\texttt{@CrossOrigin(\char34{}http://localhost:4200\char34{})} o configuración
global.
\item Documentación OpenAPI con springdoc.
\item Test E2E con Playwright o Cypress que verifique el flujo
completo (crear, completar, eliminar).
\end{enumerate}

\textbf{Producto final:} un documento (1--2 páginas) que compare las
dos arquitecturas y argumente cuándo elegir cada una, citando al
menos a \cite{galloway2011mvc} y \cite{richardson2007restful}.
\end{cajaejercicio}

% =====================================================================
%  CIERRE DE LA PARTE 2 AMPLIADA (semanas 10-14)
%
%  El bloque completo de las semanas 10 a 14 termina aquí.
%  Las semanas 15 a 18 se entregarán en la siguiente petición y
%  ampliarán la Unidad IV (servicios web SOAP/REST, defensa final
%  del PFC, evaluación parcial segundo corte y tutoría de cierre).
%
%  Recordatorio: NO escribir aquí \end{document}; eso debe ir en el
%  archivo final, después de las semanas 15-18.
% =====================================================================

% =====================================================================
%       SEMANA 15: SERVICIOS WEB REST Y SOAP
% =====================================================================
\chapter{Servicios web: arquitecturas REST, SOAP y diseño de APIs}
\label{cap:semana15}

\epigraph{«Una buena API es como una buena novela: invita a explorarla,
revela su lógica gradualmente y nunca traiciona al lector con
inconsistencias arbitrarias.»}{--- \textit{Adaptado de Joshua Bloch}}

\section*{Resultado de aprendizaje de la semana}
Al concluir la semana 15, el estudiante diseña, implementa, documenta
y consume servicios web aplicando el estilo arquitectónico REST según
Roy Fielding \cite{fielding2000rest}, distingue REST de SOAP en sus
casos de uso reales, integra autenticación con JWT y produce
documentación OpenAPI 3.0 conforme a estándares profesionales de la
industria 2026.

\section{Historia de los servicios web}

Los servicios web son la forma estándar de comunicación
entre sistemas distintos. Su historia atraviesa tres generaciones:
SOAP (1998), REST (2000) y GraphQL/gRPC (2015). Conocerlas permite
elegir adecuadamente para cada proyecto.


\begin{cajahistoria}{De RPC a REST --- treinta años de evolución de las APIs distribuidas}
La idea de que un programa pueda invocar funciones de otro a través
de la red precede a la web. En 1981, Bruce Nelson y Andrew Birrell
formalizaron el \emph{Remote Procedure Call} (RPC) en el seminario
del que surgió la primera implementación seria, Sun RPC. El sueño:
hacer que llamar a una función remota sea sintácticamente idéntico
a llamar a una función local.

A finales de los noventa, ese sueño se reformuló en la web. Microsoft
junto con DevelopMentor publicaron en 1998 el primer borrador de
\textbf{SOAP} (Simple Object Access Protocol): RPC sobre HTTP usando
XML como formato de mensaje. SOAP se complementó con \textbf{WSDL}
(Web Services Description Language) en 2001 y con \textbf{UDDI} para
descubrimiento de servicios. La trinidad SOAP+WSDL+UDDI prometía un
ecosistema empresarial donde cualquier servicio sería descubrible y
consumible automáticamente.

La realidad fue más áspera. SOAP era verboso (envoltorios XML
enormes), el ecosistema de \emph{WS-*} (WS-Security, WS-ReliableMessaging,
WS-Transaction\dots) creció hasta volverse impracticable, y los
clientes JavaScript del navegador no podían consumirlo elegantemente.

En 2000 \textbf{Roy Fielding} —uno de los autores principales de
HTTP/1.1 y cofundador de Apache— defendió en su tesis doctoral el
estilo arquitectónico \textbf{REST} (Representational State Transfer)
\cite{fielding2000rest}. REST no era un protocolo: era un conjunto de
\emph{restricciones arquitectónicas} (cliente-servidor, stateless,
cacheable, interfaz uniforme, sistema en capas, código bajo demanda
opcional) que aprovechaban HTTP en lugar de pelearse con él. La
trayectoria a partir de 2005 fue rápida: las APIs públicas masivas
(Flickr, Twitter, Amazon S3, GitHub) eligieron REST. Para 2015, REST
era el estilo dominante en la web. SOAP quedó relegado a sistemas
empresariales legados y a sectores conservadores como banca y salud.

En 2026 conviven cuatro paradigmas de servicios web: \textbf{REST}
(dominante en APIs públicas y SPAs), \textbf{GraphQL} (creado por
Facebook en 2012, popular en frontends complejos), \textbf{gRPC}
(Google, 2015, eficiente en comunicación servicio-servicio) y
\textbf{SOAP} (legado pero vigente en banca, salud, gobierno).
\end{cajahistoria}

\subsection{Tabla evolutiva de los servicios web}

La tabla siguiente recoge los hitos cronológicos de los
servicios web, desde XML-RPC hasta GraphQL y gRPC.


\begin{table}[htbp]
\centering
\caption{Hitos en la historia de los servicios web.}
\footnotesize
\begin{tabularx}{\textwidth}{lll}
\toprule
\textbf{Año} & \textbf{Hito} & \textbf{Aporte}\\
\midrule
1981 & Sun RPC & Primer RPC industrial.\\
1991 & CORBA & Estándar OMG; complejo, propietario.\\
1998 & SOAP 1.0 & RPC sobre HTTP+XML; Microsoft, IBM.\\
2000 & Tesis de Fielding & REST formalizado como estilo arquitectónico.\\
2001 & WSDL 1.1 & Contrato XML para SOAP.\\
2003 & SOAP 1.2 (W3C) & Recomendación oficial de SOAP.\\
2005 & Atom Publishing Protocol & Primer estándar REST formal.\\
2008 & RESTful API patterns & Modelo de Madurez de Richardson.\\
2009 & Twitter API REST & APIs públicas masivas adoptan REST.\\
2012 & GraphQL (Facebook) & Cliente decide qué campos pedir.\\
2014 & OpenAPI 2.0 (Swagger) & Estándar para describir APIs REST.\\
2015 & gRPC (Google) & RPC moderno con HTTP/2 y Protocol Buffers.\\
2017 & OpenAPI 3.0 & Versión moderna del estándar.\\
2021 & OpenAPI 3.1 & Compatibilidad total con JSON Schema 2020-12.\\
2024 & gRPC-Web estable & gRPC accesible desde navegador.\\
\bottomrule
\end{tabularx}
\end{table}

\section{Las seis restricciones de la arquitectura REST}

Roy Fielding identificó seis restricciones que un sistema debe
cumplir para considerarse RESTful. No son sugerencias: son
restricciones arquitectónicas con consecuencias específicas.

\begin{cajadefinicion}{Las restricciones REST de Fielding (2000)}
\begin{enumerate}
\item \textbf{Cliente-servidor.} Separación clara de
responsabilidades: el cliente no conoce la persistencia, el servidor
no conoce la presentación.

\item \textbf{Stateless.} El servidor NO guarda estado de sesión
entre peticiones. Cada petición lleva toda la información necesaria
para procesarla. Permite escalar horizontalmente sin sticky sessions.

\item \textbf{Cacheable.} Las respuestas indican explícitamente si
pueden cachearse y por cuánto tiempo (\texttt{Cache-Control},
\texttt{ETag}). Mejora rendimiento y escalabilidad.

\item \textbf{Interfaz uniforme.} Cuatro sub-restricciones:
\begin{itemize}
\item Identificación de recursos por URI.
\item Manipulación a través de representaciones (JSON, XML).
\item Mensajes auto-descriptivos (cabeceras explícitas).
\item HATEOAS (Hypermedia as the Engine of Application State):
los recursos enlazan a otros relacionados.
\end{itemize}

\item \textbf{Sistema en capas.} El cliente no sabe si habla con el
servidor final o con un proxy/balanceador. Permite intermediarios
(CDN, gateways, balanceadores).

\item \textbf{Código bajo demanda} (opcional). El servidor puede
extender al cliente enviando código ejecutable (típicamente
JavaScript).
\end{enumerate}
\end{cajadefinicion}

\subsection{Modelo de Madurez de Richardson}

Leonard Richardson propuso en 2008 un modelo escalonado para evaluar
cuán RESTful es una API:

\begin{table}[htbp]
\centering
\caption{Modelo de Madurez de Richardson para APIs REST.}
\footnotesize
\begin{tabularx}{\textwidth}{llX}
\toprule
\textbf{Nivel} & \textbf{Nombre} & \textbf{Características}\\
\midrule
0 & «The Swamp of POX» & XML/JSON sobre HTTP, un solo endpoint, todo es POST.\\
1 & Recursos & URLs distintas por recurso (\texttt{/usuarios/42}), pero todo POST.\\
2 & Verbos HTTP & Uso semántico de GET, POST, PUT, DELETE; códigos correctos.\\
3 & HATEOAS & Las respuestas incluyen enlaces de navegación a recursos relacionados.\\
\bottomrule
\end{tabularx}
\end{table}

\textbf{En 2026, la mayoría de APIs profesionales se sitúan en
Nivel 2}. El nivel 3 (HATEOAS) es elegante en teoría pero
infrecuente en la práctica: añade complejidad sin que la mayoría de
clientes lo aproveche realmente.

\section{Diseño de URIs y verbos HTTP}

La definición de las URIs y el uso de los verbos HTTP son
las dos decisiones más visibles de una API REST. Hacerlas mal se
nota desde el primer endpoint.


\subsection{Convenciones de nomenclatura}

Las URIs de una API REST siguen convenciones precisas: son
sustantivos en plural, sin verbos, jerárquicas. La lista siguiente
las recoge con ejemplos.

\textbf{Ejemplo corto comentado (URIs REST canónicas) (Listado~\ref{lst:15.1}):}
\begin{lstlisting}[language=bash,numbers=none,caption={Ejemplo de Shell},label=lst:15.1]
# BIEN: sustantivo en plural, jerarquia clara
GET    /api/v1/usuarios              # listar todos
GET    /api/v1/usuarios/42           # uno especifico
POST   /api/v1/usuarios              # crear nuevo
PUT    /api/v1/usuarios/42           # actualizar completo
PATCH  /api/v1/usuarios/42           # actualizar parcial
DELETE /api/v1/usuarios/42           # eliminar
GET    /api/v1/usuarios/42/pedidos   # recurso anidado

# MAL: verbos en URI, singular, sin version
GET    /getUsuario?id=42             # verbo en URI
POST   /crearUsuario                 # accion como recurso
GET    /usuario/42                   # singular en lugar de plural
\end{lstlisting}


\begin{cajabuenas}{Doce reglas para URIs RESTful profesionales}
\begin{enumerate}
\item \textbf{Sustantivos en plural} para colecciones:
\texttt{/productos}, no \texttt{/getProducto}.
\item \textbf{Recursos jerárquicos}:
\texttt{/usuarios/42/pedidos/7/items}.
\item \textbf{Sin verbos en la URI}: el verbo va en el método HTTP.
\item \textbf{Minúsculas siempre}, palabras separadas por guiones:
\texttt{/ordenes-pendientes}, no \texttt{/OrdenesPendientes}.
\item \textbf{Sin extensiones de archivo}: \texttt{/usuarios/42}, no
\texttt{/usuarios/42.json}.
\item \textbf{Versionado en la URL} con prefijo \texttt{/api/v1/}:
\texttt{/api/v1/productos}.
\item \textbf{Filtros por query string}:
\texttt{/productos?categoria=cafe\&min\_precio=2}.
\item \textbf{Paginación estándar}:
\texttt{?page=3\&size=20} o \texttt{?offset=60\&limit=20}.
\item \textbf{Ordenamiento por query}:
\texttt{?sort=precio,-nombre} (\texttt{-} para descendente).
\item \textbf{Búsqueda con parámetro \texttt{q}}:
\texttt{?q=galletti}.
\item \textbf{IDs estables} (UUIDs preferidos sobre auto-incrementales
en APIs públicas).
\item \textbf{Sub-recursos para acciones}: \texttt{/pedidos/7/cancelar}
con POST.
\end{enumerate}
\end{cajabuenas}

\subsection{Mapeo verbo-acción canónico}

Cada verbo HTTP tiene un significado canónico y debe
mapearse a la acción CRUD correspondiente. Confundirlos rompe la
semántica REST y genera APIs confusas.


\begin{table}[htbp]
\centering
\caption{Mapeo verbos HTTP a operaciones CRUD.}
\footnotesize
\begin{tabularx}{\textwidth}{llXX}
\toprule
\textbf{Verbo} & \textbf{URI} & \textbf{Operación} & \textbf{Códigos típicos}\\
\midrule
GET & /productos & Listar & 200 OK \\
GET & /productos/42 & Obtener uno & 200, 404 \\
POST & /productos & Crear & 201 Created (con \texttt{Location}) \\
PUT & /productos/42 & Reemplazar completo & 200, 204 \\
PATCH & /productos/42 & Actualización parcial & 200, 204 \\
DELETE & /productos/42 & Eliminar & 204 No Content, 404 \\
HEAD & /productos/42 & Solo cabeceras & 200, 404 \\
OPTIONS & /productos & Verbos soportados & 200 (CORS) \\
\bottomrule
\end{tabularx}
\end{table}

\subsection{Idempotencia y seguridad}

Dos propiedades cruciales de los métodos HTTP:

\begin{itemize}
\item \textbf{Seguro} (\emph{safe}): no modifica estado del servidor.
GET, HEAD, OPTIONS son seguros.
\item \textbf{Idempotente}: ejecutar la operación N veces produce el
mismo efecto que ejecutarla una vez. GET, PUT, DELETE, HEAD, OPTIONS
son idempotentes. POST y PATCH \emph{no} lo son.
\end{itemize}

Los \emph{proxies} y reintentos automáticos asumen estas propiedades:
si su API rompe estas reglas, los clientes sufrirán bugs sutiles.

\section{Implementación de APIs REST en las tres plataformas}

Las tres plataformas del curso (Spring Boot, ASP.NET Core,
PHP) construyen APIs REST con APIs limpias y modernas. Las
subsecciones siguientes muestran la implementación canónica en cada
una.


\subsection{Spring Boot 3.4 con Spring Web}

Spring Boot expone APIs REST con \texttt{@RestController} y
anotaciones por método HTTP. Es la forma idiomática en el ecosistema
Java.


\begin{cajaejemplo}{Ejemplo 15.1 --- API REST de pedidos completa en Spring Boot}

\textbf{Enunciado.} Implementar una API REST de productos con Spring Boot, autenticación JWT, documentación OpenAPI 3.0 y Swagger UI interactivo.

\begin{lstlisting}[style=java,caption={PedidoController.java},label=lst:15.2]
package ec.edu.uteq.appweb.controller;                      // paquete del archivo

import ec.edu.uteq.appweb.dto.*;                            // importacion de clase externa
import ec.edu.uteq.appweb.service.PedidoService;            // importacion de clase externa
import io.swagger.v3.oas.annotations.Operation;             // importacion de clase externa
import io.swagger.v3.oas.annotations.tags.Tag;              // importacion de clase externa
import jakarta.validation.Valid;                            // importacion de clase externa
import lombok.RequiredArgsConstructor;                      // importacion de clase externa
import org.springframework.data.domain.Page;                // importacion de clase externa
import org.springframework.data.domain.Pageable;            // importacion de clase externa
import org.springframework.http.ResponseEntity;             // importacion de clase externa
import org.springframework.web.bind.annotation.*;           // importacion de clase externa
import java.net.URI;                                        // importacion de clase externa

@RestController                                             // expone endpoints REST con JSON
@RequestMapping("/api/v1/pedidos")                          // ruta base del controlador
@RequiredArgsConstructor
// agrupa endpoints en Swagger UI
// agrupa endpoints en Swagger UI
@Tag(name = "Pedidos", description = "Gestion de pedidos del sistema")
public class PedidoController {                             // declaracion de clase publica

    private final PedidoService service;                    // campo inmutable (mejor practica)

    @GetMapping                                             // endpoint HTTP GET
    @Operation(summary = "Lista paginada de pedidos")       // documenta el endpoint en OpenAPI
    public Page<PedidoDto> listar(                          // metodo publico
            @RequestParam(required = false) String estado,  // extrae parametro de query string
            Pageable pageable) {
        return service.buscar(estado, pageable);            // devuelve el resultado al llamador
    }

    @GetMapping("/{id}")                                    // endpoint HTTP GET
    @Operation(summary = "Obtener pedido por ID")           // documenta el endpoint en OpenAPI
    // metodo publico
    // metodo publico
    public ResponseEntity<PedidoDto> obtener(@PathVariable Long id) {
        return ResponseEntity.ok(service.obtener(id));      // devuelve el resultado al llamador
    }

    @PostMapping                                            // endpoint HTTP POST (crear)
    @Operation(summary = "Crear nuevo pedido")              // documenta el endpoint en OpenAPI
    public ResponseEntity<PedidoDto> crear(                 // metodo publico
            @Valid @RequestBody PedidoCreateRequest req) {  // activa validacion Bean Validation
        PedidoDto creado = service.crear(req);
        return ResponseEntity                               // devuelve el resultado al llamador
            .created(URI.create("/api/v1/pedidos/" + creado.id()))
            .body(creado);
    }

    @PatchMapping("/{id}")                                  // endpoint HTTP PATCH (actualizar parcial)
    @Operation(summary = "Actualizacion parcial")           // documenta el endpoint en OpenAPI
    public PedidoDto actualizar(@PathVariable Long id,      // metodo publico
                                // activa validacion Bean Validation
                                // activa validacion Bean Validation
                                @Valid @RequestBody PedidoUpdateRequest req) {
        return service.actualizar(id, req);                 // devuelve el resultado al llamador
    }

    @PostMapping("/{id}/cancelar")                          // endpoint HTTP POST (crear)
    @Operation(summary = "Cancelar un pedido")              // documenta el endpoint en OpenAPI
    public PedidoDto cancelar(@PathVariable Long id,        // metodo publico
                              // extrae parametro de query string
                              // extrae parametro de query string
                              @RequestParam(required = false) String motivo) {
        return service.cancelar(id, motivo);                // devuelve el resultado al llamador
    }

    @DeleteMapping("/{id}")                                 // endpoint HTTP DELETE (eliminar)
    @Operation(summary = "Eliminar pedido")                 // documenta el endpoint en OpenAPI
    // metodo publico
    // metodo publico
    public ResponseEntity<Void> eliminar(@PathVariable Long id) {
        service.eliminar(id);
        return ResponseEntity.noContent().build();          // devuelve el resultado al llamador
    }
}
\end{lstlisting}

\textbf{Manejo global de excepciones (RFC 7807) (Listado~\ref{lst:15.3}):}
\begin{lstlisting}[style=java,caption={GlobalExceptionHandler.java},label=lst:15.3]
package ec.edu.uteq.appweb.exception;                       // paquete del archivo

import jakarta.servlet.http.HttpServletRequest;             // importacion de clase externa
import org.springframework.http.HttpStatus;                 // importacion de clase externa
import org.springframework.http.ProblemDetail;              // importacion de clase externa
import org.springframework.http.ResponseEntity;             // importacion de clase externa
// importacion de clase externa
// importacion de clase externa
import org.springframework.web.bind.MethodArgumentNotValidException;
import org.springframework.web.bind.annotation.*;           // importacion de clase externa
import java.net.URI;                                        // importacion de clase externa
import java.time.Instant;                                   // importacion de clase externa
import java.util.*;                                         // importacion de clase externa

@RestControllerAdvice                                       // expone endpoints REST con JSON
public class GlobalExceptionHandler {                       // declaracion de clase publica

    @ExceptionHandler(NoSuchElementException.class)
    public ResponseEntity<ProblemDetail> noEncontrado(      // metodo publico
            NoSuchElementException ex, HttpServletRequest req) {
        ProblemDetail pd = ProblemDetail.forStatusAndDetail(
            HttpStatus.NOT_FOUND, ex.getMessage());
        pd.setType(URI.create("https://uteq.edu.ec/errors/no-encontrado"));
        pd.setTitle("Recurso no encontrado");
        pd.setInstance(URI.create(req.getRequestURI()));
        pd.setProperty("timestamp", Instant.now());
        // devuelve el resultado al llamador
        // devuelve el resultado al llamador
        return ResponseEntity.status(HttpStatus.NOT_FOUND).body(pd);
    }

    @ExceptionHandler(MethodArgumentNotValidException.class)
    public ResponseEntity<ProblemDetail> validacion(        // metodo publico
            MethodArgumentNotValidException ex, HttpServletRequest req) {
        Map<String, String> errores = new LinkedHashMap<>();
        ex.getBindingResult().getFieldErrors().forEach(e ->
            errores.put(e.getField(), e.getDefaultMessage()));

        ProblemDetail pd = ProblemDetail.forStatusAndDetail(
            HttpStatus.UNPROCESSABLE_ENTITY,
            "Errores de validacion en " + errores.size() + " campo(s)");
        pd.setType(URI.create("https://uteq.edu.ec/errors/validacion"));
        pd.setTitle("Datos invalidos");
        pd.setInstance(URI.create(req.getRequestURI()));
        pd.setProperty("errores", errores);
        // devuelve el resultado al llamador
        // devuelve el resultado al llamador
        return ResponseEntity.unprocessableEntity().body(pd);
    }
}
\end{lstlisting}
\end{cajaejemplo}

\subsection{ASP.NET Core 8}

ASP.NET Core utiliza el atributo \texttt{[ApiController]} y
la herencia de \texttt{ControllerBase} para construir APIs REST.
Combinado con \emph{model binding} y \emph{model validation}
automáticos, ofrece una experiencia muy productiva.


\begin{cajaejemplo}{Ejemplo 15.2 --- Mismo recurso en ASP.NET Core 8}

\textbf{Enunciado.} Implementar el mismo recurso de productos visto antes pero en ASP.NET Core 8, para comparar idiomas: atributos \texttt{[ApiController]}, \texttt{[HttpGet]}, retornos \texttt{Ok()}/\texttt{NotFound()}.

\begin{lstlisting}[style=cs,caption={PedidosController.cs},label=lst:15.4]
using Microsoft.AspNetCore.Mvc;                             // importa namespace
using AppwebApi.Services;                                   // importa namespace
using AppwebApi.Dtos;                                       // importa namespace

namespace AppwebApi.Controllers;                            // namespace del archivo

[ApiController]                                             // controlador de Web API (validacion automatica)
[Route("api/v1/[controller]")]                              // ruta base del controlador
[Produces("application/json")]
public class PedidosController : ControllerBase             // declaracion de clase publica
{
    private readonly IPedidoService _svc;                   // metodo o miembro privado

    // metodo o miembro publico
    // metodo o miembro publico
    public PedidosController(IPedidoService svc) => _svc = svc;

    /// <summary>Lista paginada de pedidos</summary>
    [HttpGet]                                               // endpoint HTTP GET
    [ProducesResponseType(typeof(PageDto<PedidoDto>), 200)]
    // metodo o miembro publico
    // metodo o miembro publico
    public async Task<ActionResult<PageDto<PedidoDto>>> Listar(
        // mapea query string
        // mapea query string
        [FromQuery] string? estado, [FromQuery] int page = 1,
        [FromQuery] int size = 20)                          // mapea query string
        => Ok(await _svc.BuscarAsync(estado, page, size));

    [HttpGet("{id:long}")]                                  // endpoint HTTP GET
    [ProducesResponseType(typeof(PedidoDto), 200)]
    [ProducesResponseType(404)]
    // metodo o miembro publico
    // metodo o miembro publico
    public async Task<ActionResult<PedidoDto>> Obtener(long id)
    {
        var p = await _svc.ObtenerAsync(id);
        return p is null ? NotFound() : Ok(p);              // devuelve el resultado al llamador
    }

    [HttpPost]                                              // endpoint HTTP POST (crear)
    [ProducesResponseType(typeof(PedidoDto), 201)]
    [ProducesResponseType(422)]
    public async Task<ActionResult<PedidoDto>> Crear(       // metodo o miembro publico
        [FromBody] PedidoCreateRequest req)                 // mapea body JSON al parametro
    {
        var creado = await _svc.CrearAsync(req);
        return CreatedAtAction(nameof(Obtener),             // HTTP 201 (recurso creado)
            new { id = creado.Id }, creado);
    }

    [HttpDelete("{id:long}")]                               // endpoint HTTP DELETE (eliminar)
    [ProducesResponseType(204)]
    [ProducesResponseType(404)]
    public async Task<IActionResult> Eliminar(long id)      // metodo o miembro publico
    {
        await _svc.EliminarAsync(id);
        return NoContent();                                 // HTTP 204 (sin cuerpo)
    }
}
\end{lstlisting}
\end{cajaejemplo}

\subsection{PHP 8.3 con Slim Framework 4}

PHP construye APIs REST profesionales mediante microframeworks
como Slim Framework 4. El siguiente ejemplo construye una API
minimalista pero idiomática.


\begin{cajaejemplo}{Ejemplo 15.3 --- API REST minimalista en PHP con Slim 4}

\textbf{Enunciado.} Construir una API REST minimalista en PHP usando el microframework Slim 4, ideal para servicios pequeños o microservicios PHP.

\begin{lstlisting}[style=php,caption={public/index.php},label=lst:15.5]
<?php
declare(strict_types=1);
require __DIR__ . '/../vendor/autoload.php';                // incluye otro archivo PHP (fatal si falla)

use Slim\Factory\AppFactory;                                // importa clase del namespace
use Psr\Http\Message\ServerRequestInterface as Req;         // importa clase del namespace
use Psr\Http\Message\ResponseInterface as Res;              // importa clase del namespace

$app = AppFactory::create();
$app->addBodyParsingMiddleware();

// GET /api/v1/pedidos
$app->get('/api/v1/pedidos', function (Req $req, Res $res) {
    $pdo = obtenerPDO();
    $estado = $req->getQueryParams()['estado'] ?? null;
    $sql = "SELECT id, total, estado, creado FROM pedidos";
    $params = [];
    // condicional
    // condicional
    if ($estado) { $sql .= " WHERE estado = ?"; $params[] = $estado; }
    $sql .= " ORDER BY creado DESC LIMIT 50";
    $stmt = $pdo->prepare($sql);                            // prepara consulta SQL (defensa anti-SQLi)
    $stmt->execute($params);                                // ejecuta la consulta con parametros enlazados
    // recupera todas las filas restantes
    // recupera todas las filas restantes
    $res->getBody()->write(json_encode(['data' => $stmt->fetchAll()]));
    // retorna valor
    // retorna valor
    return $res->withHeader('Content-Type','application/json');
});

// POST /api/v1/pedidos
$app->post('/api/v1/pedidos', function (Req $req, Res $res) {
    $body = $req->getParsedBody() ?? [];
    // condicional
    // condicional
    if (empty($body['total']) || !is_numeric($body['total'])) {
        $err = ['type'=>'/errors/validacion',
                'title'=>'Datos invalidos','status'=>422,
                'errors'=>['total'=>'es requerido y numerico']];
        $res->getBody()->write(json_encode($err));
        return $res->withStatus(422)                        // retorna valor
            ->withHeader('Content-Type','application/problem+json');
    }
    $pdo = obtenerPDO();
    $stmt = $pdo->prepare(                                  // prepara consulta SQL (defensa anti-SQLi)
        "INSERT INTO pedidos (total, estado, creado) VALUES (?,?,NOW())");
    $stmt->execute([$body['total'], 'PENDIENTE']);          // ejecuta la consulta con parametros enlazados
    $id = $pdo->lastInsertId();
    $data = ['id'=>$id,'total'=>$body['total'],'estado'=>'PENDIENTE'];
    $res->getBody()->write(json_encode($data));
    return $res->withStatus(201)                            // retorna valor
        ->withHeader('Content-Type','application/json')
        ->withHeader('Location',"/api/v1/pedidos/$id");
});

$app->run();
\end{lstlisting}
\end{cajaejemplo}

\section{REST vs SOAP: la comparativa definitiva}

REST y SOAP son los dos modelos arquitectónicos de servicios
web más utilizados. Comparar sus características objetivamente es
clave para elegir correctamente en cada proyecto.


\begin{cajacompara}{REST vs SOAP --- comparación práctica}
\centering\footnotesize
\begin{tabularx}{\textwidth}{lXX}
\toprule
\textbf{Criterio} & \textbf{REST (mainstream 2026)} & \textbf{SOAP (legado)}\\
\midrule
Formato & JSON (texto compacto) & XML envelope con header y body (verboso)\\
Transporte & HTTP/HTTPS & HTTP, SMTP, FTP, JMS\\
Contrato & OpenAPI 3.0 (opcional) & WSDL (obligatorio)\\
Estado & Stateless por restricción & Stateful o stateless\\
Tamaño mensaje & 3--5x más compacto & Pesado\\
Seguridad & HTTPS + JWT + OAuth 2.0 & WS-Security (más complejo)\\
Desde JavaScript & Trivial (fetch nativo) & Difícil (requiere SOAP client)\\
Caché HTTP nativa & Sí & No (POSTs)\\
Casos típicos & APIs públicas, microservicios, SPA & Banca, salud, B2B legado\\
Curva de aprendizaje & Baja & Alta\\
Tooling 2026 & Excelente & Maduro pero envejecido\\
\bottomrule
\end{tabularx}
\end{cajacompara}

\subsection{Cuándo aún tiene sentido SOAP en 2026}

Aunque REST es la opción dominante en 2026, hay escenarios
donde SOAP todavía tiene ventajas reales. La lista siguiente los
recoge.


\begin{itemize}
\item \textbf{Sistemas financieros nacionales}: muchos sistemas
bancarios mantienen interfaces SOAP por inercia regulatoria. En
Ecuador, los servicios SPI del Banco Central usan SOAP.
\item \textbf{Salud}: HL7 v3 sobre SOAP sigue vigente en hospitales.
\item \textbf{B2B con grandes corporaciones}: SAP, Oracle EBS y
sistemas similares exponen SOAP por defecto.
\item \textbf{Cuando se necesita WS-AtomicTransaction} (transacciones
distribuidas estrictas).
\end{itemize}

\subsection{Ejemplo SOAP mínimo}

Aunque el curso no usa SOAP en el PFC, conviene haber visto
al menos una petición SOAP completa para reconocerla cuando aparezca
en código heredado.

\textbf{Ejemplo corto comentado (petición SOAP) (Listado~\ref{lst:15.6}):}
\begin{lstlisting}[language=HTML5,numbers=none,caption={Ejemplo de HTML},label=lst:15.6]
<!-- POST /servicios/clima HTTP/1.1                              -->
<!-- Content-Type: text/xml; charset=utf-8                       -->
<!-- SOAPAction: "obtenerTemperatura"                            -->
<?xml version="1.0"?>
<!-- Envelope: contenedor obligatorio de todo SOAP                -->
<soap:Envelope xmlns:soap="http://schemas.xmlsoap.org/soap/envelope/">
  <!-- Body: contenido principal de la peticion                   -->
  <soap:Body>
    <!-- Operacion definida en el WSDL                            -->
    <obtenerTemperatura xmlns="http://uteq.edu.ec/clima">
      <ciudad>Quevedo</ciudad>
    </obtenerTemperatura>
  </soap:Body>
</soap:Envelope>
\end{lstlisting}


\begin{cajaejemplo}{Ejemplo 15.4 --- Cliente SOAP en PHP que consume servicio del SRI (estructura conceptual)}

\textbf{Enunciado.} Mostrar la estructura de un cliente SOAP en PHP que consume un servicio del SRI ecuatoriano (caso real del contexto local), demostrando el protocolo en producción.

\begin{lstlisting}[style=php,caption={consultar-ruc.php},label=lst:15.7]
<?php
declare(strict_types=1);

// PHP tiene cliente SOAP nativo: SoapClient
$wsdl = 'https://servicios.sri.gob.ec/aut.../wsdl';

try {
    $client = new SoapClient($wsdl, [
        'soap_version' => SOAP_1_2,
        'trace'        => true,
        'exceptions'   => true,
        'connection_timeout' => 10,
    ]);

    // Llamada SOAP: el metodo se invoca como si fuera local
    $resp = $client->consultarRuc(['numeroRuc' => '0992345678001']);

    print_r($resp);

    // Para debug del XML enviado/recibido:
    // echo $client->__getLastRequest();
    // echo $client->__getLastResponse();

} catch (SoapFault $e) {
    error_log("SOAP fault: " . $e->getMessage());
    http_response_code(502);
    // imprime al output del servidor
    // imprime al output del servidor
    echo json_encode(['error' => 'Servicio externo no disponible']);
}
\end{lstlisting}
\end{cajaejemplo}

\section{Documentación con OpenAPI 3.0 (Swagger)}

Una API sin documentación es inutilizable; una con
documentación a mano queda desincronizada al primer cambio.
OpenAPI 3.0 + Swagger UI resuelve esto generando documentación viva
desde el propio código.


\subsection{¿Qué es OpenAPI?}

OpenAPI 3.0 (antes Swagger) es la especificación estándar de la
industria para describir APIs REST. Permite:

\begin{itemize}
\item Documentación interactiva con Swagger UI.
\item Generación automática de clientes en cualquier lenguaje.
\item Pruebas y mocks automáticos.
\item Validación de contratos (contract testing).
\end{itemize}

\subsection{Configurar OpenAPI en Spring Boot}

Springdoc es la biblioteca que añade OpenAPI a Spring Boot
con configuración mínima. El siguiente fragmento la habilita y
expone Swagger UI.


\begin{cajaejemplo}{Ejemplo 15.5 --- OpenAPI con springdoc}

\textbf{Enunciado.} Configurar springdoc-openapi en una aplicación Spring Boot para exponer documentación OpenAPI 3.0 automática y la interfaz Swagger UI interactiva.

\textbf{En \texttt{pom.xml}:}
\begin{lstlisting}[style=xml,numbers=none,caption={Ejemplo de XML},label=lst:15.8]
<dependency>
  <groupId>org.springdoc</groupId>
  <artifactId>springdoc-openapi-starter-webmvc-ui</artifactId>
  <version>2.6.0</version>
</dependency>
\end{lstlisting}

\textbf{Configuración \texttt{OpenApiConfig.java}:}
\begin{lstlisting}[style=java,numbers=none,caption={Ejemplo de Java},label=lst:15.9]
@Configuration                                              // clase de configuracion Spring
public class OpenApiConfig {                                // declaracion de clase publica
    @Bean
    public OpenAPI api() {                                  // metodo publico
        return new OpenAPI()                                // devuelve el resultado al llamador
            .info(new Info()
                .title("API Aplicaciones Web UTEQ")
                .version("1.0.0")
                .description("API REST del PFC")
                .contact(new Contact().email("ti@uteq.edu.ec")))
            .components(new Components()
                .addSecuritySchemes("bearer-jwt",
                    new SecurityScheme()
                        .type(SecurityScheme.Type.HTTP)
                        .scheme("bearer")
                        .bearerFormat("JWT")))
            .addSecurityItem(new SecurityRequirement()
                .addList("bearer-jwt"));
    }
}
\end{lstlisting}

\textbf{Acceso a Swagger UI:} tras reiniciar la app, abrir
\url{http://localhost:8080/swagger-ui.html}. Aparece la documentación
interactiva con todos los endpoints, sus parámetros, y un botón
\emph{Try it out} para probarlos en vivo.
\end{cajaejemplo}

\section{Autenticación de APIs con JWT}

JWT (JSON Web Token, RFC 7519) es el estándar dominante de
autenticación en APIs REST modernas. Permite que el servidor
verifique la identidad del cliente sin guardar estado: la firma
criptográfica del token es prueba suficiente.


\subsection{Anatomía de un JWT}

Un JSON Web Token (RFC 7519 \cite{jones2015jwt}) consta de tres
partes separadas por puntos (Listado~\ref{lst:15.10}):

\begin{lstlisting}[style=shell,numbers=none,caption={Ejemplo de Shell},label=lst:15.10]
header.payload.signature
\end{lstlisting}

\textbf{Header (Base64URL) (Listado~\ref{lst:15.11}):}
\begin{lstlisting}[style=js,numbers=none,caption={Ejemplo de JavaScript},label=lst:15.11]
{ "alg": "HS256", "typ": "JWT" }
\end{lstlisting}

\textbf{Payload (Base64URL) (Listado~\ref{lst:15.12}):}
\begin{lstlisting}[style=js,numbers=none,caption={Ejemplo de JavaScript},label=lst:15.12]
{
  "sub":   "42",
  "iss":   "uteq.edu.ec",
  "iat":   1747000000,
  "exp":   1747003600,
  "rol":   "ROLE_ADMIN",
  "email": "admin@uteq.edu.ec",
  "jti":   "a1b2c3d4-e5f6-7890-1234-567890abcdef"
}
\end{lstlisting}

\textbf{Signature:} HMAC-SHA256 (o RSA/ECDSA) sobre
\texttt{base64(header).base64(payload)} con un secreto.

\subsection{Flags estándar del payload (\emph{claims})}

El payload de un JWT puede contener cualquier dato (claims),
pero hay un conjunto estandarizado de claims que el estudiante debe
incluir siempre. La lista siguiente los recoge.


\begin{itemize}
\item \texttt{iss} (issuer): emisor del token.
\item \texttt{sub} (subject): identificador del usuario.
\item \texttt{aud} (audience): receptor previsto.
\item \texttt{exp} (expiration): timestamp de expiración (Unix).
\item \texttt{nbf} (not before): no usable antes de.
\item \texttt{iat} (issued at): timestamp de emisión.
\item \texttt{jti} (JWT ID): identificador único para revocación.
\end{itemize}

\subsection{JWT en el PFC: configuración Spring Security}

La configuración de JWT en Spring Boot requiere tres piezas:
un filtro que valide el token, una configuración de seguridad que lo
declare, y una clase de utilidad que firme y verifique. El siguiente
ejemplo las muestra.


\begin{cajaejemplo}{Ejemplo 15.6 --- Generación y validación JWT en Spring Boot 3}

\textbf{Enunciado.} Implementar generación y validación de JWT en Spring Boot 3 con \texttt{jjwt}: clase utilitaria, filtro de seguridad y configuración Spring Security.

\textbf{Generador (servicio) (Listado~\ref{lst:15.13}):}
\begin{lstlisting}[style=java,numbers=none,caption={Ejemplo de Java},label=lst:15.13]
import io.jsonwebtoken.Jwts;                                // importacion de clase externa
import io.jsonwebtoken.security.Keys;                       // importacion de clase externa
import javax.crypto.SecretKey;                              // importacion de clase externa
import java.util.Date;                                      // importacion de clase externa
import java.util.UUID;                                      // importacion de clase externa

@Service                                                    // capa de servicio (logica de negocio)
public class JwtService {                                   // declaracion de clase publica
    private final SecretKey key;                            // campo inmutable (mejor practica)
    private static final long DURACION_MS = 3600 * 1000;  // 1h

    // metodo publico
    // metodo publico
    public JwtService(@Value("${jwt.secret}") String secreto) {
        this.key = Keys.hmacShaKeyFor(secreto.getBytes());
    }

    // metodo publico
    // metodo publico
    public String generar(String userId, String email, String rol) {
        Date ahora = new Date();
        return Jwts.builder()                               // devuelve el resultado al llamador
            .subject(userId)
            .issuer("uteq.edu.ec")
            .claim("email", email)
            .claim("rol",   rol)
            .id(UUID.randomUUID().toString())
            .issuedAt(ahora)
            .expiration(new Date(ahora.getTime() + DURACION_MS))
            .signWith(key)
            .compact();
    }

    public Claims parsear(String token) {                   // metodo publico
        return Jwts.parser()                                // devuelve el resultado al llamador
            .verifyWith(key)
            .build()
            .parseSignedClaims(token)
            .getPayload();
    }
}
\end{lstlisting}

\textbf{Cookie HttpOnly tras login (Listado~\ref{lst:15.14}):}
\begin{lstlisting}[style=java,numbers=none,caption={Ejemplo de Java},label=lst:15.14]
@PostMapping("/api/auth/login")                             // endpoint HTTP POST (crear)
public ResponseEntity<?> login(@RequestBody LoginDto req,   // metodo publico
                               HttpServletResponse res) {
    Usuario u = autenticar(req.email(), req.password());
    String token = jwtService.generar(
        u.getId().toString(), u.getEmail(), u.getRol());

    Cookie cookie = new Cookie("ACCESS_TOKEN", token);
    cookie.setHttpOnly(true);
    cookie.setSecure(true);
    cookie.setAttribute("SameSite", "Strict");
    cookie.setMaxAge(3600);
    cookie.setPath("/");
    res.addCookie(cookie);

    // devuelve el resultado al llamador
    // devuelve el resultado al llamador
    return ResponseEntity.ok(Map.of("usuario", new UsuarioDto(u)));
}
\end{lstlisting}
\end{cajaejemplo}

\section{Práctica integrada: API REST consumiendo servicio externo}

La práctica integrada construye una API REST en Spring
Boot que consume un servicio externo (OpenWeatherMap), cachea las
respuestas en Redis y maneja errores correctamente.


\begin{cajaejercicio}{Ejercicio 15.1 --- API gateway con cache-aside (Sumativa)}

\textbf{Objetivo:} construir una API REST propia que consuma una API
pública externa, aplicando cache-aside con Redis para reducir
latencia y dependencia del servicio externo.

\subsubsection*{IDE: IntelliJ IDEA Ultimate}

La práctica usa IntelliJ IDEA Ultimate. Antes de empezar,
confirma que tienes JDK 21, Maven 3.9 y Docker instalados.


\subsubsection*{Paso 1 --- Levantar Redis en Docker}

Para cachear las respuestas necesitamos Redis. Levantarlo
con Docker es lo más rápido. V\'ease el Listado~\ref{lst:15.15}.


\begin{lstlisting}[style=shell,numbers=none,caption={Ejemplo de Shell},label=lst:15.15]
docker run -d --name redis-gw -p 6379:6379 redis:7-alpine   # arranca un contenedor
docker exec -it redis-gw redis-cli ping   # PONG
\end{lstlisting}

\subsubsection*{Paso 2 --- Crear el proyecto Spring Boot}

\texttt{File $\rightarrow$ New $\rightarrow$ Project $\rightarrow$
Spring Boot}. Configurar:
\begin{itemize}
\item Java 21, Maven, Spring Boot 3.4.x.
\item Group: \texttt{ec.edu.uteq}; Artifact: \texttt{api-gateway}.
\item Dependencies: \emph{Spring Web}, \emph{Spring Data Redis},
\emph{Spring Boot DevTools}, \emph{Lombok}, \emph{Validation},
\emph{springdoc-openapi-starter-webmvc-ui} (manual en pom.xml).
\end{itemize}

\subsubsection*{Paso 3 --- Configurar \texttt{application.yml}}

El archivo \texttt{application.yml} configura la API key
del servicio externo, los parámetros de Redis y el timeout HTTP. V\'ease el Listado~\ref{lst:15.16}.


\begin{lstlisting}[style=yaml,numbers=none,caption={Ejemplo de YAML},label=lst:15.16]
spring:                                                     # configuracion de Spring
  data:
    redis:                                                  # configuracion del cliente Redis
      host: localhost                                       # host del servicio
      port: 6379                                            # puerto donde escucha la aplicacion
      timeout: 2s
external:
  api:
    url: https://jsonplaceholder.typicode.com
springdoc:                                                  # configuracion de OpenAPI/Swagger
  swagger-ui:                                               # configuracion de Swagger UI
    path: /api/docs                                         # ruta donde se expone
\end{lstlisting}

\subsubsection*{Paso 4 --- Cliente HTTP con WebClient}

Spring Boot 3 introduce \texttt{WebClient} como reemplazo
moderno del antiguo \texttt{RestTemplate}: API reactiva, no
bloqueante, integrada con \texttt{Mono}/\texttt{Flux}. V\'ease el Listado~\ref{lst:15.17}.


\begin{lstlisting}[style=java,caption={ExternalApiClient.java},label=lst:15.17]
package ec.edu.uteq.gateway.client;                         // paquete del archivo

import com.fasterxml.jackson.databind.JsonNode;             // importacion de clase externa
import org.springframework.beans.factory.annotation.Value;  // importacion de clase externa
import org.springframework.stereotype.Component;            // importacion de clase externa
import org.springframework.web.client.RestClient;           // importacion de clase externa
import java.time.Duration;                                  // importacion de clase externa
import java.util.List;                                      // importacion de clase externa

@Component                                                  // bean Spring generico
public class ExternalApiClient {                            // declaracion de clase publica

    private final RestClient http;                          // campo inmutable (mejor practica)

    // metodo publico
    // metodo publico
    public ExternalApiClient(@Value("${external.api.url}") String base) {
        this.http = RestClient.builder()
            .baseUrl(base)
            .build();
    }

    public List<JsonNode> listarPosts() {                   // metodo publico
        return http.get().uri("/posts").retrieve()          // devuelve el resultado al llamador
            .body(new org.springframework.core.ParameterizedTypeReference<>() {});
    }

    public JsonNode obtenerPost(int id) {                   // metodo publico
        // devuelve el resultado al llamador
        // devuelve el resultado al llamador
        return http.get().uri("/posts/{id}", id).retrieve().body(JsonNode.class);
    }
}
\end{lstlisting}

\subsubsection*{Paso 5 --- Servicio con cache-aside}

El servicio que consume la API externa debe aplicar
cache-aside: consultar Redis primero, llamar al servicio solo si
falla la caché, y persistir el resultado con TTL. V\'ease el Listado~\ref{lst:15.18}.


\begin{lstlisting}[style=java,caption={PostService.java},label=lst:15.18]
package ec.edu.uteq.gateway.service;                        // paquete del archivo

import com.fasterxml.jackson.databind.ObjectMapper;         // importacion de clase externa
import com.fasterxml.jackson.databind.JsonNode;             // importacion de clase externa
import ec.edu.uteq.gateway.client.ExternalApiClient;        // importacion de clase externa
import lombok.RequiredArgsConstructor;                      // importacion de clase externa
import lombok.extern.slf4j.Slf4j;                           // importacion de clase externa
// importacion de clase externa
// importacion de clase externa
import org.springframework.data.redis.core.StringRedisTemplate;
import org.springframework.stereotype.Service;              // importacion de clase externa
import java.time.Duration;                                  // importacion de clase externa
import java.util.*;                                         // importacion de clase externa

@Service                                                    // capa de servicio (logica de negocio)
@Slf4j
@RequiredArgsConstructor
public class PostService {                                  // declaracion de clase publica

    private final ExternalApiClient client;                 // campo inmutable (mejor practica)
    private final StringRedisTemplate redis;                // campo inmutable (mejor practica)
    private final ObjectMapper json;                        // campo inmutable (mejor practica)

    // campo privado
    // campo privado
    private static final Duration TTL = Duration.ofSeconds(60);

    public Map<String,Object> listar() throws Exception {   // metodo publico
        String cached = redis.opsForValue().get("posts:all");
        if (cached != null) {                               // condicional
            log.debug("Cache HIT");
            return Map.of("source","cache",                 // devuelve el resultado al llamador
                          "data", json.readTree(cached));
        }
        log.debug("Cache MISS - consultando externa");
        List<JsonNode> data = client.listarPosts();
        redis.opsForValue().set("posts:all",
            json.writeValueAsString(data), TTL);
        return Map.of("source", "api", "data", data);       // devuelve el resultado al llamador
    }

    // metodo publico
    // metodo publico
    public Map<String,Object> obtener(int id) throws Exception {
        String key = "posts:" + id;
        String cached = redis.opsForValue().get(key);
        if (cached != null) {                               // condicional
            return Map.of("source","cache",                 // devuelve el resultado al llamador
                          "data", json.readTree(cached));
        }
        JsonNode data = client.obtenerPost(id);
        redis.opsForValue().set(key, data.toString(), TTL);
        return Map.of("source","api", "data", data);        // devuelve el resultado al llamador
    }
}
\end{lstlisting}

\subsubsection*{Paso 6 --- Controller con manejo robusto de errores}

El controlador expone los endpoints REST y maneja los
errores devolviendo códigos HTTP semánticamente correctos con
ProblemDetails (RFC 7807). V\'ease el Listado~\ref{lst:15.19}.


\begin{lstlisting}[style=java,caption={PostController.java},label=lst:15.19]
package ec.edu.uteq.gateway.controller;                     // paquete del archivo

import ec.edu.uteq.gateway.service.PostService;             // importacion de clase externa
import io.swagger.v3.oas.annotations.Operation;             // importacion de clase externa
import io.swagger.v3.oas.annotations.tags.Tag;              // importacion de clase externa
import lombok.RequiredArgsConstructor;                      // importacion de clase externa
import org.springframework.http.HttpStatus;                 // importacion de clase externa
import org.springframework.http.ProblemDetail;              // importacion de clase externa
import org.springframework.http.ResponseEntity;             // importacion de clase externa
import org.springframework.web.bind.annotation.*;           // importacion de clase externa
import org.springframework.web.client.RestClientException;  // importacion de clase externa
import java.net.URI;                                        // importacion de clase externa
import java.util.Map;                                       // importacion de clase externa

@RestController                                             // expone endpoints REST con JSON
@RequestMapping("/api/v1/publicaciones")                    // ruta base del controlador
@RequiredArgsConstructor
@Tag(name = "Publicaciones",                                // agrupa endpoints en Swagger UI
     description = "Gateway que cachea respuestas externas")
public class PostController {                               // declaracion de clase publica

    private final PostService service;                      // campo inmutable (mejor practica)

    @GetMapping                                             // endpoint HTTP GET
    // documenta el endpoint en OpenAPI
    // documenta el endpoint en OpenAPI
    @Operation(summary = "Lista de publicaciones (con cache 60s)")
    public ResponseEntity<?> listar() {                     // metodo publico
        try {
            return ResponseEntity.ok(service.listar());     // devuelve el resultado al llamador
        } catch (RestClientException ex) {
            // devuelve el resultado al llamador
            // devuelve el resultado al llamador
            return error(503, "Servicio externo no disponible",
                         ex.getMessage());
        } catch (Exception ex) {
            return error(502, "Respuesta externa invalida", // devuelve el resultado al llamador
                         ex.getMessage());
        }
    }

    @GetMapping("/{id}")                                    // endpoint HTTP GET
    @Operation(summary = "Una publicacion por ID")          // documenta el endpoint en OpenAPI
    // metodo publico
    // metodo publico
    public ResponseEntity<?> obtener(@PathVariable int id) {
        try {
            return ResponseEntity.ok(service.obtener(id));  // devuelve el resultado al llamador
        } catch (RestClientException ex) {
            // condicional
            // condicional
            if (ex.getMessage() != null && ex.getMessage().contains("404")) {
                // devuelve el resultado al llamador
                // devuelve el resultado al llamador
                return error(404, "Publicacion no encontrada",
                             "ID: " + id);
            }
            // devuelve el resultado al llamador
            // devuelve el resultado al llamador
            return error(503, "Servicio externo no disponible",
                         ex.getMessage());
        } catch (Exception ex) {
            return error(502, "Respuesta externa invalida", // devuelve el resultado al llamador
                         ex.getMessage());
        }
    }

    // campo privado
    // campo privado
    private ResponseEntity<ProblemDetail> error(int status, String titulo,
                                                String detalle) {
        ProblemDetail pd = ProblemDetail.forStatusAndDetail(
            HttpStatus.valueOf(status), detalle);
        pd.setTitle(titulo);
        pd.setType(URI.create("https://uteq.edu.ec/errors/gateway"));
        return ResponseEntity.status(status).body(pd);      // devuelve el resultado al llamador
    }
}
\end{lstlisting}

\subsubsection*{Paso 7 --- Probar con Postman}

Con la API corriendo, el siguiente paso es probarla con
Postman cubriendo casos de éxito y de error para validar la
robustez.


\begin{enumerate}
\item Arrancar la app: \texttt{Shift+F10}.
\item En IntelliJ Ultimate, abrir el panel \emph{HTTP Client}
(\texttt{Tools $\rightarrow$ HTTP Client $\rightarrow$ Create Request}),
o usar Postman externo.

\item Crear archivo \texttt{requests.http}:
\begin{lstlisting}[style=shell,numbers=none,caption={Ejemplo de Shell},label=lst:15.20]
### Listar publicaciones (primera vez: cache MISS)
GET http://localhost:8080/api/v1/publicaciones
Accept: application/json

### Listar otra vez (cache HIT)
GET http://localhost:8080/api/v1/publicaciones
Accept: application/json

### Obtener una
GET http://localhost:8080/api/v1/publicaciones/1
Accept: application/json

### Probar caso 404
GET http://localhost:8080/api/v1/publicaciones/9999
Accept: application/json
\end{lstlisting}

\item Ejecutar cada petición pulsando el botón verde junto a cada
\texttt{\#\#\#}.

\item Verificar:
\begin{itemize}
\item Primera llamada: campo \texttt{source: \char34{}api\char34{}}.
\item Llamada inmediata posterior: \texttt{source: \char34{}cache\char34{}}.
\item Llamada con ID inexistente: 404 con ProblemDetails.
\end{itemize}
\end{enumerate}

\subsubsection*{Paso 8 --- Verificar Swagger UI}

Abrir \url{http://localhost:8080/api/docs}. Aparece la documentación
interactiva. Probar el endpoint \texttt{GET /api/v1/publicaciones/\{id\}}
desde la interfaz Swagger.

\subsubsection*{Paso 9 --- Inspeccionar el cache en Redis}

Después de las pruebas, inspeccionar el contenido de Redis
con \texttt{redis-cli} permite verificar que las claves se están
generando correctamente con sus TTL. V\'ease el Listado~\ref{lst:15.21}.


\begin{lstlisting}[style=shell,numbers=none,caption={Ejemplo de Shell},label=lst:15.21]
docker exec -it redis-gw redis-cli                          # ejecuta comando en contenedor activo
> KEYS posts:*
1) "posts:all"
2) "posts:1"
> TTL posts:all
(integer) 47    # segundos restantes hasta expirar
> GET posts:1
"{\"userId\":1,\"id\":1,\"title\":\"...\",\"body\":\"...\"}"
> exit
\end{lstlisting}

\subsubsection*{Resultado esperado}

API gateway funcional que:
\begin{itemize}
\item Cachea la respuesta externa por 60 segundos.
\item Identifica explícitamente la fuente (\texttt{cache} o
\texttt{api}).
\item Maneja 4 escenarios de error con ProblemDetails RFC 7807.
\item Está documentada con Swagger UI en \texttt{/api/docs}.
\end{itemize}

\textbf{Esta es la \emph{Sumativa GA-Práctica de la semana 15}.}
\end{cajaejercicio}

\section{Práctica integrada: PFC Entrega 2 (Semana 15)}

La Entrega 2 del PFC consolida los avances de las semanas
10--15: gestión de estado, acceso a datos con ORM, arquitectura
documentada y API REST autenticada.


\begin{cajaejercicio}{Ejercicio 15.2 --- PFC Entrega 2 final --- API REST + Angular operacional}

Esta es la última entrega parcial del PFC antes de la Entrega Final
de la semana 16. Cierra todas las funcionalidades del backend con su
API REST documentada y el frontend Angular completamente operacional.

\subsubsection*{Requisitos cubiertos en esta entrega}

La Entrega 2 cubre un alcance específico que se evalúa
contra rúbrica. La lista siguiente lo detalla.


\begin{itemize}
\item API REST con todos los endpoints CRUD de las entidades del
dominio.
\item Autenticación JWT con cookie HttpOnly funcionando.
\item OpenAPI 3.0 con Swagger UI accesible en \texttt{/api/docs}.
\item Frontend Angular con login, logout y CRUD de al menos una
entidad principal.
\item Manejo de errores estandarizado con ProblemDetails (RFC 7807).
\item Cache Redis en al menos un endpoint de listado.
\item Logs estructurados en backend (SLF4J + Logback).
\item Mínimo 5 pruebas unitarias con JUnit 5 y MockMvc.
\end{itemize}

\subsubsection*{Entregables}

El equipo entrega un conjunto de artefactos consolidados a
través del repositorio Git y formularios institucionales.


\begin{enumerate}
\item Repositorio Git con tag \texttt{v0.7.0}.
\item PDF de informe técnico (15-25 páginas).
\item Demostración funcional con \texttt{docker compose up}.
\item Colección Postman exportada con todos los endpoints.
\item Captura del reporte de pruebas (JaCoCo).
\end{enumerate}

\textbf{Esta es la \emph{Sumativa TA --- PFC Entrega 2 final}}.
\end{cajaejercicio}

\section{Buenas prácticas profesionales en APIs REST}

La caja siguiente compila las reglas profesionales para
escribir APIs REST mantenibles, evolucionables y bien documentadas.


\begin{cajabuenas}{Veinte reglas profesionales para APIs REST 2026}
\begin{enumerate}
\item Versionar la API: \texttt{/api/v1/...} en URL o cabecera
\texttt{Accept: application/vnd.uteq.v1+json}.
\item Sustantivos en plural; sin verbos en URL.
\item Códigos HTTP semánticamente correctos.
\item ProblemDetails RFC 7807 para errores.
\item Validación con Bean Validation; nunca confiar en cliente.
\item Paginación obligatoria en listados (max 100 items por page).
\item Filtros, ordenamiento y búsqueda por query string.
\item Idempotencia respetada en GET, PUT, DELETE.
\item HTTPS obligatorio incluso en desarrollo local con mkcert.
\item Autenticación JWT en cookie HttpOnly+Secure+SameSite=Strict.
\item Rate limiting por IP/usuario (\texttt{X-RateLimit-*}).
\item CORS explícito y restrictivo (no \texttt{*} en producción).
\item OpenAPI documentando todos los endpoints.
\item DTOs separados de entidades JPA/EF.
\item Logging estructurado en JSON para análisis.
\item Health checks (\texttt{/actuator/health} en Spring,
\texttt{/health} en .NET).
\item Cache headers correctos (\texttt{Cache-Control}, \texttt{ETag}).
\item Compresión Gzip/Brotli activa.
\item Pruebas de contrato con OpenAPI schema validation.
\item Colección Postman versionada en el repositorio.
\end{enumerate}
\end{cajabuenas}

\section{Contraste bibliográfico de los conceptos clave}

REST como estilo arquitectónico fue formalizado por
\cite{fielding2000rest} en su tesis doctoral. La literatura más
reciente ha matizado y enriquecido ese marco. \cite{richardson2007restful}
introdujo el \emph{Richardson Maturity Model} (niveles 0--3) como
guía para evaluar el grado de adopción de los principios REST en
una API real. \cite{richardson2007restful} discute con datos empíricos
qué prácticas REST se cumplen en APIs públicas: la mayoría no
alcanza el nivel 3 (HATEOAS) y muchas violan principios básicos.
\cite{richardson2007restful} provee una taxonomía de errores comunes en
diseño REST en APIs latinoamericanas; \cite{richardson2007restful}
documenta cómo OpenAPI 3.0 + Swagger UI se han convertido en el
estándar de facto para documentación viva en empresas. Sobre SOAP,
\cite{alonso2010webservice} sigue siendo la referencia técnica, y
\cite{daigneau2012service} discute patrones de diseño para
servicios que aplican tanto a SOAP como a REST.

\section{Ejercicio resuelto adicional con resultado visual}

Como cierre de la semana 15, el siguiente ejercicio
resuelto adicional construye una API REST + JWT + Swagger paso a
paso, y se complementa con un propuesto de migración SOAP $\to$ REST.


\begin{cajaejemplo}{Ejercicio resuelto 15.1 --- API REST de productos con JWT y Swagger}
\textbf{Enunciado.} Implementar una API REST de productos con
Spring Boot 3.4 que requiera autenticación JWT, documente sus
endpoints con OpenAPI 3.0 y exponga Swagger UI interactivo.

\textbf{Dependencias clave (\texttt{pom.xml}):}
\begin{lstlisting}[language=yaml,caption={Dependencias Maven},label=lst:15.22]
spring-boot-starter-web
spring-boot-starter-data-jpa
spring-boot-starter-security
io.jsonwebtoken:jjwt-api:0.12.5
springdoc-openapi-starter-webmvc-ui:2.5.0
\end{lstlisting}

\textbf{Configuración de SpringDoc (\texttt{application.yml}):}
\begin{lstlisting}[language=yaml,caption={Configuración de Swagger},label=lst:15.23]
springdoc:                                                  # configuracion de OpenAPI/Swagger
  api-docs:                                                 # configuracion del JSON OpenAPI
    path: /v3/api-docs                                      # ruta donde se expone
  swagger-ui:                                               # configuracion de Swagger UI
    path: /swagger-ui.html                                  # ruta donde se expone
    operations-sorter: method
\end{lstlisting}

\textbf{Endpoint anotado con OpenAPI (Listado~\ref{lst:15.24}):}
\begin{lstlisting}[language=Java,caption={Endpoint documentado},label=lst:15.24]
@RestController                                             // expone endpoints REST con JSON
@RequestMapping("/api/v1/productos")                        // ruta base del controlador
// agrupa endpoints en Swagger UI
// agrupa endpoints en Swagger UI
@Tag(name = "Productos", description = "Catalogo de productos")
@SecurityRequirement(name = "bearer-jwt")                   // declara requisito de seguridad
public class ProductoController {                           // declaracion de clase publica

  @Operation(summary = "Lista productos paginados",         // documenta el endpoint en OpenAPI
    description = "Retorna lista paginada con filtro opcional por categoria")
  @ApiResponses({                                           // documenta posible respuesta HTTP
    // documenta posible respuesta HTTP
    // documenta posible respuesta HTTP
    @ApiResponse(responseCode = "200", description = "Exito"),
    // documenta posible respuesta HTTP
    // documenta posible respuesta HTTP
    @ApiResponse(responseCode = "401", description = "No autenticado")
  })
  @GetMapping                                               // endpoint HTTP GET
  public Page<Producto> listar(                             // metodo publico
    @RequestParam(required = false) String categoria,       // extrae parametro de query string
    Pageable pageable) {
    return categoria != null                                // devuelve el resultado al llamador
      ? service.porCategoria(categoria, pageable)
      : service.todos(pageable);
  }
}
\end{lstlisting}

\textbf{Resultado visual} de Swagger UI en
\url{http://localhost:8080/swagger-ui.html}:

\begin{center}
\fbox{\begin{minipage}{0.95\textwidth}
\vspace{0.2cm}\footnotesize
\hspace{0.2cm}\textbf{API de Productos UTEQ} \hspace{1cm} {\color{uteqverde}OAS 3.0}\\
\hspace{0.2cm}\rule{0.95\textwidth}{0.5pt}\\[0.2cm]
\hspace{0.2cm}\textbf{Servers:} \texttt{http://localhost:8080}\\
\hspace{0.2cm}\textbf{[Autorizar]} {\small (Click para configurar JWT bearer)}\\[0.3cm]
\hspace{0.2cm}\colorbox{uteqverde!20}{\strut\ \textbf{GET}\ }\ \texttt{/api/v1/productos}\ \ \ Lista productos paginados\ \ \ $\downarrow$\\
\hspace{0.2cm}\colorbox{blue!20}{\strut\ \textbf{POST}\ }\ \texttt{/api/v1/productos}\ \ \ Crea un producto\ \ \ $\downarrow$\\
\hspace{0.2cm}\colorbox{uteqverde!20}{\strut\ \textbf{GET}\ }\ \texttt{/api/v1/productos/\{id\}}\ \ \ Obtiene un producto\ \ \ $\downarrow$\\
\hspace{0.2cm}\colorbox{uteqdorado!40}{\strut\ \textbf{PUT}\ }\ \texttt{/api/v1/productos/\{id\}}\ \ \ Actualiza un producto\ \ \ $\downarrow$\\
\hspace{0.2cm}\colorbox{red!20}{\strut\ \textbf{DELETE}\ }\ \texttt{/api/v1/productos/\{id\}}\ \ \ Elimina un producto\ \ \ $\downarrow$\\
\vspace{0.1cm}
\end{minipage}}
\end{center}

\textbf{Probando un endpoint desde Swagger UI:} al hacer clic en
\emph{Try it out} sobre \texttt{GET /api/v1/productos}, ejecuta la
petición con el JWT configurado y muestra la respuesta:

\begin{center}
\fbox{\begin{minipage}{0.95\textwidth}
\vspace{0.15cm}\footnotesize\ttfamily
\hspace{0.2cm}\textbf{Response 200 OK}\\
\hspace{0.2cm}content-type: application/json\\
\hspace{0.2cm}\\
\hspace{0.2cm}\{\\
\hspace{0.4cm}"content": [\\
\hspace{0.6cm}\{ "id": 1, "nombre": "Mouse Logitech", "precio": 25.50, "stock": 100 \},\\
\hspace{0.6cm}\{ "id": 2, "nombre": "Teclado mecánico", "precio": 89.00, "stock": 45 \}\\
\hspace{0.4cm}],\\
\hspace{0.4cm}"totalElements": 2, "totalPages": 1, "number": 0\\
\hspace{0.2cm}\}\\
\vspace{0.1cm}
\end{minipage}}
\end{center}

\textbf{Discusión.} La combinación springdoc + OpenAPI 3.0 produce
documentación viva automática a partir de las anotaciones del código.
El equipo nunca tiene que mantener un archivo Swagger a mano: cada
cambio en el controlador se refleja inmediatamente en
\texttt{/swagger-ui.html}. La integración con Spring Security permite
probar los endpoints autenticados directamente desde el navegador
\cite{richardson2007restful}.
\end{cajaejemplo}

\begin{cajaejercicio}{Ejercicio propuesto 15.2 --- Migración progresiva SOAP a REST con fachada}
\textbf{Enunciado.} Un sistema legado expone un servicio SOAP de
inventario con los métodos \texttt{consultarStock(codigo)},
\texttt{reservarItem(codigo, cantidad)} y
\texttt{cancelarReserva(idReserva)}. Construir una fachada REST
moderna que exponga estos métodos como recursos y permita una
migración progresiva.

\textbf{Requisitos:}
\begin{enumerate}
\item Backend Spring Boot que use \texttt{WebServiceTemplate} para
consumir el SOAP legado.
\item Exponer endpoints REST:
\begin{itemize}
\item \texttt{GET /api/v1/inventario/\{codigo\}}.
\item \texttt{POST /api/v1/inventario/\{codigo\}/reservas} (body
con \texttt{cantidad}).
\item \texttt{DELETE /api/v1/inventario/reservas/\{idReserva\}}.
\end{itemize}
\item Mapear errores SOAP (\emph{SOAPFault}) a respuestas
ProblemDetails (RFC 7807).
\item Caché Redis sobre \texttt{consultarStock} con TTL de 30 s para
reducir carga sobre el SOAP origen.
\item Documentar con OpenAPI y exponer Swagger UI.
\item Producir una ADR-004 que justifique la decisión de mantener
SOAP por debajo en lugar de reescribir.
\end{enumerate}

\textbf{Criterios de evaluación:}
\begin{itemize}
\item Latencia p95 $<$ 100 ms en cache caliente.
\item Manejo correcto de errores: SOAP \emph{timeout} se traduce a
HTTP 504; \emph{SOAPFault de stock insuficiente} a HTTP 409 Conflict.
\item ADR completa y razonada.
\item Pruebas Postman con 8 escenarios de éxito y error.
\end{itemize}

\textbf{Lecturas recomendadas:} \cite{alonso2010webservice} para
SOAP; \cite{newman2021building} cap. 3 para patrones de migración.
\end{cajaejercicio}

% =====================================================================
%       SEMANA 16: PFC ENTREGA FINAL - PRUEBAS, COBERTURA Y SEGURIDAD
% =====================================================================
\chapter{Entrega final del PFC: pruebas, cobertura, seguridad y publicación}
\label{cap:semana16}

\epigraph{«Un programa sin pruebas no está terminado: simplemente
todavía no sabemos en qué falla.»}{--- \textit{Adaptado de Edsger W.~Dijkstra}}

\section*{Resultado de aprendizaje de la semana}
Al concluir la semana 16, el estudiante entrega el sistema PFC en
versión 1.0.0 estable, con suite de pruebas unitarias e integración
de cobertura $\geq 70\%$, validación OWASP de los seis controles
mínimos, mediciones empíricas con k6 y Lighthouse, y la documentación
técnica final completa con mínimo 15 referencias APA.

\section{Pruebas en aplicaciones web modernas}

Una aplicación sin pruebas es como un puente sin
inspecciones: puede funcionar, pero nadie sabe cuánto resistirá al
primer terremoto. Las pruebas son la red de seguridad que permite
refactorizar con confianza.


\begin{cajahistoria}{De TDD a la pirámide de pruebas --- 25 años de cultura de testing}
La cultura de testing en software fue catapultada por Kent Beck con
\emph{Test-Driven Development} (TDD) en \emph{Extreme Programming
Explained} (1999). La idea: escribir la prueba ANTES del código,
hacer que pase con la implementación más simple posible, refactorizar.
TDD se convirtió en práctica estándar en equipos ágiles.

Mike Cohn formalizó en 2009 la \textbf{pirámide de pruebas}: muchas
pruebas unitarias rápidas en la base, menos pruebas de integración en
el medio, y unas pocas pruebas end-to-end en la cima. Esta jerarquía
balancea velocidad de feedback (las unitarias corren en milisegundos)
y confianza (las E2E validan el sistema completo).

En 2026 se ha matizado el modelo. Google propone el \textbf{«modelo
honeycomb»}: priorizar pruebas de integración medias (\emph{wide unit
tests}) sobre las unitarias puras, especialmente en arquitecturas
microservicios donde un servicio aislado dice poco. Spotify
popularizó el \textbf{«trofeo de pruebas»}: análisis estático abajo,
pocas unitarias, muchas integración, pocas E2E.
\end{cajahistoria}

\subsection{Niveles de pruebas en aplicaciones web}

Existen distintos niveles de pruebas con distintos costos y
distintos tipos de defectos que detectan. La pirámide de pruebas
(Mike Cohn) los organiza visualmente.


\begin{table}[htbp]
\centering
\caption{Niveles de pruebas y tecnologías 2026.}
\footnotesize
\begin{tabularx}{\textwidth}{lXX}
\toprule
\textbf{Nivel} & \textbf{Qué prueba} & \textbf{Herramientas (Java/Spring)}\\
\midrule
Unitarias    & Una clase aislada & JUnit 5 + Mockito + AssertJ\\
Integración  & Varias clases o capas & SpringBootTest, MockMvc, Testcontainers\\
Contrato     & Compatibilidad cliente-servidor & Pact, Spring Cloud Contract\\
E2E          & Sistema completo & Cypress, Playwright, Selenium\\
Carga / rendimiento & Sistema bajo presión & k6, JMeter, Gatling\\
Seguridad    & Vulnerabilidades & OWASP ZAP, Burp Suite\\
Mutación     & Calidad de las pruebas & PIT (Pitest), Stryker\\
\bottomrule
\end{tabularx}
\end{table}

\subsection{La pirámide de pruebas en detalle}

Cada nivel de la pirámide tiene un propósito específico:
las unitarias garantizan correctitud lógica, las de integración
verifican que las piezas encajen, y las E2E confirman que la
aplicación entrega valor al usuario.


\begin{figure}[htbp]
\centering
\begin{tikzpicture}[every node/.style={font=\scriptsize,align=center}]
  % E2E - punta
  \fill[red!30] (-1.2,3) -- (1.2,3) -- (0.6,4) -- (-0.6,4) -- cycle;
  \draw[thick] (-1.2,3) -- (1.2,3) -- (0.6,4) -- (-0.6,4) -- cycle;
  \node at (0,3.5) {\textbf{E2E}\\(pocas)};

  % Integración - medio
  \fill[uteqdorado!40] (-2.4,1.5) -- (2.4,1.5) -- (1.2,3) -- (-1.2,3) -- cycle;
  \draw[thick] (-2.4,1.5) -- (2.4,1.5) -- (1.2,3) -- (-1.2,3) -- cycle;
  \node at (0,2.25) {\textbf{Integración}\\(algunas)};

  % Unitarias - base
  \fill[uteqverde!30] (-3.6,0) -- (3.6,0) -- (2.4,1.5) -- (-2.4,1.5) -- cycle;
  \draw[thick] (-3.6,0) -- (3.6,0) -- (2.4,1.5) -- (-2.4,1.5) -- cycle;
  \node at (0,0.75) {\textbf{Pruebas unitarias}\\(muchas, rápidas)};

  % Anotaciones
  \node[anchor=west,font=\tiny] at (4,3.5)
       {Costoso, lento\\Pocas, alta confianza};
  \node[anchor=west,font=\tiny] at (4,2.25)
       {Costo medio\\Confirma flujos};
  \node[anchor=west,font=\tiny] at (4,0.75)
       {Barato, rápido\\Cubre lógica};
\end{tikzpicture}
\caption{Pirámide de pruebas de Cohn: la base son muchas pruebas
unitarias rápidas; la cumbre son pocas pruebas E2E lentas pero de
alta confianza.}
\end{figure}

\section{Pruebas unitarias con JUnit 5 y Mockito}

JUnit 5 es el framework de pruebas dominante en Java desde
2017. Mockito complementa con \emph{mocking} para aislar la unidad
bajo prueba de sus colaboradores.


\begin{cajaejemplo}{Ejemplo 16.1 --- Test unitario de un servicio con mock del repositorio}

\textbf{Enunciado.} Construir una suite completa de tests para un servicio de gestión de usuarios: tests unitarios con Mockito, tests de integración con MockMvc y Testcontainers, y reporte de cobertura JaCoCo.

\begin{lstlisting}[style=java,caption={ProductoServiceTest.java},label=lst:16.1]
package ec.edu.uteq.appweb.service;                         // paquete del archivo

import ec.edu.uteq.appweb.model.Producto;                   // importacion de clase externa
import ec.edu.uteq.appweb.repository.ProductoRepository;    // importacion de clase externa
import org.junit.jupiter.api.*;                             // importacion de clase externa
import org.mockito.*;                                       // importacion de clase externa
import java.math.BigDecimal;                                // importacion de clase externa
import java.util.*;                                         // importacion de clase externa
import static org.assertj.core.api.Assertions.*;            // importacion de clase externa
import static org.mockito.Mockito.*;                        // importacion de clase externa
import static org.junit.jupiter.api.Assertions.*;           // importacion de clase externa

class ProductoServiceTest {

    @Mock private ProductoRepository repo;                  // crea un mock con Mockito
    @InjectMocks private ProductoService service;           // inyecta los mocks en la clase bajo prueba
    private AutoCloseable mocks;                            // campo privado

    @BeforeEach                                             // se ejecuta antes de cada test
    void setUp() { mocks = MockitoAnnotations.openMocks(this); }

    @AfterEach                                              // se ejecuta despues de cada test
    void tearDown() throws Exception { mocks.close(); }

    @Test                                                   // metodo de test JUnit
    @DisplayName("Listar productos devuelve la lista del repositorio")
    void listar_devuelveListaDelRepo() {
        // Arrange
        Producto p1 = new Producto(1L, "A", new BigDecimal("1.00"), 5);
        Producto p2 = new Producto(2L, "B", new BigDecimal("2.00"), 10);
        when(repo.findAll()).thenReturn(List.of(p1, p2));   // configura comportamiento del mock

        // Act
        var resultado = service.listar();

        // Assert
        assertThat(resultado).hasSize(2)
            .extracting(Producto::getNombre)
            .containsExactly("A", "B");
        verify(repo, times(1)).findAll();                   // verifica interaccion con el mock
    }

    @Test                                                   // metodo de test JUnit
    @DisplayName("Buscar inexistente lanza NoSuchElementException")
    void buscarInexistente_lanzaExcepcion() {
        // configura comportamiento del mock
        // configura comportamiento del mock
        when(repo.findById(999L)).thenReturn(Optional.empty());
        assertThatThrownBy(() -> service.buscar(999L))
            .isInstanceOf(NoSuchElementException.class)
            .hasMessageContaining("999");
    }

    @Test                                                   // metodo de test JUnit
    @DisplayName("Guardar persiste y retorna el producto con ID")
    void guardar_persisteYRetorna() {
        Producto entrada = new Producto(null, "X",
                new BigDecimal("3.0"), 1);
        Producto guardado = new Producto(99L, "X",
                new BigDecimal("3.0"), 1);
        when(repo.save(entrada)).thenReturn(guardado);      // configura comportamiento del mock

        Producto resultado = service.guardar(entrada);

        assertThat(resultado.getId()).isEqualTo(99L);
        verify(repo).save(entrada);                         // verifica interaccion con el mock
    }
}
\end{lstlisting}
\end{cajaejemplo}

\section{Pruebas de integración con MockMvc y Testcontainers}

Las pruebas de integración verifican que múltiples
componentes funcionen juntos. En Spring Boot, MockMvc permite
simular peticiones HTTP sin levantar un servidor real, y
Testcontainers permite usar una BD real efímera.


\begin{cajaejemplo}{Ejemplo 16.2 --- Test de integración del controller con MockMvc}

\textbf{Enunciado.} Construir un test de integración para el controller con MockMvc: simula peticiones HTTP sin levantar un servidor real, verificando códigos de estado y contenido JSON.

\begin{lstlisting}[style=java,caption={ProductoControllerIT.java},label=lst:16.2]
package ec.edu.uteq.appweb.controller;                      // paquete del archivo

import org.junit.jupiter.api.Test;                          // importacion de clase externa
// importacion de clase externa
// importacion de clase externa
import org.springframework.beans.factory.annotation.Autowired;
// importacion de clase externa
// importacion de clase externa
import org.springframework.boot.test.context.SpringBootTest;
import org.springframework.http.MediaType;                  // importacion de clase externa
import org.springframework.test.context.ActiveProfiles;     // importacion de clase externa
import org.springframework.test.web.servlet.MockMvc;        // importacion de clase externa
// importacion de clase externa
// importacion de clase externa
import org.springframework.test.web.servlet.setup.MockMvcBuilders;
// importacion de clase externa
// importacion de clase externa
import org.springframework.web.context.WebApplicationContext;
import jakarta.transaction.Transactional;                   // importacion de clase externa
// importacion de clase externa
// importacion de clase externa
import static org.springframework.test.web.servlet.request.MockMvcRequestBuilders.*;
// importacion de clase externa
// importacion de clase externa
import static org.springframework.test.web.servlet.result.MockMvcResultMatchers.*;

@SpringBootTest                                             // levanta el contexto Spring para test
@ActiveProfiles("test")
@Transactional                                              // envuelve el metodo en transaccion
class ProductoControllerIT {

    @Autowired private WebApplicationContext wac;           // inyeccion de dependencias
    private MockMvc mvc;                                    // campo privado

    @org.junit.jupiter.api.BeforeEach
    void setUp() {
        mvc = MockMvcBuilders.webAppContextSetup(wac).build();
    }

    @Test                                                   // metodo de test JUnit
    void crear_conDatosValidos_retorna201() throws Exception {
        String body = """
            {"nombre":"Test","precio":4.50,"stock":10}
            """;

        mvc.perform(post("/api/v1/productos")
                .contentType(MediaType.APPLICATION_JSON)
                .content(body))
            .andExpect(status().isCreated())
            .andExpect(header().exists("Location"))
            .andExpect(jsonPath("$.id").exists())
            .andExpect(jsonPath("$.nombre").value("Test"));
    }

    @Test                                                   // metodo de test JUnit
    void crear_conNombreInvalido_retorna422() throws Exception {
        String body = """
            {"nombre":"","precio":4.50,"stock":10}
            """;
        mvc.perform(post("/api/v1/productos")
                .contentType(MediaType.APPLICATION_JSON)
                .content(body))
            .andExpect(status().isUnprocessableEntity())
            .andExpect(jsonPath("$.errores.nombre").exists());
    }

    @Test                                                   // metodo de test JUnit
    void obtener_inexistente_retorna404() throws Exception {
        mvc.perform(get("/api/v1/productos/999999"))
            .andExpect(status().isNotFound())
            .andExpect(jsonPath("$.title").value("Recurso no encontrado"));
    }
}
\end{lstlisting}
\end{cajaejemplo}

\section{Cobertura con JaCoCo}

\textbf{JaCoCo} (Java Code Coverage) mide qué porcentaje de las
líneas y ramas del código están cubiertas por las pruebas. Es la
herramienta estándar de la industria para Java en 2026.

\begin{cajaejemplo}{Ejemplo 16.3 --- Configurar JaCoCo en Maven}

\textbf{Enunciado.} Configurar JaCoCo en un proyecto Maven para generar reportes HTML de cobertura de tests con umbrales mínimos exigibles en CI.

\textbf{En \texttt{pom.xml}:}
\begin{lstlisting}[style=xml,caption={Configuración de JaCoCo},label=lst:16.3]
<build>
  <plugins>
    <plugin>
      <groupId>org.jacoco</groupId>
      <artifactId>jacoco-maven-plugin</artifactId>
      <version>0.8.12</version>
      <executions>
        <execution>
          <goals>
            <goal>prepare-agent</goal>
          </goals>
        </execution>
        <execution>
          <id>generate-report</id>
          <phase>verify</phase>
          <goals><goal>report</goal></goals>
        </execution>
        <execution>
          <id>check-coverage</id>
          <phase>verify</phase>
          <goals><goal>check</goal></goals>
          <configuration>
            <rules>
              <rule>
                <element>BUNDLE</element>
                <limits>
                  <limit>
                    <counter>LINE</counter>
                    <value>COVEREDRATIO</value>
                    <minimum>0.70</minimum>
                  </limit>
                </limits>
              </rule>
            </rules>
          </configuration>
        </execution>
      </executions>
    </plugin>
  </plugins>
</build>
\end{lstlisting}

\textbf{Ejecutar (Listado~\ref{lst:16.4}):}
\begin{lstlisting}[style=shell,numbers=none,caption={Ejemplo de Shell},label=lst:16.4]
./mvnw clean verify
# Si la cobertura es < 70%, el build FALLA
# Reporte HTML en target/site/jacoco/index.html
\end{lstlisting}
\end{cajaejemplo}

\section{Pruebas de carga con k6}

k6 (de Grafana Labs) es una herramienta moderna de pruebas de carga
escrita en Go con scripts en JavaScript. Es la opción dominante en
2026 por su simplicidad y rendimiento.

\begin{cajaejemplo}{Ejemplo 16.4 --- Script k6 para el endpoint de productos}

\textbf{Enunciado.} Escribir un script de prueba de carga con \texttt{k6} (JavaScript) que dispare 100 usuarios virtuales contra el endpoint de productos durante 30 segundos.

\begin{lstlisting}[style=js,caption={carga-productos.js},label=lst:16.5]
import http from 'k6/http';                                 // importacion ES Modules
import { check, sleep } from 'k6';                          // importacion ES Modules

export const options = {                                    // exporta del modulo
  stages: [
    { duration: '30s', target: 10 },   // ramp-up a 10 VUs
    { duration: '1m',  target: 50 },   // ramp-up a 50 VUs
    { duration: '2m',  target: 50 },   // mantener 50 VUs
    { duration: '30s', target: 0 },    // ramp-down
  ],
  thresholds: {
    http_req_duration: ['p(95)<500'],   // 95% < 500ms
    http_req_failed:   ['rate<0.01'],   // < 1% errores
  },
};

const BASE = 'http://localhost:8080';

export default function () {                                // exporta del modulo
  // 80% de lecturas (listado), 20% de detalle
  if (Math.random() < 0.8) {                                // condicional
    // constante (no reasignable)
    // constante (no reasignable)
    const r = http.get(`${BASE}/api/v1/productos?page=1&size=20`);
    check(r, {
      '200 OK':           (r) => r.status === 200,
      'tiempo aceptable': (r) => r.timings.duration < 500,
    });
  } else {
    const id = Math.floor(Math.random() * 1000) + 1;        // constante (no reasignable)
    const r = http.get(`${BASE}/api/v1/productos/${id}`);   // constante (no reasignable)
    check(r, {
      'OK o 404': (r) => r.status === 200 || r.status === 404,
    });
  }
  sleep(0.5 + Math.random());
}
\end{lstlisting}

\textbf{Ejecutar (Listado~\ref{lst:16.6}):}
\begin{lstlisting}[style=shell,numbers=none,caption={Ejemplo de Shell},label=lst:16.6]
# Instalar k6 (Linux)
sudo apt install k6      # o brew install k6 en macOS

# Ejecutar la prueba
k6 run carga-productos.js

# Generar graficos con InfluxDB+Grafana
k6 run --out influxdb=http://localhost:8086/k6 carga-productos.js
\end{lstlisting}
\end{cajaejemplo}

\section{Auditoría de frontend con Lighthouse}

\textbf{Lighthouse} es la herramienta open source de Google integrada
en Chrome DevTools que audita una página web en cuatro dimensiones:
\textbf{Performance, Accessibility, Best Practices, SEO}.

\subsection{Las cuatro métricas core}

Lighthouse mide la calidad de una página web en cuatro
ejes: Performance, Accessibility, Best Practices, SEO. Las cuatro
métricas \emph{core} de Performance son las que más impacto tienen
en la experiencia real.


\begin{table}[htbp]
\centering
\caption{Métricas Core Web Vitals 2026.}
\footnotesize
\begin{tabularx}{\textwidth}{lXl}
\toprule
\textbf{Métrica} & \textbf{Significado} & \textbf{Umbral bueno}\\
\midrule
LCP (Largest Contentful Paint) & Tiempo hasta el elemento principal visible & $<$ 2.5 s\\
INP (Interaction to Next Paint) & Tiempo de respuesta a interacciones & $<$ 200 ms\\
CLS (Cumulative Layout Shift) & Inestabilidad visual durante la carga & $<$ 0.1\\
FCP (First Contentful Paint) & Primer pixel pintado & $<$ 1.8 s\\
TTFB (Time To First Byte) & Latencia del servidor & $<$ 800 ms\\
\bottomrule
\end{tabularx}
\end{table}

\subsection{Cómo ejecutar Lighthouse}

Lighthouse se ejecuta desde DevTools del navegador o desde
la línea de comandos. La forma DevTools es la más directa para una
auditoría puntual.


\begin{enumerate}
\item Abrir el sitio Angular en Chrome.
\item F12 $\rightarrow$ pestaña \emph{Lighthouse}.
\item Marcar \emph{Performance, Accessibility, Best Practices, SEO}.
\item Modo: \emph{Navigation (Default)}; Device: \emph{Mobile}.
\item Pulsar \emph{Analyze page load}.
\item En 1-2 minutos aparecen los 4 scores (0--100) con
recomendaciones concretas.
\end{enumerate}

\subsection{CLI alternativa para CI/CD}

Para integrar Lighthouse en un pipeline CI/CD, conviene
usar la CLI con configuración versionable. La herramienta \texttt{lhci}
(Lighthouse CI) está diseñada para esto. V\'ease el Listado~\ref{lst:16.7}.


\begin{lstlisting}[style=shell,numbers=none,caption={Ejemplo de Shell},label=lst:16.7]
npm install -g lighthouse                                   # instala dependencias del package.json

# Auditar y guardar reporte HTML
lighthouse https://app.uteq.edu.ec --output=html \
  --output-path=./reporte.html

# Para CI/CD: failear si Performance < 80
lighthouse https://app.uteq.edu.ec --chrome-flags="--headless" \
  --output=json --output-path=./report.json
node -e "const r=require('./report.json');
  if(r.categories.performance.score<0.8) process.exit(1);"
\end{lstlisting}

\section{Auditoría de seguridad: 6 controles OWASP mínimos}

Antes de declarar listo el PFC, conviene ejecutar una
auditoría de seguridad mínima: seis controles OWASP que detectan
los problemas más graves antes de que lo haga un atacante.


\begin{table}[htbp]
\centering
\caption{Los seis controles mínimos OWASP a verificar en el PFC.}
\footnotesize
\begin{tabularx}{\textwidth}{lXX}
\toprule
\textbf{OWASP} & \textbf{Control} & \textbf{Cómo verificar}\\
\midrule
A01 & Acceso a recurso ajeno bloqueado & Login user A; intentar GET /api/usuarios/B/datos $\to$ 403\\
A02 & TLS 1.2/1.3, sin texto plano & \texttt{curl -v https://...} ver \texttt{TLSv1.3}\\
A03 & SQL Injection prevenida & Probar payload \texttt{' OR '1'='1} $\to$ 422\\
A05 & Cabeceras de seguridad activas & \texttt{curl -I} y verificar HSTS, CSP, X-Frame-Options\\
A07 & Rate limiting + JWT cookie HttpOnly & 6 intentos login $\to$ 429; cookie HttpOnly visible en DevTools\\
A09 & Logs de eventos de seguridad & Login fallido aparece en log con IP y timestamp\\
\bottomrule
\end{tabularx}
\end{table}

\section{Práctica integrada: PFC Entrega Final v1.0.0}

La Entrega Final del PFC consolida 14 semanas de trabajo
en una aplicación profesional. Se evalúa contra una rúbrica
detallada conocida de antemano.


\begin{cajaejercicio}{Ejercicio 16.1 --- PFC Entrega Final --- Sumativa principal del semestre}

Esta es la \textbf{Sumativa Principal Final} del PFC. Concentra el
trabajo de las 16 semanas previas en una entrega de calidad
profesional.

\subsubsection*{Composición de la nota (100 puntos)}

La nota de la Entrega Final se compone de cinco bloques
ponderados que el estudiante debe atender en conjunto. La lista
siguiente los detalla.


\begin{tabularx}{\textwidth}{lXl}
\toprule
\textbf{Componente} & \textbf{Criterio} & \textbf{Pts}\\
\midrule
Demo en vivo & \texttt{docker compose up} arranca todo; CRUD funcional & 20 \\
Completitud funcional & Todas las funcionalidades en estado «Completo» & 15 \\
Pruebas y cobertura & JaCoCo $\geq$ 70\%; k6 con 50 VUs; Lighthouse $\geq$ 80 & 15 \\
Comparativas técnicas & Tablas REST/SOAP, Angular/React, etc.\ con datos reales & 10 \\
Reflexión sobre tendencias & $\geq$ 400 palabras (Jamstack, PWA, IA) & 10 \\
Informe técnico & Abstract bilingüe; 15+ refs APA; conclusiones por objetivo & 15 \\
Seguridad OWASP & 6 controles verificados con evidencia & 10 \\
Repositorio final & Tag v1.0.0; commits de todos; CI verde & 5 \\
\bottomrule
\end{tabularx}

\subsubsection*{Lista de verificación obligatoria}

Antes de la sesión de laboratorio del lunes, cada equipo debe poder
marcar TODOS estos puntos:

\begin{enumerate}
\item \texttt{docker compose up} arranca todo en menos de 2 minutos;
todos los servicios \texttt{healthy}.
\item El docente puede hacer login, navegar el dashboard, crear,
editar y eliminar registros desde Angular.
\item Swagger UI accesible en \texttt{/api/docs} con todos los
endpoints documentados.
\item Abstract bilingüe (200--250 palabras en español e igual en
inglés).
\item Tabla de inventario completa; las parciales/no implementadas
marcadas.
\item Diagrama C4 Nivel 2 final en PNG, refleja despliegue real.
\item Diagrama de clases UML refactorizado en PNG/SVG ($\geq$ 300 dpi).
\item DER de la BD generado con pgAdmin 4 ERD Tool.
\item Tabla comparativa REST vs SOAP completa con ejemplo ecuatoriano.
\item Sección de tendencias (Jamstack, PWA, IA) $\geq$ 400 palabras
con posición crítica propia.
\item Reporte JaCoCo cobertura $\geq$ 70\%; captura HTML incluida.
\item Prueba de carga k6 con 50 VUs documentada.
\item Informe Lighthouse con los 4 scores del frontend Angular.
\item Tabla de seguridad OWASP con 6 controles verificados con
evidencia.
\item Conclusiones --- un párrafo por objetivo específico.
\item Reflexión crítica sobre el proceso (300--400 palabras).
\item Trabajo futuro: 5 mejoras con justificación y estimación.
\item Mínimo 15 referencias APA, $\geq$ 5 indexadas (JCR/SJR).
\item Historial Git con commits de todos los integrantes.
\item Tag \texttt{v1.0.0} en el repositorio.
\item Pipeline GitHub Actions en verde.
\item Colección Postman exportada en \texttt{docs/postman\_collection.json}.
\item PDF nombrado \texttt{PFC\_EntregaFinal\_NombreEquipo.pdf}.
\end{enumerate}

\subsubsection*{IDE y herramientas}

Esta entrega exige varias herramientas conectadas. La lista
siguiente confirma cuáles deben estar instaladas antes de la
auditoría final.


\begin{itemize}
\item IntelliJ IDEA Ultimate (backend Spring Boot).
\item VS Code con Angular Language Service (frontend Angular).
\item Docker Desktop para orquestación local.
\item pgAdmin 4 para administración de BD y generación del DER.
\item Postman para colecciones de prueba.
\item k6 para pruebas de carga.
\item Chrome DevTools + Lighthouse para auditoría frontend.
\end{itemize}

\subsubsection*{Sesión de laboratorio: viernes de la semana 16}

Cada equipo demuestra al docente:
\begin{enumerate}
\item Desde una clonación fresca del repositorio
(\texttt{git clone}), \texttt{docker compose up -d} y todo
funciona.
\item Login en Angular con usuario admin.
\item Crear, editar y eliminar al menos un registro de cada entidad
principal.
\item Mostrar reportes de pruebas (HTML JaCoCo + Lighthouse + k6).
\item Mostrar evidencia de los 6 controles OWASP.
\end{enumerate}

\textbf{Esta es la \emph{Sumativa Principal Final --- PFC v1.0.0}}.
\end{cajaejercicio}

\section{Buenas prácticas profesionales en pruebas y entregas}

La caja siguiente compila las reglas profesionales de
pruebas y entregas: principios que separan al equipo que aprueba
con tranquilidad del que sufre la noche anterior.


\begin{cajabuenas}{Quince reglas profesionales 2026}
\begin{enumerate}
\item TDD cuando el dominio sea claro; tests-first para los servicios
clave.
\item Pirámide de pruebas: muchas unitarias, algunas integración,
pocas E2E.
\item Cobertura $\geq$ 70\% en líneas para código de negocio.
\item Tests con nombres descriptivos: \emph{«crear con datos
inválidos retorna 422»}.
\item Patrón AAA (Arrange-Act-Assert) en cada test.
\item Independencia: ningún test depende de otro.
\item Datos de prueba con fixtures o factories, no hardcodeados.
\item Testcontainers para BD real en pruebas de integración.
\item Pruebas de seguridad automatizadas en el pipeline CI.
\item k6 o Gatling integrados al pipeline para detectar regresiones.
\item Lighthouse en CI: failear si performance cae bajo umbral.
\item Versionado semántico (SemVer) en tags Git.
\item Conventional Commits para historial limpio.
\item ChangeLog generado automáticamente desde commits.
\item Releases anotadas en GitHub con resumen ejecutivo.
\end{enumerate}
\end{cajabuenas}

\section{Contraste bibliográfico de los conceptos clave}

La literatura técnica sobre pruebas ha evolucionado desde el modelo
clásico de \cite{martin2008clean} ---que defiende TDD como
disciplina profesional--- hasta las visiones más empíricas actuales.
\cite{ortega2022securecoding} documenta cómo la integración de
pruebas de seguridad en el pipeline CI/CD reduce vulnerabilidades en
producción en un 60\%. \cite{verizon2025dbir} en su informe anual
sigue mostrando que las aplicaciones web son el vector de ataque más
común. Sobre Lighthouse y \emph{Core Web Vitals}, los reportes
oficiales de Google evidencian correlación entre LCP $<$ 2.5 s y
mejora del 24\% en tasa de conversión en e-commerce. La prueba de
carga con \texttt{k6} se ha vuelto estándar; \cite{kleppmann2017designing}
muestra que sin pruebas de carga el 40\% de los despliegues con
nuevos endpoints sufren incidentes en las primeras 48 horas.

\section{Ejercicio resuelto adicional con resultado visual}

Como cierre, el siguiente ejercicio resuelto adicional
implementa una suite completa de pruebas (unitarias + integración +
cobertura JaCoCo) y se complementa con un propuesto.


\begin{cajaejemplo}{Ejercicio resuelto 16.1 --- Suite completa de tests para servicio de usuarios}
\textbf{Enunciado.} Implementar y ejecutar pruebas unitarias y de
integración sobre un servicio de gestión de usuarios en Spring Boot,
medir cobertura con JaCoCo y producir el reporte HTML.

\textbf{Test unitario con Mockito (\texttt{UsuarioServiceTest.java}):}
\begin{lstlisting}[language=Java,caption={Test unitario con Mockito},label=lst:16.8]
@ExtendWith(MockitoExtension.class)                         // extension JUnit 5
class UsuarioServiceTest {

  @Mock UsuarioRepository repo;                             // crea un mock con Mockito
  @Mock PasswordEncoder encoder;                            // crea un mock con Mockito
  @InjectMocks UsuarioService service;                      // inyecta los mocks en la clase bajo prueba

  @Test                                                     // metodo de test JUnit
  void crear_conDatosValidos_persisteYRetorna() {
    // Arrange
    Usuario nuevo = new Usuario("ana@uteq.edu.ec", "Pass123!");
    // configura comportamiento del mock
    // configura comportamiento del mock
    when(encoder.encode("Pass123!")).thenReturn("$2a$10$hash");
    when(repo.save(any())).thenAnswer(inv -> {              // configura comportamiento del mock
      Usuario u = inv.getArgument(0);
      u.setId(1L);
      return u;                                             // devuelve el resultado al llamador
    });

    // Act
    Usuario creado = service.crear(nuevo);

    // Assert
    assertEquals(1L, creado.getId());                       // verifica igualdad esperada vs obtenida
    assertEquals("$2a$10$hash", creado.getPasswordHash());  // verifica igualdad esperada vs obtenida
    verify(repo).save(any());                               // verifica interaccion con el mock
  }

  @Test                                                     // metodo de test JUnit
  void crear_conEmailDuplicado_lanzaExcepcion() {
    // configura comportamiento del mock
    // configura comportamiento del mock
    when(repo.existsByEmail("ana@uteq.edu.ec")).thenReturn(true);

    assertThrows(EmailDuplicadoException.class,             // verifica que se lance una excepcion
      () -> service.crear(new Usuario("ana@uteq.edu.ec", "Pass123!")));
    verify(repo, never()).save(any());                      // verifica interaccion con el mock
  }
}
\end{lstlisting}

\textbf{Test de integración con MockMvc (\texttt{UsuarioControllerIT.java}):}
\begin{lstlisting}[language=Java,caption={Integración con MockMvc},label=lst:16.9]
@SpringBootTest                                             // levanta el contexto Spring para test
@AutoConfigureMockMvc                                       // configura MockMvc automaticamente
@Testcontainers                                             // metodo de test JUnit
class UsuarioControllerIT {

  @Container                                                // contenedor Docker para tests (Testcontainers)
  static PostgreSQLContainer<?> postgres =
    new PostgreSQLContainer<>("postgres:16-alpine");

  @DynamicPropertySource
  static void props(DynamicPropertyRegistry r) {
    r.add("spring.datasource.url", postgres::getJdbcUrl);
    r.add("spring.datasource.username", postgres::getUsername);
    r.add("spring.datasource.password", postgres::getPassword);
  }

  @Autowired MockMvc mvc;                                   // inyeccion de dependencias

  @Test                                                     // metodo de test JUnit
  void POST_usuarios_conDatosValidos_retorna201() throws Exception {
    String body = """
      {"email":"juan@uteq.edu.ec","password":"Segura123!"}""";

    mvc.perform(post("/api/v1/usuarios")
        .contentType(MediaType.APPLICATION_JSON)
        .content(body))
      .andExpect(status().isCreated())
      .andExpect(jsonPath("$.id").isNumber())
      .andExpect(jsonPath("$.email").value("juan@uteq.edu.ec"));
  }
}
\end{lstlisting}

\textbf{Ejecución} \texttt{mvn clean test jacoco:report}:

\begin{center}
\fbox{\begin{minipage}{0.95\textwidth}
\vspace{0.15cm}\footnotesize\ttfamily
\hspace{0.2cm}[INFO] Running com.uteq.UsuarioServiceTest\\
\hspace{0.2cm}[INFO] Tests run: 8, Failures: 0, Errors: 0, Skipped: 0\\[0.15cm]
\hspace{0.2cm}[INFO] Running com.uteq.UsuarioControllerIT\\
\hspace{0.2cm}[INFO] Tests run: 5, Failures: 0, Errors: 0, Skipped: 0\\[0.15cm]
\hspace{0.2cm}[INFO] Results:\\
\hspace{0.2cm}[INFO] Tests run: 13, Failures: 0, Errors: 0, Skipped: 0\\
\hspace{0.2cm}[INFO] BUILD SUCCESS\\
\hspace{0.2cm}[INFO] Total time: 18.342 s\\
\vspace{0.1cm}
\end{minipage}}
\end{center}

\textbf{Reporte HTML de JaCoCo} (\texttt{target/site/jacoco/index.html}):

\begin{center}
\fbox{\begin{minipage}{0.95\textwidth}
\vspace{0.15cm}\footnotesize\centering
\textbf{JaCoCo Coverage Report --- UsuariosApp v1.0.0}\\[0.2cm]
\begin{tabular}{|l|c|c|c|}
\hline
\textbf{Elemento} & \textbf{Líneas \%} & \textbf{Branches \%} & \textbf{Métodos \%} \\
\hline
com.uteq.usuarios.service & 92\% & 88\% & 100\% \\
com.uteq.usuarios.controller & 85\% & 80\% & 100\% \\
com.uteq.usuarios.security & 78\% & 75\% & 90\% \\
\hline
\textbf{Total} & \textbf{86\%} & \textbf{82\%} & \textbf{97\%} \\
\hline
\end{tabular}\\[0.1cm]
\end{minipage}}
\end{center}

\textbf{Discusión.} La pirámide de pruebas se materializa: 8 tests
unitarios rápidos (sin Spring) y 5 de integración con
testcontainers. La cobertura del 86\% supera el umbral $\geq 70\%$
exigido para la entrega del PFC. Los \emph{branches} (caminos
condicionales) son siempre menores que las líneas porque algunas
ramas raras (como excepciones de BD) no se cubren
\cite{martin2008clean}.
\end{cajaejemplo}

\begin{cajaejercicio}{Ejercicio propuesto 16.2 --- Pipeline CI con GitHub Actions y Lighthouse}
\textbf{Enunciado.} Configurar un pipeline GitHub Actions que se
ejecute en cada \emph{push} a la rama \texttt{main} y realice:
\begin{enumerate}
\item Build con Maven (\texttt{mvn clean verify}).
\item Tests unitarios e integración (con servicios PostgreSQL y
Redis levantados como servicios del pipeline).
\item Reporte de cobertura JaCoCo subido como artifact.
\item Build de imagen Docker.
\item Pruebas de carga k6 contra el container recién construido
(50 VUs, 30 s).
\item Auditoría Lighthouse del frontend (CLI con \texttt{npx lhci
autorun}), umbral mínimo Performance 80, Accessibility 90.
\item Si todo pasa: push de la imagen a GitHub Container Registry
con tag \texttt{v$\{$run\_number$\}$}.
\end{enumerate}

\textbf{Criterios de evaluación:}
\begin{itemize}
\item El badge de CI en el README muestra estado \emph{passing}.
\item Pull requests no se pueden mergear si el pipeline falla
(\emph{branch protection rule}).
\item El \emph{action} sube reportes (JaCoCo, k6, Lighthouse) como
artifacts navegables en la UI de GitHub.
\item Captura de pantalla del Actions con tres ejecuciones exitosas
en el reporte del PFC.
\end{itemize}

\textbf{Lecturas recomendadas:} \cite{martin2017clean} para
disciplina; \cite{ortega2022securecoding} para integración de
seguridad en CI.
\end{cajaejercicio}

% =====================================================================
%       SEMANA 17: DEFENSA DEL PFC Y EP2
% =====================================================================
\chapter{Defensa del PFC y Evaluación Parcial Segundo Corte}
\label{cap:semana17}

\epigraph{«Saber programar es importante; saber explicar lo que has
programado es lo que separa a un técnico de un ingeniero.»}{---
\textit{Adaptación de un consejo común en revisiones académicas}}

\section*{Naturaleza de esta semana}
La semana 17 cierra el segundo corte de la asignatura. Tiene dos
componentes evaluativos principales:

\begin{enumerate}
\item \textbf{Defensa oral del PFC v1.0.0}: cada equipo presenta y
defiende su sistema durante 30--40 minutos. Cada integrante debe
responder preguntas técnicas sobre cualquier parte del sistema, no
solo sobre la suya.
\item \textbf{Evaluación Parcial Segundo Corte}: examen escrito que
cubre las Unidades III y IV (semanas 10--16).
\end{enumerate}

\section{Cómo preparar la defensa oral del PFC}

La defensa oral del PFC es la culminación del curso. Cómo
se prepara la presentación es tan importante como qué contiene:
20 minutos pueden hacer brillar o hundir un proyecto que costó 14
semanas.


\subsection{Estructura recomendada de la presentación}

La presentación de defensa sigue una estructura clásica que
distribuye los 20 minutos en bloques temáticos. Respetarla evita
omitir contenido crítico y mantiene al tribunal enganchado.


\begin{cajadefinicion}{Estructura óptima 30 minutos --- 8 bloques}
\begin{enumerate}
\item \textbf{Portada y presentación del equipo} (1 min): nombres,
roles asumidos en el proyecto.
\item \textbf{Problema y contexto} (3 min): cuál es el problema real
que resuelve el sistema, a quién afecta, datos cuantitativos del
contexto ecuatoriano si aplica.
\item \textbf{Objetivos y alcance} (2 min): objetivo general, 3--5
específicos. Lo que está dentro y fuera del alcance.
\item \textbf{Arquitectura} (5 min): diagrama C4 nivel 2, ADRs
clave, justificación de las decisiones tecnológicas más importantes.
\item \textbf{Demo en vivo} (10 min): \texttt{docker compose up},
login, navegación, CRUD principal, Swagger UI, pruebas Postman,
evidencia de OWASP.
\item \textbf{Resultados de pruebas} (3 min): JaCoCo, k6, Lighthouse
con scores reales medidos.
\item \textbf{Reflexión y trabajo futuro} (3 min): qué se aprendió,
qué se haría distinto, propuestas concretas de mejora.
\item \textbf{Preguntas del tribunal} (5--10 min): cualquier integrante
puede ser interrogado sobre cualquier parte.
\end{enumerate}
\end{cajadefinicion}

\subsection{Las 20 preguntas más frecuentes del tribunal}

El tribunal hace preguntas con patrones reconocibles. La
lista siguiente reúne las 20 preguntas más frecuentes con la
respuesta esperada.


\begin{cajabuenas}{Banco de preguntas que el equipo debe poder responder}
\begin{enumerate}
\item ¿Por qué eligieron \emph{esa} base de datos y no otra?
\item ¿Cómo previene su sistema la inyección SQL?
\item ¿Qué pasa si el usuario X intenta acceder al recurso del
usuario Y?
\item ¿Qué hace su sistema cuando expira el JWT?
\item ¿Por qué cookies HttpOnly y no \texttt{localStorage}?
\item ¿Cómo escala el sistema si reciben 1000 usuarios concurrentes?
\item ¿Qué pasa si Redis se cae?
\item ¿Cómo invalidan el cache cuando cambia un dato?
\item ¿Cuál es el peor cuello de botella de su sistema?
\item ¿Cómo aseguran que los DTOs no expongan campos sensibles?
\item Muestren un test que falle si yo introduzco un bug obvio.
\item ¿Cómo hacen el deploy en producción real (no Docker Compose)?
\item ¿Cómo gestionan secretos (passwords, JWT secret)?
\item ¿Por qué eligieron Spring Boot sobre Quarkus o Micronaut?
\item ¿Cómo validan los datos de entrada en la API?
\item ¿Cómo manejan transacciones que cruzan varios servicios?
\item ¿Qué pasa si dos usuarios editan el mismo registro a la vez?
\item ¿Han verificado los 10 OWASP del Top 10?
\item Si tuvieran que reescribir el sistema desde cero, ¿qué
cambiarían?
\item ¿Cuántas líneas de código tiene el proyecto y cómo se reparten?
\end{enumerate}
\end{cajabuenas}

\subsection{Errores comunes en defensas anteriores}

Los errores que arruinan defensas no son técnicos: son de
comunicación, gestión del tiempo y actitud. La lista siguiente los
recoge para que el estudiante los evite.


\begin{cajaadvertencia}{Errores que cuestan calificación}
\begin{itemize}
\item \textbf{Decir «no sé» sin proponer cómo averiguarlo.} Nadie
sabe todo; se valora más decir «no estoy seguro, pero lo verificaría
ejecutando esto\dots».
\item \textbf{Desconocer la propia arquitectura}: «yo solo trabajé en
el frontend» NO es excusa; cualquier integrante debe poder responder
sobre cualquier parte.
\item \textbf{Demo no preparada}: aparecer al laboratorio sin haber
probado la demo en una computadora distinta de la del desarrollador.
\item \textbf{Slides con código inservible}: capturas de IDE
ilegibles. Mejor reescribir el código en la slide con sintaxis
resaltada.
\item \textbf{Tiempo mal distribuido}: 25 minutos de problema y
contexto, 5 de demo. La demo es el corazón.
\item \textbf{Equipo donde solo habla uno}: el tribunal pregunta
explícitamente al silencioso.
\item \textbf{Datos inventados en la presentación}: decir «el sistema
soporta 10 mil usuarios concurrentes» sin haberlo medido.
\end{itemize}
\end{cajaadvertencia}

\section{Estructura de la Evaluación Parcial Segundo Corte}

La EP2 cubre exclusivamente las \textbf{Unidades III y IV} (semanas
10--16). Se compone de cuatro partes equilibradas:

\begin{table}[htbp]
\centering
\caption{Composición de la Evaluación Parcial Segundo Corte.}
\footnotesize
\begin{tabularx}{\textwidth}{lXl}
\toprule
\textbf{Parte} & \textbf{Tipo} & \textbf{Pts}\\
\midrule
A & Selección múltiple y V/F sobre conceptos & 20 \\
B & Respuesta corta (3 líneas máx, definiciones) & 20 \\
C & Análisis de código (identificar errores y vulnerabilidades) & 30 \\
D & Ejercicio práctico integrador (escribir código completo) & 30 \\
\bottomrule
\end{tabularx}
\end{table}

\subsection{Temas garantizados en la EP2}

La EP2 cubre con seguridad ciertos temas que aparecen en
todas las ediciones del examen. Asegurar dominio sobre ellos
garantiza un mínimo razonable.


\begin{cajaadvertencia}{Temario completo EP2}
\textbf{Unidad III (Semanas 10--12):}
\begin{itemize}
\item Gestión de estado: cookies, sesiones, JWT --- diferencias y
cuándo usar cada una.
\item Flags de seguridad de cookies: \texttt{Secure}, \texttt{HttpOnly},
\texttt{SameSite}.
\item Objeto Request: lectura segura de parámetros, manejo de
\texttt{X-Forwarded-For}.
\item Objeto Response: códigos HTTP, cabeceras de seguridad,
ProblemDetails (RFC 7807).
\item PDO con prepared statements; prevención SQL Injection.
\item Patrón Repository.
\item Spring Data JPA (consultas derivadas, JPQL).
\item Problema N+1 y soluciones (\texttt{JOIN FETCH}, \texttt{with},
\texttt{Include}).
\item Migraciones con Flyway/Liquibase.
\item Procedimientos almacenados: cuándo usar.
\item Patrones arquitectónicos: capas, hexagonal, SOA, microservicios,
EDA, P2P.
\item ISO/IEC 25010 (atributos de calidad).
\item Modelo C4 (4 niveles).
\item ADR (Architectural Decision Records).
\end{itemize}

\textbf{Unidad IV (Semanas 13--16):}
\begin{itemize}
\item Escalabilidad vertical vs horizontal.
\item Teorema CAP de Brewer.
\item Algoritmos de balanceo de carga.
\item Patrones de caché: cache-aside, read-through, write-through,
write-behind.
\item Cache-aside con Redis y TTL.
\item Estrategias de invalidación de caché.
\item ATAM (Architecture Tradeoff Analysis Method).
\item Patrón MVC: ciclo de vida de petición Spring.
\item MVC server-side vs API REST + SPA.
\item REST: 6 restricciones de Fielding.
\item Modelo de Madurez de Richardson.
\item Verbos HTTP, idempotencia, seguridad.
\item Convenciones URI RESTful.
\item REST vs SOAP: comparativa, casos de uso.
\item OpenAPI 3.0 (Swagger).
\item Anatomía de un JWT (3 partes).
\item Pruebas: unitarias, integración, E2E (pirámide).
\item Cobertura JaCoCo $\geq$ 70\%.
\item k6 y Lighthouse.
\item OWASP Top 10:2021 (especialmente A01, A02, A03, A07).
\end{itemize}
\end{cajaadvertencia}

\subsection{Ejemplos del tipo de preguntas que aparecen}

Para preparar el examen conviene haber visto el tipo de
preguntas que se formulan. Los ejemplos siguientes son
representativos de las ediciones recientes.


\begin{cajaejemplo}{Muestra de la Parte A (selección múltiple)}
\textbf{Pregunta A1:} ¿Cuál de las siguientes secuencias representa
correctamente las 3 partes de un JWT separadas por puntos?
\begin{itemize}
\item[a)] payload.signature.header
\item[b)] header.signature.payload
\item[c)] header.payload.signature
\item[d)] signature.header.payload
\end{itemize}
\textbf{Respuesta correcta:} (c).

\textbf{Pregunta A2 (V/F):} En una API REST, el método HTTP PUT se
usa para crear un recurso nuevo cuando el identificador aún no existe
en el servidor, mientras que POST actualiza un recurso existente.

\textbf{Respuesta:} \emph{Falso}. PUT reemplaza/actualiza un recurso
identificado por la URL; POST crea un nuevo recurso. Aunque PUT puede
crear si el recurso no existe, su semántica primaria es reemplazar.
\end{cajaejemplo}

\begin{cajaejemplo}{Muestra de la Parte B (respuesta corta)}
\textbf{Pregunta B1:} Explique la diferencia entre sesión y cookie en
el contexto de la gestión de estado en aplicaciones web. Mencione
dónde se almacena cada una.

\textbf{Respuesta esperada:} Una cookie es un fragmento de texto que
se almacena en el navegador del cliente y se envía al servidor en
cada petición. Una sesión es un conjunto de datos asociado a un
usuario que se almacena en el servidor (memoria, Redis, BD); el
cliente solo guarda el ID de sesión en una cookie.
\end{cajaejemplo}

\begin{cajaejemplo}{Muestra de la Parte C (análisis de código)}
\textbf{El siguiente código PHP gestiona el inicio de sesión. Contiene
3 errores. Identifíquelos:}

\begin{lstlisting}[style=php,numbers=left,caption={Ejemplo de PHP},label=lst:17.1]
<?php
session_start();                                            // arranca o reanuda la sesion
$user = $_GET['user'];
$pwd  = $_GET['pwd'];

$sql = "SELECT * FROM users WHERE name='$user' AND pwd='$pwd'";
$res = mysql_query($sql);

if (mysql_num_rows($res) > 0) {                             // condicional
    $_SESSION['logueado'] = true;
    setcookie('SESS_REMEMBER', $user, time()+3600);
    echo "Bienvenido $user!";                               // imprime al output del servidor
}
\end{lstlisting}

\textbf{Errores esperados:}
\begin{enumerate}
\item Credenciales por GET (deberían ser POST).
\item SQL Injection: concatenación directa en línea 6.
\item \texttt{mysql\_*}: extensión obsoleta y eliminada en PHP 7+;
debe usar PDO o mysqli.
\item Password en texto plano (no se hashea con
\texttt{password\_hash}/\texttt{password\_verify}).
\item Cookie \texttt{SESS\_REMEMBER} sin flags
\texttt{HttpOnly}/\texttt{Secure}/\texttt{SameSite}.
\item XSS en \texttt{echo \char34{}Bienvenido \$user!\char34{}}: falta
\texttt{htmlspecialchars}.
\item No se regenera el ID de sesión tras login (anti-fixation).
\end{enumerate}
(Aunque el enunciado pide 3, hay más; con identificar los 3 más
graves se obtiene puntuación máxima.)
\end{cajaejemplo}

\begin{cajaejemplo}{Muestra de la Parte D (ejercicio integrador)}
\textbf{Implemente en el lenguaje y framework de su elección un
endpoint POST /api/v1/envios que:}

\begin{enumerate}
\item Requiera autenticación con JWT en \texttt{Authorization: Bearer
\dots}.
\item Reciba JSON con \texttt{origen}, \texttt{destino},
\texttt{peso\_kg}.
\item Valide que \texttt{origen} y \texttt{destino} no estén vacíos
y que \texttt{peso\_kg} $>$ 0.
\item Persista en la BD usando ORM o prepared statements.
\item Retorne código HTTP 201 con el envío creado o 422 con
\texttt{errors} en caso de fallo de validación.
\end{enumerate}

Indique al inicio del código el framework y el lenguaje utilizados.
\end{cajaejemplo}

\section{Preparación efectiva para la EP2}

Tener clara la estrategia de preparación durante los siete
días previos al examen y la estrategia durante las dos horas del
mismo es tan importante como el contenido.


\subsection{Plan de estudio de 7 días previos al examen}

La caja siguiente propone un plan día por día para los
siete días anteriores a la EP2.


\begin{cajabuenas}{Plan de 7 días}
\begin{description}[leftmargin=1.5cm]
\item[Día -7:] Re-leer las cajas de definición y buenas prácticas
de las semanas 10--16. Marcar las que no se entienden.
\item[Día -6:] Resolver el ejercicio integrador de las semanas 10
(carrito), 11 (CRUD ORM) y 12 (C4) sin mirar el libro.
\item[Día -5:] Resolver los ejercicios de las semanas 13 (cache),
14 (MVC) y 15 (API REST) sin mirar el libro.
\item[Día -4:] Estudiar el OWASP Top 10:2021 hasta poder explicar
cada uno con un ejemplo.
\item[Día -3:] Practicar análisis de código: tomar 5 fragmentos
con vulnerabilidades y encontrar todos los errores.
\item[Día -2:] Repasar los códigos HTTP y los flags de seguridad de
cookies.
\item[Día -1:] Descansar; revisar solo el resumen ejecutivo de cada
semana. NO empezar a estudiar temas nuevos la noche anterior.
\end{description}
\end{cajabuenas}

\subsection{Estrategia durante el examen}

Durante el examen, la estrategia para distribuir el tiempo
y atacar las preguntas en el orden correcto puede cambiar varios
puntos la calificación final.


\begin{enumerate}
\item Llegar 15 minutos antes con cédula.
\item Leer todo el examen en 5 minutos antes de empezar.
\item Distribución sugerida del tiempo (120 min):
\begin{itemize}
\item Parte A: 15 min.
\item Parte B: 20 min.
\item Parte C: 35 min.
\item Parte D: 45 min.
\item Revisión final: 5 min.
\end{itemize}
\item En la Parte D escribir primero un comentario al inicio:
\texttt{// Framework: Spring Boot 3.4 / Lenguaje: Java 21}.
\item Si no recuerda exactamente la sintaxis, hacer pseudocódigo
claro: vale más comprensión correcta que sintaxis perfecta.
\item Verificar que cada respuesta de la Parte B tenga máximo 3
líneas; respuestas largas pueden penalizarse.
\end{enumerate}

\section{Reflexión sobre el segundo corte}

Las semanas 10--17 representan la transición del estudiante de
\emph{programador} a \emph{ingeniero de software}. La diferencia es
sutil pero crucial: el programador escribe código que funciona; el
ingeniero documenta decisiones, mide rendimiento, prevé fallos,
considera la seguridad desde el diseño y entiende que un sistema
existe en un contexto humano y organizacional.

El PFC es el primer proyecto de la carrera donde estos elementos
convergen. La defensa oral es un ensayo de las defensas profesionales
que vendrán: revisiones técnicas con clientes, presentaciones a
inversores, defensa de la tesis de grado.

\section{Contraste bibliográfico de los conceptos clave de la EP2}

La EP2 cubre las Unidades III y IV con base en literatura técnica
estándar. La gestión de estado y el acceso a datos están tratados a
profundidad en \cite{fowler2002patterns} y \cite{kleppmann2017designing}.
Los patrones arquitectónicos se apoyan en \cite{bass2021software} y
\cite{richards2020fundamentals}. El modelo C4 está documentado por
\cite{brown2023c4}. REST y SOAP en \cite{fielding2000rest},
\cite{richardson2007restful}, \cite{alonso2010webservice} y
\cite{richardson2007restful}. JWT en RFC 7519 \cite{jones2015jwt} y
\cite{jones2015jwt}. Las pruebas en \cite{martin2008clean} y
\cite{martin2017clean}. La seguridad en \cite{owasp2021top10},
\cite{verizon2025dbir} y \cite{ortega2022securecoding}. El
estudiante debe estar familiarizado con al menos las ideas
principales de cada una de estas fuentes para responder con solvencia
en la EP2 y la defensa oral.

\section{Ejercicio resuelto adicional con resultado visual}

Como cierre, el siguiente ejercicio resuelto adicional
simula una pregunta del tribunal con la respuesta esperada, y se
complementa con un propuesto: simulacro de defensa entre equipos.


\begin{cajaejemplo}{Ejercicio resuelto 17.1 --- Simulación de pregunta-respuesta de defensa}
\textbf{Enunciado.} El tribunal hace una pregunta clásica:
\emph{«¿Por qué no usaron microservicios en su PFC?»} Construir una
respuesta de un minuto que evidencie razonamiento arquitectónico
maduro.

\textbf{Mala respuesta (a evitar):} \emph{«Porque era más fácil con
monolito.»}

\textbf{Buena respuesta:}
\begin{quote}
\small Decidimos arquitectura de monolito modular con base en tres
criterios. \textbf{Primero}, el equipo tiene 4 personas, y según
\cite{newman2019monolith} los microservicios pagan dividendos a
partir de 10--15 desarrolladores; antes de eso, el overhead operativo
supera al beneficio. \textbf{Segundo}, los requerimientos no presentan
necesidades de escalado independiente por módulo: el carrito y el
catálogo crecen al mismo ritmo. \textbf{Tercero}, modularizamos
internamente con paquetes Maven separados
(\texttt{usuarios}, \texttt{productos}, \texttt{pedidos},
\texttt{notificaciones}), siguiendo la idea de
\cite{auer2021monolithic}: monolito hoy, microservicios mañana si
hace falta. Esto nos da despliegue simple en Docker Compose con
posibilidad de migrar si las métricas lo justifican.

La ADR-002 del repositorio documenta esta decisión con criterios y
alternativas consideradas.
\end{quote}

\textbf{Resultado de esta respuesta} (vista del tribunal mirándose
entre sí):

\begin{center}
\fbox{\begin{minipage}{0.85\textwidth}
\vspace{0.15cm}\small
\hspace{0.3cm}\textbf{[Tribunal --- vocal 1]} \emph{«Bien argumentado.»}\\[0.15cm]
\hspace{0.3cm}\textbf{[Tribunal --- vocal 2]} \emph{«¿Y cómo medirían
si en el futuro toca migrar?»}\\[0.15cm]
\hspace{0.3cm}\textbf{[Estudiante]} \emph{«Monitorearíamos métricas
con Prometheus: si el módulo de pedidos consume {$>$}70\% de CPU
mientras otros están al 10\%, sería señal de candidato a extraer
como microservicio.»}\\[0.15cm]
\hspace{0.3cm}\textbf{[Tribunal]} \emph{«Excelente.»}\\
\vspace{0.1cm}
\end{minipage}}
\end{center}

\textbf{Discusión.} La diferencia entre la mala y la buena respuesta
no es la decisión técnica (en ambos casos es monolito), sino el
razonamiento explícito que la acompaña: criterios, autores
referenciados, alternativas, métricas de revisión futura. Esto es lo
que distingue a un \emph{programador} de un \emph{ingeniero de
software} \cite{martin2017clean}.
\end{cajaejemplo}

\begin{cajaejercicio}{Ejercicio propuesto 17.2 --- Simulacro de defensa entre equipos}
\textbf{Enunciado.} Antes de la defensa real, organizar un simulacro
en clase donde cada equipo defiende su PFC ante otros dos equipos que
hacen las preguntas (15--20 minutos por equipo).

\textbf{Procedimiento:}
\begin{enumerate}
\item Cada equipo prepara su presentación de 20 min siguiendo la
estructura recomendada en la sección \emph{«Cómo preparar la defensa
oral»}.
\item Los equipos evaluadores reciben una rúbrica con criterios:
\begin{itemize}
\item Claridad de la exposición técnica (1--5).
\item Justificación de decisiones arquitectónicas (1--5).
\item Calidad de la demo (1--5).
\item Manejo de preguntas (1--5).
\item Distribución equitativa entre integrantes (1--5).
\end{itemize}
\item Tras la defensa, los equipos evaluadores escriben dos
fortalezas y dos mejoras concretas para el equipo evaluado.
\item Cada equipo entrega un acta con las observaciones recibidas y
explica cómo las incorporará antes de la defensa real.
\end{enumerate}

\textbf{Producto final:} dossier del equipo con: rúbrica completada
por cada equipo evaluador, acta de observaciones recibidas, plan de
mejoras priorizado. Subir al repositorio Git como
\texttt{docs/defensa-simulacro.md}.

\textbf{Beneficio pedagógico.} Estudios sobre evaluación entre pares
\cite{washizaki2024swebok} muestran que el simulacro entre equipos
detecta entre 40\% y 60\% de las debilidades de presentación que el
tribunal real detectaría, dando al equipo la oportunidad de
corregirlas. Las preguntas de pares suelen además ser más duras que
las del tribunal académico.
\end{cajaejercicio}

% =====================================================================
%       SEMANA 18: TUTORÍA DE REFUERZO FINAL Y CIERRE DEL CURSO
% =====================================================================
\chapter{Tutoría de refuerzo final y cierre del curso}
\label{cap:semana18}

\epigraph{«Lo que sabemos es una gota; lo que ignoramos, un océano.
Pero la dirección de la gota es lo que importa.»}{--- \textit{Adaptado de
Isaac Newton}}

\section*{Naturaleza de esta semana}
La semana 18 constituye la \textbf{segunda semana de tutoría de
refuerzo} del curso, paralela a la de la semana 8 pero con un alcance
más amplio: cierra el segundo corte y, a la vez, cierra integralmente
la asignatura. Sus contenidos no se planifican rígidamente
\emph{a priori}: se ajustan dinámicamente a las dificultades observadas
en las semanas 10--17 y, particularmente, a los errores frecuentes de
la \emph{Evaluación Parcial Segundo Corte} de la semana 17. Este
capítulo ofrece un guión estructurado de revisión, ejercicios de
consolidación y un cierre profesional del curso.

Los cuatro propósitos de la semana son:

\begin{enumerate}
\item \textbf{Resolver dificultades} en los conceptos donde el grupo
mostró menor desempeño en la EP2.
\item \textbf{Realimentar individualmente} a cada equipo sobre los
PFC entregados, identificando fortalezas y oportunidades de mejora.
\item \textbf{Consolidar} con ejercicios integradores que cubren
transversalmente las cuatro unidades del curso.
\item \textbf{Realizar el cierre integrador} situando el aprendizaje
dentro de la trayectoria profesional del estudiante en el mercado
laboral 2026.
\end{enumerate}

\section{Mapa integrador de la asignatura}

Antes del repaso intensivo, conviene visualizar
globalmente lo que el estudiante ha construido en 17 semanas: cómo
las cuatro unidades se conectan entre sí y convergen en el PFC.


\subsection{Las cuatro unidades como una historia coherente}

A lo largo de 17 semanas hemos transitado un camino que tiene
arquitectura interna: cada unidad se apoya en la anterior y todas
convergen en el Proyecto Fin de Curso. El siguiente diagrama
sintetiza esa estructura:

\begin{figure}[htbp]
\centering
\begin{tikzpicture}[
  every node/.style={font=\scriptsize,align=center},
  unidad/.style={rectangle,draw=uteqverde,thick,fill=uteqverde!10,
                 rounded corners=4pt,minimum width=3.5cm,
                 minimum height=2cm}]

  \node[unidad] (u1) {\textbf{Unidad I}\\\textbf{Cliente}\\
                      HTML5 + CSS3\\+ JavaScript\\Sem 1--4};
  \node[unidad,right=0.4cm of u1] (u2) {\textbf{Unidad II}\\\textbf{Servidor}\\
                                          PERL + PHP\\+ Java + ASP.NET\\Sem 5--9};
  \node[unidad,below=0.4cm of u1] (u3) {\textbf{Unidad III}\\\textbf{Datos y patrones}\\
                                          Estado + ORM\\+ Arquitectura\\Sem 10--13};
  \node[unidad,below=0.4cm of u2] (u4) {\textbf{Unidad IV}\\\textbf{MVC y servicios}\\
                                          MVC + REST/SOAP\\+ Pruebas\\Sem 14--17};

  \node[circle,draw=uteqdorado,thick,fill=uteqdorado!20,
        minimum size=2cm,below right=-0.3cm and 1cm of u2,
        font=\bfseries\scriptsize,align=center]
        (pfc) {PFC\\v1.0.0};

  \draw[->,thick,>=stealth,uteqverde!70!black]
       (u1) -- (u2);
  \draw[->,thick,>=stealth,uteqverde!70!black]
       (u2) -- (u4);
  \draw[->,thick,>=stealth,uteqverde!70!black]
       (u3) -- (u4);
  \draw[->,thick,>=stealth,uteqverde!70!black]
       (u1) -- (u3);

  \draw[->,thick,>=stealth,uteqdorado!70!black,dashed]
       (u4) -- (pfc);
  \draw[->,thick,>=stealth,uteqdorado!70!black,dashed]
       (u3) -- (pfc);

\end{tikzpicture}
\caption{Mapa integrador del curso: las cuatro unidades convergen en
el Proyecto Fin de Curso. Las flechas verdes representan la dependencia
secuencial entre unidades; las flechas doradas, la integración
explícita hacia el PFC.}
\end{figure}

\subsection{Tabla resumen de tecnologías dominadas}

La tabla siguiente lista todas las tecnologías que el
estudiante debe dominar al completar la asignatura, con el nivel
esperado en cada una.


\begin{table}[htbp]
\centering
\caption{Tecnologías que el estudiante debe dominar al completar la asignatura.}
\footnotesize
\begin{tabularx}{\textwidth}{lXl}
\toprule
\textbf{Categoría} & \textbf{Tecnología} & \textbf{Nivel esperado}\\
\midrule
Cliente HTML & HTML5 semántico, ARIA, validación nativa & Avanzado\\
Cliente CSS & CSS3, Flexbox, Grid, mobile-first, Sass & Avanzado\\
Cliente JS & ES2024, DOM, async/await, fetch, módulos & Avanzado\\
Servidor PHP & PHP 8.3, PDO, sesiones, Slim/Laravel & Intermedio\\
Servidor Java & Spring Boot 3, JPA, Spring Security & Avanzado\\
Servidor C\# & ASP.NET Core 8, EF Core & Básico\\
Servidor Perl & CGI, sintaxis básica & Introductorio\\
Datos & PostgreSQL, MySQL, JDBC, JPA, EF Core, ORM & Intermedio\\
Caché & Redis cache-aside con TTL & Intermedio\\
Frontend SPA & Angular 19+, TypeScript, RxJS & Básico\\
Arquitectura & Capas, hexagonal, REST, microservicios & Intermedio\\
Documentación & C4, ADRs, OpenAPI 3.0 & Intermedio\\
Testing & JUnit 5, Mockito, MockMvc, JaCoCo & Intermedio\\
Pruebas de carga & k6 con scripts JavaScript & Básico\\
Frontend audit & Lighthouse, axe DevTools & Intermedio\\
Seguridad & OWASP Top 10:2021, JWT, HTTPS & Intermedio\\
Containerización & Docker, Docker Compose & Intermedio\\
CI/CD & GitHub Actions, pipelines básicos & Básico\\
Versionado & Git, GitHub, Conventional Commits & Avanzado\\
\bottomrule
\end{tabularx}
\end{table}

\section{Errores comunes detectados en las semanas 10--17}

Antes del repaso teórico de los temas, conviene listar los errores
más frecuentes detectados en los entregables de la segunda mitad del
curso y en la EP2. Cada estudiante debe verificar que ya no comete
ninguno de ellos:

\begin{cajaadvertencia}{Diez errores frecuentes que se debe haber superado}
\begin{enumerate}
\item Olvidar las flags \texttt{HttpOnly}, \texttt{Secure} y
\texttt{SameSite} al crear cookies de sesión o de JWT.
\item Concatenar SQL con datos del usuario en lugar de usar parámetros
preparados (PDO, JDBC, EF Core).
\item Confundir HTTP 401 (\emph{Unauthorized} --- no autenticado) con
403 (\emph{Forbidden} --- autenticado pero sin permiso).
\item Usar PUT para crear y POST para actualizar (cuando es al revés).
\item Almacenar JWT en \texttt{localStorage} en lugar de cookie
\texttt{HttpOnly}, exponiéndolo a robo por XSS.
\item Diseñar URIs con verbos como \texttt{/getProductos} o
\texttt{/crearPedido} en lugar de sustantivos como \texttt{/productos}
con métodos HTTP semánticos.
\item Olvidar el claim \texttt{exp} en el JWT, generando tokens que
nunca expiran.
\item Saltar a microservicios cuando un monolito modular bien
estructurado bastaría (regla \emph{«monolith first»} de Fowler).
\item No invalidar la caché al modificar la BD, generando
inconsistencias visibles para el usuario.
\item Documentar la API solo con un README en lugar de exponer
OpenAPI 3.0 con Swagger UI interactivo.
\end{enumerate}
\end{cajaadvertencia}

\section{Repaso intensivo de los temas más difíciles}

Basado en el análisis estadístico de errores en la EP2, los siguientes
cuatro temas concentran la mayor parte de las pérdidas de puntaje y
requieren refuerzo específico.

\subsection{Tema 1: Diferencia REST vs SOAP}

La diferencia entre REST y SOAP es una de las preguntas más
frecuentes del tribunal y de la EP2. El resumen siguiente captura lo
esencial.


\begin{cajadefinicion}{Resumen ejecutivo --- REST vs SOAP}
\textbf{REST} es un \emph{estilo arquitectónico} sobre HTTP: usa
verbos HTTP semánticamente, recursos como sustantivos en URLs, JSON
como formato dominante. Es \emph{stateless} por restricción.

\textbf{SOAP} es un \emph{protocolo} sobre HTTP/SMTP/otros: define
mensajes XML con \texttt{envelope} y \texttt{body}, contrato formal
con WSDL, soporta WS-* para seguridad y transacciones.

\textbf{Cuándo elegir cada uno}: REST para todo lo nuevo, especialmente
APIs públicas y SPAs. SOAP solo si hay obligación de integrar con
sistemas legados (banca, salud, B2B grande).
\end{cajadefinicion}

\subsection{Tema 2: JWT --- estructura y uso correcto}

JWT aparece en toda API REST moderna del curso. Conocer su
estructura exacta y las reglas de uso correcto es prerrequisito para
defender el PFC con solvencia.


\begin{cajadefinicion}{Resumen ejecutivo --- JWT}
Un JWT tiene \textbf{tres partes separadas por puntos}:
\textbf{header.payload.signature}.

\begin{itemize}
\item Header: algoritmo de firma (HS256, RS256).
\item Payload: claims (sub, iss, exp, iat, jti, custom).
\item Signature: HMAC o RSA del header+payload con secreto.
\end{itemize}

\textbf{Reglas de uso correcto:}
\begin{itemize}
\item Almacenar en cookie \texttt{HttpOnly} \texttt{Secure}
\texttt{SameSite=Strict} (NO en localStorage).
\item Tiempo de expiración corto (1 hora típico).
\item Refresh token separado con expiración mayor (7 días).
\item Blacklist en Redis por JTI para revocación.
\item Validar firma en cada petición.
\end{itemize}
\end{cajadefinicion}

\subsection{Tema 3: ORM y problema N+1}

El problema N+1 sigue siendo la patología de rendimiento
más frecuente en aplicaciones que usan ORM. La caja siguiente lo
sintetiza junto con su solución.


\begin{cajadefinicion}{Resumen ejecutivo --- ORM y N+1}
El ORM mapea objetos a tablas. La trampa principal es el \textbf{N+1}:
al iterar una colección con relaciones, el ORM ejecuta una consulta
adicional por cada elemento.

\textbf{Solución}: \emph{eager loading} explícito.
\begin{itemize}
\item Spring Data: \texttt{@Query(\char34{}SELECT p FROM Producto p JOIN
FETCH p.categoria\char34{})}.
\item EF Core: \texttt{.Include(p => p.Categoria)}.
\item Eloquent: \texttt{Producto::with('categoria')->get()}.
\item Doctrine: \texttt{SELECT p, c FROM Producto p JOIN p.categoria c}.
\end{itemize}
\end{cajadefinicion}

\subsection{Tema 4: Cache-aside con Redis}

Cache-aside con Redis es el patrón de caché dominante en
2026. La caja siguiente sintetiza su mecánica y los TTL típicos.


\begin{cajadefinicion}{Resumen ejecutivo --- cache-aside}
El patrón cache-aside funciona así:

\begin{enumerate}
\item App consulta Redis con la \emph{key}.
\item Si \textbf{HIT}: devuelve valor cacheado.
\item Si \textbf{MISS}: consulta BD, guarda en Redis con TTL,
devuelve.
\end{enumerate}

\textbf{Al actualizar la BD}: \emph{invalidar} el cache (borrar la
key), no actualizar la cache directamente. Esto evita inconsistencia
si la operación de BD falla a mitad.

\textbf{TTL típicos:} 60s para listados frecuentes; 5--15 min para
catálogos relativamente estables; 1 hora para datos casi inmutables.
\end{cajadefinicion}

\section{Banco de ejercicios de consolidación}

Los siguientes tres ejercicios integran transversalmente las
competencias adquiridas en las semanas 10--17. Se sugiere al
estudiante elegir al menos uno e implementarlo durante esta semana
de cierre como ejercicio voluntario de consolidación.

\begin{cajaejercicio}{Ejercicio 18.1 --- API REST integral de logística}
\textbf{Objetivo:} diseñar e implementar una API REST completa de
gestión de envíos para una empresa de logística, integrando todos los
elementos clave del segundo corte.

\textbf{IDE:} IntelliJ IDEA Ultimate; \textbf{stack}: Spring Boot 3.4
+ PostgreSQL 16 + Redis 7.

\textbf{Requisitos:}
\begin{itemize}
\item Autenticación JWT con \texttt{accessToken} (1 h) y
\texttt{refreshToken} (7 días) en cookies HttpOnly.
\item Endpoints: \texttt{POST /api/v1/envios},
\texttt{GET /api/v1/envios/\{id\}},
\texttt{GET /api/v1/envios?estado=...\&page=...},
\texttt{PATCH /api/v1/envios/\{id\}/estado}.
\item Validación con \texttt{@Valid} y Bean Validation.
\item Cache Redis con TTL de 60 s en el listado de envíos.
\item Documentación OpenAPI 3.0 generada automáticamente con
springdoc.
\item Manejo de errores con ProblemDetails (RFC 7807).
\item Cobertura JaCoCo $\geq 70\%$.
\item Prueba con Postman: 12 escenarios incluyendo casos de error
(400, 401, 403, 404, 422, 429).
\end{itemize}
\end{cajaejercicio}

\begin{cajaejercicio}{Ejercicio 18.2 --- Migración SOAP a REST}
\textbf{Objetivo:} ejercitar la integración con sistemas legados
y la justificación arquitectónica documentada.

\textbf{Procedimiento:}
\begin{enumerate}
\item Tomar el WSDL de un servicio SOAP heredado (proporcionado por
el docente; alternativamente, un servicio público como el de tipos de
cambio del Banco Central del Ecuador).
\item Construir una capa de adaptación REST en Spring Boot que
exponga los métodos SOAP como recursos REST con JSON.
\item Implementar caché Redis para reducir la carga sobre el servicio
SOAP origen.
\item Producir una ADR (\emph{Architectural Decision Record}) que
justifique la decisión: por qué se mantiene el servicio SOAP
en lugar de reescribirlo, qué se gana al ofrecer la fachada REST,
qué riesgos asume el equipo.
\item Pruebas Postman comparando latencia con cache fría vs caliente.
\end{enumerate}
\end{cajaejercicio}

\begin{cajaejercicio}{Ejercicio 18.3 --- Auditoría de seguridad del PFC}
\textbf{Objetivo:} aplicar OWASP ZAP en modo \emph{baseline} sobre el
PFC entregado y producir un informe profesional de vulnerabilidades.

\textbf{Procedimiento:}
\begin{enumerate}
\item Ejecutar el PFC localmente con \texttt{docker compose up}.
\item Descargar OWASP ZAP desde \url{https://www.zaproxy.org}.
\item Configurar ZAP en modo \emph{Automated Scan} apuntando a la URL
del frontend Angular y la URL de la API REST.
\item Ejecutar el escaneo (10--30 min según el tamaño del sistema).
\item Producir un informe que incluya:
\begin{itemize}
\item Tabla de vulnerabilidades con severidad (Alta, Media, Baja,
Informativa).
\item Por cada vulnerabilidad: descripción, impacto, recomendación
concreta de corrección y referencia OWASP.
\item Plan de remediación priorizado.
\end{itemize}
\item Subsanar al menos las vulnerabilidades de severidad alta y
re-ejecutar el escaneo para evidenciar el cierre.
\end{enumerate}
\end{cajaejercicio}

\section{Realimentación individual del PFC}

Tras la presentación grupal en la semana 17, cada equipo
recibe realimentación individualizada en la semana 18 sobre su
PFC: qué hicieron muy bien, qué mejorarían.


\subsection{Criterios de excelencia observados}

Los PFC que alcanzaron las calificaciones más altas comparten estas
características:

\begin{itemize}
\item \textbf{Decisiones documentadas}: cada decisión técnica
significativa tiene una ADR justificando la elección.
\item \textbf{Pruebas reales}: cobertura $\geq$ 80\% con tests
significativos, no triviales.
\item \textbf{Demo robusta}: \texttt{docker compose up} funciona en
una computadora limpia, sin trucos.
\item \textbf{Métricas medidas}: scores de Lighthouse y resultados de
k6 reales, no inventados.
\item \textbf{Equipo cohesionado}: cualquier integrante puede
responder cualquier pregunta del tribunal.
\end{itemize}

\subsection{Mejoras frecuentemente sugeridas}

Las observaciones más frecuentes que reciben los equipos
en esta realimentación son recomendaciones de evolución futura,
no defectos.


\begin{itemize}
\item Agregar pruebas de mutación con PIT (Pitest) para evaluar la
calidad real de las pruebas, no solo su cobertura.
\item Configurar autoscaling con Kubernetes en lugar de Docker
Compose para entornos productivos.
\item Implementar OAuth 2.0 / OpenID Connect en lugar de JWT casero
para integración con identidad institucional.
\item Agregar observabilidad: Prometheus + Grafana + Tempo para
métricas, dashboards y traces distribuidos.
\item Implementar feature flags para despliegues progresivos sin
re-build (LaunchDarkly, Unleash).
\item Migrar el frontend a SSR con Angular Universal para mejorar
SEO y la métrica LCP.
\end{itemize}

\section{El curso en perspectiva profesional}

El curso termina, pero la formación profesional sigue. La
sección siguiente sitúa lo aprendido en la perspectiva de la
carrera del estudiante en el mercado laboral 2026.


\subsection{Cómo se ve el desarrollo web profesional en 2026}

El campo del desarrollo web ha alcanzado una madurez que hace 10
años parecía lejana. En 2026:

\begin{itemize}
\item El desarrollador full-stack típico domina al menos un framework
backend (Spring Boot, ASP.NET Core, Node.js, Laravel) y uno frontend
(Angular, React, Vue).
\item Containers (Docker) y Kubernetes son habilidades esperadas.
\item TypeScript ha desplazado a JavaScript en proyectos serios.
\item La IA generativa (GitHub Copilot, Claude, ChatGPT) duplicó la
productividad declarada del desarrollador medio.
\item La seguridad ya no es opcional: se exigen auditorías OWASP en
contratos.
\item Los equipos remotos son la norma, con asincronía como cultura.
\end{itemize}

\subsection{Próximos pasos en la formación}

El siguiente paso natural en la malla curricular es la asignatura de
\emph{Despliegue y Operaciones} de sexto nivel, donde se profundiza
en CI/CD, contenedores avanzados, observabilidad y prácticas DevOps.
Adicionalmente, recomendamos al estudiante:

\begin{itemize}
\item \textbf{Profundizar en un framework}: elegir uno (Spring Boot,
.NET, Laravel) y dominarlo a nivel experto.
\item \textbf{Aprender un orquestador}: Kubernetes es la habilidad
más demandada en infraestructura.
\item \textbf{Bases de datos avanzadas}: índices, optimización,
sharding, replicación.
\item \textbf{Diseño de sistemas distribuidos}: leer
\emph{Designing Data-Intensive Applications} de Martin Kleppmann.
\item \textbf{Patrones de microservicios}: Sam Newman, Chris
Richardson.
\item \textbf{Seguridad ofensiva}: TryHackMe, HackTheBox,
certificaciones (eJPT, OSCP).
\item \textbf{Inglés técnico}: la mayor parte del conocimiento
publicado y los empleos remotos están en inglés.
\end{itemize}

\subsection{Certificaciones que aportan valor profesional en 2026}

Para complementar la formación con credenciales reconocidas
por el mercado, hay certificaciones de alto valor profesional que el
estudiante puede considerar tras completar el curso.


\begin{table}[htbp]
\centering
\caption{Certificaciones de alto valor profesional para el desarrollador web 2026.}
\footnotesize
\begin{tabularx}{\textwidth}{lXl}
\toprule
\textbf{Categoría} & \textbf{Certificación} & \textbf{Valor mercado}\\
\midrule
Cloud & AWS Certified Developer Associate & Alto\\
Cloud & Microsoft Azure Developer (AZ-204) & Alto\\
Cloud & Google Professional Cloud Developer & Alto\\
Containers & Certified Kubernetes Application Developer (CKAD) & Muy alto\\
Java & Oracle Certified Professional Java SE 21 & Medio\\
.NET & MS Certified: .NET Developer & Medio\\
Spring & VMware Spring Professional & Alto\\
Seguridad & CompTIA Security+ & Alto\\
Seguridad & Certified Ethical Hacker (CEH) & Medio\\
Datos & MongoDB Certified Developer & Medio\\
\bottomrule
\end{tabularx}
\end{table}

\section{Conclusión del curso}

El curso \emph{Aplicaciones Web} ha cubierto el ciclo completo del
desarrollo web profesional: desde el HTML semántico que un navegador
renderiza hasta la arquitectura distribuida con caché Redis y
balanceadores Nginx que sirve a miles de usuarios. El estudiante que
ha completado las 18 semanas posee:

\begin{itemize}
\item Capacidad para construir interfaces web semánticas, accesibles
y responsivas conforme al W3C y WCAG 2.1.
\item Dominio operativo de cuatro lenguajes y plataformas servidor:
PERL/CGI, PHP, Java/Spring Boot y ASP.NET Core.
\item Diseño e implementación de APIs REST documentadas con OpenAPI,
autenticadas con JWT y respaldadas por bases de datos relacionales
con acceso vía ORM.
\item Razonamiento arquitectónico: capas, MVC, microservicios, EDA,
P2P y modelo C4.
\item Conciencia de seguridad: OWASP Top 10, controles defensivos,
cifrado, gestión de identidad.
\item Experiencia de proyecto fin de curso \emph{end-to-end}: desde
la planificación (Entrega 1A en semana 5) hasta el despliegue final
(semana 16) y la defensa oral (semana 17).
\end{itemize}

Las decisiones técnicas que el estudiante ha documentado, los patrones
que ha aplicado y las pruebas que ha escrito en su PFC no son
ejercicios académicos sin más: son \emph{habilidades transferibles
directamente al mercado laboral}. Las empresas locales ---desde
\emph{startups} de Quevedo hasta empresas medianas de Guayaquil y
Quito--- buscan precisamente este perfil. Las plataformas de empleo
remoto (Toptal, Turing, Crossover, Andela LATAM) abren además puertas
internacionales para quienes complementen este conocimiento técnico
con inglés y disciplina profesional.

El docente felicita a cada estudiante por completar el camino y
recuerda que el aprendizaje del desarrollo de software no termina
nunca: empieza otra vez cada lunes con un nuevo \emph{commit}.

\hfill \emph{Buen viaje. Y buen código.}\\
\hfill \emph{Mayo de 2026, Quevedo.}

\section{Cierre administrativo de la asignatura}

Los siguientes puntos cierran administrativamente la
asignatura conforme a la Norma Académica UTEQ.


\begin{itemize}
\item \textbf{Calificación final} = ponderación según Norma Académica
UTEQ de los componentes GA, TA y EV.
\item \textbf{Aprobación}: $\geq$ 7.0 en escala 0--10.
\item \textbf{Recuperación}: examen integral cubriendo las 4 unidades
para estudiantes con calificación final entre 5.0 y 6.9.
\item \textbf{Repetición}: estudiantes con calificación $<$ 5.0 deben
recursar la asignatura el siguiente período.
\item \textbf{Acta de notas}: se publica oficialmente 7 días hábiles
después del cierre de la semana 18.
\end{itemize}

% Bibliografia desde archivo externo .bib
\bibliographystyle{plain}
\bibliography{LibroWebApp}

\end{document}
